% Para utilizar este template siga o tutorial disponível em http://www.biblioteca.ufc.br/wp-content/uploads/2015/09/tutorial-sharelatex.pdf

%%%%%%%%%%%%%%%%%%%%%%%%%%%%%%%%%%%%%%%%%%%%%%%%%%%%%%%
%% Você deve criar uma conta no Overleaf. Depois,    %%
%% vá nas opções no canto esquerdo superior da tela  %%
%% e clique em "Copiar Projeto". Dê um novo nome pa- %%
%% ra o projeto.                                     %%
%%                                                   %%
%% Os principais desenvolvedores deste template são: %%
%%                                                   %%
%%            Ednardo Moreira Rodrigues              %%
%%       (Doutor em Engenharia Elétrica - UFC)       %%
%%                      &                            %%
%%            Alan Batista de Oliveira               %%
%%           (Engenheiro Eletricista - UFC)          %%
%%                                                   %%
%% Consultoria Bibliotecária                         %%
%%                                                   %%
%%  Versão 2016 - ShareLaTeX:                        %% 
%%                                                   %%
%% - Francisco Edvander Pires Santos;                %%
%% - Juliana Soares Lima;                            %%
%% - Izabel Lima dos Santos;                         %%
%% - Kalline Yasmin Soares Feitosa;                  %%
%% - Eliene Maria Vieira de Moura.                   %%
%%
%%  Versão 2019 - Overleaf:
%%  
%%  Biblioteca de Ciências Humanas: 
%% - Francisco Edvander Pires Santos;                %%
%% - Juliana Soares Lima;                            %%
%% - Eliene Maria Vieira de Moura;                   %%
%% - Edmundo Moreira de Sousa Filho.                 %%
%%                                                   %%
%% Biblioteca da FEAAC:                              %%
%% - Izabel Lima dos Santos;                         %%
%% - Kalline Yasmin Soares Feitosa;                  %%
%% - Kleber Lima dos Santos.                         %%
%%                                                   %%
%%  Biblioteca do Curso de Física:                   %%
%% - Aline Rodrigues de Lima Mendes;                 %%
%% - Maria de Jesus Silva dos Santos.                %%
%%                                                   %%
%%  Biblioteca Central do Campus do Pici:            %%
%% - Raquel da Silva Nascimento.                     %%
%%                                                   %%
%% Colaboradores                                     %%
%%                                                   %%
%% -Andrei Bosco Bezerra Torres                      %% 
%% (Professor - Sistemas e Mídias Digitais -         %%
%% Instituto Universidade Virtual - UFC)             %%
%% Tiago ALves Lima                                  %% 
%% (Aluno de Mestrado em Eng. Elétrica)              %%
%%                                                   %%
%% Grande parte do trabalho foi adaptado do template %%
%% da UECE elaborado por:                            %%
%% Thiago Nascimento  (UECE)                         %%
%% Project available on:                             %%
%% https://github.com/thiagodnf/uecetex2             %%
%%                                                   %%
%% "Dúvidas, esclarecimentos ou sugestões podem ser  %%
%% enviadas para o seguinte e-mail:                  %%
%%                                                   %%
%%             atendimentobch@ufc.br                 %%
%%                                                   %%
%% As últimas atualizações estão descritas no inicio %%
%% do arquivo "README.md".                           %%
%%                                                   %%
%%%%%%%%%%%%%%%%%%%%%%%%%%%%%%%%%%%%%%%%%%%%%%%%%%%%%%%

\documentclass[        
    a4paper,          % Tamanho da folha A4
    12pt,             % Tamanho da fonte 12pt
    chapter=TITLE,    % Todos os capitulos devem ter caixa alta
    section=Title,    % Todas as secoes devem ter caixa alta somente na primeira letra
    subsection=Title, % Todas as subsecoes devem ter caixa alta somente na primeira letra
    oneside,          % Usada para impressao em apenas uma face do papel
    english,          % Hifenizacoes em ingles
    spanish,          % Hifenizacoes em espanhol
    brazil,           % Ultimo idioma eh o idioma padrao do documento
    fleqn,             % Comente esta linha se quiser centralizar as equacoes. Comente também a linha 65 abaixo
]{abntex2}

% Para utilizar este template siga o tutorial disponível em http://www.biblioteca.ufc.br/wp-content/uploads/2015/09/tutorial-sharelatex.pdf

%%%%%%%%%%%%%%%%%%%%%%%%%%%%%%%%%%%%%%%%%%%%%%%%%%%%%%%
%% Você deve criar uma conta no Overleaf. Depois,    %%
%% vá nas opções no canto esquerdo superior da tela  %%
%% e clique em "Copiar Projeto". Dê um novo nome pa- %%
%% ra o projeto.                                     %%
%%                                                   %%
%% Os principais desenvolvedores deste template são: %%
%%                                                   %%
%%            Ednardo Moreira Rodrigues              %%
%%       (Doutor em Engenharia Elétrica - UFC)       %%
%%                      &                            %%
%%            Alan Batista de Oliveira               %%
%%           (Engenheiro Eletricista - UFC)          %%
%%                                                   %%
%% Revisão:                                          %%
%%                                                   %%
%% - Francisco Edvander Pires Santos;                %%
%% - Juliana Soares Lima;                            %%
%% - Izabel Lima dos Santos;                         %%
%% - Kalline Yasmin Soares Feitosa.                  %%
%% - Eliene Maria Vieira de Moura;                   %%
%%                                                   %%
%% Colaboradores                                     %%
%%                                                   %%
%% -Andrei Bosco Bezerra Torres                      %% 
%% (Professor - Sistemas e Mídias Digitais -         %%
%% Instituto Universidade Virtual - UFC)             %%
%% Tiago ALves Lima                                  %% 
%% (Aluno de Mestrado em Eng. Elétrica)              %%
%%                                                   %%
%% Grande parte do trabalho foi adaptado do template %%
%% da UECE elaborado por:                            %%
%% Thiago Nascimento  (UECE)                         %%
%% Project available on:                             %%
%% https://github.com/thiagodnf/uecetex2             %%
%%                                                   %%
%% "Dúvidas, esclarecimentos ou sugestões podem ser  %%
%% enviadas para o seguinte e-mail:                  %%
%%                                                   %%
%%             atendimentobch@ufc.br                 %%
%%                                                   %%
%% As últimas atualizações estão descritas no inicio %%
%% do arquivo "README.md".                           %%
%%                                                   %%
%%%%%%%%%%%%%%%%%%%%%%%%%%%%%%%%%%%%%%%%%%%%%%%%%%%%%%%

% Importações de pacotes
\usepackage[utf8]{inputenc}                         % Acentuação direta
\usepackage[T1]{fontenc}                            % Codificação da fonte em 8 bits
\usepackage{graphicx}                               % Inserir figuras
\usepackage{amsfonts, amssymb, amsmath}             % Fonte e símbolos matemáticos
\usepackage{booktabs}                               % Comandos para tabelas
\usepackage{verbatim}                               % Texto é interpretado como escrito no documento
\usepackage{multirow, array}                        % Múltiplas linhas e colunas em tabelas
\usepackage{indentfirst}                            % Endenta o primeiro parágrafo de cada seção.
\usepackage{listings}                               % Utilizar codigo fonte no documento
\usepackage{xcolor}
\usepackage{microtype}                              % Para melhorias de justificação?
\usepackage[portuguese,ruled,lined]{algorithm2e}    % Escrever algoritmos
\usepackage{algorithmic}                            % Criar Algoritmos  
%\usepackage{float}                                 % Utilizado para criação de floats
\usepackage{amsgen}
\usepackage{lipsum}                                 % Usar a simulação de texto Lorem Ipsum
%\usepackage{titlesec}                              % Permite alterar os títulos do documento
\usepackage{tocloft}                                % Permite alterar a formatação do Sumário
\usepackage{etoolbox}                               % Usado para alterar a fonte da Section no Sumário
\usepackage[nogroupskip,nonumberlist]{glossaries}   % Permite fazer o glossario

\usepackage[font=singlespacing,format=plain,justification=justified,skip=0pt,singlelinecheck = false]{caption}            % Altera o comportamento da tag caption

\usepackage[alf, abnt-emphasize=bf, recuo=0cm, abnt-etal-cite=2, abnt-etal-list=0, abnt-etal-text=it]{abntex2cite}  % Citações padrão ABNT
%\usepackage[bottom]{footmisc}                      % Mantém as notas de rodapé sempre na mesma posição
%\usepackage{times}                                 % Usa a fonte Times
%%%%%%%%%%%%%%%%%%% AVISO %%%%%%%%%%%%%%%%%%%%%%%%%%%%%%%%%%%%%%%%
%descomente as duas linhas abaixo para alterar o texto de Times New Roman para Arial:

%\usepackage{helvet}
%\renewcommand{\familydefault}{\sfdefault}  % Usa a fonte Arial              
%%%%%%%%%%%%%%%%%%%%%%%%%%%%%%%%%%%%%%%%%%%%%%%%%%%%%%%%%%%%%%%%%%

\usepackage{mathptmx}         % Usa a fonte Times New Roman			%\usepackage{lmodern}         % Usa a fonte Latin Modern
%\usepackage{subfig}          % Posicionamento de figuras
%\usepackage{scalefnt}        % Permite redimensionar tamanho da fonte
%\usepackage{color, colortbl} % Comandos de cores
%\usepackage{lscape}          % Permite páginas em modo "paisagem"
%\usepackage{ae, aecompl}     % Fontes de alta qualidade
%\usepackage{picinpar}        % Dispor imagens em parágrafos
%\usepackage{latexsym}        % Símbolos matemáticos
%\usepackage{upgreek}         % Fonte letras gregas
\usepackage{appendix}         % Gerar o apendice no final do documento
\usepackage{paracol}          % Criar paragrafos sem identacao
\usepackage{lib/ufctex}	      % Biblioteca com as normas da UFC para trabalhos academicos
\usepackage{pdfpages}         % Incluir pdf no documento
\usepackage{amsmath}          % Usar equacoes matematicas

\makeglossaries % Organiza e gera a lista de abreviaturas, simbolos e glossario
\makeindex      % Gera o Indice do documento         


% -------------------------------------

\usepackage{tikz}
\usetikzlibrary{positioning,shapes,fit}
\usepackage{xspace}


\setlength{\mathindent}{0pt} %Complementa o alinhamento de equações para totalmente a esquerda.

%%%%%%%%%%%%%%%%%%%%%%%%%%%%%%%%%%%%%%%%%%%%%%%%%%%%%
%%                     ATENCAO                     %%
%%%%%%%%%%%%%%%%%%%%%%%%%%%%%%%%%%%%%%%%%%%%%%%%%%%%%
%  Qual e o nivel do trabalho academico que voce esta 
% escrevendo? Retire o simbolo "%" apenas de um dos 
% quatro topicos abaixo refente ao nível do seu traba
% -lho.

\trabalhoacademico{tccgraduacao}
% \trabalhoacademico{tccespecializacao}
%\trabalhoacademico{dissertacao}
%\trabalhoacademico{tese}

%%%%%%%%%%%%%%%%%%%%%%%%%%%%%%%%%%%%%%%%%%%%%%%%%%%%%

% Define se o trabalho e uma qualificacao
% Coloque 'nao' para versao final do trabalho

\ehqualificacao{nao}

% Remove as bordas vermelhas e verdes do PDF gerado
% Coloque 'sim' pare remover

\removerbordasdohyperlink{sim} 

% Adiciona a cor Azul a todos os hyperlinks

\cordohyperlink{nao}

\usepackage{xcolor}
\usepackage[table]{xcolor}

%%%%%%%%%%%%%%%%%%%%%%%%%%%%%%%%%%%%%%%%%%%%%%%%%%%%%
%%         Informacao sobre a instituicao          %%
%%%%%%%%%%%%%%%%%%%%%%%%%%%%%%%%%%%%%%%%%%%%%%%%%%%%%

\ies{Universidade Federal do Ceará}
\iessigla{UFC}
\centro{Campus de Quixadá}
\departamento{Bacharelado em Inteligência Artificial}

%%%%%%%%%%%%%%%%%%%%%%%%%%%%%%%%%%%%%%%%%%%%%%%%%%%%%
%%        Informacao para TCC de Graduacao         %%
%%%%%%%%%%%%%%%%%%%%%%%%%%%%%%%%%%%%%%%%%%%%%%%%%%%%%

\graduacaoem{Inteligência Artificial}
\habilitacao{bacharel} % Ou licenciado(a)

% AVISO: Caso necessario alterar o texto de apresenta-
% cao da Especializacao, ir a pasta "lib", arquivo 
% "ufctex.sty" na linha 502.


%%%%%%%%%%%%%%%%%%%%%%%%%%%%%%%%%%%%%%%%%%%%%%%%%%%%%
%%     Informacao para TCC de Especializacao       %%
%%%%%%%%%%%%%%%%%%%%%%%%%%%%%%%%%%%%%%%%%%%%%%%%%%%%%

\especializacaoem{Yyyyyyyyy}

% AVISO: Caso necessario alterar o texto de apresenta-
% cao da Especializacao, ir a pasta "lib", arquivo 
% "ufctex.sty" na linha 507.

%%%%%%%%%%%%%%%%%%%%%%%%%%%%%%%%%%%%%%%%%%%%%%%%%%%%%
%%         Informacao para Dissertacao             %%
%%%%%%%%%%%%%%%%%%%%%%%%%%%%%%%%%%%%%%%%%%%%%%%%%%%%%

\programamestrado{Programa de Pós-Graduação em Xxxxxxx}
\nomedomestrado{Mestrado Acadêmico em Xxxxxxx}
\mestreem{Engenharia Xxxxxx}
\areadeconcentracaomestrado{Engenharia Xxxxxx}

% AVISO: Caso necessario alterar o texto de apresenta-
% cao da dissertacao, ir a pasta "lib", arquivo 
% "ufctex.sty" na linha 511.

%%%%%%%%%%%%%%%%%%%%%%%%%%%%%%%%%%%%%%%%%%%%%%%%%%%%%
%%               Informação para Tese              %%
%%%%%%%%%%%%%%%%%%%%%%%%%%%%%%%%%%%%%%%%%%%%%%%%%%%%%

\programadoutorado{Programa de Pós-Graduação em Xxxxxx}
\nomedodoutorado{Doutorado em Xxxxxxx}
\doutorem{Engenharia Xxxxxx}
\areadeconcentracaodoutorado{Engenharia Xxxxxxx}

% AVISO: Caso necessario alterar o texto de apresenta-
% cao da tese, ir a pasta "lib", arquivo "ufctex.sty" 
% na linha 515.

%%%%%%%%%%%%%%%%%%%%%%%%%%%%%%%%%%%%%%%%%%%%%%%%%%%%%
%%      Informacoes relacionadas ao trabalho       %%
%%%%%%%%%%%%%%%%%%%%%%%%%%%%%%%%%%%%%%%%%%%%%%%%%%%%%

\autor{ }
\titulo{Projeto Pedagógico do Curso}
\data{2025}
\local{Quixadá}

% Exemplo: \dataaprovacao{01 de Janeiro de 2012}
\dataaprovacao{}

%%%%%%%%%%%%%%%%%%%%%%%%%%%%%%%%%%%%%%%%%%%%%%%%%%%%%
%%           Informação sobre o Orientador         %%
%%%%%%%%%%%%%%%%%%%%%%%%%%%%%%%%%%%%%%%%%%%%%%%%%%%%%

\orientador{Prof. Dr. Xxxxxxx Xxxxxx Xxxxxxx}
\orientadories{Universidade Federal do Ceará (UFC)}
\orientadorcentro{Centro de Ciências e Tecnologia (CCT)}
\orientadorfeminino{nao} % Coloque 'sim' se for do sexo feminino

%%%%%%%%%%%%%%%%%%%%%%%%%%%%%%%%%%%%%%%%%%%%%%%%%%%%%
%%          Informação sobre o Coorientador        %%
%%%%%%%%%%%%%%%%%%%%%%%%%%%%%%%%%%%%%%%%%%%%%%%%%%%%%

% Deixe o nome do coorientador em branco para remover do documento

\coorientador{}
\coorientadories{Universidade Coorientador (SIGLA)}
\coorientadorcentro{Centro do Coorientador (SIGLA)}
\coorientadorfeminino{nao} % Coloque 'sim' se for do sexo feminino

%%%%%%%%%%%%%%%%%%%%%%%%%%%%%%%%%%%%%%%%%%%%%%%%%%%%%
%%              Informação sobre a banca           %%
%%%%%%%%%%%%%%%%%%%%%%%%%%%%%%%%%%%%%%%%%%%%%%%%%%%%%

% Atenção! Deixe em branco o nome do membro da banca para remover da folha de aprovacao

% Exemplo de uso:
% \membrodabancadois{Prof. Dr. Fulano de Tal}
% \membrodabancadoisies{Universidade Federal do Ceará - UFC}


\membrodabancadois{Prof. Dr. Xxxxxxx Xxxxxx Xxxxxxx}
\membrodabancadoiscentro{Faculdade de Filosofia Dom Aureliano Matos (FAFIDAM)}
\membrodabancadoisies{Universidade do Membro da Banca Dois (SIGLA)}
\membrodabancatres{Prof. Dr. Xxxxxxx Xxxxxx Xxxxxxx}
\membrodabancatrescentro{Centro de Ciências e Tecnologia (CCT)}
\membrodabancatresies{Universidade do Membro da Banca Três (SIGLA)}
\membrodabancaquatro{Prof. Dr. Xxxxxxx Xxxxxx Xxxxxxx}
\membrodabancaquatrocentro{Centro de Ciências e Tecnologia (CCT)}
\membrodabancaquatroies{Universidade do Membro da Banca Quatro (SIGLA)}
\membrodabancacinco{Prof. Dr. Xxxxxxx Xxxxxx Xxxxxxx}
\membrodabancacincocentro{Teste}
\membrodabancacincoies{Universidade do Membro da Banca Cinco (SIGLA)}
\membrodabancaseis{Prof. Dr. Xxxxxxx Xxxxxx Xxxxxxx}
\membrodabancaseiscentro{}
\membrodabancaseisies{Universidade do Membro da Banca Seis (SIGLA)}


\renewcommand{\imprimircapa}{%
	\begin{capa}%
		% Todos os elementos da capa devem esta centralizado
		\center
		% Todos os elementos da capa devem esta em negrito
		\bfseries

        \includegraphics[scale=0.10]{lib/logo-ufc} \\
        \vspace{-0.4cm}
		\MakeUppercase{\imprimiries} \\
		\MakeUppercase{\imprimircentro} \\
		\MakeUppercase{\imprimirdepartamento} \\
		% \MakeUppercase{Curso de Graduação em \imprimirgraduacaoem} \\

		\vspace{3\onelineskip}

		\vspace{2.0cm}

		\MakeUppercase{\imprimirautor}
		
		\vspace{3.0cm}
		
		\MakeUppercase{\imprimirtitulo}
		
		\vfill
		\normalsize\MakeUppercase\imprimirlocal\par	
		\normalsize\imprimirdata
	\end{capa}	
}

\def\nomedocurso{\imprimirgraduacaoem\xspace}

% \usepackage{array}
% \usepackage{longtable}

\begin{document}	

	% Elementos pré-textuais
        \imprimircapa
	% \imprimircapa
	% \imprimirfolhaderosto{}
	% \imprimirfichacatalografica{1-pre-textuais/ficha-catalografica}
	%\imprimirerrata{elementos-pre-textuais/errata}
	% \imprimirfolhadeaprovacao
	% \imprimirdedicatoria{1-pre-textuais/dedicatoria}
	% \imprimiragradecimentos{1-pre-textuais/agradecimentos}
	% \imprimirepigrafe{1-pre-textuais/epigrafe}
	% \imprimirresumo{1-pre-textuais/resumo}
	% \imprimirabstract{1-pre-textuais/abstract}
	% \renewcommand*\listfigurename{Lista de Figuras} %Se você comentar esta linha o título da lista fica: LISTA DE ILUSTRAÇÕES
	% \imprimirlistadeilustracoes
	% \imprimirlistadetabelas
	%\imprimirlistadequadros
	%\imprimirlistadealgoritmos
	%\imprimirlistadecodigosfonte
	% \imprimirlistadeabreviaturasesiglas
	% \imprimirlistadesimbolos{1-pre-textuais/lista-de-simbolos}   
	% \imprimirsumario
	
	\setcounter{table}{0}% Deixe este comando antes da primeira tabela.
	
	%Elementos textuais
	\textual
	\chapter{APRESENTAÇÃO}
\label{cap:apresentacao}

A Inteligência Artificial (IA) é a subárea da Ciência da Computação que busca desenvolver sistemas e algoritmos capazes de realizar tarefas que, quando realizadas por humanos, requerem inteligência. Ela envolve a criação de máquinas que podem aprender com dados, tomar decisões, resolver problemas complexos e simular comportamentos inteligentes. De acordo com \citeonline{poole2017ai}, a IA é ``o estudo de agentes computacionais inteligentes, ou seja, sistemas que percebem seu ambiente e tomam decisões para alcançar objetivos específicos''. As abordagens da IA tradicionalmente dividem-se entre paradigmas simbólicos (baseados em regras e lógica formal) e subsimbólicos (baseados em redes neurais e outras técnicas conexionistas), sendo muitas vezes complementares. Tais abordagens são aplicadas a problemas como percepção (visão computacional, compreensão de linguagem natural), raciocínio (busca heurística, planejamento, escalonamento, inferência), ação (robótica, agentes inteligentes) e interação (sistemas conversacionais, interfaces adaptativas), muitas vezes combinadas em arquiteturas complexas como  agentes autônomos e sistemas multiagentes \cite{russell2022ia, poole2017ai} 

Nos últimos anos, a IA tornou-se um dos pilares da transformação digital com impactos econômicos, sociais, políticos e culturais. O \textit{AI Index Report} de 2025 \cite{aiindex2025}, produzido pela \textit{Stanford University}, evidencia esse crescimento da área, desde a pesquisa até sua aplicação na indústria, educação, saúde, justiça e políticas públicas \cite{russell2022ia}. Tal avanço refletiu-se no surgimento de cursos de nível superior de formação de profissionais especializados nesta área tais como: o \textit{B.S. in Artificial Intelligence}, ofertado a partir de 2018 pela \textit{Carnegie Mellon University}, EUA \cite{cmu_bsai}; o Bacharelado em Inteligência Artificial, ofertado pela Universidade Federal de Goiás \cite{ufg_bia} a partir de 2020 e; o Bacharelado em Inteligência Artificial, ofertado pela Universidade Federal do Rio Grande do Norte \cite{ufrn_bia} a partir de 2025. 


A Inteligência Artificial já nasceu interdisciplinar. Desde a conferência de \textit{Dartmouth} \cite{mccarthy2006proposal}, em 1956, considerada o marco inaugural da IA, o desenvolvimento da área ocorreu por contribuições da filosofia, matemática, neurociência, psicologia, teoria do controle e cibernética, linguística, estatística, entre outras \cite{russell2022ia}.  A formação em IA, portanto, requer uma sólida base conceitual e prática em áreas como lógica, algoritmos, teoria da computação, estatística, otimização, sistemas de representação do conhecimento, raciocínio automatizado e aprendizagem de máquina. %(Russell \& Norvig, 2022; Poole \& Mackworth, 2017).
Além disso, o(a) bacharel em Inteligência Artificial deve ser capaz de refletir criticamente sobre os fundamentos, impactos e limites da área, compreendendo os aspectos éticos, sociais, legais e filosóficos associados à automação inteligente, ao viés algorítmico, à explicabilidade dos modelos e à regulação do uso da IA \cite{floridi2018ai4people,mittelstadt2016ethics, eu_ai_act_2021}.



%Modelos computacionais fundamentais para a IA surgiram dessas interações interdisciplinares. Por exemplo, o \textit{perceptron} \cite{rosenblatt1958perceptron}, inspirado em modelos neurobiológicos de aprendizado, estabeleceu as bases para as redes neurais artificiais. Décadas depois, a redescoberta do algoritmo de \textit{backpropagation} \cite{rumelhart1986learning} permitiu avanços no treinamento de redes profundas. Ainda, a evolução do poder computacional e do grande volume de dados disponíveis, impulsionaram o desenvolvimento da técnica de aprendizado profundo, revolucionando áreas como visão computacional, processamento de linguagem natural e reconhecimento de fala \cite{lecun2015deep}.
%(LeCun, Bengio & Hinton, 2015; Wang & Benning, 2020). %Marvin Minsky, John McCarthy, Allen Newell e Herbert Simon, figuras centrais na fundação da IA, já defendiam que o estudo da inteligência exigiria métodos e teorias oriundas de múltiplas disciplinas \cite{russell2022ia}.


%(Floridi et al., 2018; Mittelstadt et al., 2016).


Dessa forma, o curso de Bacharelado em Inteligência Artificial da Universidade Federal do Ceará visa formar profissionais aptos a desenvolver e aplicar soluções baseadas em IA em diversos setores da sociedade, com responsabilidade, criatividade e pensamento crítico. Trata-se de uma formação estratégica para o avanço científico, tecnológico e social do país, considerando o papel central que a IA desempenha na promoção da inovação, na competitividade das empresas e na melhoria da qualidade de vida dos cidadãos.

%Nas últimas décadas, a Inteligência Artificial (IA) passou de uma área de interesse puramente acadêmico para uma área de grande interesse econômico e da sociedade em geral como pode ser visto no AHI Stanford University, 2023, \url{https://aiindex.stanford.edu/report/}. Tal crescimento se refletiu no surgimento de vários cursos de nível superior de formação de profissionais mais especializados na área.  

%A IA já nasce multi e interdisciplinar, em 1958. Ela sempre esteve muito ligada à ciência da computação, mas não só. Desde sua origem ela foi sendo influenciada por outras ciências, por exemplo, a psicologia e a filosofia, com os modelos de conhecimento e as arquiteturas intencionais. Recentemente, passou também a influenciar e ser influenciada por outras áreas do conhecimento, por exemplo, a neurociência (hoje se fala em máquinas que imitam o cérebro e máquinas que aumentam o cérebro). 
%Estudos interdisciplinares entre essas áreas renderam à IA os modelos matemáticos como o \textit{perceptron} (Rosenblatt,1958) e as redes neurais com \textit{backpropagation} (Rumelhart, D.E., Hinton, G.E., Williams, R.J, 1986). 
%Ainda, com a evolução do poder de computação, do volume de dados disponíveis e da pesquisa em neurociência e em ciência cognitiva, impulsionaram o 
%desenvolvimento da técnica de aprendizado profundo - esta denominação de profundo refere-se ao fato de que o algoritmo utiliza uma rede neural com grande quantidade de camadas ocultas e neurônios digitais interconectados, 
%processando milhões de parâmetros – uma arquitetura complexa que alavancou o desenvolvimento da IA nos últimos 10 anos (Wang; Benning, 2020).   

%O bacharel em Inteligência Artificial precisa, além de saber como usar a IA para resolver problemas, também precisa pensar sobre a IA, isto é: entender como a IA resolve os problemas, quais são seus riscos e vantagens e quais são seus limites e potencialidades para a sociedade.   

%Tradicionalmente, a IA inclui uma mistura de abordagens simbólicas e subsimbólicas. As soluções que a IA fornece baseiam-se num amplo conjunto de esquemas de representação de conhecimento geral e especializado, mecanismos de resolução de problemas e técnicas de otimização. Essas abordagens lidam com percepção (por exemplo, reconhecimento de fala, compreensão e geração de linguagem natural, visão computacional), resolução de problemas (por exemplo, busca, planejamento automatizado, otimização), ação (por exemplo, robótica, automação de tarefas, controle) e as arquiteturas necessárias para apoiar estes sistemas (por exemplo, agentes únicos ou sistemas multiagentes). O aprendizado de máquina pode ser usado em cada um desses aspectos e pode até ser empregado de ponta a ponta em todos eles. Ou seja, a IA congrega várias tecnologias para resolver problemas.   

%A formação sólida de bacharéis em Inteligência Artificial influenciará decisivamente na melhoria e na evolução do país e da sociedade como um todo, no que se refere ao atendimento das demandas de inovação, na evolução das empresas e dos cidadãos. 

% >> Ver PPC de Eng. de Software:
% >> https://es.quixada.ufc.br/wp-content/uploads/2023/05/PPC-ES-2023.pdf
%
% 3.1. REALIDADE LOCAL
% 3.1.1. CENÁRIO EDUCACIONAL
% 3.1.2. ASPECTOS SOCIOAMBIENTAIS
% 3.1.3. ASPECTOS REGIONAIS
% 3.2. MISSÃO, VISÃO, VALORES E VOCAÇÃO DO CURSO
% 1.1. JUSTIFICATIVA PARA A CRIAÇÃO DO CURSO
% 1.2. PRINCIPAIS DOCUMENTOS QUE SUBSIDIARAM A ELABORAÇÃO DO PPC
% 1.3. ORGANIZAÇÃO DO DOCUMENTO

        \chapter{HISTÓRICO DA UNIVERSIDADE FEDERAL DO CEARÁ (UFC)}
\label{cap:historico-da-ufc}

\def\ufcAnoReferencia{2024\xspace}

\def\ufcCurGraduacao{118\xspace}
\def\ufcCurGradPresencial{110\xspace}
\def\ufcCurGradEAD{8\xspace}
\def\ufcCurEspecializacao{9\xspace}
\def\ufcCurMestrado{85\xspace}
\def\ufcCurDoutorado{56\xspace}
\def\ufcResidMedica{53\xspace}

\def\ufcVagasGrad{6.278\xspace}

\def\ufcPolosEAD{17\xspace}
\def\ufcMatriculasEAD{412\xspace}

\def\ufcDiscentesPos{8.338\xspace}

\def\ufcCartaPatente{22\xspace}
\def\ufcPedidoPatente{45\xspace}
\def\ufcRegSoftware{9\xspace}
\def\ufcPesqCNPQ{336\xspace}
\def\ufcAcoesExtensao{1.148\xspace}
% \def\ufc{}

% Textos do PDI: https://pdi.ufc.br/wp-content/uploads/2025/01/pdi-2023-2027-3a-revisao-16.12.2024.pdf

A Universidade Federal do Ceará (UFC) foi criada em 16 de dezembro de 1954, pela Lei n.º 2.373, e instalada em 25 de junho do ano seguinte. Com
o lema ``O universal pelo regional'', o reitor Antônio Martins Filho, idealizador e fomentador da fundação da instituição, afirmou o compromisso com o progresso sustentável da nação, associando a importância de uma rede de universidades públicas e gratuitas comprometidas com um projeto de desenvolvimento para o Brasil.

Nessa perspectiva, a UFC, instituição federal de ensino superior, estabelecida como autarquia educacional de regime especial, vinculada ao Ministério da Educação, vem desde então formando profissionais de excelência, gerando e difundindo conhecimentos, preservando e divulgando valores éticos, científicos, artísticos e culturais.%, em conformidade com a sua missão institucional.

Constituída inicialmente a partir de faculdades e de escolas de nível superior já existentes, como a Escola de Agronomia do Ceará, a Faculdade de Direito do Ceará, a Faculdade de Medicina do Ceará e a Faculdade de Farmácia e Odontologia do Ceará, a UFC hoje possui outra dimensão. Com a sua expansão e interiorização ao longo dos últimos anos, a UFC está presente em quase todas as regiões do estado do Ceará, por meio de seus 8 (oito) campi, 3 (três) dos quais localizados no município de Fortaleza, sede da UFC: Campus do Benfica, Campus do Pici e Campus do Porangabuçu; e 5 (cinco) dos quais distribuídos no interior do Ceará, nas cidades Sobral, Quixadá, Crateús, Russas e Itapajé.

São, ao todo, 19 unidades acadêmicas, divididas em centros, faculdades, institutos e campi no interior do estado. São eles: I. Centros – Centro de Ciências; Centro de Ciências Agrárias; Centro de Humanidades e Centro de Tecnologia; II. Faculdades – Faculdade de Direito; Faculdade de Economia, Administração, Atuária e Contabilidade; Faculdade de Educação; Faculdade de Farmácia, Odontologia e Enfermagem; Faculdade de Medicina; III. Institutos – Instituto de Ciências do Mar; Instituto de Cultura e Arte; Instituto de Educação Física e Esportes; e Instituto Universidade Virtual; e Instituto de Arquitetura e Urbanismo e Design; IV. Campi no interior do estado – Campus da UFC em Crateús; Campus da UFC em Quixadá; Campus da UFC em Russas; Campus da UFC em Sobral; e Campus da UFC em Itapajé.

A UFC conta ainda com 1 (um) Centro de Estudos em Aquicultura (CEAC / Labomar Eusébio); 4 (quatro) Fazendas Experimentais: Fazenda Experimental Vale do Curu (Pentecoste), Fazenda Raposa (Maracanaú), Sítio São José (Maracanaú) e Fazenda Lavoura Seca (Quixadá); além dos seguintes equipamentos científicos, tecnológicos e culturais: a Seara da Ciência e o Condomínio de Empreendedorismo e Inovação, localizados no Campus do Pici, a Casa de José de Alencar, localizada no Sítio Alagadiço Novo, a Rádio Universitária FM, o Museu de Arte da UFC (MAUC), a Casa Amarela Eusélio Oliveira (CAEO) e o Teatro Universitário Paschoal Carlos Magno, localizados no Campus do Benfica.

Ao todo a universidade disponibiliza \ufcCurGraduacao cursos de graduação, sendo \ufcCurGradPresencial presenciais e \ufcCurGradEAD na modalidade a distância, oferecidos por meio do Programa Universidade Aberta do Brasil (UAB). Anualmente, a UFC oferta mais de 6 mil novas vagas de ingresso para seus cursos de graduação. Em \ufcAnoReferencia, foram \ufcVagasGrad vagas ofertadas. Os cursos a distância, disponibilizados nos graus de licenciatura, bacharelado e tecnólogo, contaram, no ano de \ufcAnoReferencia, com \ufcMatriculasEAD alunos matriculados e \ufcPolosEAD polos de atendimento distribuídos nos municípios cearenses. Os polos possuem salas de aula, biblioteca, laboratórios de informática, laboratórios para práticas de física, química, dentre outros espaços. Eles também dão suporte aos alunos para sanar dúvidas e resolver problemas.\footnote{A abrangência geográfica dos polos de apoio com alunos ativos ou em matrícula institucional pode ser acessada no endereço \url{https://ead.virtual.ufc.br/}.}

Na pós-graduação, no ano de \ufcAnoReferencia, foram ofertados \ufcCurDoutorado cursos de doutorado, \ufcCurMestrado cursos de mestrado, \ufcCurEspecializacao cursos de especialização, contemplando mais de \ufcDiscentesPos discentes matriculados nesse período, além das \ufcResidMedica residências médicas. 

No âmbito da pesquisa, em \ufcAnoReferencia, a UFC atingiu a marca de \ufcCartaPatente cartas-patente concedidas, \ufcPedidoPatente patentes solicitadas e \ufcRegSoftware registros de softwares depositados no Instituto Nacional da Propriedade Industrial (INPI). Ainda na esfera da pesquisa, a UFC contava, em \ufcAnoReferencia, com \ufcPesqCNPQ pesquisadores CNPq. 

Já na extensão universitária, a UFC realizou \ufcAcoesExtensao ações de extensão em \ufcAnoReferencia, em cinco modalidades: Programas, Projetos, Cursos, Eventos e Prestação de Serviços. Estas ações articulam o conhecimento científico do ensino e da pesquisa com as necessidades da comunidade, interagindo e transformando a realidade social.

A Universidade Federal do Ceará tem a inovação como seu quarto pilar, junto ao ensino, pesquisa e extensão. A partir da Resolução n.º 38/CONSUNI, de 18 de agosto de 2017, a UFC instituiu a sua Política de Inovação, a qual dispõe sobre a definição, geração e gestão de direitos relativos à propriedade intelectual e à inovação tecnológica no âmbito da universidade. 

A estrutura do aparelho administrativo da universidade é composta por seus órgãos colegiados deliberativos – o Conselho Universitário (CONSUNI), o Conselho de Ensino, Pesquisa e Extensão (CEPE) e o Conselho de Curadores –, pelos órgãos de assistência direta e de assessoramento ao reitor, pelos órgãos de planejamento e administração, pelos órgãos de atividades específicas e pelos órgãos suplementares.  
\begin{itemize}
    \item  Órgão de assistência direta e imediata ao reitor: Gabinete e Procuradoria Geral;
 \item  Órgãos de assessoramento ao reitor: Ouvidoria Geral; Coordenadoria Geral de Auditoria, Secretaria Executiva da Comissão de Ética e Comissão Permanente de Processo Administrativo Disciplinar;
 \item  Órgãos de planejamento e administração: Pró-Reitoria de Planejamento e Administração, Pró-Reitoria de Gestão de Pessoas, Superintendência de Infraestrutura e Superintendência de Hospitais Universitários;
 \item   Órgãos de atividades específicas: Pró-Reitoria de Graduação, Pró-Reitoria de Pesquisa e Pós-Graduação, Pró-Reitoria de Extensão, Pró-Reitoria de Assistência Estudantil, PróReitoria de Cultura, Pró-Reitoria de Inovação e Relações Interinstitucionais.
 \item  Órgãos suplementares: Biblioteca Universitária, Superintendência de Tecnologia da Informação, Secretaria de Acessibilidade, Secretaria de Governança, Secretaria de Esportes, Escola Integrada de Desenvolvimento e Inovação Acadêmica, Centro de Excelência em Políticas Educacionais, Centro Estratégico de Excelência em Políticas de Água e Secas, Editora UFC, Secretaria de Comunicação e Marketing, Secretaria de Meio Ambiente, Central Analítica, Colégio de Estudos Avançados.
\end{itemize}

O Conselho Universitário (CONSUNI) é o órgão superior deliberativo e consultivo que concebe a política universitária e arbitra em matéria administrativa, inclusive na gestão econômico-financeira. O Conselho de Ensino, Pesquisa e Extensão (CEPE) é o órgão superior deliberativo e consultivo da universidade, em matéria de ensino, pesquisa e extensão. Já o Conselho de Curadores atua na fiscalização econômico-financeira da universidade. Tais conselhos são compostos por servidores docentes e técnico-administrativos, discentes e representantes da comunidade.

De forma análoga, a administração acadêmica na UFC é exercida por intermédio dos órgãos colegiados deliberativos e por órgãos executivos pertencentes às unidades acadêmicas. Os órgãos executivos que compõem a administração acadêmica na UFC são os centros, as faculdades, os institutos e os campi localizados no interior do estado. Já os órgãos colegiados deliberativos responsáveis pela administração acadêmica são os conselhos das unidades acadêmicas, os departamentos acadêmicos (nos casos de centros e faculdades), as coordenações de cursos de graduação e as coordenações de cursos de pós-graduação (quando houver). Nota-se que os institutos (Fortaleza) e os campi  (Sobral, Quixadá, Russas, Crateús e Itapajé) possuem uma estrutura acadêmica diferente da dos centros e faculdades, dispensando a presença dos departamentos acadêmicos.

Com uma trajetória sólida e em constante expansão, a Universidade Federal do Ceará consolida-se como uma das principais instituições de ensino superior do país, destacando-se pelo compromisso com a excelência acadêmica, a produção científica de relevância e o impacto social de suas ações. Sua presença em diferentes regiões do estado, aliada à diversidade de cursos, programas e iniciativas, reflete o papel estratégico que exerce no desenvolvimento regional e nacional.
        \chapter{HISTÓRICO DO CURSO E POLÍTICAS INSTITUCIONAIS NO ÂMBITO DO CURSO}
\label{cap:historico-do-curso}

O Bacharelado em \nomedocurso da Universidade Federal do Ceará (UFC), ofertado no Campus de Quixadá, foi criado com o propósito de formar profissionais altamente capacitados para atuar em uma das áreas mais promissoras e estratégicas da atualidade. Sua criação foi aprovada pelo Conselho Universitário da Universidade Federal do Ceará , por meio da Resolução n$^\circ$ \textcolor{red}{XXX}, de \textcolor{red}{XX} de \textcolor{red}{XXXXXXX} de \textcolor{red}{XXXX}, como parte do compromisso institucional da UFC com a inovação acadêmica e a formação orientada aos desafios contemporâneos.

Vinculado academicamente e administrativamente ao Campus da UFC em Quixadá, o curso se insere em um ambiente que se destaca por sua vocação tecnológica e pela concentração de graduações na área da Computação. A presença do curso reforça o papel estratégico do campus como polo de formação em TI no interior do Ceará, contribuindo para o desenvolvimento científico, econômico e social da região.

A criação do curso está fortemente alinhada às políticas institucionais estabelecidas no Plano de Desenvolvimento Institucional da UFC (PDI 2023–2027), especialmente no que se refere à ampliação da oferta de cursos inovadores, ao estímulo à pesquisa aplicada e ao fortalecimento da extensão universitária. Nesse contexto, o curso promove uma formação interdisciplinar e sólida, integrando ensino, pesquisa e extensão de maneira orgânica e voltada para a realidade social e tecnológica do país.

Entre os objetivos estratégicos destacados no PDI com os quais o curso se articula em seu projeto pedagógico, destacam-se:
\begin{itemize}
    % OE1 Aprimorar a formação discente
    \item {Aprimorar a formação discente}, mediante o uso de metodologias de ensino ativas, projetos interdisciplinares e conteúdos atualizados;

    % OE2 Destacar-se, nacional e internacionalmente, pelo desenvolvimento da ciência, tecnologia, inovação e empreendedorismo.
    \item {Estimular o desenvolvimento da ciência, tecnologia e inovação}, incentivando o envolvimento dos estudantes em projetos de pesquisa e em ecossistemas de inovação;

    % OE3 Fortalecer a extensão universitária na UFC.
    \item {Ampliar e consolidar a extensão universitária}, promovendo a inclusão digital, o desenvolvimento regional e a aplicação prática do conhecimento acadêmico em benefício da sociedade.
    
\end{itemize}

Essas diretrizes são colocadas em prática pelo curso, seguindo programas institucionais que fortalecem a missão institucional, como ``Propostas formativas flexíveis e arranjos curriculares modernos'', ``Inovação tecnológica'' e ``Democratização da extensão universitária''. A incorporação dessas políticas no cotidiano do curso contribui diretamente para a formação de um egresso crítico, tecnicamente preparado e socialmente comprometido para atuar de forma inovadora em diversos contextos.

Dessa forma, o Bacharelado em \nomedocurso representa uma iniciativa estratégica, não apenas por sua aderência a um campo de conhecimento promissor, mas também por reafirmar o compromisso institucional com a excelência acadêmica, a relevância social e com o desenvolvimento sustentável baseado no conhecimento e na inovação.
        \chapter{IDENTIFICAÇÃO DO CURSO}
\label{cap:identificacao-do-curso}

\section*{Nome do curso}
\nomedocurso

\section*{Grau acadêmico do curso}
Bacharelado

\section*{Carga horária total}
3.200 horas

\section*{Duração do curso}
Integralização mínima em 4 anos (8 semestres) e máxima em 6 anos (12 semestres).

\section*{Regime do curso}
Regime semestral, com o ano letivo dividido em dois períodos.

\section*{Turnos previstos para ofertas}
Integral (matutino e vespertino), com disciplinas obrigatórias ofertadas	preferencialmente no turno matutino.

\section*{Ano e semestre de início de funcionamento do curso}
2026.1

\section*{Ato de Autorização}
Resolução n$^\circ$ \textcolor{red}{XXX}/CONSUNI, de \textcolor{red}{XX} de \textcolor{red}{XXXXXXX} de \textcolor{red}{XXXX}.


\section*{Número de vagas oferecidas por semestre/ano}
50 vagas de ingresso, com uma única entrada prevista por ano.

\section*{Processo de ingresso}
Sistema de Seleção Unificada (SISU), com base no desempenho no Exame Nacional do Ensino Médio (ENEM), e por meio de editais institucionais de transferência.

\section*{Titulação conferida em diplomas}
Bacharel(a) em Inteligência Artificial.


        \chapter{PRINCÍPIOS NORTEADORES}
\label{cap:principios-norteadores}

Os princípios norteadores definidos para o curso de \nomedocurso procuram estabelecer um equilíbrio entre as necessidades do mercado e as demandas do indivíduo e da  sociedade. Assim, esses princípios foram estabelecidos em conformidade com os princípios institucionais da UFC, expressos no PDI (Plano de Desenvolvimento Institucional) da instituição. Neste capítulo, são apresentados os princípios norteadores do curso.

\section{ÉTICA E CIDADANIA}

As Diretrizes Curriculares Nacionais para a Educação em Direitos Humanos – Parecer CNE/CP 1/2012 \cite{brasil2012dcnedh} destacam a responsabilidade das IES com a formação de cidadãos éticos, comprometidos com a construção da paz, da defesa dos direitos humanos e dos valores da democracia, além da responsabilidade de gerar conhecimento mundial, visando atender aos atuais desafios dos direitos humanos, como a erradicação da pobreza, do preconceito e da discriminação. Assim, o curso assume a ética e a cidadania como princípios fundamentais que orientam a atuação profissional e social de seus estudantes e egressos. %Assim, o curso defende a ética e a cidadania como norteadoras do comportamento profissional e social de seus alunos e egressos.

\section{RESPEITO ÀS DIFERENÇAS E À DIVERSIDADE HUMANA}

As Diretrizes Curriculares Nacionais para a Educação em Direitos Humanos – Parecer CNE/CP 1/2012 \cite{brasil2012dcnedh} recomendam a transversalidade curricular das temáticas relativas aos direitos humanos. O documento define, como princípios da educação em direitos, dentre outros: a dignidade humana, a igualdade de direitos, o reconhecimento e valorização das diferenças e das diversidades, a democracia na educação e a transversalidade. Assim, as atividades do curso de \nomedocurso pautam-se em combater a indiferença, a discriminação, o preconceito, a injustiça e os rótulos em relação a todo e qualquer indivíduo.

\section{EQUILÍBRIO NAS ATIVIDADES DE ENSINO, PESQUISA E EXTENSÃO}

O curso de \nomedocurso segue o princípio da indissociabilidade entre o ensino, a pesquisa e a extensão, estabelecido no Estatuto da UFC \cite{ufc2025estatuto}. Nas ações de ensino, pesquisa e extensão, os indivíduos devem ser sempre considerados como sujeitos integrados e integradores da sociedade. Além das atividades de ensino, os alunos são incentivados e têm a oportunidade de se envolver em atividades de pesquisa e extensão ao longo do curso. Essas experiências permitem a aplicação dos conhecimentos adquiridos em sala de aula na resolução de problemas reais da sociedade. Tais atividades enriquecem o aprendizado, ampliando a produção de conhecimento para além dos muros da universidade, promovendo a participação coletiva na comunidade e fomentando uma interação permanente com a sociedade.
%Além das atividades de ensino, ao longo do curso, os alunos são estimulados e têm a oportunidade de participar ativamente de projetos de pesquisa e extensão, de modo a aplicarem os conhecimentos adquiridos em sala de aula em problemas reais da sociedade, atividades que potencializam o conhecimento que se produz fora do ambiente universitário, bem como dão oportunidade ao aluno adquirir conhecimentos que vão além dos muros da universidade, estimulando a participação coletiva na comunidade e possibilitando a interação permanente com a sociedade. %(UFC, 2018c)

\section{FLEXIBILIDADE NA ESTRUTURAÇÃO CURRICULAR}

A flexibilização curricular é considerada parte essencial na organização dos projetos pedagógicos dos cursos de graduação. O Plano Nacional de Educação \cite{brasil2014pne} define em seus objetivos que se devem estabelecer, em nível nacional, diretrizes curriculares que assegurem a necessária flexibilidade e diversidade nos programas oferecidos pelas diferentes IES, de forma a melhor atender às necessidades distintas de suas clientelas e às peculiaridades das regiões nas quais estão inseridas.

Ao se construir currículos flexíveis, evidencia-se a importância de uma estrutura curricular que permita incorporar outras formas de aprendizagem e formação presentes na realidade social. Segundo \citeonline{cabralneto2004flexibilizacao}, a flexibilização curricular possibilita ao aluno participar do processo de formação profissional; rompe com o enfoque unicamente disciplinar e sequenciado; cria novos espaços de aprendizagem; busca a articulação entre teoria e prática; possibilita ao aluno ampliar os horizontes do conhecimento e a aquisição de uma visão crítica que lhe permita extrapolar a aptidão específica de seu campo de atuação profissional e propicia a diversidade de experiências aos alunos. %Cabral Neto (2004)

Nesse sentido, a organização curricular do curso de \nomedocurso compreende uma quantidade limitada de pré-requisitos entre suas componentes curriculares, além de permitir que o aluno construa seu percurso próprio no curso, estruturando seu currículo de acordo com suas necessidades e interesses pessoais e profissionais, a partir da escolha entre a grande variedade de disciplinas optativas e livres ofertadas pela instituição. Admite-se assim, com o esforço pela construção de um currículo flexível e abrangente, que o aluno é responsável direto na construção de seu próprio itinerário formativo.

\section{DESENVOLVIMENTO DA CAPACIDADE CRÍTICA E DA PROATIVIDADE DO EDUCANDO}

As atividades de ensino, pesquisa e extensão do curso refletem a preponderância da educação sobre a instrução, ou seja, há uma preocupação com a aprendizagem baseada na construção do saber a partir da experiência, prévia ou induzida, do próprio indivíduo, a despeito da simples passagem de informações unidirecionais do professor para o aluno. Os discentes são constantemente estimulados a desenvolver trabalhos e projetos críticos e criativos em que apresentam suas próprias visões a partir do que foi aprendido e discutido nas aulas, e não apenas reproduções mecânicas dos conhecimentos adquiridos.

O curso procura estimular uma postura empreendedora e proativa, de modo que este seja pensado não apenas em nível operacional, como resolução de problemas, mas a partir de uma prática estratégica de gerência dos projetos. %Seguindo a visão de Nogueira e Portinari (2016), o foco não está unicamente na resolução de problemas práticos, mas sim em compreender os problemas inseridos em contextos complexos, que merecem respostas à altura dessa realidade. \textcolor{red}{aqui não achei a referência pessoal..dissertação de 2025 orientadora nogueira....ver no PPC de ES....MUDAR A REFERÊNCIA}.
%

\section{INTERDISCIPLINARIDADE}

A \nomedocurso possui sua base na computação, porém sua aplicabilidade é virtualmente ilimitada, com a necessidade de o desenvolvedor de \nomedocurso entender sobre a área de aplicação. Ao longo do curso, as disciplinas se complementam e mostram o caráter transversal do conhecimento de \nomedocurso.
%O desenvolvimento de software integra o emprego de dois domínios: 1) o de \nomedocurso ligado à computação e; 2) o domínio onde está inserido o problema que motiva a construção do software. No curso de \nomedocurso, a interdisciplinaridade é assegurada desde a concepção do projeto pedagógico, já que são inseridas disciplinas integradoras para permitir um diálogo mais coeso entre as várias disciplinas. %Existe também a possibilidade de participação de projetos no Núcleo de Práticas em Informática, que serve como um ambiente mais próximo do mercado de trabalho do Engenheiro de Software através de diversos projetos com diferentes domínios, aproximando a Universidade ao “mundo real”. 
Além disso, por ser um campus temático, o aluno tem a possibilidade de interagir com docentes, alunos e disciplinas de áreas correlatas ao seu curso, propiciando um ambiente maior de teoria e prática.

\section{INTEGRAÇÃO ENTRE TEORIA E PRÁTICA}

Conforme apresentado no Parecer nº CNE/CP 009/2001 , a integração entre teoria e prática está alinhada com a concepção da prática como componente curricular, que ``implica vê-la como uma dimensão do conhecimento, que tanto está presente nos cursos de formação, nos momentos em que se trabalha na reflexão sobre a atividade profissional, como durante o estágio nos momentos em que se exercita a atividade profissional'' \cite[p. 23]{cne2001parecer9}. Nessa visão, busca-se superar a ideia de que ``o estágio é o espaço reservado à prática, enquanto, na sala de aula se dá conta da teoria'' \cite[p. 23]{cne2001parecer9}. %(BRASIL, 2001b, p. 23)

Desta forma, as atividades do curso buscam contemplar a integração entre teoria-prática, visando proporcionar ao estudante uma educação baseada na reflexão crítica e no fazer. Ao longo do curso, os alunos são desafiados a desenvolver projetos práticos. %a maior parte convergindo para o contexto das disciplinas de Projeto Integrado, cujo objetivo é integrar a participação de alunos e professores nas diversas disciplinas ofertadas em um mesmo semestre letivo, objetivando-se uma maior contextualização do conteúdo a ser aprendido bem como ressaltando a importância do inter-relacionamento dos saberes e dos profissionais envolvidos.

Outra iniciativa de integração teoria-prática é a realização de atividades formativas ao mesmo tempo transversais e paralelas ao curso, como oficinas, exposições, palestras e debates, estabelecidas a partir de parcerias entre alunos e professores de diversas áreas.

Considerando os elementos em referência, o Projeto Pedagógico do Curso de \nomedocurso busca a consolidação de uma identidade própria, orientado por princípios que compreendem que a formação profissional em \nomedocurso envolve uma prática específica, que pressupõe saberes e competências coerentes. Para isso, é preciso que o currículo seja flexível e possibilite não só a formação de competências técnicas, como também o compromisso da ciência com as transformações sociais.
        \chapter{OBJETIVOS DO CURSO}
\label{cap:objetivos-do-curso}

% Os objetivos do curso, constantes no PPC, estão implementados, considerando
% o perfil profissional do egresso, a estrutura curricular, o contexto educacional, características locais e regionais e novas práticas emergentes no campo do conhecimento relacionado ao curso.

O estudo da IA prepara os alunos para determinar quando uma abordagem de IA é apropriada para um determinado problema, identificar modelos e representações e mecanismos de raciocínio apropriados, implementá-los e avaliá-los em relação ao desempenho e ao seu impacto social mais amplo.   

Segundo a proposta da  UNESCO \cite{unesco2024ai}, o letramento em IA consiste nas seguintes vertentes: letramento em dados, letramento em algoritmos e letramento em modelos.  O uso de dados, de algoritmos e de modelos, estatísticos e/ou simbólicos, permite que as máquinas, além de coletar e processar dados durante tarefas de aprendizado, possam também raciocinar sobre os modelos aprendidos. O letramento em IA pode ser alcançado através dos eixos de conteúdos definidos neste referencial curricular.

Os três componentes do letramento em IA apontam as competências específicas da IA que serão abordadas nesta proposta de referencial curricular, que são: 

\begin{itemize}
    \item percepção/atuação: como as máquinas percebem e atuam no mundo. Agentes inteligentes possuem formas de interagir com humanos - ou com outros sistemas - de maneira mais natural, através de textos, da fala, de imagens ou de movimentos. Desta forma, a coleta de dados é feita de maneira automática; 
    \item representação e raciocínio (máquinas usam representações do mundo - projetadas por desenvolvedores - e as utilizam para raciocinar e tomada de decisão); 
    \item aprendizado de máquina: os computadores conseguem aprender modelos para raciocínio e tomada de decisão a partir de dados. Nos últimos anos a IA tem focado ainda mais em novas formas de representação do mundo, isto é, passaram de modelos projetados pelos desenvolvedores para modelos aprendidos a partir dos dados ou híbridos; e
    \item impactos sociais: a IA tem impactos que assumem formas positivas e negativas na sociedade, distinguindo seu uso e desenvolvimento ético e responsável, incluindo ainda questões legais, confiança e explicabilidade. Pensar sobre a IA ajuda a construir o pensamento crítico dos alunos.   
\end{itemize}

Há uma grande intersecção da IA com a computação e a ciência de dados. Enquanto a mineração de dados desenvolve metodologias para explorar os diferentes tipos de dados coletados nos diferentes ambientes, o \textit{learning analytics} refere-se diretamente à medição, à coleta, à análise e à produção de relatórios e \textit{insights} sobre os contextos, permitindo que se possa melhor compreender e otimizar os cenários estudados. Ou seja, a área de conhecimento de IA possui intersecção com a ciência de dados por meio de conexões cruzadas com a área de conhecimento de gestão de dados.

Da mesma forma, existem objetivos explícitos para desenvolver a literacia básica em IA e o pensamento crítico em todos os estudantes de ciências da computação, dada a amplitude das interconexões entre a IA e outras áreas de conhecimento, na prática. É mais fácil falar sobre o que a IA tem para além da computação ao invés do que ela compartilha. A IA vai além da perspectiva da lógica e dos algoritmos e enfatiza aspectos como um conjunto de representações, nem sempre estruturadas, para contextualizar a resolução de problemas diversos, raciocinar por padrões permitindo o processamento de semânticas e contextos e a capacidade de aprendizado dos algoritmos.   Ou   seja,   da   mesma   forma   que   a   IA   dialoga   com   várias ciências, ela possui intimidade com muitas disciplinas da Computação.  

Esses aspectos têm como objetivo garantir uma formação mais holística, levando o egresso a refletir sobre o mundo, a entender e resolver problemas computacionais aplicados em diversas áreas e sabendo agir de forma consciente, ética, empreendedora e inovadora, contribuindo para a evolução e melhoria da sociedade. Para isso, durante o curso de Inteligência Artificial, é importante que o estudante tenha a oportunidade de desenvolver os seguintes aspectos complementares à sua formação:

\begin{enumerate}
    \item Atuar com diferentes profissionais de diferentes áreas para identificar oportunidades do mercado e atender às necessidades da sociedade, sabendo trabalhar em equipe.
    \item Praticar a interdisciplinaridade para atuar em diferentes domínios de sistemas computacionais.
    \item Ter a capacidade de resiliência frente à evolução acelerada da IA.
    \item Realizar ações empreendedoras na busca de soluções mais eficazes, incluindo novas tecnologias, produtos e serviços.
    \item Aprender de forma contínua e autônoma sobre métodos, instrumentos, tecnologias de infraestrutura e domínios de aplicação da computação, além de se adequar rapidamente às mudanças tecnológicas e aos novos ambientes de trabalho.
    \item Exercitar a inovação em IA, por meio de conhecimentos científicos e tecnológicos que vão além dos necessários para suas aplicações tradicionais  (incluindo a possibilidade de criar o seu próprio negócio).
    \item Participar de intercâmbio e internacionalização da ciência e tecnologia.
    \item Envolver-se em pesquisa científica.
    \item Interagir com empresas, por meio, por exemplo, de estágio, laboratórios-empresa e empresa júnior.
    \item As aplicações de IA podem ter impactos significativos na sociedade, afetando tanto indivíduos como a coletividade. Isto traz a necessidade de que os alunos compreendam as implicações do trabalho em IA para a sociedade e que façam escolhas responsáveis sobre quando e como aplicar técnicas de IA.
\end{enumerate}


        \chapter{PERFIL PROFISSIONAL DO EGRESSO}
\label{cap:perfil-profissional}
% regis

% IACG: O perfil profissional do egresso consta no PPC, está de acordo com as DCN (quando houver), expressa as competências a serem desenvolvidas pelo
% discente e as articula com necessidades locais e regionais, sendo ampliado em função de novas demandas apresentadas pelo mundo do trabalho.

O avanço acelerado da Inteligência Artificial (IA) nas últimas décadas tem impulsionado a criação de formações acadêmicas específicas voltadas à preparação de profissionais aptos a lidar com os desafios técnicos, éticos e sociais associados ao desenvolvimento e à aplicação de sistemas inteligentes. Nesse contexto, os cursos de Bacharelado em Inteligência Artificial têm como objetivo formar egressos com um conjunto de competências e habilidades que vão além da base tradicional da Computação, integrando conhecimentos interdisciplinares e uma visão crítica sobre os impactos da tecnologia.

O perfil específico de egressos dos cursos de Bacharelado em Inteligência Artificial em termos de competências e habilidades específicas está definido a partir dos Referenciais de Formação para Cursos de Bacharelado em Inteligência Artificial \cite{sbc2024}. Esses elementos conferem características distintivas à formação em IA e estão alinhados com o que estabelecem os Artigos 4º e 5º das Diretrizes Curriculares Nacionais (DCN16) para a área de Computação \cite{DCN16}.

O egresso do curso de Bacharelado em Inteligência Artificial será um profissional com formação sólida e interdisciplinar, preparado para atuar com competência técnica, responsabilidade ética e visão crítica sobre os impactos sociais, econômicos e ambientais das tecnologias que desenvolve e utiliza. Espera-se que este profissional:

\begin{itemize}

    \item Possua base consistente em Ciência da Computação, Matemática, Estatística e Ciência de Dados, que o capacite a identificar, modelar e resolver problemas complexos por meio da construção de sistemas computacionais, incluindo soluções inovadoras em contextos diversos como sistemas autônomos, embarcados e distribuídos;

    \item Domine os principais paradigmas, técnicas e ferramentas da Inteligência Artificial, como representação do conhecimento, raciocínio automático, aprendizado de máquina, otimização, mineração de dados e visualização interativa, sendo capaz de selecionar, integrar e aplicar tais métodos de forma eficiente e fundamentada;

    \item Seja capaz de analisar e explorar grandes volumes de dados, extraindo conhecimento relevante que subsidie a tomada de decisões em ambientes organizacionais, industriais, científicos ou sociais;
    
    \item Atue de maneira empreendedora e inovadora na criação de soluções computacionais, considerando aspectos técnicos, econômicos e de viabilidade de mercado, além de reconhecer oportunidades de negócio emergentes no campo da Inteligência Artificial;
    
    \item Seja capaz de criar soluções, individualmente ou em equipe, para problemas complexos; reconhecer o caráter fundamental da inovação e da criatividade e compreender as perspectivas de negócios e oportunidades comerciais relevantes na área de Inteligência Artificial;

    \item Seja capaz de agir de forma consciente na construção e no uso de sistemas de IA, compreendendo seus impactos diretos e indiretos sobre a sociedade e os indivíduos no que tange a questões de privacidade, transparência, equanimidade e preconceito implícito em sistemas de software e bases de dados, entre outras.
    
    \item Compreenda e considere os impactos éticos, sociais e legais da Inteligência Artificial, incluindo questões de privacidade, viés algorítmico, transparência, justiça e segurança, adotando uma postura crítica e consciente no desenvolvimento e uso dessas tecnologias;
    
    \item Demonstre habilidades interpessoais e de trabalho em equipe, sendo capaz de colaborar com profissionais de diferentes áreas do conhecimento em projetos multidisciplinares de forma proativa, ética e comunicativa;
    
    \item Seja capaz de comunicar ideias, projetos e resultados de forma clara e objetiva, tanto na forma escrita quanto oral, adaptando a linguagem ao público-alvo;
    
    \item Mantenha atitude investigativa, curiosidade científica e disposição para o aprendizado contínuo, sendo capaz de acompanhar a evolução tecnológica e de atuar na área de Inteligência Artificial;
    
    \item Enfrente com resiliência os desafios impostos por problemas de alta complexidade, demonstrando capacidade de adaptação, análise crítica e pensamento estratégico para propor soluções robustas e contextualizadas na área de IA.

\end{itemize}

        \chapter{ÁREAS DE ATUAÇÃO DO FUTURO PROFISSIONAL}
\label{cap:areas-de-atuacao}
% regis


O egresso do curso de Bacharelado em Inteligência Artificial estará apto a atuar em uma ampla gama de setores e contextos que demandam soluções inteligentes, baseadas em dados, automação de processos, aprendizado de máquina e raciocínio computacional. A formação proposta pelo curso visa preparar profissionais com sólida base teórica e prática para enfrentar desafios reais por meio do desenvolvimento, aplicação e avaliação de sistemas inteligentes, sempre com responsabilidade ética e compromisso social.

As áreas de atuação para esse profissional são diversas e em constante expansão, destacando-se:

\begin{itemize}
    \item \textbf{Tecnologia da Informação e Comunicação (TIC):} desenvolvimento de sistemas autônomos, softwares inteligentes, assistentes virtuais, soluções de recomendação e personalização, além de aplicações em computação em nuvem e computação embarcada.

    \item \textbf{Indústria e Manufatura:} aplicação de IA em processos de automação, manutenção preditiva, controle de qualidade, robótica industrial e otimização de cadeias produtivas.
    
    \item \textbf{Saúde e Biotecnologia:} uso de técnicas de IA para diagnóstico assistido por computador, análise de imagens médicas, descoberta de fármacos, modelagem de sistemas biológicos e monitoramento inteligente de pacientes.
    
    \item \textbf{Agronegócio e Meio Ambiente:} implementação de soluções para agricultura de precisão, monitoramento ambiental, previsão climática, detecção de pragas, otimização de recursos naturais e sustentabilidade.
    
    \item \textbf{Serviços Financeiros e Seguros:} desenvolvimento de modelos preditivos para análise de risco, detecção de fraudes, algoritmos de investimento, scoring de crédito e automação de processos bancários.
    
    \item \textbf{Transporte e Logística:} aplicação de IA em sistemas de roteamento inteligente, veículos autônomos, gerenciamento de frotas, previsão de demanda e logística urbana.
    
    \item \textbf{Educação e Tecnologias Educacionais:} criação de sistemas adaptativos de aprendizagem, tutores inteligentes, análise de desempenho acadêmico e personalização de conteúdos.
    
    \item \textbf{Setor Público e Políticas Públicas:} uso de IA para apoio à tomada de decisão, melhoria na prestação de serviços públicos, análise de grandes volumes de dados governamentais e formulação de políticas públicas baseadas em evidências.
    
    \item \textbf{Pesquisa Científica e Desenvolvimento:} atuação em centros de pesquisa, universidades e laboratórios de inovação, contribuindo para o avanço do conhecimento e desenvolvimento de novas técnicas, algoritmos e aplicações em IA.
\end{itemize}

Além disso, o egresso poderá empreender, criar startups ou liderar iniciativas inovadoras que envolvam a aplicação ética e eficiente da Inteligência Artificial nos mais diversos setores. Sua formação o habilita também a continuar os estudos em programas de pós-graduação, tanto acadêmicos quanto profissionais, ampliando ainda mais suas possibilidades de atuação.

        \chapter{ESTRUTURA CURRICULAR}
\label{cap:estrutura-curricular}

As Diretrizes Curriculares Nacionais para cursos da área de Computação \cite{cne2001parecer9} e a Referenciais de Formação para o Curso de Bacharelado em Inteligência Artificial \cite{sbc2024} serviram como base para a elaboração e organização do currículo do Bacharelado em \nomedocurso. O projeto pedagógico do curso valoriza o desenvolvimento de competências tecnológicas, a capacidade de aprender a aprender, a adaptabilidade às transformações do mundo contemporâneo e a flexibilidade curricular.

A definição dos conteúdos dos componentes curriculares, das metodologias de ensino e aprendizagem, bem como dos instrumentos de avaliação, está alinhada aos princípios norteadores do curso, aos seus objetivos formativos e ao perfil do egresso. Neste capítulo, evidencia-se como as unidades curriculares articulam-se no interior da integralização curricular, favorecendo o desenvolvimento do pensamento crítico e a integração entre teoria e prática, com ênfase na flexibilidade e no uso de metodologias ativas de aprendizagem.

Destaca-se também, no âmbito do curso de \nomedocurso e do Campus da UFC em Quixadá, o compromisso com a acessibilidade. As ações desenvolvidas estão em sintonia com a atuação da \textit{Secretaria de Acessibilidade UFC Inclui}, setor institucional criado em 2010, que elabora e coordena estratégias para promover a inclusão de pessoas com deficiência em seis dimensões: atitudinal, arquitetônica, comunicacional, instrumental, metodológica e programática. A acessibilidade é compreendida de forma ampla, não se restringindo apenas a aspectos físicos, mas considerando um conjunto de elementos complementares essenciais para garantir um processo educacional inclusivo e equitativo.

Outro aspecto relevante é a articulação do curso com o setor produtivo. Parcerias com empresas da área de Tecnologia da Informação e Comunicação (TIC) possibilitam que a coordenação e o corpo docente mantenham uma escuta ativa sobre as demandas locais. Essa interação contínua subsidia a atualização do currículo e a adequação da formação às exigências do mercado de trabalho e da sociedade.

Adicionalmente, as ações de extensão do curso são estruturadas em consonância com as diretrizes da Resolução CNE/CES nº 7/2018. As ações devem necessariamente promover uma interação dialógica com a sociedade, inserindo os estudantes em projetos junto a comunidades externas para a troca de saberes e o enfrentamento de questões complexas. Essa vivência, integrada à estrutura curricular, é um pilar para a formação cidadã dos discentes, que aplicam seus conhecimentos de modo interprofissional e interdisciplinar. Ao construir e aplicar soluções de Inteligência Artificial em cenários reais, busca-se a produção de mudanças que impactem positivamente a comunidade e a própria instituição, materializando a articulação entre ensino, extensão e pesquisa em um processo pedagógico único e indissociável, conforme os princípios deste projeto.

A proposta curricular foi construída com base nas necessidades regionais de formação de profissionais qualificados para atuar em diversas áreas da informática. Considerando o crescimento do mercado de TIC no estado do Ceará, o projeto pedagógico prioriza a formação cidadã e técnica, voltada à promoção do desenvolvimento regional, à inclusão social e à inovação tecnológica.


\section{Conteúdos curriculares}

O currículo do Curso de \nomedocurso da UFC está estruturado para promover o desenvolvimento integral do perfil profissional do egresso, por meio de uma formação sólida em fundamentos da computação, inteligência artificial, ética e inovação. Está organizado de forma a fomentar o comportamento autodidata, a atualização contínua, a criatividade, a experimentação e a reflexão crítica sobre os impactos sociais e ambientais das tecnologias emergentes.

A proposta curricular valoriza a interdisciplinaridade, o que contribui diretamente para a integração entre ensino, pesquisa e extensão, conforme previsto nas Diretrizes Curriculares Nacionais e nas diretrizes institucionais da UFC. A matriz inclui componentes obrigatórios, optativos e de extensão, com cargas horárias expressas em horas-relógio, respeitando os parâmetros normativos e assegurando um equilíbrio entre teoria, prática e formação cidadã.

Além dos componentes que compõem a integralização curricular, o curso contempla temas transversais voltados às políticas de educação ambiental, de educação em direitos humanos, de educação das relações étnico-raciais e ao ensino de história e cultura afro-brasileira, africana e indígena. Tais conteúdos são inseridos especialmente nas disciplinas integradoras, como os projetos extensionistas, nos componentes de formação humanística como Sociedade, Culturas e Tecnologia e Ética em Inteligência Artificial, em componentes específicos sobre cada temática, como Relações Étnico-Raciais e Africanidades, Educação Ambiental e Educação em Direitos Humanos, bem como em atividades complementares promovidas pelo campus.

A educação ambiental é um eixo de atuação presente no Campus da UFC em Quixadá desde 2013, com iniciativas como palestras temáticas, ações de reflorestamento e visitas técnicas. O curso de Inteligência Artificial integra essas ações e oferece a disciplina de Educação Ambiental como optativa, por meio da qual foram realizadas campanhas de conscientização como “Menos um copo” (2019, 2022), além de concursos, palestras e visitas técnicas. Projetos como o EcoPET, promovido pelo PET-TI, integram os discentes em ações contínuas de preservação e educação ambiental, com uso de tecnologias móveis desenvolvidas pelos próprios estudantes do campus.

As questões étnico-raciais e de direitos humanos emergem tanto nos debates em sala de aula quanto em atividades culturais e científicas promovidas no campus. Um exemplo é a “I Mostra de Cinema Africano – Mamma África”, realizada em parceria com a Casa de Saberes Cego Aderaldo, que buscou fomentar o contato com a produção cultural africana e reflexões sobre diversidade, história e identidade. Essas temáticas também surgem em disciplinas Sociedade, Culturas e Tecnologia e Ética em Inteligência Artificial, e nos componentes voltados à atuação social e comunitária.

%A formação oferecida também valoriza a educação empreendedora e a inovação tecnológica. O currículo inclui a disciplina obrigatória de Empreendedorismo em Inteligência Artificial e estimula a participação discente em programas como o INOVE, voltado à pesquisa aplicada, à geração de soluções tecnológicas e à criação de empreendimentos inovadores com impacto regional.

O curso promove ao longo de seu percurso o contato dos estudantes com o fundamento da área, enfatizando o método e o rigor científico que orientam a produção de conhecimento em diferentes temáticas, como Mineração de Dados, Sistemas Multiagentes, Processamento de Linguagem Natural e Visão Computacional. A disciplina Projeto de Pesquisa Científico-Tecnológico desenvolve formalmente a metodologia de elaboração de projetos de pesquisa, fornecendo aos discentes ferramentas conceituais e técnicas necessárias para a definição de problemas, construção de hipóteses, escolha de métodos adequados e estruturação rigorosa de investigações científicas na área do curso. Essa competência se consolida nas disciplinas de Trabalho de Conclusão de Curso I e II, em que os alunos desenvolvem e apresentam pesquisas com base em problemas atuais e aplicação de técnicas aprendidas ao longo do curso. Essa estrutura foi reforçada com a implementação do Mestrado em Ciência da Computação do campus, aprovado pelo MEC em 2019, o que fortalece o vínculo entre graduação e pós-graduação.

A extensão universitária é formalizada no currículo por meio da curricularização prevista na Resolução CEPE nº 28/2017. Diversas disciplinas contemplam ações extensionistas que envolvem comunidades, instituições públicas, escolas e organizações da sociedade civil, ampliando o impacto social da formação acadêmica e a interação com saberes diversos.

Por fim, a estrutura curricular do curso promove um diálogo entre conteúdos atualizados e inovadores, metodologias acessíveis, bibliografias de qualidade e políticas de inclusão e diversidade. Essa articulação qualifica a formação dos egressos, diferenciando o curso na área da Computação e contribuindo para a formação de profissionais críticos, comprometidos com o desenvolvimento tecnológico e com os desafios sociais contemporâneos.


\section{Unidades curriculares}
\label{sec:unidades-curriculares}

As unidades curriculares (UC) do curso de \nomedocurso são áreas de conhecimento que agrupam componentes curriculares afins. Além de cumprirem uma função administrativa — já que um representante de cada unidade compõe o Colegiado do curso — exercem, sobretudo, uma função pedagógica, constituindo-se como espaços específicos de discussão sobre questões didático-pedagógicas de cada área.

Essas unidades visam formar o futuro profissional em \nomedocurso com uma base sólida e interdisciplinar, capacitando-o para atuar em diferentes subáreas da Inteligência Artificial, sempre com consciência crítica, responsabilidade social e comprometimento ético.

Os componentes curriculares estão organizados em seis Unidades Curriculares que cobrem três grandes áreas: Formação Básica; Formação Tecnológica; e Formação Complementar e Humanística.

A formação básica é subdividida nas seguintes Unidades Curriculares: Formação Básica em Matemática; Formação Básica em Ciência da Computação; Formação Básica em Inteligência Artificial; Formação Básica em Aprendizado de Máquina e Ciência de Dados. 
%
A formação tecnológica e realizada por meio da Unidade Curricular Formação Tecnológica em Inteligência Artificial e Aprendizado de Máquina, que contempla os saberes técnicos e práticos para o desenvolvimento e aplicação de soluções em IA. 
%
A formação complementar e humanística é promovida pela Unidade Curricular de mesmo nome, que tem como foco o desenvolvimento das dimensões ética, social, crítica e comunicacional do futuro profissional em \nomedocurso.

De forma mais específica, recomenda-se o seguinte perfil docente para o curso:

\begin{enumerate}
    \item Os docentes da Formação Básica em Matemática e da Formação Complementar e Humanística devem possuir formação nas áreas específicas das disciplinas que ministram. É desejável que tenham vivência acadêmica ou profissional que lhes permita articular os conteúdos desenvolvidos com sua aplicação na área da Inteligência Artificial.

    \item Para a Formação Básica em Ciência da Computação, recomenda-se que os docentes tenham formação em Computação ou áreas afins, com domínio de fundamentos teóricos e práticos. Experiência na aplicação desses conhecimentos à modelagem de problemas e à construção de soluções computacionais é altamente desejável.

    \item Os professores da Formação Básica em Inteligência Artificial e da Formação Básica em Aprendizado de Máquina e Ciência de Dados devem ter formação e experiência em áreas como Inteligência Artificial, Ciência de Dados, Aprendizado de Máquina, Estatística Aplicada ou áreas correlatas. Espera-se que possam proporcionar uma sólida base conceitual e prática, alinhada às tendências contemporâneas da área.

    \item Os docentes da Formação Tecnológica em Inteligência Artificial e Aprendizado de Máquina devem possuir experiência comprovada em projetos de pesquisa, desenvolvimento e inovação envolvendo técnicas e aplicações de IA, como visão computacional, processamento de linguagem natural, sistemas autônomos, entre outros. É recomendável que atuem também na orientação de trabalhos de conclusão de curso e em projetos.

    \item Na Formação Complementar e Humanística, os docentes devem estar preparados para promover reflexões sobre as implicações éticas, sociais, políticas e econômicas do uso da Inteligência Artificial, contribuindo para uma formação cidadã crítica e responsável. %Deve ainda enfatizar a necessidade da realização de estágio profissional e/ou de trabalhos de conclusão de curso.

\end{enumerate}

% Cada Unidade Curricular elege um representante para o Colegiado do curso, juntamente com seu respectivo suplente, com mandato de três anos, permitida uma recondução. As regras de composição e funcionamento das UC seguem as diretrizes da Resolução nº 07/CEPE de 1994 (UFC, 1994), com alterações da Resolução nº 03/CEPE de 2016 (UFC, 2016a).

A distribuição dos componentes curriculares por Unidade Curricular revela o caráter híbrido e inovador do curso de \nomedocurso, evidenciando sua ênfase em fundamentos sólidos, competências técnicas avançadas e formação humanística.  
A seguir, estão descritas em mais detalhes as Unidades Curriculares definidas para o curso, bem como os componentes curriculares que as compõem.

\subsection{Formação Básica em Matemática}

Esta unidade curricular reúne disciplinas obrigatórias e optativas, de caráter prioritariamente teórico, relacionadas com a base matemática necessária para o entendimento e desenvolvimento de técnicas em Inteligência Artificial. A Tabela \ref{tab:ucfbm} apresenta as disciplinas que compõem a \textit{Unidade Curricular Formação Básica em Matemática}.


\begin{table}[h]
  \caption{Componentes curriculares da Unidade Curricular de Formação Básica em Matemática}
    \centering
    \resizebox{0.95\textwidth}{!}{%
    \footnotesize
    \begin{tabular}{|
        >{\centering\arraybackslash}m{1.5cm}|
        >{\centering\arraybackslash}m{6.5cm}|
        >{\centering\arraybackslash}m{1.5cm}|
        >{\centering\arraybackslash}m{1.7cm}|
        >{\centering\arraybackslash}m{1.5cm}|
        >{\centering\arraybackslash}m{1.8cm}| 
    }
    \hline
    \textbf{Código} & \textbf{Componente} & \textbf{Tipo} & \textbf{Caráter} & \textbf{Regime} & \textbf{Unidade responsável} \\ \hline 
QXD0116 & Álgebra Linear / \textit{Linear Algebra} & Disciplina & Obrigatório & Semestral & Campus Quixadá \\ \hline
QXD0006 & Cálculo Diferencial e Integral I / \textit{Calculus I} & Disciplina & Obrigatório & Semestral & Campus Quixadá \\ \hline
QXD0134 & Cálculo Diferencial e Integral II / \textit{Calculus II} & Disciplina & Optativo & Semestral & Campus Quixadá \\ \hline
QXD0135 & Cálculo Diferencial e Integral III / \textit{Calculus III} & Disciplina & Optativo & Semestral & Campus Quixadá \\ \hline
QXD0179 & Estatística Multivariada / \textit{Multivariate Statistics} & Disciplina & Optativo & Semestral & Campus Quixadá \\ \hline
QXD0017 & Lógica para Computação / \textit{Logic for Computer Science} & Disciplina & Obrigatório & Semestral & Campus Quixadá \\ \hline
QXD0056 & Matemática Básica / \textit{Basic Mathematics} & Disciplina & Obrigatório & Semestral & Campus Quixadá \\ \hline
QXD0120 & Matemática Computacional / \textit{Computational Mathematics} & Disciplina & Optativo & Semestral & Campus Quixadá \\ \hline
QXD0008 & Matemática Discreta / \textit{Discrete Mathematics} & Disciplina & Obrigatório & Semestral & Campus Quixadá \\ \hline
QXD0109 & Pré-Cálculo / \textit{Pre-Calculus} & Disciplina & Obrigatório & Semestral & Campus Quixadá \\ \hline
QXD0112 & Probabilidade e Estatística / \textit{Probability and Statistics} & Disciplina & Obrigatório & Semestral & Campus Quixadá \\ \hline
QXD0152 & Teoria dos Grafos / \textit{Graph Theory} & Disciplina & Optativo & Semestral & Campus Quixadá \\ \hline
<27> & Teoria dos Jogos / \textit{Game Theory} & Disciplina & Optativo & Semestral & Campus Quixadá \\ \hline
<32> & Tópicos Especiais em Matemática Computacional / \textit{Special Topics in Computational Mathematics} & Disciplina & Optativo & Semestral & Campus Quixadá \\ \hline
    \end{tabular}
}
  \label{tab:ucfbm}
\end{table}

\subsection{Formação Básica em Ciência da Computação}

Esta unidade curricular reúne disciplinas obrigatórias e optativas, de caráter prioritariamente teórico-prático, relacionadas com a área de formação básica em Ciência da Computação. A Tabela \ref{tab:ucfbcc} apresenta as disciplinas que compõem a \textit{Unidade Curricular Formação Básica em Ciência da Computação}.

\begin{table}[h!]
    \caption{Componentes curriculares da Unidade Curricular de Formação Básica em Ciência da Computação}
    \centering
    \resizebox{0.95\textwidth}{!}{%
    \footnotesize
    \begin{tabular}{|
        >{\centering\arraybackslash}m{1.5cm}|
        >{\centering\arraybackslash}m{6.5cm}|
        >{\centering\arraybackslash}m{1.5cm}|
        >{\centering\arraybackslash}m{1.7cm}|
        >{\centering\arraybackslash}m{1.5cm}|
        >{\centering\arraybackslash}m{1.8cm}| 
    }
    \hline
    \textbf{Código} & \textbf{Componente} & \textbf{Tipo} & \textbf{Caráter} & \textbf{Regime} & \textbf{Unidade responsável} \\ \hline
QXD0005 & Arquitetura de Computadores / \textit{Computer Architecture} & Disciplina & Optativo & Semestral & Campus Quixadá \\ \hline
<07> & Busca e Otimização / \textit{Search and Optimization} & Disciplina & Obrigatório & Semestral & Campus Quixadá \\ \hline
QXD0025 & Compiladores / \textit{Compilers} & Disciplina & Optativo & Semestral & Campus Quixadá \\ \hline
QXD0079 & Computação em Nuvem / \textit{Cloud Computing} & Disciplina & Optativo & Semestral & Campus Quixadá \\ \hline
QXD0183 & Computação Paralela / \textit{Parallel Computing} & Disciplina & Optativo & Semestral & Campus Quixadá \\ \hline
QXD0153 & Desafios de Programação / \textit{Programming Challenges} & Disciplina & Optativo & Semestral & Campus Quixadá \\ \hline
QXD0276 & Desenvolvimento de Software para Dispositivos Móveis / \textit{Software Development for Mobile Devices} & Disciplina & Optativo & Semestral & Campus Quixadá \\ \hline
QXD0253 & Desenvolvimento de Software para Web / \textit{Software Development for the Web} & Disciplina & Optativo & Semestral & Campus Quixadá \\ \hline
QXD0019 & Engenharia de Software / \textit{Software Engineering} & Disciplina & Obrigatório & Semestral & Campus Quixadá \\ \hline
QXD0010 & Estrutura de Dados / \textit{Data Structure} & Disciplina & Obrigatório & Semestral & Campus Quixadá \\ \hline
QXD0115 & Estrutura de Dados Avançada / \textit{Advanced Data Structure} & Disciplina & Obrigatório & Semestral & Campus Quixadá \\ \hline
QXD0011 & Fundamentos de Banco de Dados / \textit{Database Fundamentals} & Disciplina & Obrigatório & Semestral & Campus Quixadá \\ \hline
QXD0001 & Fundamentos de Programação / \textit{Programming Fundamentals} & Disciplina & Obrigatório & Semestral & Campus Quixadá \\ \hline
QXD0254 & Gerência de Projetos / \textit{Project Management} & Disciplina & Optativo & Semestral & Campus Quixadá \\ \hline
QXD0256 & Interação Humano-Computador / \textit{Computer-Human Interaction} & Disciplina & Optativo & Semestral & Campus Quixadá \\ \hline
QXD0016 & Linguagens de Programação / \textit{Programming Languages} & Disciplina & Optativo & Semestral & Campus Quixadá \\ \hline
QXD0040 & Linguagens Formais e Autômatos / \textit{Formal Languages and Automata} & Disciplina & Optativo & Semestral & Campus Quixadá \\ \hline
QXD0181 & Pesquisa Operacional / \textit{Operations Research} & Disciplina & Optativo & Semestral & Campus Quixadá \\ \hline
QXD0188 & Processamento de Imagens / \textit{Image Processing} & Disciplina & Optativo & Semestral & Campus Quixadá \\ \hline
QXD0114 & Programação Funcional / \textit{Functional Programming} & Disciplina & Optativo & Semestral & Campus Quixadá \\ \hline
QXD0007 & Programação Orientada a Objetos / \textit{Object-Oriented Programming} & Disciplina & Obrigatório & Semestral & Campus Quixadá \\ \hline
    \end{tabular}
    }
    \label{tab:ucfbcc}
\end{table}

\begin{table}[h!]
    % \captionsetup{labelformat=empty}
    % \caption{Componentes curriculares da Unidade Curricular de Formação Básica em Ciência da Computação (\textit{continuação})}
    \centering
    \resizebox{0.95\textwidth}{!}{%
    \footnotesize
    \begin{tabular}{|
        >{\centering\arraybackslash}m{1.5cm}|
        >{\centering\arraybackslash}m{6.5cm}|
        >{\centering\arraybackslash}m{1.5cm}|
        >{\centering\arraybackslash}m{1.7cm}|
        >{\centering\arraybackslash}m{1.5cm}|
        >{\centering\arraybackslash}m{1.8cm}| 
    }
    \hline
    \textbf{Código} & \textbf{Componente} & \textbf{Tipo} & \textbf{Caráter} & \textbf{Regime} & \textbf{Unidade responsável} \\ \hline
QXD0041 & Projeto e Análise de Algoritmos / \textit{Design and Analysis of Algorithms} & Disciplina & Obrigatório & Semestral & Campus Quixadá \\ \hline
QXD0021 & Redes de Computadores / \textit{Computer Networks} & Disciplina & Optativo & Semestral & Campus Quixadá \\ \hline
QXD0069 & Segurança / \textit{Security} & Disciplina & Optativo & Semestral & Campus Quixadá \\ \hline
QXD0043 & Sistemas Distribuídos / \textit{Distributed Systems} & Disciplina & Optativo & Semestral & Campus Quixadá \\ \hline
QXD0013 & Sistemas Operacionais / \textit{Operating Systems} & Disciplina & Optativo & Semestral & Campus Quixadá \\ \hline
QXD0046 & Teoria da Computação / \textit{Theory of Computation} & Disciplina & Optativo & Semestral & Campus Quixadá \\ \hline
    \end{tabular}
    }
\end{table}

\subsection{Formação Básica em Aprendizado de Máquina e Ciência de Dados}

Esta unidade curricular abrange disciplinas que apresentam os princípios estatísticos e computacionais do aprendizado de máquina e da análise de dados, proporcionando uma base para a modelagem e interpretação de dados em diferentes contextos. A Tabela \ref{tab:ucfbamcd} apresenta as disciplinas que compõem a \textit{Unidade Curricular Formação Básica em Aprendizado de Máquina e Ciência de Dados}.

\begin{table}[h]
  \caption{Componentes curriculares da Unidade Curricular de Formação Básica em Aprendizado de Máquina e Ciência de Dados}
    \centering
    \resizebox{0.95\textwidth}{!}{%
    \footnotesize
    \begin{tabular}{|
        >{\centering\arraybackslash}m{1.5cm}|
        >{\centering\arraybackslash}m{6.5cm}|
        >{\centering\arraybackslash}m{1.5cm}|
        >{\centering\arraybackslash}m{1.7cm}|
        >{\centering\arraybackslash}m{1.5cm}|
        >{\centering\arraybackslash}m{1.8cm}| 
    }
    \hline
<21> & Análise Exploratória de Dados / \textit{Exploratory Data Analysis} & Disciplina & Optativo & Semestral & Campus Quixadá \\ \hline
<05> & Aprendizado Não Supervisionado / \textit{Unsupervised Learning} & Disciplina & Obrigatório & Semestral & Campus Quixadá \\ \hline
<12> & Aprendizado por Reforço / \textit{Reinforcement Learning} & Disciplina & Obrigatório & Semestral & Campus Quixadá \\ \hline
<09> & Aprendizado Profundo / \textit{Deep Learning} & Disciplina & Obrigatório & Semestral & Campus Quixadá \\ \hline
<06> & Aprendizado Supervisionado / \textit{Supervised Learning} & Disciplina & Obrigatório & Semestral & Campus Quixadá \\ \hline
<22> & Inteligência Artificial Explicável / \textit{Explainable Artificial Intelligence} & Disciplina & Optativo & Semestral & Campus Quixadá \\ \hline
QXD0178 & Mineração de Dados / \textit{Data Mining} & Disciplina & Obrigatório & Semestral & Campus Quixadá \\ \hline
<10> & Processamento de Dados em Larga Escala / \textit{Large-Scale Data Processing} & Disciplina & Obrigatório & Semestral & Campus Quixadá \\ \hline
QXD0177 & Recuperação de Informação / \textit{Information Retrieval} & Disciplina & Optativo & Semestral & Campus Quixadá \\ \hline
<28> & Tópicos Especiais em Aprendizado de Máquina / \textit{Special Topics in Machine Learning} & Disciplina & Optativo & Semestral & Campus Quixadá \\ \hline
<29> & Tópicos Especiais em Ciência de Dados / \textit{Special Topics in Data Science} & Disciplina & Optativo & Semestral & Campus Quixadá \\ \hline
<33> & Visualização de Dados / \textit{Data Visualization} & Disciplina & Optativo & Semestral & Campus Quixadá \\ \hline
    \end{tabular}
    }
  \label{tab:ucfbamcd}
\end{table}

\subsection{Formação Básica em Inteligência Artificial}

Esta unidade curricular compreende disciplinas que introduzem os fundamentos da Inteligência Artificial, suas técnicas clássicas e modernas, abordando conceitos como busca, raciocínio, representação de conhecimento e tomada de decisão. A Tabela \ref{tab:ucfbia} apresenta as disciplinas que compõem a \textit{Unidade Curricular Formação Básica em Inteligência Artificial}.

\begin{table}[h]
  \caption{Componentes curriculares da Unidade Curricular de Formação Básica em Inteligência Artificial}
    \centering
    \resizebox{0.95\textwidth}{!}{%
    \footnotesize
    \begin{tabular}{|
        >{\centering\arraybackslash}m{1.5cm}|
        >{\centering\arraybackslash}m{6.5cm}|
        >{\centering\arraybackslash}m{1.5cm}|
        >{\centering\arraybackslash}m{1.7cm}|
        >{\centering\arraybackslash}m{1.5cm}|
        >{\centering\arraybackslash}m{1.8cm}| 
    }
    \hline
<03> & Agentes Inteligentes / \textit{Intelligent Agents} & Disciplina & Obrigatório & Semestral & Campus Quixadá \\ \hline
<02> & Introdução à Inteligência Artificial / \textit{Introduction to Artificial Intelligence} & Disciplina & Obrigatório & Semestral & Campus Quixadá \\ \hline
<15> & Planejamento Automatizado / \textit{Automated Planning} & Disciplina & Obrigatório & Semestral & Campus Quixadá \\ \hline
<26> & Raciocínio sob Incerteza / \textit{Reasoning Under Uncertainty} & Disciplina & Optativo & Semestral & Campus Quixadá \\ \hline
<04> & Representação do Conhecimento e Raciocínio / \textit{Knowledge Representation and Reasoning} & Disciplina & Obrigatório & Semestral & Campus Quixadá \\ \hline
QXD0076 & Sistemas Multiagentes / \textit{Multiagent Systems} & Disciplina & Obrigatório & Semestral & Campus Quixadá \\ \hline
<31> & Tópicos Especiais em Inteligência Artificial / \textit{Special Topics in Artificial Intelligence} & Disciplina & Optativo & Semestral & Campus Quixadá \\ \hline
    \end{tabular}
    }
  \label{tab:ucfbia}
\end{table}

\subsection{Formação Tecnológica em Inteligência Artificial e Aprendizado de Máquina}

Esta unidade curricular contempla disciplinas voltadas à aplicação prática e desenvolvimento de sistemas baseados em Inteligência Artificial, com foco em tecnologias emergentes, frameworks computacionais e soluções para problemas reais. A Tabela \ref{tab:uctia} apresenta as disciplinas que compõem a \textit{Unidade Curricular Formação Tecnológica em Inteligência Artificial e Aprendizado de Máquina}.

\begin{table}[h]
  \caption{Componentes curriculares da Unidade Curricular de Formação Tecnológica em Inteligência Artificial e Aprendizado de Máquina}
    \centering
    \resizebox{0.95\textwidth}{!}{%
    \footnotesize
    \begin{tabular}{|
        >{\centering\arraybackslash}m{1.5cm}|
        >{\centering\arraybackslash}m{6.5cm}|
        >{\centering\arraybackslash}m{1.5cm}|
        >{\centering\arraybackslash}m{1.7cm}|
        >{\centering\arraybackslash}m{1.5cm}|
        >{\centering\arraybackslash}m{1.8cm}| 
    }
    \hline
    \textbf{Código} & \textbf{Componente} & \textbf{Tipo} & \textbf{Caráter} & \textbf{Regime} & \textbf{Unidade responsável} \\ \hline
<13> & Inteligência Artificial Generativa / \textit{Generative Artificial Intelligence} & Disciplina & Obrigatório & Semestral & Campus Quixadá \\ \hline
<14> & MLOPs / \textit{MLOps} & Disciplina & Obrigatório & Semestral & Campus Quixadá \\ \hline
<23> & Percepção e Ação Robótica / \textit{Robotic Perception and Action} & Disciplina & Optativo & Semestral & Campus Quixadá \\ \hline
<25> & Processamento de Áudio e Voz / \textit{Audio and Speech Processing} & Disciplina & Optativo & Semestral & Campus Quixadá \\ \hline
<11> & Processamento de Linguagem Natural / \textit{Natural Language Processing} & Disciplina & Obrigatório & Semestral & Campus Quixadá \\ \hline
QXD0182 & Visão Computacional / \textit{Computer Vision} & Disciplina & Obrigatório & Semestral & Campus Quixadá \\ \hline
    \end{tabular}
    }
  \label{tab:uctia}
\end{table}

\subsection{Formação Complementar e Humanística}

Esta unidade curricular reúne componentes voltados ao desenvolvimento de competências éticas, sociais, comunicacionais e interdisciplinares, contribuindo para a formação crítica e cidadã do bacharel em Inteligência Artificial. 
Nesta unidade encontram-se também os componentes curriculares de extensão, como as Atividades de Extensão e os Projetos Extensionistas, que atuam como um eixo articulador do currículo. Eles permitem que os estudantes apliquem os saberes técnicos e teóricos, desenvolvidos nas demais unidades curriculares, como a Formação Básica em Ciência da Computação, Aprendizado de Máquina e a Formação Tecnológica em IA, em contextos sociais reais, promovendo a integração entre a formação tecnológica e a humanística e consolidando a indissociabilidade entre ensino, pesquisa e extensão. 
A Tabela \ref{tab:uccfh} apresenta as disciplinas e atividades que compõem a \textit{Unidade Curricular Formação Complementar e Humanística}.

\begin{table}[h]
  \caption{Componentes curriculares da Unidade Curricular de Formação Complementar e Humanística}
    \centering
    \resizebox{0.95\textwidth}{!}{%
    \footnotesize
    \begin{tabular}{|
        >{\centering\arraybackslash}m{1.5cm}|
        >{\centering\arraybackslash}m{6.5cm}|
        >{\centering\arraybackslash}m{1.5cm}|
        >{\centering\arraybackslash}m{1.7cm}|
        >{\centering\arraybackslash}m{1.5cm}|
        >{\centering\arraybackslash}m{1.8cm}| 
    }
    \hline
    \textbf{Código} & \textbf{Componente} & \textbf{Tipo} & \textbf{Caráter} & \textbf{Regime} & \textbf{Unidade responsável} \\ \hline
<01> & Atividades Complementares / \textit{Complementary Activities} & Atividade & Obrigatório & Semestral & Campus Quixadá \\ \hline
<AE> & Atividades de Extensão / \textit{Extension Activities} & Atividade & Obrigatório & Semestral & Campus Quixadá \\ \hline
QXD0232 & Educação Ambiental / \textit{Environmental Education} & Disciplina & Optativo & Semestral & Campus Quixadá \\ \hline
QXD0245 & Educação em Direitos Humanos / \textit{Human Rights Education} & Disciplina & Optativo & Semestral & Campus Quixadá \\ \hline
QXD0029 & Empreendedorismo / \textit{Enterpreneurship} & Disciplina & Optativo & Semestral & Campus Quixadá \\ \hline
<16> & Estágio Curricular Supervisionado I / \textit{Supervised Internship} & Atividade & Obrigatório & Semestral & Campus Quixadá \\ \hline
QXD0508 & Ética e Legislação / \textit{Ethics and Legislation} & Disciplina & Optativo & Semestral & Campus Quixadá \\ \hline
<08> & Ética em Inteligência Artificial / \textit{Ethics in Artificial Intelligence} & Disciplina & Obrigatório & Semestral & Campus Quixadá \\ \hline
QXD0035 & Inglês Instrumental I / \textit{Instrumental English I} & Disciplina & Optativo & Semestral & Campus Quixadá \\ \hline
QXD0036 & Inglês Instrumental II / \textit{Instrumental English II} & Disciplina & Optativo & Semestral & Campus Quixadá \\ \hline
QXD0113 & Língua Brasileira de Sinais - LIBRAS / \textit{Brazilian Sign Language - Libras} & Disciplina & Optativo & Semestral & Campus Quixadá \\ \hline
<24> & Pesquisa e Inovação / \textit{Research and Innovation} & Disciplina & Optativo & Semestral & Campus Quixadá \\ \hline
QXD0110 & Projeto de Pesquisa Científico-Tecnológico / \textit{Scientific and Technological Research Project} & Disciplina & Obrigatório & Semestral & Campus Quixadá \\ \hline
<19> & Projeto Extensionista I / \textit{Extension Project I} & Disciplina & Obrigatório & Semestral & Campus Quixadá \\ \hline
<20> & Projeto Extensionista II / \textit{Extension Project II} & Disciplina & Obrigatório & Semestral & Campus Quixadá \\ \hline
QXD0126 & Psicologia e Percepção / \textit{Psychology and Perception} & Disciplina & Optativo & Semestral & Campus Quixadá \\ \hline
QXD0246 & Relações Étnico-Raciais e Africanidades / \textit{Ethnic-Racial Relationships and Africanities} & Disciplina & Optativo & Semestral & Campus Quixadá \\ \hline
QXD0162 & Sociedade, Culturas e Tecnologia / \textit{Society, Cultures and Technology} & Disciplina & Obrigatório & Semestral & Campus Quixadá \\ \hline
<30> & Tópicos Especiais em Filosofia, Ética e Humanidades / \textit{Special Topics in Philosophy, Ethics and Humanities} & Disciplina & Optativo & Semestral & Campus Quixadá \\ \hline
<17> & Trabalho de Conclusão de Curso I / \textit{Final Project I} & Atividade & Obrigatório & Semestral & Campus Quixadá \\ \hline
<18> & Trabalho de Conclusão de Curso II / \textit{Final Project II} & Atividade & Obrigatório & Semestral & Campus Quixadá \\ \hline
    \end{tabular}
    }
  \label{tab:uccfh}
\end{table}

\clearpage

\section{Integralização curricular} \label{sec:integralizacao}

A estrutura curricular do curso de \nomedocurso foi elaborada de forma a contemplar os objetivos do curso e atingir o perfil profissional proposto, alicerçado nos seus princípios norteadores. A organização do currículo permite a compreensão, o entendimento e o conhecimento necessários para aplicar e desenvolver modelos e métodos baseados em Inteligência Artificial, utilizando tecnologias e metodologias atuais, assegurando as inter-relações com outras áreas do conhecimento. 

Os componentes curriculares podem ser disciplinas ou atividades. As disciplinas do curso são de dois tipos: obrigatórias e optativas. Entre as optativas, os(as) estudantes podem escolher componentes que fazem parte da integralização curricular do curso e também disciplinas livres. As disciplinas livres são aquelas escolhidas fora do elenco específico de componentes curriculares do curso, podendo ser cursadas em qualquer outro curso da universidade. Assim, qualquer código de componente que não pertença ao rol de obrigatórios e optativos elencados aqui, quando cursado pelo(a) estudante, será integralizado como ``livre''. %Dessa forma, a carga horária de disciplinas livres é, necessariamente, parte da carga horária optativa do curso.
%
Na integralização curricular do curso de \nomedocurso, o estudante poderá contabilizar até 192 horas em disciplinas optativas livres, dentro do total de 448 horas exigidas em disciplinas optativas. 

Entre os componentes curriculares do tipo ``atividade'', estão previstas para o curso as seguintes: Estágio Supervisionado, Trabalho de Conclusão de Curso, Atividades de Extensão e Atividades Complementares. Essas atividades são regulamentadas por manuais específicos, anexos a este documento. 

A Tabela~\ref{tab:organizacao_curricular} apresenta a organização curricular por semestre, com informações sobre carga horária, pré-requisitos, correquisitos e equivalências. A Tabela \ref{tab:disciplinas_optativas} apresenta o rol de disciplinas optativas que compõem a integralização curricular do curso. As disciplinas optativas são uma das formas de trabalhar a flexibilidade do currículo.%, onde o aluno poderá escolher 8 disciplinas dentre as disciplinas listadas e que forem ofertadas durante o percurso do aluno no curso.


\begin{table}[h]
\centering
\caption{{Integralização Curricular}}
\renewcommand{\arraystretch}{1.5}
\resizebox{\textwidth}{!}{%
\small
\begin{tabular}{|
>{\centering\arraybackslash}m{1.5cm}|
>{\centering\arraybackslash}m{10.4cm}|
>{\centering\arraybackslash}m{1.4cm}|
>{\centering\arraybackslash}m{1.4cm}|
>{\centering\arraybackslash}m{1.4cm}|
>{\centering\arraybackslash}m{1.5cm}|
>{\centering\arraybackslash}m{1.4cm}|
% >{\centering\arraybackslash}m{1.4cm}|
>{\centering\arraybackslash}m{1.8cm}|
>{\centering\arraybackslash}m{2.4cm}|
>{\centering\arraybackslash}m{2.4cm}|
}
\hline
\centering{\textbf{Código}} & 
\centering{\textbf{Nome do Componente Curricular}\\ (em português e inglês)} & 
\textbf{Carga Horária Teórica} & 
\textbf{Carga Horária Prática} & 
\textbf{Carga Horária EAD} & 
\textbf{Carga Horária Extensão} & 
\textbf{Carga Horária TOTAL} & 
% \textbf{Carga Horária PCC} & 
\textbf{Pré-requisito(s)} & 
\textbf{Correquisito(s)}&  
\textbf{Equivalência(s)} \\ \hline

\multicolumn{10}{|>{\centering\arraybackslash}c|}{\normalsize\textbf{Semestre 1}} \\ \hline
QXD0001 & Fundamentos de Programação / \textit{Programming Fundamentals} & 32 & 64 & 0 & 0 & 96 &  &  &  \\ \hline
<02> & Introdução à Inteligência Artificial / \textit{Introduction to Artificial Intelligence} & 64 & 0 & 0 & 0 & 64 &  &  &  \\ \hline
QXD0056 & Matemática Básica / \textit{Basic Mathematics} & 64 & 0 & 0 & 0 & 64 &  &  &  \\ \hline
QXD0109 & Pré-Cálculo / \textit{Pre-Calculus} & 32 & 0 & 0 & 0 & 32 &  &  &  \\ \hline
QXD0162 & Sociedade, Culturas e Tecnologia / \textit{Society, Cultures and Technology} & 64 & 0 & 0 & 0 & 64 &  &  &  \\ \hline
<AC> & Atividades Complementares / \textit{Complementary Activities} & 32 & 0 & 0 & 0 & 32 &  &  &  \\ \hline

\multicolumn{10}{|>{\centering\arraybackslash}c|}{\normalsize\textbf{Semestre 2}} \\ \hline
QXD0006 & Cálculo Diferencial e Integral I / \textit{Calculus I} & 64 & 0 & 0 & 0 & 64 & QXD0109 &  &  \\ \hline
QXD0010 & Estrutura de Dados / \textit{Data Structure} & 32 & 32 & 0 & 0 & 64 & QXD0001 &  &  \\ \hline
QXD0017 & Lógica para Computação / \textit{Logic for Computer Science} & 64 & 0 & 0 & 0 & 64 & QXD0056 &  &  \\ \hline
QXD0008 & Matemática Discreta / \textit{Discrete Mathematics} & 64 & 0 & 0 & 0 & 64 & QXD0056 &  &  \\ \hline
QXD0007 & Programação Orientada a Objetos / \textit{Object-Oriented Programming} & 32 & 32 & 0 & 0 & 64 & QXD0001 &  &  \\ \hline
<AC> & Atividades Complementares / \textit{Complementary Activities} & 32 & 0 & 0 & 0 & 32 &  &  &  \\ \hline

\multicolumn{10}{|>{\centering\arraybackslash}c|}{\normalsize\textbf{Semestre 3}} \\ \hline
<03> & Agentes Inteligentes / \textit{Intelligent Agents} & 48 & 16 & 0 & 0 & 64 & <02> &  &  \\ \hline
QXD0116 & Álgebra Linear / \textit{Linear Algebra} & 64 & 0 & 0 & 0 & 64 & QXD0056 &  &  \\ \hline
QXD0115 & Estrutura de Dados Avançada / \textit{Advanced Data Structure} & 32 & 32 & 0 & 0 & 64 & QXD0010 &  &  \\ \hline
QXD0011 & Fundamentos de Banco de Dados / \textit{Database Fundamentals} & 32 & 32 & 0 & 0 & 64 & QXD0001 &  &  \\ \hline
QXD0112 & Probabilidade e Estatística / \textit{Probability and Statistics} & 64 & 0 & 0 & 0 & 64 & QXD0056 &  &  \\ \hline
<AE> & Atividades de Extensão / \textit{Extension Activities} & 0 & 0 & 0 & 32 & 32 &  &  &  \\ \hline

\multicolumn{10}{|>{\centering\arraybackslash}c|}{\normalsize\textbf{Semestre 4}} \\ \hline
QXD0019 & Engenharia de Software / \textit{Software Engineering} & 64 & 0 & 0 & 0 & 64 & QXD0007 &  &  \\ \hline
QXD0178 & Mineração de Dados / \textit{Data Mining} & 48 & 16 & 0 & 0 & 64 & QXD0011 &  &  \\ \hline
QXD0041 & Projeto e Análise de Algoritmos / \textit{Design and Analysis of Algorithms} & 32 & 32 & 0 & 0 & 64 & QXD0008, QXD0010 &  &  \\ \hline
<04> & Representação do Conhecimento e Raciocínio / \textit{Knowledge Representation and Reasoning} & 64 & 0 & 0 & 0 & 64 & QXD0017 &  &  \\ \hline
QXD0076 & Sistemas Multiagentes / \textit{Multiagent Systems} & 32 & 32 & 0 & 0 & 64 & <03> &  &  \\ \hline
<AE> & Atividades de Extensão / \textit{Extension Activities} & 0 & 0 & 0 & 32 & 32 &  &  &  \\ \hline

\multicolumn{10}{|>{\centering\arraybackslash}c|}{\normalsize\textbf{Semestre 5}} \\ \hline
<05> & Aprendizado Não Supervisionado / \textit{Unsupervised Learning} & 32 & 32 & 0 & 0 & 64 & QXD0112, QXD0116 &  &  \\ \hline
<06> & Aprendizado Supervisionado / \textit{Supervised Learning} & 32 & 32 & 0 & 0 & 64 & QXD0112, QXD0116 &  &  \\ \hline
<07> & Busca e Otimização / \textit{Search and Optimization} & 32 & 32 & 0 & 0 & 64 & QXD0041 &  &  \\ \hline
<08> & Ética em Inteligência Artificial / \textit{Ethics in Artificial Intelligence} & 64 & 0 & 0 & 0 & 64 & QXD0162 &  &  \\ \hline
<$\ast$> & Optativas / \textit{Elective Courses} & 64 & 0 & 0 & 0 & 64 &  &  &  \\ \hline
<19> & Projeto Extensionista I / \textit{Extension Project I} & 16 & 0 & 0 & 16 & 32 & QXD0178 &  &  \\ \hline
<AE> & Atividades de Extensão / \textit{Extension Activities} & 0 & 0 & 0 & 96 & 96 &  &  &  \\ \hline

\multicolumn{10}{|>{\centering\arraybackslash}c|}{\normalsize\textbf{Semestre 6}} \\ \hline
<09> & Aprendizado Profundo / \textit{Deep Learning} & 32 & 32 & 0 & 0 & 64 & <06> &  &  \\ \hline
<$\ast$> & Optativas / \textit{Elective Courses} & 64 & 0 & 0 & 0 & 64 &  &  &  \\ \hline
<10> & Processamento de Dados em Larga Escala / \textit{Large-Scale Data Processing} & 48 & 16 & 0 & 0 & 64 & QXD0011 &  &  \\ \hline
<11> & Processamento de Linguagem Natural / \textit{Natural Language Processing} & 48 & 16 & 0 & 0 & 64 & <06> &  &  \\ \hline
QXD0182 & Visão Computacional / \textit{Computer Vision} & 48 & 16 & 0 & 0 & 64 & <06> &  &  \\ \hline
<20> & Projeto Extensionista II / \textit{Extension Project II} & 16 & 0 & 0 & 16 & 32 & <19> &  &  \\ \hline
<AE> & Atividades de Extensão / \textit{Extension Activities} & 0 & 0 & 0 & 96 & 96 &  &  &  \\ \hline

\multicolumn{10}{|>{\centering\arraybackslash}c|}{\normalsize\textbf{Semestre 7}} \\ \hline
<12> & Aprendizado por Reforço / \textit{Reinforcement Learning} & 32 & 32 & 0 & 0 & 64 & <09> &  &  \\ \hline
<13> & Inteligência Artificial Generativa / \textit{Generative Artificial Intelligence} & 32 & 32 & 0 & 0 & 64 & <09> &  &  \\ \hline
<14> & MLOPs / \textit{MLOps} & 48 & 16 & 0 & 0 & 64 & <06> &  &  \\ \hline
<$\ast$> & Optativas / \textit{Elective Courses} & 64 & 0 & 0 & 0 & 64 &  &  &  \\ \hline
<15> & Planejamento Automatizado / \textit{Automated Planning} & 48 & 16 & 0 & 0 & 64 & QXD0017 &  &  \\ \hline
QXD0110 & Projeto de Pesquisa Científico-Tecnológico / \textit{Scientific and Technological Research Project} & 16 & 16 & 0 & 0 & 32 & <09> & <17> &  \\ \hline
<17> & Trabalho de Conclusão de Curso I / \textit{Final Project I} & 32 & 0 & 0 & 0 & 32 &  & QXD0110 &  \\ \hline
<AE> & Atividades de Extensão / \textit{Extension Activities} & 0 & 0 & 0 & 32 & 32 &  &  &  \\ \hline

\multicolumn{10}{|>{\centering\arraybackslash}c|}{\normalsize\textbf{Semestre 8}} \\ \hline
<$\ast$> & Optativas / \textit{Elective Courses} & 256 & 0 & 0 & 0 & 256 &  &  &  \\ \hline
<AE> & Atividades de Extensão / \textit{Extension Activities} & 0 & 0 & 0 & 32 & 32 &  &  &  \\ \hline
<16> & Estágio Curricular Supervisionado I / \textit{Supervised Internship} & 0 & 160 & 0 & 0 & 160 & <06> &  &  \\ \hline
<18> & Trabalho de Conclusão de Curso II / \textit{Final Project II} & 64 & 0 & 0 & 0 & 64 & <17> &  &  \\ \hline
\end{tabular}
}
\label{tab:organizacao_curricular}
\end{table}





%Na coluna ``Componente Curricular'', todas as disciplinas nomeadas são obrigatórias, enquanto as disciplinas optativas são representadas pela palavra ``Optativa'' seguida de um número, indicando sua posição na contagem. No curso, recomenda-se que o(a) estudante curse: do primeiro ao terceiro semestre, cinco disciplinas obrigatórias por semestre; no quarto semestre, seis disciplinas obrigatórias; no quinto e sexto semestres, cinco disciplinas obrigatórias e uma optativa; no sétimo semestre, três optativas, além do Estágio Supervisionado I e da primeira parte do TCC; e no oitavo semestre, quatro optativas, o Estágio Supervisionado II e a finalização do TCC. Além disso, ao longo de todo o curso, o(a) estudante deverá cumprir 192 horas de Atividades Complementares e 320 horas de Atividades de Extensão, distribuídas entre disciplinas e a Unidade Curricular Especial de Extensão.



\begin{table}[h]
\centering
\caption{Disciplinas optativas do curso}
\renewcommand{\arraystretch}{1.5}
\resizebox{\textwidth}{!}{%
\small
\begin{tabular}{|
>{\centering\arraybackslash}m{1.5cm}|
>{\centering\arraybackslash}m{10.4cm}|
>{\centering\arraybackslash}m{1.4cm}|
>{\centering\arraybackslash}m{1.4cm}|
>{\centering\arraybackslash}m{1.4cm}|
>{\centering\arraybackslash}m{1.5cm}|
>{\centering\arraybackslash}m{1.4cm}|
>{\centering\arraybackslash}m{1.8cm}|
>{\centering\arraybackslash}m{2.4cm}|
>{\centering\arraybackslash}m{2.4cm}|
}
\hline
% \multicolumn{11}{|>{\centering\arraybackslash}l|}{\normalsize\textbf{Disciplinas optativass}} \\ \hline
\centering{\textbf{Código}} & 
\centering{\textbf{Nome do Componente Curricular}\\ (em português e inglês)} & 
\textbf{Carga Horária Teórica} & 
\textbf{Carga Horária Prática} & 
\textbf{Carga Horária EAD} & 
\textbf{Carga Horária Extensão} & 
\textbf{Carga Horária TOTAL} & 
\textbf{Pré-requisito(s)} & 
\textbf{Correquisito(s)}&  
\textbf{Equivalência(s)} \\ \hline

<21> & Análise Exploratória de Dados / \textit{Exploratory Data Analysis} & 48 & 16 & 0 & 0 & 64 & QXD0112, QXD0178 &  &  \\ \hline
QXD0005 & Arquitetura de Computadores / \textit{Computer Architecture} & 64 & 0 & 0 & 0 & 64 &  &  &  \\ \hline
QXD0134 & Cálculo Diferencial e Integral II / \textit{Calculus II} & 64 & 0 & 0 & 0 & 64 & <01> &  &  \\ \hline
QXD0135 & Cálculo Diferencial e Integral III / \textit{Calculus III} & 64 & 0 & 0 & 0 & 64 & QXD0134 &  &  \\ \hline
QXD0025 & Compiladores / \textit{Compilers} & 32 & 32 & 0 & 0 & 64 & QXD0040 &  &  \\ \hline
QXD0079 & Computação em Nuvem / \textit{Cloud Computing} & 32 & 32 & 0 & 0 & 64 & QXD0011, QXD0021 &  &  \\ \hline
QXD0183 & Computação Paralela / \textit{Parallel Computing} & 48 & 16 & 0 & 0 & 64 & QXD0013, QXD0041 &  &  \\ \hline
QXD0153 & Desafios de Programação / \textit{Programming Challenges} & 32 & 32 & 0 & 0 & 64 & QXD0008, QXD0010 &  &  \\ \hline
QXD0276 & Desenvolvimento de Software para Dispositivos Móveis / \textit{Software Development for Mobile Devices} & 32 & 16 & 0 & 16 & 64 & QXD0007 &  &  \\ \hline
QXD0253 & Desenvolvimento de Software para Web / \textit{Software Development for the Web} & 32 & 16 & 0 & 16 & 64 & QXD0007 &  &  \\ \hline
QXD0232 & Educação Ambiental / \textit{Environmental Education} & 16 & 48 & 0 & 0 & 64 &  &  &  \\ \hline
QXD0245 & Educação em Direitos Humanos / \textit{Human Rights Education} & 64 & 0 & 0 & 0 & 64 &  &  &  \\ \hline
QXD0029 & Empreendedorismo / \textit{Enterpreneurship} & 32 & 0 & 0 & 32 & 64 &  &  &  \\ \hline
QXD0179 & Estatística Multivariada / \textit{Multivariate Statistics} & 48 & 16 & 0 & 0 & 64 & QXD0112, QXD0116 &  &  \\ \hline
QXD0508 & Ética e Legislação / \textit{Ethics and Legislation} & 48 & 0 & 0 & 16 & 64 &  &  &  \\ \hline
QXD0254 & Gerência de Projetos / \textit{Project Management} & 32 & 16 & 0 & 16 & 64 & QXD0019 &  &  \\ \hline
QXD0035 & Inglês Instrumental I / \textit{Instrumental English I} & 64 & 0 & 0 & 0 & 64 &  &  &  \\ \hline
QXD0036 & Inglês Instrumental II / \textit{Instrumental English II} & 64 & 0 & 0 & 0 & 64 & QXD0035 &  &  \\ \hline
<22> & Inteligência Artificial Explicável / \textit{Explainable Artificial Intelligence} & 64 & 0 & 0 & 0 & 64 & <04>, <06> &  &  \\ \hline
QXD0256 & Interação Humano-Computador / \textit{Computer-Human Interaction} & 32 & 16 & 0 & 16 & 64 &  &  &  \\ \hline
QXD0113 & Língua Brasileira de Sinais - LIBRAS / \textit{Brazilian Sign Language - Libras} & 32 & 32 & 0 & 0 & 64 &  &  &  \\ \hline
QXD0016 & Linguagens de Programação / \textit{Programming Languages} & 64 & 0 & 0 & 0 & 64 & QXD0007 &  &  \\ \hline
QXD0040 & Linguagens Formais e Autômatos / \textit{Formal Languages and Automata} & 64 & 0 & 0 & 0 & 64 & QXD0008 &  &  \\ \hline
QXD0120 & Matemática Computacional / \textit{Computational Mathematics} & 48 & 16 & 0 & 0 & 64 & QXD0116 &  &  \\ \hline
<23> & Percepção e Ação Robótica / \textit{Robotic Perception and Action} & 48 & 16 & 0 & 0 & 64 & <03> &  &  \\ \hline
<24> & Pesquisa e Inovação / \textit{Research and Innovation} & 64 & 0 & 0 & 0 & 64 & QXD0162 &  &  \\ \hline
QXD0181 & Pesquisa Operacional / \textit{Operations Research} & 48 & 16 & 0 & 0 & 64 & QXD0116, QXD0041 &  &  \\ \hline
<25> & Processamento de Áudio e Voz / \textit{Audio and Speech Processing} & 48 & 16 & 0 & 0 & 64 & <11> &  &  \\ \hline
QXD0188 & Processamento de Imagens / \textit{Image Processing} & 48 & 16 & 0 & 0 & 64 & QXD0116 &  &  \\ \hline
QXD0114 & Programação Funcional / \textit{Functional Programming} & 32 & 32 & 0 & 0 & 64 &  &  &  \\ \hline
QXD0126 & Psicologia e Percepção / \textit{Psychology and Perception} & 64 & 0 & 0 & 0 & 64 &  &  &  \\ \hline
<26> & Raciocínio sob Incerteza / \textit{Reasoning Under Uncertainty} & 64 & 0 & 0 & 0 & 64 & QXD0017, QXD0112 &  &  \\ \hline
QXD0177 & Recuperação de Informação / \textit{Information Retrieval} & 48 & 16 & 0 & 0 & 64 & QXD0010, QXD0011 &  &  \\ \hline
QXD0021 & Redes de Computadores / \textit{Computer Networks} & 48 & 16 & 0 & 0 & 64 &  &  &  \\ \hline
QXD0246 & Relações Étnico-Raciais e Africanidades / \textit{Ethnic-Racial Relationships and Africanities} & 64 & 0 & 0 & 0 & 64 &  &  &  \\ \hline
QXD0069 & Segurança / \textit{Security} & 32 & 32 & 0 & 0 & 64 & QXD0021 &  &  \\ \hline
QXD0043 & Sistemas Distribuídos / \textit{Distributed Systems} & 32 & 32 & 0 & 0 & 64 & QXD0013 &  &  \\ \hline
QXD0013 & Sistemas Operacionais / \textit{Operating Systems} & 48 & 16 & 0 & 0 & 64 & QXD0005 &  &  \\ \hline
QXD0046 & Teoria da Computação / \textit{Theory of Computation} & 64 & 0 & 0 & 0 & 64 & QXD0040 &  &  \\ \hline
QXD0152 & Teoria dos Grafos / \textit{Graph Theory} & 64 & 0 & 0 & 0 & 64 & QXD0008 &  &  \\ \hline
<27> & Teoria dos Jogos / \textit{Game Theory} & 64 & 0 & 0 & 0 & 64 & QXD0112 &  &  \\ \hline
<28> & Tópicos Especiais em Aprendizado de Máquina / \textit{Special Topics in Machine Learning} & 64 & 0 & 0 & 0 & 64 & <05>, <06> &  &  \\ \hline
<29> & Tópicos Especiais em Ciência de Dados / \textit{Special Topics in Data Science} & 64 & 0 & 0 & 0 & 64 & QXD0178 &  &  \\ \hline
<30> & Tópicos Especiais em Filosofia, Ética e Humanidades / \textit{Special Topics in Philosophy, Ethics and Humanities} & 64 & 0 & 0 & 0 & 64 & <08> &  &  \\ \hline
<31> & Tópicos Especiais em Inteligência Artificial / \textit{Special Topics in Artificial Intelligence} & 64 & 0 & 0 & 0 & 64 & <03> &  &  \\ \hline
<32> & Tópicos Especiais em Matemática Computacional / \textit{Special Topics in Computational Mathematics} & 64 & 0 & 0 & 0 & 64 & QXD0116 &  &  \\ \hline
<33> & Visualização de Dados / \textit{Data Visualization} & 48 & 16 & 0 & 0 & 64 & QXD0178 &  &  \\ \hline

\end{tabular}
}
\label{tab:disciplinas_optativas}
\end{table}

% No âmbito da UFC em Quixadá, é comum a existência de disciplinas semelhantes ofertadas por diferentes cursos, que podem ser cursadas por estudantes de \nomedocurso e aproveitadas conforme sua equivalência registrada.

Um dos princípios norteadores do curso de \nomedocurso é a integração entre teoria e prática. Essa articulação é evidenciada na distribuição da carga horária teórica e prática dos componentes curriculares. As atividades práticas, no âmbito das disciplinas, incluem tanto aulas realizadas em laboratório quanto a execução de trabalhos e projetos experimentais desenvolvidos pelos(as) estudantes. Dessa forma, a contagem das horas teóricas e práticas atribuídas a cada disciplina fornece uma visão concreta de como ocorre essa integração entre teoria e prática ao longo da formação.

%As disciplinas com previsão de atividades em laboratório contam com alocação apropriada em horários e espaços físicos. Além disso, os(as) estudantes têm à sua disposição um laboratório de informática destinado a estudos extraclasse, conforme será detalhado no Capítulo~\ref{cap:infraestrutura-do-curso}. 

\clearpage

A Tabela~\ref{tab:carga_horaria} apresenta a distribuição da carga horária total do curso, abrangendo disciplinas obrigatórias e optativas, bem como atividades complementares e de extensão. %Esses dados são apresentados em créditos e equivalentes em horas. Observa-se que a integralização curricular do curso de \nomedocurso prevê a conclusão de um total de 200 créditos, correspondentes a 3200 horas.

\def\ufcCHObrigatoria{2.144h\xspace}
\def\ufcCHTeorica{1.536h\xspace}
\def\ufcCHPratica{576h\xspace}
\def\ufcCHEAD{0h\xspace}
\def\ufcCHExtensao{32h\xspace}
\def\ufcHorExtensao{288h\xspace}

\begin{table}[h]
\centering
\caption{Distribuição da Carga Horária Total do curso}
\label{tab:carga_horaria}
\renewcommand{\arraystretch}{1.5}
\resizebox{.9\textwidth}{!}{%
\small
\begin{tabular}{|m{4.5cm}|m{7cm}|m{3cm}|}
\hline
\textbf{Tipo do Componente} & \textbf{Componente Curricular} & \textbf{Carga horária} \\ \hline

\multirow{5}{*}{Componentes obrigatórios} 
& Disciplinas obrigatórias & \ufcCHObrigatoria (distribuidas em \ufcCHTeorica teóricas, \ufcCHPratica práticas,  \ufcCHExtensao extensão, \ufcCHEAD EAD)  \\ \cline{2-3}
% & Unidade Curricular Especial de Extensão & \\ \cline{2-3}
& Atividades Complementares &  64h \\ \cline{2-3}
& Atividades Extensionistas & \ufcHorExtensao \\ \cline{2-3}
& Estágio Supervisionado & 160h \\ \cline{2-3}
& Trabalho de Conclusão de Curso & 96h \\ \cline{2-3}
% & Atividades Complementares & \\ 
\hline
\multirow{2}{*}{Componentes optativos}
& Disciplinas optativas & 256h \\ \cline{2-3}
& Disciplinas optativas livres & 192h \\ \hline

\multicolumn{2}{|c|}{\textbf{Total}} & 3200h \\ \hline

\end{tabular}
}
% \vspace{0.5cm}
\end{table}

%Na Universidade Federal do Ceará, cada crédito equivale a 16 horas. O tempo ideal para a conclusão do curso é estimado em 4 (quatro) anos, ou 8 (oito) semestres letivos. Sendo assim, o(a) estudante do curso de \nomedocurso, modalidade Bacharelado, deverá observar o tempo máximo para sua conclusão, estipulado em 6 (seis) anos, ou 12 (doze) semestres letivos.

As cargas horárias mínima, média e máxima por semestre do curso foram definidas seguindo as orientações da Portaria n.º 31/2022, de 20 de abril de 2022. A carga horária mínima é de 192 horas por semestre, correspondendo ao mínimo de 3 disciplinas de 64 horas. A carga horária semestral média corresponde a 400 horas e a carga horária semestral máxima, a 544 horas.

\begin{table}[h]
\centering
\caption{Carga horária por semestre do curso e prazos de conclusão}
\vspace{1em}
\renewcommand{\arraystretch}{1.5}
\resizebox{.75\textwidth}{!}{%
\small
\begin{tabular}{|p{10cm}|c|}
\hline
\textbf{Carga horária por semestre} & \textbf{Horas} \\
\hline
{Carga horária semestral mínima do currículo} & 192h \\
\hline
{Carga horária semestral média do currículo} & 400h \\
\hline
{Carga horária semestral máxima do currículo} & 544h \\ 
\hline
\multicolumn{2}{c}{} \\
\hline
\textbf{Prazos} & \textbf{Semestres} \\
\hline
Mínimo & 7\\
\hline
Médio & 8\\
\hline
Máximo & 12\\
\hline
\end{tabular}
\label{tab:carga_horaria_semestral}
}
\end{table}


\clearpage
\section{Ementário e bibliografias}

As ementas de todos os componentes curriculares que compõem a estrutura curricular do curso, juntamente com as respectivas bibliografias básica e complementar, estão disponíveis no Apêndice \ref{ap:ementario_bibliografia} deste documento. 

Cada bibliografia básica contém pelo menos três títulos, e cada bibliografia complementar inclui pelo menos cinco títulos. As ementas e bibliografias dos componentes curriculares devem revisadas e atualizadas regularmente, levando em consideração os avanços na área de conhecimento de cada componente. Isso visa garantir que o material bibliográfico adotado no curso contribua para a formação definida neste PPC, assegurando abrangência, aprofundamento e coerência teórica.

Conforme detalhado na Seção~\ref{sec:biblioteca}, a Biblioteca do Campus Quixadá (BCQ), em parceria com a Secretaria de Acessibilidade da UFC, oferece serviços especializados, recursos e tecnologia assistiva para atender usuários com deficiência, permitindo-lhes acesso ao material bibliográfico recomendado nas disciplinas. Além disso, todos os títulos físicos indicados possuem exemplares em quantidade adequada no acervo da BCQ, conforme relatório de adequação assinado pelo Núcleo Docente Estruturante (NDE) do curso.

        \chapter{ESTÁGIO CURRICULAR SUPERVISIONADO}
\label{cap:estagio-curricular-supervisionado}
% regis

O estágio supervisionado é parte de grande importância na estrutura curricular do curso ao inserir os alunos, de forma supervisionada, em contato direto com as práticas do mercado de trabalho. Os alunos têm a oportunidade de observar as técnicas, métodos, processos e afins, vistos ao longo do curso, sendo aplicados no dia a dia das empresas, além de observarem as dificuldades e alternativas que lhes são inerentes, contribuindo assim para a prática das competências e habilidades previstas no perfil do egresso. No âmbito do curso de \nomedocurso esse componente curricular é disciplinado e regimentado pelo Regulamento de Estágio Curricular Supervisionado, elaborado e aprovado pelo Colegiado do Curso.

O estágio supervisionado é firmado através da assinatura de um Termo de Compromisso de Estágio por parte da UFC, do estudante estagiário e da empresa concedente do estágio, e é acompanhado pelo professor orientador de estágio e por um profissional supervisor, conforme disposição da Lei n.º 11.788, de 29 de setembro de 2008 \cite{brasil2008lei11788}, da Resolução n.º 32/CEPE, de 30 de outubro de 2009 \cite{ufc2009resolucao32}. Além disso, as experiências individuais devem ser relatadas pelo estudante estagiário. Também é possível realizar o aproveitamento de atividade de estágio por parte do estudante que já possua experiência profissional como trabalhador formal ou que tenha iniciado estágio fora do período regular de matrícula, observadas as disposições vigentes.

É de responsabilidade do professor orientador de estágio acompanhar e avaliar o Plano de Trabalho fornecido pelo estudante, a Avaliação do Rendimento Discente realizada pelo profissional supervisor, o Seminário de Relato de Experiência e o Relatório Final de Estágio apresentados pelo estudante, a partir da experiência do estágio, da inter-relação teoria e prática, e do desenvolvimento de ações exitosas e inovadoras das quais participou. Esses dispositivos serão também constituintes da avaliação do Estágio Supervisionado, sendo considerado aprovado o aluno que obtiver nota superior ou igual a 7,0, com frequência mínima de 90\%, tendo respeitado todos os requisitos dispostos no Regulamento.

Ao final do estágio, são avaliadas as experiências que poderão ser sistematizadas para publicações e apresentações em eventos da comunidade acadêmica. A carga horária do estágio no curso de \nomedocurso é de 160 horas (obrigatórias), correspondendo a 5\% da carga horária do curso. O estágio é normalmente realizado nos dois últimos anos do curso e é implementado com a atividade Estágio Supervisionado, realizada na área de \nomedocurso.
%Apesar de o estágio supervisionado normalmente ser cumprido no fim do curso, muitos dos alunos de \nomedocurso, que já realizaram ou estão realizando estágio, o fizeram antes do último ano do curso, o que reflete o nível de interesse dos alunos para ingressarem no mercado e também a demanda por profissionais com o perfil do egresso, mesmo antes de se formarem.

% No que diz respeito ao aproveitamento da atuação profissional do discente como atividade de estágio, poderá ser pleiteado para atividades realizadas em uma das seguintes modalidades: estágio prévio, trabalhador formal, atuação como pessoa jurídica, empresa sediada no exterior, programa de imersão profissional e Projeto de Pesquisa e Desenvolvimento (P\&D). Essas modalidades de aproveitamento de estágio estão previstas no documento de Regulamento de Estágio Curricular Supervisionado do Curso de \nomedocurso do Campus Quixadá.

Dentro do contexto do estágio, é importante que o curso desenvolva estratégias para a gestão da integração entre ensino e mundo do trabalho. Uma das formas de se fazer isso é através da promoção dos Seminários de Oportunidades em Pesquisa, Desenvolvimento e Inovação, cujo objetivo é aproximar o Campus da UFC em Quixadá do meio empresarial local e regional. Nesses seminários, os docentes apresentam suas áreas e projetos de pesquisa, e os representantes das empresas apresentam lacunas e necessidades de pesquisa e implementação de soluções, momento em que se busca identificar oportunidades de cooperação entre o campus e as empresas. A partir desses seminários e outros contatos com profissionais, por exemplo, com a participação deles em eventos do campus, é possível identificar demandas por conhecimentos e tecnologias que podem ser discutidos e incorporados nas atividades de ensino, a fim de capacitar os alunos ao mundo do trabalho, durante o estágio e após se formarem.

A regulamentação de estágio curricular supervisionado do curso de \nomedocurso define os procedimentos, responsabilidades, obrigações, formulários, objetivos e avaliação do estágio do curso de \nomedocurso.

As atividades desenvolvidas pelos discentes durante o Estágio Supervisionado devem ser compatíveis com as áreas de atuação previstas no perfil do egresso do curso (Capítulo~\ref{cap:perfil-profissional}), e podem se enquadrar em diferentes modalidades, conforme descrito a seguir.

\section{Estágio em Empresas Conveniadas à UFC}

O estágio pode ser realizado em empresas públicas ou privadas que possuam convênio com a UFC. A formalização desses convênios e de todo o trâmite legal entre empresa, aluno e universidade é de responsabilidade da Agência de Estágios da UFC\footnote{Agência de Estágios da UFC: \url{https://estagios.ufc.br/}}, vinculada à Pró-Reitoria de Extensão (PREx). Essa agência coordena o Programa de Estágio Curricular Supervisionado nas unidades acadêmicas da UFC, sendo responsável pela articulação, agenciamento e formalização dos estágios obrigatórios e não obrigatórios com organizações conveniadas. Atua ainda como elo institucional com os ambientes de estágio e fornece subsídios para a atualização das práticas de estágio.

O Termo de Compromisso de Estágio (TCE), regulamentado pela Procuradoria Geral da União, é o documento que oficializa o vínculo entre a Instituição de Ensino (UFC), a Concedente (empresa) e o Estagiário (aluno), conferindo segurança jurídica às partes envolvidas\footnote{A lista de empresas conveniadas à UFC pode ser acessada no site da agência: \url{https://si3.ufc.br/sigaa/public/estagio/lista.jsf}}.

Além disso, a Agência de Estágios realiza a divulgação de oportunidades de estágio, emprego e trainee, e colabora com empresas no processo de recrutamento, orientando sobre os cursos mais adequados ao perfil da vaga.

\section{Aproveitamento de Experiências Profissionais e Acadêmicas}

Atividades atuais ou anteriores desenvolvidas pelos discentes na área de IA podem ser aproveitadas como carga horária total do Estágio Supervisionado, desde que estejam em conformidade com o Regulamento de Estágio Curricular Supervisionado do Curso de \nomedocurso do Campus Quixadá. As modalidades reconhecidas incluem:

\begin{itemize}
    \item \textbf{Contrato Formal de Trabalho}: Atuação na área do curso, com vínculo empregatício formal.
    \item \textbf{Atuação como Pessoa Jurídica}: Prestação de serviços na área do curso por meio de empresa própria, no Brasil ou no exterior.
    \item \textbf{Projetos de Pesquisa}: Participação em projetos de Pesquisa e Desenvolvimento (P\&D) em inteligência artificial, formalizados em instituições de ensino superior.
    \item \textbf{Programas de Imersão Profissional}: Participação em programas estruturados voltados à inserção no mercado de trabalho.
    \item \textbf{Estágio Não Obrigatório Anterior}: Solicitação de aproveitamento de atividades previamente realizadas como estágio não obrigatório.
\end{itemize}


% O discente também pode pleitear o aproveitamento de suas atividades anteriores de estágio não obrigatório, como carga horária total das atividades de Estágio Supervisionado, de acordo com as regras estabelecidas no Regulamento de Estágio Curricular Supervisionado do Curso de \nomedocurso do Campus Quixadá.

% desde que o aproveitamento seja solicitado até o semestre letivo seguinte ao desenvolvimento das atividades pleiteadas, e as atividades desenvolvidas pelo discente tenham sido realizadas em áreas de atuação afins com o perfil de egresso previsto no Projeto Pedagógico do Curso.

% \section{CONTRATO FORMAL DE TRABALHO}

% O discente que está atuando ou que já atuou na área do curso como trabalhador formal, pode pleitear o aproveitamento de suas atividades como carga horária total das atividades de Estágio Supervisionado, de acordo com as regras estabelecidas no Regulamento de Estágio Curricular Supervisionado do Curso de \nomedocurso do Campus Quixadá.

% \section{ATUAÇÃO COMO PESSOA JURÍDICA}

% O discente que está desempenhando ou que já desempenhou atividades como pessoa jurídica para empresas do Brasil ou do exterior poderá pleitear o aproveitamento de suas atividades como carga horária total de Estágio Supervisionado, de acordo com as regras estabelecidas no Regulamento de Estágio Curricular Supervisionado do Curso de \nomedocurso do Campus Quixadá.

% \section{NÚCLEO DE PRÁTICAS EM INFORMÁTICA (NPI)}

% O Núcleo de Práticas em Informática (NPI) foi criado com o objetivo de atender à comunidade acadêmica e à sociedade do sertão central com soluções de Tecnologia da Informação. As atividades do núcleo foram iniciadas em 2009, à época com o nome ``Escritório de Projetos'', através de projetos executados por docentes e alunos do grupo PET-SI, atendendo às demandas da comunidade acadêmica. Posteriormente, perceberam-se outras possibilidades para este, como, por exemplo, uma alternativa de provimento de estágio para estudantes dos cursos de graduação do campus.

% Com o estabelecimento oficial do NPI em 2011, foram iniciados projetos com alunos concludentes do curso de Sistemas de Informação. Em 2013, alunos concludentes do curso de \nomedocurso também começaram a atuar nos projetos. Também em 2013, o núcleo passou a operar em infraestrutura própria, contando com três salas equipadas com estações de trabalho. Em 2016, os alunos de Ciência da Computação também começaram a atuar nos projetos e, em 2017, foi a vez dos alunos do curso de Design Digital também começarem a participar dos projetos do NPI. A integração dos alunos de vários cursos, atuando nos mesmos projetos, permite que os projetos explorem habilidades e competências específicas de cada curso.

% O NPI é estruturado para funcionar nos moldes de uma fábrica de software, com processo definido e projetos com clientes reais. Um processo de desenvolvimento foi elaborado e implantado como forma de padronizar as atividades dos alunos no desenvolvimento de software e incorporar melhores práticas de \nomedocurso de acordo com metodologias, métodos e modelos de maturidade de processo que já são largamente utilizados na indústria de software e na academia \cite{goncalves2013}. Os alunos participantes são orientados por uma equipe composta por professores e servidores técnico-administrativos, além de fazerem registro de ponto para o acompanhamento da frequência pelo coordenador de estágio. Esta frequência deve ser, no mínimo, de 90\% para o aluno ser considerado aprovado.

% O NPI é registrado como um programa de extensão com múltiplas perspectivas, já que atende necessidades da comunidade (fornecimento de soluções de TI), dos alunos (provimento de estágio) e dos docentes, servindo a estes últimos como espaço para ampliação da experiência profissional. Em termos de infraestrutura, o núcleo conta, no início de 2025, com 27 computadores distribuídos em três salas.

% \section{PROJETOS DE PESQUISA}

% O discente que participa ou participou de projeto de Pesquisa e Desenvolvimento (P\&D) na área de inteligência artificial, formalizado em uma instituição de ensino superior, pode utilizar suas atividades como carga horária das atividades de Estágio Supervisionado, de acordo com as regras estabelecidas no Regulamento de Estágio Curricular Supervisionado do Curso de \nomedocurso do Campus Quixadá.

% Para isso, as seguintes regras devem ser observadas:

% \begin{itemize}
%     \item O discente deverá apresentar Plano de Trabalho / Lista de Tarefas a serem desenvolvidas para análise do supervisor de estágio;
%     \item O discente deverá, durante o estágio, participar do desenvolvimento de uma ou mais soluções de IA;
%     \item O discente deve entregar, até o final de estágio, artefatos como: documentação técnica, pipelines de dados, scripts de treinamento e avaliação de modelos de IA, código de API ou aplicação que faz uso de IA, Jupyter Notebooks com experimentos, roteiros de teste, testes automatizados, modelos de IA treinados e análises de interpretabilidade.
% \end{itemize}


% \section{ATUAÇÃO EM PROGRAMA DE IMERSÃO PROFISSIONAL}

% O discente que está participando ou que participou de programa(s) de imersão profissional poderá pleitear o aproveitamento de suas atividades como carga horária total das atividades de Estágio Supervisionado, de acordo com as regras estabelecidas no Regulamento de Estágio Curricular Supervisionado do Curso de \nomedocurso do Campus Quixadá.

% \section{INICIATIVA EMPREENDEDORA (IE)}

% O estágio em Iniciativa Empreendedora (IE)\footnote{Iniciativa empreendedora (IE) na UFC em Quixadá: \url{https://www.quixada.ufc.br/iniciativa-empreendedora/}} tem o objetivo de fomentar iniciativas empreendedoras voltadas para tecnologia no Campus da UFC em Quixadá. Esse tipo de estágio é possível desde que aprovado pelo professor orientador de estágio do curso. Para ser aprovado, o aluno deve submeter uma proposta de modelo de negócio que apresente um nível adequado de inovação e que possa também ser interdisciplinar.

% Assim como o estágio tradicional em empresas, a IE do aluno é acompanhada durante todo o período vigente, de acordo com as seguintes regras:
% \begin{itemize}
%     \item Encontros quinzenais de acompanhamento das atividades realizadas e planejamento das atividades da próxima quinzena;
%     \item As atividades devem ser passíveis de auditoria para comprovação de horas trabalhadas;
%     \item A carga horária mínima é de 12 horas semanais;
%     \item Acompanhamento através de ferramentas de gerenciamento de projetos / sistemas de controle de versão.
%     \item Cumprir os itens de avaliação do estágio supervisionado.
% \end{itemize}

        \chapter{TRABALHO DE CONCLUSÃO DE CURSO}
\label{cap:trabalho-de-conclusao-de-curso}

O curso de \nomedocurso possui um Trabalho de Conclusão de Curso (TCC) que envolve todos os procedimentos de uma investigação técnico-científica nos parâmetros acadêmicos. O TCC de \nomedocurso deverá ser desenvolvido pelo estudante ao longo dos dois últimos semestres do curso e envolve as seguintes disciplinas e atividades\footnote{A Seção \ref{sec:integralizacao} apresenta o detalhamento de correquisitos e pré-requisitos relacionados aos três componentes curriculares listados.}:
\begin{itemize}
    \item Disciplina ``Projeto de Pesquisa Científica e Tecnológica'' (PPCT), com 2 créditos (32h), ofertada no 7º semestre;
    \item Atividade ``Trabalho de Conclusão de Curso I'' (TCC I), com 2 créditos (32h), ofertada simultaneamente à disciplina acima citada, no 7º semestre;
    \item Atividade ``Trabalho de Conclusão de Curso II'' (TCC II), com 4 créditos (64h), ofertada no 8º semestre.
\end{itemize}



Cabe ao professor responsável pela disciplina ``Projeto de Pesquisa Científica e Tecnológica'' ministrar o conteúdo de metodologia científica nos seguintes termos:
\begin{itemize}
    \item No início do semestre, apresentar Plano de Ensino, contemplando o cronograma de entrega das versões parciais e da versão final do trabalho realizado pelo discente; no caso do TCC I, auxiliar os alunos na escolha dos temas e do professor-orientador;

    \item Explanar detalhadamente sobre a elaboração, estrutura, redação e apresentação de Trabalhos de Conclusão de Curso, orientando os alunos nessas questões;
    
    \item Reunir-se semanalmente com o grupo de alunos para a construção gradual e colaborativa dos projetos de pesquisa, buscando referencial bibliográfico que fundamente a metodologia empregada em cada projeto;

    \item Verificar a conformidade do TCC com as normas de elaboração de trabalhos acadêmicos da UFC;
    
   % \item Organizar a agenda de defesas e auxiliar na composição das bancas;

   % \item Providenciar, junto à Coordenação do curso, os encaminhamentos administrativos necessários.

\end{itemize}

Enquanto o professor de PPCT contribui com técnicas para elaboração do projeto do TCC, apresentando cada uma de suas etapas de forma conceitual e aplicada, cabe ao orientador as responsabilidades relacionadas ao conteúdo do trabalho e acompanhamento das atividades do aluno em TCC I e TCC II. O TCC I e o TCC II deverão ser entregues em formato acadêmico e defendidos perante uma banca de pelo menos três membros, de acordo com os critérios gerais da UFC, detalhados no Regimento Geral da UFC \cite{ufc_regimento_geral_2019}, serão aprovados com nota igual ou superior a 7,0 e frequência igual ou superior a 90\%.

Assim como acontece com os outros cursos do Campus, para os alunos de \nomedocurso, há a disponibilização dos Trabalhos de Conclusão de Curso em repositórios institucionais próprios da UFC, acessíveis pela internet, para sua devida exposição e consulta pela comunidade em geral\footnote{Os TCCs são armazenados no Repositório Institucional da UFC: \url{http://www.repositorio.ufc.br/} e podem ser acessados também pelo site do campus: \url{https://www.quixada.ufc.br/monografias/}}.

Além disso, para auxiliar a normalização de trabalhos acadêmicos, a Biblioteca Universitária elaborou o Guia de Normalização de Trabalhos Acadêmicos da UFC, tomando como base as normas da ABNT. No site da Biblioteca Universitária constam ainda \textit{templates} em três formatos distintos, já contemplando as recomendações das normas da ABNT: Word, Libre Office e Overleaf\footnote{Todos os documentos sobre normalização de trabalhos acadêmicos estão disponíveis no site da Biblioteca Universitária: \url{https://biblioteca.ufc.br/pt/servicos-e-produtos/normalizacao-de-trabalhos-academicos/}}. Para mais detalhes, observar o Regulamento dos Trabalhos de Conclusão de Curso.

A atividade de conclusão de curso é  intrinsecamente interdisciplinar, onde os conhecimentos adquiridos ao longo do curso são aplicados na criação de soluções científicas ou tecnológicas em \nomedocurso. O desenvolvimento dessa atividade colabora  no alcance das competências previstas no perfil do egresso, como o desenvolvimento do pensamento sistêmico, que permite analisar e entender os problemas e os processos organizacionais, e o desenvolvimento e a gestão de soluções baseadas em tecnologia da informação para os processos de negócio de organizações.
        \chapter{ATIVIDADES COMPLEMENTARES}
\label{cap:atividades-complementares}

As Atividades Complementares constituem um conjunto de estratégias pedagógico-didáticas que permitem, no âmbito do currículo, a articulação entre teoria e prática e a complementação, por parte do estudante, dos saberes e habilidades necessárias à sua formação. O Programa de Atividades Complementares (PAC) do Campus da UFC em Quixadá, busca qualificar o aluno e desenvolver de forma complementar, nos futuros profissionais, competências valorizadas em sua área de formação, tais como perfil empreendedor, iniciativa, liderança, autoconhecimento, perseverança e habilidade em lidar com obstáculos, mudanças e transformações, além de prestar serviços à comunidade.

O programa possibilita que o aluno realize atividades práticas ligadas à profissão que escolheu, a partir do primeiro semestre do curso, criando um diferencial na formação universitária, oferecendo uma variedade de atividades complementares, agrupadas em sete categorias: atividades de iniciação à docência, à pesquisa e/ou à extensão, atividades artístico-culturais e esportivas, atividades de participação e/ou organização de eventos, experiências ligadas à formação profissional e/ou correlatas, produção técnica e/ou científica, vivências de gestão, outras atividades.

Entre as sete categorias, são permitidos mais de cinquenta tipos diferentes de atividades complementares, abrangendo uma grande diversidade de opções para o aluno, indo desde “publicação de artigos científicos” a “doação de sangue”. As possibilidades de atividades complementares permitidas são aderentes à formação geral do discente, por meio de iniciativas de socialização com colegas e demais membros da comunidade acadêmica, atividades em prol da sociedade, atividades relacionadas à saúde, esporte, cultura, etc.; são também aderentes à formação específica do discente, como experiências técnicas em minicursos, palestras, projetos de pesquisa, etc.

Diversas atividades e eventos de extensão, atualmente, podem ser contabilizados como horas complementares. Entretanto, a partir do atendimento à Resolução CEPE nº 28, de 1º de dezembro de 2017 (UFC, 2017c), que regulamenta a Curricularização da Extensão nos cursos de graduação da UFC, poderão ser integralizadas, como “Atividades Complementares”, apenas as horas excedentes das ações extensionistas.
A forma de aproveitamento das horas varia dependendo da atividade. Em alguns casos, o aluno contabiliza horas complementares com valores diferentes das horas reais. Por exemplo, a participação em congressos nacionais gera uma pontuação de seis horas complementares por dia de evento (hora complementar provavelmente menor que hora real), enquanto a participação como ministrante de um minicurso gera duas horas complementares para cada hora ministrada (hora complementar maior que hora real). Outra variação considerada no aproveitamento das horas é que, em alguns casos, a mesma atividade gera horas complementares diferentes dependendo da natureza da participação do aluno. Por exemplo, para um mesmo minicurso com duração de quatro horas, o aluno que ministrou contabiliza oito horas complementares, o aluno que organizou contabiliza duas horas complementares e o aluno que assistiu contabiliza uma hora complementar.

O cumprimento das atividades complementares é obrigatório à colação de grau. O seu acompanhamento e validação ficam a cargo da Coordenação do Curso. O registro e acompanhamento das Atividades Complementares é realizado através do sistema de acompanhamento SISAC (Sistema de Atividades Complementares) pela Coordenação do Curso, auxiliada pela Secretaria do Curso. O SISAC foi desenvolvido no campus para auxiliar no registro das atividades e no acompanhamento e controle por parte da Coordenação e do próprio aluno, que pode se organizar e planejar o cumprimento das horas complementares durante todo o curso.

Semestralmente, os alunos são orientados a protocolar as atividades complementares realizadas, a serem validadas e lançadas no sistema pela Coordenação e/ou secretaria, mediante a devida comprovação. As atividades protocoladas serão, então, distribuídas entre os itens presentes em cada categoria, de acordo com o Regulamento de Atividades Complementares do curso, observadas as equivalências e os limites de aproveitamento. Recomenda-se que o estudante conclua suas atividades complementares até o semestre anterior àquele em que ele pretende colar grau. O estudante de \nomedocurso que protocolar atividades que contabilizem o mínimo de 64 horas será considerado aprovado nesse componente.

O elenco específico de atividades complementares aproveitáveis, bem como os limites de aproveitamento máximo e mínimo de cada atividade são definidos no Regulamento de Atividades Complementares que atende à Resolução No.07/CEPE, de 17 de junho de 2005, que dispõe sobre estas atividades nos cursos de graduação da UFC (UFC, 2005).
        \chapter{EXTENSÃO}
\label{cap:extensao}

O conceito de Extensão Universitária, instituído no I Encontro Nacional de Pró-Reitores de Extensão das Universidades Públicas Brasileiras, em 1987, e ratificado pelo Plano Nacional de Extensão Universitária em 2012, pode ser entendido como ``o processo educativo, cultural e científico que articula o ensino e a pesquisa de forma indissociável e viabiliza a relação transformadora entre a universidade e a sociedade''. Assim, a Extensão Universitária consiste em um conjunto de ações de natureza interdisciplinar e multidisciplinar, articulando os saberes produzidos na universidade com a sociedade de modo geral.

A Extensão Universitária na Universidade Federal do Ceará (UFC) é definida pela Resolução Nº 05/CEPE, de 07 de março de 2025, como um processo de educação superior que objetiva promover uma relação mutuamente transformadora entre a universidade e a sociedade, integrando ensino e pesquisa por meio da cultura, arte, ciência, tecnologia e inovação, tendo em vista a articulação e o desenvolvimento social. 
Neste sentido, o processo de desenvolvimento da Extensão na UFC segue os seguintes os princípios norteadores:% \cite{ufc2025resolucao05}: 

\begin{enumerate}[label=\Roman* --, leftmargin=*]
    \item alinhamento da ciência, da arte e da tecnologia às prioridades do local, região e país; sensibilidade da universidade à necessidade de interação com outros saberes e necessidades, adquiridos nas atividades de ensino, pesquisa e extensão, sem pretender ser proprietária de um saber universal e pronto; 
    \item participação nos movimentos sociais para desenvolvimento de ações que busquem a superação das desigualdades e processos excludentes; 
    \item democratização de conhecimentos acadêmicos e científicos como parte do processo de interação e inclusão social desenvolvidos pela Extensão; 
    \item prestação de serviços como desenvolvimento de trabalho social; 
    \item atuar junto aos sistemas públicos de ensino para o fortalecimento da educação básica. 
\end{enumerate}


A Resolução Nº 05/CEPE, de 07 de março de 2025, define como diretrizes gerais para orientar a formulação, execução e avaliação das atividades de extensão a interação dialógica universidade-sociedade, a indissociabilidade entre o ensino, a pesquisa e a extensão, a interdisciplinaridade e a interprofissionalidade, a formação cidadã de estudantes e o impacto e transformação social, considerando que:

\begin{enumerate}[label=\Roman* --, leftmargin=*]
    \item a interação dialógica orienta o desenvolvimento das relações entre a universidade e setores sociais, superando o discurso da hegemonia acadêmica;
    \item a indissociabilidade entre o ensino, a pesquisa e a extensão pressupõe que as atividades de extensão são mais efetivas quando vinculadas ao processo de formação, pessoal e profissional, e de geração de conhecimento, tendo como princípio o desenvolvimento de práticas que promovam ações articuladoras da interdisciplinaridade, da transdisciplinaridade entre as áreas do conhecimento, bem como da interprofissionalidade;
    \item a interdisciplinaridade e a interprofissionalidade contribuem para o desenvolvimento da cultura extensionista na universidade, atendendo à complexidade das comunidades de saberes, aqui entendidas como grupos de pessoas que se articulam com as atividades de extensão Universitária por meio de uma relação dialógica. Estas comunidades não se limitam a grupos vulneráveis, podem ser movimentos sociais, entidades profissionais e outras formas de organização;
    \item a formação cidadã de estudantes implica no desenvolvimento de uma consciência cidadã, agregada à sua contribuição para a transformação social e integrada à matriz curricular de seu curso de graduação e/ou pós-graduação;
    \item a transformação social diz respeito à forma de relacionamento da universidade com a sociedade, com objetivo de, por meio da interação dialógica, trabalhar nos interesses e necessidades da população, possibilitando melhor compreensão das problemáticas e para o aprimoramento das políticas públicas.
\end{enumerate}

A concepção de currículo estabelecida na Lei nº 9.394/96, em seu capítulo IV – Da Educação Superior, expressa em seu Artigo 43, incisos VII, as seguintes finalidades da educação superior: ``VII - promover a extensão, aberta à participação da população, visando à difusão das conquistas e benefícios resultantes da criação cultural e da pesquisa científica e tecnológica geradas na instituição''. 

A promoção da extensão na UFC e no curso ocorre por meio de ações de extensão, sendo classificadas na Resolução Nº 05/CEPE, de 07 de março de 2025, em cinco modalidades: Programa, Projeto, Curso e Oficina, Evento ou Prestação de Serviço. 

Um Programa consiste em um conjunto articulado de ações com um objetivo comum, executado em um prazo de dois a cinco anos e podendo ser de natureza estratégica, estruturante ou setorial. Já um Projeto é uma ação de caráter processual com objetivo específico e prazo determinado de quatro meses a dois anos, que dialoga com as comunidades de saberes e promove a formação dos estudantes, estabelecendo um vínculo com atividades de pesquisa científica.

Cursos e Oficinas são ações pedagógicas, teóricas ou práticas, com carga horária mínima de duas horas, podendo ser classificadas como iniciação, atualização, treinamento, qualificação profissional ou aperfeiçoamento. Os Eventos, por sua vez, são ações de curta duração que implicam na apresentação ou exibição pública de conhecimento ou produto cultural, artístico, esportivo ou científico, como congressos, seminários e exposições.

A Prestação de Serviço refere-se à aplicação do conhecimento universitário para a solução de problemas de parceiros externos, incluindo serviços técnicos, assistência e consultoria. Dentre suas classificações, destacam-se o desenvolvimento tecnológico e inovação, que abrange a criação de softwares, aplicativos e sistemas inovadores, e as atividades de empresa júnior, ambas diretamente alinhadas à área de tecnologia. Todas essas modalidades podem gerar produtos de extensão, como publicações, relatórios técnicos, softwares, sistemas tecnológicos e jogos educativos, que materializam o conhecimento produzido.

Conforme a Resolução Nº 05/CEPE, de 07 de março de 2025, todas as atividades de extensão devem ser obrigatoriamente classificadas em Áreas do Conhecimento, Áreas Temáticas e Linhas de Extensão. A classificação em uma das Áreas do Conhecimento toma por base as definições do Conselho Nacional de Desenvolvimento Científico e Tecnológico (CNPq). Adicionalmente, as atividades devem ser classificadas conforme uma área temática principal e, opcionalmente, uma secundária. Por fim, as ações também precisam ser categorizadas segundo uma linha de extensão principal e, opcionalmente, uma secundária, sendo estas importantes para a articulação das atividades em programas. A Resolução Nº 05/CEPE, de 07 de março de 2025, apresenta as listas e os detalhamentos de cada uma dessas categorias de classificação.

Alinhando o perfil do egresso do Curso de \nomedocurso e considerando as necessidades da região, as atividades de extensão da UCEE do curso deverão ser realizadas em projetos que se enquadrem, preferencialmente, nas seguintes classificações, conforme a Resolução Nº 05/CEPE, de 07 de março de 2025:

\begin{itemize}
    \item \textbf{Áreas do Conhecimento}: Ciências Exatas e da Terra ou Multidisciplinar.
    \item \textbf{Áreas Temáticas}: Comunicação; Cultura; Direitos Humanos e Justiça; Educação; Meio Ambiente; Saúde; Tecnologia e Produção; Trabalho; ou Gestão.
    \item \textbf{Linhas de Extensão}: Comunicação estratégica; Desenvolvimento de produtos; Desenvolvimento regional; Desenvolvimento tecnológico; Desenvolvimento urbano; Educação profissional; Empreendedorismo; Espaços de ciência; Formação de professores; Gestão informacional; Gestão institucional e pública; Inovação tecnológica; Metodologias e estratégias de ensino/aprendizagem; Pessoas com deficiência; Questões ambientais;  Tecnologia da informação; ou Temas específicos/desenvolvimento humano.
\end{itemize}

A Estratégia 7 da Meta 12 do Plano Nacional de Educação 2014-2024 (Lei nº 13.005/2014), estabeleceu que as Instituições de Ensino Superior (IES) devem assegurar, no mínimo, 10\% (dez por cento) do total de créditos curriculares exigidos para a graduação em programas e projetos de extensão universitária. 
%
A Resolução Nº 05/CEPE, de 07 de março de 2025, estabelece a curricularização da extensão por meio da inserção de ações curriculares de extensão nos cursos de graduação como componentes obrigatórios para sua integralização. Por essa resolução, os projetos pedagógicos dos cursos de graduação da UFC devem prever a realização de atividades de Extensão que correspondam ao percentual mínimo de 10\% do total da carga horária do curso. %As ações de extensão inseridas nos cursos deverão reforçar as diretrizes nacionais da extensão, e, quando ofertadas por disciplinas, tais elementos devem estar descritos na ementa.

O processo de curricularização da extensão poderá ser realizado de três formas, desde que sejam atendidas as diretrizes da Extensão Universitária nas atividades realizadas, identificando a importância do ensino pela extensão, destacando-se objetivos, conhecimentos, habilidades e atitudes por ela constituídos ou desenvolvidos: 

\begin{enumerate}[label=\Roman* --, leftmargin=*]
    \item Unidade Curricular Especial de Extensão (UCEE), constituída de carga horária atribuída por meio de Programas e Projetos que identifiquem no seu planejamento a importância, o desenvolvimento e a avaliação de ações formativas de ensino de extensão; 
    \item Componentes curriculares com destinação de carga horária de extensão definidas na sua criação e regulamentação; 
    \item Ações Curriculares em Comunidades de Saberes (ACCS), integradas ao Plano Pedagógico do Curso (PPC), aprovadas pela Câmara de Extensão do CEPE e/ou pela Câmara de Graduação do CEPE; 
\end{enumerate}

Em conformidade com o Art. 4º da Resolução CNE/CES Nº 7, de 18 de dezembro de 2018, e com o parágrafo 1º do Art. 15 da Resolução Nº 05/CEPE, de 07 de março de 2025, o curso de \nomedocurso destina um total de 320 horas para as atividades de extensão, o que corresponde a 10\% da carga horária total de 3.200 horas do curso. Essa carga horária é integralizada por meio da combinação das modalidades previstas no Art. 20 da Resolução Nº 05/CEPE, de 07 de março de 2025.

A distribuição dessa carga horária é realizada em dois blocos, utilizando as modalidades previstas no Art. 20 da referida resolução:

\begin{itemize}
    \item \textbf{32 horas} são alocadas em Componentes Curriculares com destinação de carga horária de extensão, onde as práticas extensionistas compõem parte da carga horária de disciplinas da matriz curricular, conforme detalhado na seção de Integralização Curricular (Seção \ref{sec:integralizacao}).
    \item \textbf{288 horas} restantes podem ser integralizadas de forma flexível pelo discente por meio de uma das outras duas modalidades: a Unidade Curricular Especial de Extensão (UCEE) , constituída pela participação em programas e projetos de extensão; ou por meio de Ações Curriculares em Comunidades de Saberes (ACCS).
\end{itemize}

Para facilitar o processo de gestão das Atividades de Extensão e inserção das informações, foi desenvolvido o Módulo de Registro de Ação de Extensão no Sistema Integrado de Gestão de Atividades Acadêmicas (SIGAA) com o objetivo de melhorar a experiência de acesso e uso aos servidores (docentes e TAEs) e discentes.

O estudante poderá solicitar aproveitamento de carga horária das ações curriculares de extensão declaradas por setores com competência de outras instituições de ensino superior no Brasil ou no Exterior. Os pedidos serão analisados pela Câmara de Extensão do CEPE, em conformidade com o Projeto Pedagógico do Curso.

Caso o discente realize uma carga horária em atividades de extensão superior à exigida para a curricularização, as horas excedentes poderão ser contabilizadas para o cumprimento do total de horas de atividades complementares.

Para a integralização da carga horária obrigatória de extensão, é vedada a duplicidade no cômputo de horas com outras atividades curriculares. Dessa forma, as cargas horárias dedicadas ao Estágio Curricular Supervisionado, ao Trabalho de Conclusão de Curso e às Atividades Complementares não podem ser contabilizadas para este fim, por se tratarem de componentes distintos da matriz curricular. O mesmo princípio de não duplicidade se aplica a outras atividades de natureza extensionista, como a participação em Empresas Juniores e no Programa de Educação Tutorial (PET), cujas horas não podem ser simultaneamente utilizadas para o cômputo da extensão e de outra atividade formativa.
 
% As atividades de extensão são caracterizadas assim pela participação ativa de discentes, docentes e comunidade extra-acadêmica, que, atuando em conjunto, promovem a interação transformadora entre a Universidade e a sociedade. As atividades de extensão são acompanhadas no âmbito do curso por um supervisor de extensão designado pelo colegiado do curso, sendo este responsável por analisar e validar o cumprimento das ações da extensão.



O Manual de Procedimentos de Integralização de Atividades de Extensão (Anexo \ref{an:manual-extensao}) detalha as especificidades para o cumprimento da carga horária destinada à extensão pelos estudantes, as regras para o cadastro e acompanhamento das ações nas diferentes modalidades, como a Unidade Curricular Especial de Extensão (UCEE), e os critérios de avaliação e validação das atividades.
        %\chapter{ATIVIDADES PRÁTICAS DE ENSINO / ATIVIDADES DE TUTORIA}
\label{cap:atividades-praticas-de-ensino}

\section{Atividades Práticas de Ensino}

As Atividades Práticas de Ensino (APE) no Bacharelado em Inteligência Artificial têm como objetivo proporcionar aos estudantes experiências concretas de aplicação dos conhecimentos teóricos adquiridos nas disciplinas, fortalecendo a capacidade de resolução de problemas, desenvolvimento de soluções inovadoras e atuação autônoma em contextos reais e simulados.

Essas atividades estão integradas à matriz curricular por meio de componentes como laboratórios, oficinas, projetos integradores, estágios supervisionados, atividades de iniciação científica e desafios de inovação. As APE visam, ainda, estimular o raciocínio crítico e sistêmico, o trabalho em equipe, a comunicação científica e a aprendizagem baseada em projetos.

Dentre as principais formas de APE previstas neste curso, destacam-se:

\begin{itemize}
\item Desenvolvimento de sistemas inteligentes e modelos de aprendizado de máquina, incluindo tarefas como tratamento de dados, experimentação computacional, avaliação de desempenho e interpretação de resultados;
\item Construção de pipelines de engenharia de dados e MLOps, com uso de ferramentas modernas para versionamento, automação e monitoramento de modelos;
\item Elaboração e implantação de APIs com modelos embarcados, utilizando frameworks como FastAPI e Docker;
\item Participação em hackathons, competições e desafios práticos, estimulando a criatividade e o trabalho sob pressão;
\item Projetos interdisciplinares integradores, nos quais os estudantes resolvem problemas reais ou propostos por parceiros externos, utilizando abordagens baseadas em inteligência artificial;
\item Implementação de algoritmos e arquiteturas clássicas e modernas de IA, como redes neurais profundas, \textit{transformers}, sistemas de recomendação e modelos probabilísticos.
\end{itemize}

Essas atividades são registradas, avaliadas e sistematizadas por meio de relatórios técnicos, repositórios de código, apresentações públicas e autoavaliações, permitindo à coordenação pedagógica acompanhar o progresso dos estudantes e garantir a coerência formativa do curso.

\section{Atividades de Tutoria}

As atividades de tutoria desempenham papel fundamental no processo formativo ao promover um acompanhamento mais próximo e individualizado da trajetória dos estudantes, tanto no âmbito acadêmico quanto no desenvolvimento de competências pessoais e profissionais.

A tutoria é exercida por docentes do curso, monitores ou mentores vinculados a programas institucionais, e pode ocorrer em diferentes formatos, como encontros regulares, sessões de orientação, acompanhamento de projetos ou mediação de grupos de estudo.

As principais finalidades das atividades de tutoria incluem:

\begin{itemize}
\item Apoiar o processo de adaptação e permanência dos estudantes no curso, promovendo acolhimento, escuta ativa e orientação acadêmica;
\item Orientar a organização dos estudos e o aproveitamento de oportunidades de aprendizagem, como iniciação científica, estágios, eventos e cursos complementares;
\item Acompanhar o desenvolvimento de projetos práticos, TCCs e atividades de extensão, contribuindo com feedback técnico e metodológico;
\item Estimular a reflexão sobre o papel social da inteligência artificial, por meio de discussões sobre ética, privacidade, justiça algorítmica e impactos sociais das tecnologias emergentes;
\item Fornecer suporte na construção do itinerário formativo e profissional, orientando quanto às trilhas de carreira possíveis no campo da inteligência artificial, bem como à preparação de portfólios, currículos e perfis profissionais.
\end{itemize}

As atividades de tutoria complementam as dimensões técnica e acadêmica do curso, contribuindo para a formação integral dos estudantes e para o alinhamento das ações pedagógicas com os objetivos formativos do bacharelado.

        \chapter{METODOLOGIAS DE ENSINO E DE APRENDIZAGEM}
\label{cap:metodologias-de-ensino}

Os métodos de ensino e aprendizagem, efetivamente aplicados na formação proporcionada pelo curso de \nomedocurso, são definidos com base nos objetivos de aprendizagem estabelecidos nos planos de ensino de cada componente curricular, observando-se as especificidades de cada área quanto às estratégias mais adequadas e eficazes de transposição didática dos conteúdos e sua aprendizagem por parte dos alunos, visando assegurar que os correspondentes objetivos sejam alcançados.

Buscando desenvolver práticas de ensino-aprendizagem inovadoras, a escolha dos métodos é orientada por alguns princípios gerais: a) congruência entre objetivos de aprendizagem e métodos; b) factibilidade dos métodos em termos de recursos; c) observância dos aspectos de acessibilidade em todo o seu espectro.

Como forma de aprimoramento das metodologias de ensino e de aprendizagem adotadas nos cursos de graduação, a UFC investe na formação continuada dos docentes, especialmente através do projeto CASa\footnote{Site do CASa: \url{http://casa.virtual.ufc.br/}}. A Comunidade de Cooperação e Aprendizagem Significativa (CASa) é o programa de formação docente da UFC, fundado em 2009, e que hoje compõe a Escola Integrada de Desenvolvimento e Inovação Acadêmica (EIDEIA). O CASa trata-se de um espaço institucional na universidade que reúne iniciativas voltadas para a inovação de práticas de ensino e aprendizagem, e a promoção da excelência no ensino, pesquisa e extensão. Esse projeto busca, por meio de proposta dialógica e colaborativa, proporcionar aos docentes um espaço favorável à troca de experiências, de maneira que valorize o protagonismo e a trajetória de cada docente sem perder de vista os princípios de heterogeneidade, trabalho coletivo, interação, solidariedade, equidade e transformação que constituem bases sólidas a esse programa.

Adicionalmente, as metodologias de ensino e aprendizagem adotadas pelo curso, em especial as metodologias ativas como a Aprendizagem Baseada em Problemas e Projetos (PBL), estão diretamente alinhadas às diretrizes da extensão, conforme a Resolução nº 7, de 18 de dezembro de 2018, e a Resolução Nº 05/CEPE/UFC, de 07 de março de 2025. Ao propor desafios baseados em problemas reais, a metodologia exige que o discente utilize o conhecimento teórico (ensino), investigue e desenvolva soluções (pesquisa) e aplique essas soluções em um contexto comunitário (extensão). Esse ciclo promove a interação dialógica com a sociedade, a formação cidadã e a busca por transformação social, consolidando a extensão como um eixo central do processo de ensino-aprendizagem e não apenas uma atividade complementar.

As metodologias de ensino e aprendizagem estão em constante atualização, como fruto das experiências acadêmicas vividas por docentes e discentes, da formação continuada dos professores e também como resultado das Autoavaliações Institucionais semestrais. Conforme será visto no Capítulo 16, na UFC, os professores são avaliados em três dimensões com pesos diferentes. A primeira delas, e com maior peso (40\%), é “Planejamento pedagógico, didático e domínio do conteúdo”, que questiona ao aluno, entre outras coisas, o quanto ele concorda ou discorda da afirmação “Utilizou metodologias de ensino que motivaram e facilitaram o aprendizado”. O instrumento de avaliação dos docentes pelos discentes também apresenta um campo aberto, por meio do qual os alunos podem se manifestar livremente, opinando, entre outras coisas, sobre o desempenho do professor na disciplina avaliada. Todas essas informações permitem à Coordenação e ao Colegiado do curso realizarem o contínuo acompanhamento das atividades realizadas e das estratégias de ensino e aprendizagem adotadas, gerando subsídios para intervenções e melhorias sempre que necessário.

O curso de \nomedocurso, conforme listado nas subseções a seguir, adota alguns procedimentos e projetos acadêmicos visando a dar suporte a fatores como: desenvolvimento de conteúdos, estratégias de aprendizagem, acessibilidade metodológica e autonomia do discente.

\section{INTERDISCIPLINARIDADE}

Uma das principais características do curso de \nomedocurso é seu caráter multidisciplinar. Assim, como foi possível ver nas seções anteriores, sua matriz curricular é composta por formação básica, tecnológica, complementar, humanística e suplementar. Para que os alunos compreendam a relação entre essas áreas para sua formação acadêmica e atuação profissional, é primordial que, ao longo do curso, eles tenham a oportunidade de trabalhar com os conceitos de forma integrada. Por isso, a interdisciplinaridade, no curso, é praticada extensivamente como um dos princípios norteadores.

Dentro desse contexto, acreditamos que o conhecimento não se constitui apenas em uma linearidade organizada por disciplinas e seus pré-requisitos, mas sim, a partir da complexidade do real e das experiências significativas, de onde se origina a necessidade de fomento à interdisciplinaridade em eixos de formação \cite{pombo1993interdisciplinaridade}.

No curso de \nomedocurso, a interdisciplinaridade é explorada sob as perspectivas de interligação entre disciplinas e atividades intercursos, conforme apresentado a seguir.

% \subsection{PROJETOS INTEGRADORES}

% Ao longo do curso, os alunos passam por três componentes curriculares do tipo projetos integradores: Projeto integrado I, Projeto integrado II e Projeto integrado III. Tais projetos são integradores ``no sentido da integração curricular'' (SANTOS \& BARRA, 2012, p. 2).

% Considerando a importância da interdisciplinaridade na formação dos alunos, os três projetos integrados promovem o trabalho em equipe e a participação de professores com diferentes saberes de forma integrada para o desenvolvimento de projetos que aliem a teoria com estudos práticos reais. As disciplinas de Projeto integrado, distribuídas ao longo do curso, visam fornecer ao aluno a oportunidade de por em prática, de forma integrada, os assuntos estudados em disciplinas de semestres anteriores, bem como nas disciplinas cursadas simultaneamente com cada disciplina de Projeto integrado.

% Os projetos integrados são disciplinas que acontecem em sala de aula com a supervisão de professores. O conteúdo dos projetos integrados é essencialmente prático e visa aplicar de forma integrada os conhecimentos adquiridos nas outras disciplinas do semestre em que o Projeto integrado acontece. Os professores das outras disciplinas do semestre participam da formatação do que se espera dos projetos elaborados pelos alunos no semestre, a partir dos conteúdos abordados em cada disciplina que ocorre em paralelo, definindo em conjunto um projeto integrado que deve ser executado em grupo pelos alunos e com a supervisão e coordenação desses professores.

% Dessa forma, os projetos integrados criam oportunidades de aprendizado que permitem ao aluno praticar os conhecimentos teóricos e técnicos adquiridos durante sua formação acadêmica através da execução de projetos, objetivando vivenciar e assim melhor compreender a realidade em que o aluno está inserido social e profissionalmente. As disciplinas de Projeto integrado visam, ainda, exercitar o trabalho em equipe, a divisão de tarefas e suas responsabilidades decorrentes, reconhecer a diferença entre ação responsável e obrigações sociais, praticar a ética e seus processos e construir produtos e objetos da prática profissional relativos ao curso de \nomedocurso.

\subsection{INTERLIGAÇÃO ENTRE AS COMPONENTES CURRICULARES DO CURSO}

Além das disciplinas de Projeto integrado, a matriz curricular para o curso de \nomedocurso, está organizada de forma a facilitar a execução de ações de interdisciplinaridade entre as diversas disciplinas que compõem a formação do aluno e entre os demais aspectos que perpassam a sua formação geral como cidadão.

Buscou-se aproximar disciplinas com potencial de interação de forma a facilitar o surgimento de ações de mútua interferência entre conteúdos e conhecimentos. Além disso, procurou-se a existência de disciplinas com natureza intrinsecamente interdisciplinar ou que resultem dessa interação.

Além disso, o próprio agrupamento de disciplinas segundo UC (Seção \ref{sec:unidades-curriculares}) favorece o surgimento e a execução continuada de metodologias ou práticas interdisciplinares quando se observa que professores de componentes curriculares afins congregam-se em um fórum com o objetivo de pensar de forma coletiva as ações e revisões nas UC das quais participam.

\subsection{ATIVIDADES INTERCURSOS}

Dada a situação de um campus temático de TI, como o Campus da UFC em Quixadá, com seis cursos da área instalados, existe ainda a possibilidade de os alunos de \nomedocurso observarem as relações entre os conhecimentos explorados de seu curso e os conhecimentos explorados nos demais cursos na oportunidade em que estiverem cursando conjuntamente disciplinas desses cursos, situação possível na forma de disciplinas optativas em comum ou de optativas livres. Essa integração os capacita a observar onde e de que formas os conhecimentos do curso estão inseridos ou são capazes de influenciar os demais cursos do campus.

Portanto, considerando a integração do curso de \nomedocurso no Campus da UFC em Quixadá, é possibilitada aos alunos a oportunidade de cursar disciplinas e realizar atividades dos demais cursos de computação do campus. Através dessas atividades intercursos, o aluno poderá desenvolver várias competências, tais como:
\begin{itemize}
    \item Conhecer e identificar soluções de software para problemas em diversificados domínios de aplicação;
    \item Compreender aspectos das tecnologias de construção de software, suas possibilidades e limitações;
    \item Trabalhar em equipe com pessoas de outra formação técnica;
    \item Conhecer o mercado de trabalho relacionado a softwares e serviços.
\end{itemize}

Da mesma forma, alunos de outros cursos também poderão desenvolver competências semelhantes em relação à área de \nomedocurso. A matriz curricular do curso prevê uma parte dos créditos reservada para disciplinas optativas livres, que podem ser usadas para as atividades intercursos.

Outra forma de integração intercursos acontece quando professores de \nomedocurso unem-se a professores dos outros cursos do campus em disciplinas, pesquisas e projetos.

Existe ainda no campus uma discussão em andamento sobre a possibilidade da criação de uma disciplina nos moldes de Projeto integrado, mas que fosse ofertada para todos os cursos, a fim de fomentar a integração dos alunos de cursos diferentes em projetos interdisciplinares envolvendo múltiplas áreas de conhecimento e atuação. É também uma forma de concretizar a vocação do campus em formar recursos humanos na área de TI, mas com especializações complementares.

O campus também promove anualmente eventos acadêmicos (Workshops, Encontros Universitários e Empreenday), cuja programação tem sido enriquecida com palestras, cursos e oficinas ligadas a temas de interesse dos discentes de \nomedocurso.

\section{INTEGRAÇÃO ENTRE TEORIA E PRÁTICA}

A formação no curso integra teoria e prática de modo essencial e mútuo: o conhecimento teórico fundamenta a aplicação em situações reais, enquanto a prática enriquece e direciona a teoria. Essa conexão constante entre saber e fazer é central para o aprendizado.
%
O projeto pedagógico valoriza a interação entre professores e alunos, reconhecendo o aprendizado como um processo coletivo \cite{morin2016sete}, e busca ampliar a participação da comunidade em atividades relevantes e inovadoras. A proposta prioriza a construção conjunta do saber e a aprendizagem colaborativa, superando um modelo puramente transmissor do conhecimento.

O aprendizado é um ciclo contínuo de teoria e prática, que se informam e complementam desde o início do curso, e não apenas no estágio final. Disciplinas introdutórias, como Fundamentos de Programação, já combinam conceitos teóricos com aplicação prática em laboratório. Essa abordagem visa consolidar o aprendizado pela experiência desde os primeiros semestres, buscando reduzir a evasão associada à concentração teórica inicial.
%
Essa dinâmica de ensino mantém o curso atualizado e focado em temas relevantes para os estudantes, ao mesmo tempo, em que garante o aprendizado das teorias essenciais para a análise crítica dos conteúdos. O objetivo é um conteúdo moderno e pertinente, alicerçado em uma base teórica sólida para uma formação completa em IA.

Em suma, a integração teoria-prática desenvolve habilidades e competências cruciais para o futuro profissional, definindo o perfil do egresso e preparando os alunos para o mercado de trabalho. Aprender fazendo e compreendendo a teoria subjacente é fundamental nesta formação.

\section{FLEXIBILIDADE NA ESTRUTURAÇÃO CURRICULAR}

A estrutura curricular do curso é planejada para ser flexível e abrangente. Com poucos pré-requisitos entre as disciplinas, e com uma ampla variedade de disciplinas optativas e optativas livres, os alunos têm autonomia para construir seu próprio percurso formativo.
%
Além disso, dada a rápida evolução da área de Inteligência Artificial, a atualização contínua do currículo e a existência de disciplinas de ementa flexível (Tópicos Especiais) permitem a rápida incorporação dos últimos avanços da área.

\section{AS TIC NO PROCESSO ENSINO-APRENDIZAGEM} \label{sec:TICEA}

Os alunos geralmente já possuem familiaridade com diversas TIC que utilizam no seu cotidiano, como computadores, celulares, redes sociais e aplicativos de conversação online. Durante a sua experiência na universidade, novos sistemas e ambientes digitais são integrados à sua rotina acadêmica.
%
Muitas disciplinas do curso com carga horária prática têm aulas em laboratório. Nesses casos, os alunos têm à disposição sistemas e ambientes específicos, de acordo com as necessidades das disciplinas, instalados nos computadores dos laboratórios, conforme solicitação dos docentes. Em alguns casos, é possível instalá-los também nos computadores pessoais dos alunos, por meio de software livre, com acesso gratuito, ou software proprietário, com licenças acadêmicas liberadas aos alunos ou licenças de teste.

Para apoiar o processo de ensino-aprendizagem, o curso adota um conjunto de sistemas com diversos fins:

\begin{itemize}
    \item Sistema Integrado de Gestão Acadêmica (SIGAA)\footnote{SIGAA: \url{https://si3.ufc.br/sigaa/}}: É a ferramenta de tecnologia da informação que a UFC disponibiliza para sua comunidade e no qual os procedimentos da área acadêmica são informatizados através de módulos e portais específicos TIC{souza2015docentes}. Do ponto de vista de ensino-aprendizagem, é no SIGAA que os docentes disponibilizam o plano de aula e consolidam o diário de aula e as informações de frequência e notas, utilizados para definir as aprovações e reprovações dos alunos. É o sistema oficial da UFC para esse tipo de informação. O sistema também oferece um ambiente completo para a disponibilidade de material de aula e comunicação entre o docente e os alunos, através de notícias e fóruns.
    \item Modular Object-Oriented Dynamic Learning Environment (Moodle)\footnote{Moodle: \url{https://moodle2.quixada.ufc.br}}: É uma plataforma de aprendizagem a distância baseada em software livre, bastante popular no Brasil e no mundo, com servidor exclusivo no campus e utilizado por grande parte dos professores e alunos do curso. Em comparação com o SIGAA, o Moodle apresenta uma variedade maior de recursos para o desenvolvimento das atividades, tais como: notícias, fóruns, eventos, chat, enquete, glossário, tarefa, questionário, wiki e até um laboratório virtual de programação, por meio do qual o professor cadastra uma tarefa de programação e pode avaliá-la automaticamente com o uso do ambiente pelo aluno. %O uso do Moodle não é obrigatório, mas muitas disciplinas do curso são gerenciadas nele.
    \item Ferramentas adicionais: dependendo das características da disciplina e do interesse dos docentes e discentes, podem ser usadas ferramentas adicionais, como, por exemplo: Discord\footnote{\url{https://discord.com/}},  Slack\footnote{\url{https://slack.com/}} (ambiente de mensagens e trabalho em grupo), Trello\footnote{\url{https://trello.com/}} (ferramenta de controle de atividades e projetos), Google Sala de Aula\footnote{\url{https://edu.google.com/intl/pt-BR_ALL/k-12-solutions/classroom}} (ambiente virtual de aprendizagem), GitHub\footnote{\url{https://github.com/}} (plataforma colaborativa de desenvolvimento), Integrated Development Environments (IDEs) (ambientes completos de apoio ao desenvolvimento de software), redes sociais (Facebook, WhatsApp, Telegram e Instagram) e ambientes de armazenamento em nuvem (Google Drive, Dropbox, Microsoft OneDrive).
\end{itemize}


\section{METODOLOGIAS ATIVAS DE APRENDIZAGEM}
% regis

De forma geral, um aspecto fundamental que norteia as metodologias adotadas para o curso de \nomedocurso é o aprendizado ativo. O método ativo ou metodologia ativa tem como princípio o protagonismo do aluno, ou seja, há um “deslocamento da perspectiva do docente (ensino) para o estudante (aprendizagem)” \cite[p. 270]{diesel2017metodologias}. Em comparação com o método tradicional, ao invés da prioridade na transmissão de informações e na centralidade na figura do docente (abordagem tradicional), no método ativo “os estudantes ocupam o centro das ações educativas e o conhecimento é construído de forma colaborativa” \cite[p. 271]{diesel2017metodologias}.

\begin{quote}
Assim, em contraposição ao método tradicional, em que os estudantes possuem postura passiva de recepção de teorias, o método ativo propõe o movimento inverso, ou seja, passam a ser compreendidos como sujeitos históricos e, portanto, a assumir um papel ativo na aprendizagem, posto que têm suas experiências, saberes e opiniões valorizadas como ponto de partida para a construção do conhecimento. \cite[p. 271]{diesel2017metodologias}.    
\end{quote}

No curso de \nomedocurso, na relação professor-aluno, busca-se adotar práticas pedagógicas de formação intelectual, técnica e profissional do aluno visando o desenvolvimento de sua consciência crítica e autonomia, em linha com os Princípios Norteadores estabelecidos neste documento. O desenvolvimento desses processos constitui-se tanto na relação professor-aluno, nos momentos de encontros presenciais (sala de aula, laboratórios e demais espaços de uso comum do curso), quanto nas relações mediadas por ferramentas tecnológicas digitais (uso de ambientes virtuais de aprendizagem, ferramentas de comunicação, dentre outras).
Isso significa estabelecer que:

\begin{itemize}
    \item Professor e aluno são coautores dos saberes a serem desenvolvidos ao longo das disciplinas e do curso, evitando-se o posicionamento do professor como único detentor e disseminador de informações e conhecimentos;

    \item A aprendizagem coletiva, em grupos, ganha destaque, haja vista a necessidade de agregação de saberes, experiências e práticas diversas para o enfrentamento de questões complexas e atuais que se colocam aos aprendizes;

    \item O professor assume o papel de mediador de tal aprendizagem coletiva e de gestor de tais espaços coletivos de aprendizagem, sejam eles presenciais ou a distância.
\end{itemize}

No campus, há ainda práticas exitosas e experimentais de técnicas específicas e inovadoras alinhadas com o método ativo.

\subsection{APRENDIZAGEM BASEADA EM PROBLEMAS}
% regis


A forma mais recorrente no curso de uso da metodologia ativa é empregando a aprendizagem baseada em problemas ou projetos (PBL – \textit{Problem/Project Based Learning}). A PBL promove o pensamento crítico, a construção coletiva do conhecimento e estimula a participação ativa do estudante na sua formação. Segundo \citeonline[p. 9]{bender2014aprendizagem}, “a aprendizagem baseada em projetos é um modelo de ensino que consiste em permitir que os alunos confrontem as questões e os problemas do mundo real que consideram significativos, determinando como abordá-los e, então, agindo de forma cooperativa em busca de soluções.”

Essa metodologia vem sendo aplicada com o objetivo de incentivar uma articulação entre teoria e prática, ao longo de todo o curso e de forma profunda. Ao associar competência prática e conhecimento teórico, desde o início do curso, fornece-se uma base para a construção da autonomia intelectual do aluno, favorecendo o desenvolvimento do interesse tanto em buscar o uso prático de conhecimentos adquiridos, como buscar a fundamentação teórica para práticas conhecidas.

Nas disciplinas extensionistas, os alunos serão estimulados a desenvolver projetos envolvendo problemas reais, muitas vezes pertencentes ao seu próprio cotidiano.

\subsection{APRENDIZAGEM ENTRE PARES OU TIMES}
% regis


Uma prática muito comum em escolas e universidades é a realização de atividades (trabalhos, projetos, pesquisas, exercícios) coletivamente (em duplas ou grupos maiores). A partir da formação de equipes, procura-se que o aprendizado seja feito em conjunto e haja compartilhamento de ideias \cite{lyceum2017metodologias}.

Segundo essa visão, quando os alunos resolvem os desafios e trabalham juntos, podem beneficiar-se na busca pelo conhecimento. Com a ajuda mútua, “eles podem aprender e ensinar ao mesmo tempo, formando o pensamento crítico, construído por meio de discussões embasadas e levando em consideração opiniões divergentes” \cite{lyceum2017metodologias}.

No curso de \nomedocurso, além das disciplinas de extensão, muitas outras adotam essa abordagem, como, por exemplo, Fundamentos de Programação e Programação Orientada a Objetos.

\section{PROGRAMAS DE ACOMPANHAMENTO E AUXÍLIO A ALUNOS COM DIFICULDADES DE APRENDIZAGEM}

O curso de \nomedocurso, assim como os outros cursos do Campus da UFC em Quixadá, conta com dois programas básicos de acompanhamento para alunos com dificuldade de aprendizagem: a) Programa de Iniciação à Docência (PID), que incentiva o interesse do estudante de graduação por atividades docentes. Nesse projeto, vinculados a disciplinas específicas do curso e orientados por um professor da área, alunos mais experientes ministram atividades de monitoria e acompanhamento dos alunos de uma determinada disciplina; b) Programa de Orientação Acadêmica (POA), que busca favorecer a integração dos alunos à vida universitária, orientando-os quanto às suas atividades acadêmicas, prioritariamente nos dois anos iniciais do curso. Contribui, portanto, para o processo de socialização e ambientação dos alunos ao campus.

%No POA, há um acompanhamento contínuo feito por professores e servidores técnico-administrativos aos alunos participantes do programa. Por conhecer as dificuldades que a maioria dos alunos enfrenta para continuarem o curso até o fim, é reconhecida a importância do acompanhamento individual, onde aluno e orientador acadêmico têm a oportunidade de estabelecer uma relação com base no diálogo, com o objetivo de contornar algumas das dificuldades enfrentadas por esses sujeitos, proporcionando um momento de fala e de escuta. Assim, todos os envolvidos nesse processo beneficiam-se: os alunos, com a oportunidade de melhorar seu rendimento acadêmico e assim concluir o curso de modo exitoso; os orientadores, com a possibilidade de conhecerem mais de perto os alunos com quem dividem o espaço da sala de aula; e, por fim, o próprio processo ensino-aprendizagem (mais informações sobre o POA podem ser consultadas na Seção \ref{sec:POA}).

%\textcolor{red}{esse parágrafo deveria ser uma nova seção pois está falando de outro tipo de assistência} 
É importante ressaltar que é uma política do curso incentivar os professores a executar esses projetos, visando a melhoria do processo de formação do estudante. Além disso, o aluno do curso de \nomedocurso pode, caso necessário, utilizar o serviço de apoio psicopedagógico do campus, disponibilizado pelo Núcleo de Atendimento Social (NAS), que conta com corpo técnico especializado composto por psicólogo, assistente social e nutricionista (mais sobre o NAS na Seção \ref{sec:NAS}).

\section{ACESSIBILIDADE METODOLÓGICA}
\label{sec:acessibilidade}

A acessibilidade metodológica diz respeito à ausência de barreiras nas metodologias e técnicas de ensino-aprendizagem, considerando sempre o aprendiz em suas necessidades individuais, sejam elas relacionadas a deficiências ou não. Assim, no curso de \nomedocurso, busca-se estimular os docentes a refletirem sobre noções de conhecimento, aprendizagem, avaliação e inclusão educacional para remover as barreiras pedagógicas.

Há pelo menos três iniciativas já em andamento no campus que favorecem a acessibilidade metodológica: a) diversificação curricular: os alunos podem cursar disciplinas optativas e optativas livres, compartilhadas ou não com outros cursos; b) flexibilização do tempo: em geral, os professores mantêm-se disponíveis fora dos horários de aula para atendimento aos alunos que necessitam de acompanhamento especial; c) utilização de recursos de acessibilidade: por ser um campus temático de TI, professores e alunos estão habituados a trabalhar com tecnologias digitais, inclusive em sala de aula. Isso favorece o acesso aos softwares e materiais didáticos por estudantes com deficiência e para aqueles que preferem estudar extraclasse. Uma ferramenta importante nesse sentido é a disponibilização de materiais didáticos no Moodle, conforme descrito na Seção \ref{sec:TICEA}.

Outra preocupação relacionada à acessibilidade metodológica é o tamanho das turmas. Em alguns casos, a necessidade de acompanhamento constante e de atividades práticas exige que a turma seja dividida. Assim como realizado nos demais cursos do Campus, há a previsão de ocorrer, principalmente nos semestres iniciais, a divisão da turma. %com um mesmo professor ou com professores diferentes. 

A partir do SiSU 2018, conforme a Lei n.º 13.409, sancionada em dezembro de 2016 \cite{brasil2016lei13409}, o preenchimento das vagas deve levar em consideração também uma reserva em cada modalidade de cota para pessoas com deficiência, no mínimo igual à proporção da população da unidade da federação onde está instalada a instituição, de acordo com o IBGE. Assim, o curso deve estar preparado para receber alunos com deficiências diversas. Com o apoio da ``Secretaria de Acessibilidade UFC Inclui''\footnote{Secretaria de Acessibilidade UFC Inclui: \url{https://acessibilidade.ufc.br/pt/sobre/}} e em parceria com os profissionais da Biblioteca Universitária\footnote{Biblioteca Acessível: \url{http://www.biblioteca.ufc.br/biblioteca-acessivel/}}, dependendo da necessidade, serão disponibilizados recursos e tecnologias assistivas para comunicação e estudo como: pranchas de comunicação, texto impresso e ampliado, softwares ampliadores de tela, softwares de comunicação alternativa, softwares leitores de tela, intérpretes de Libras, entre outros recursos. Mais aspectos relacionados à acessibilidade são apresentados na Seção \ref{sec:AssistenciaEmAcessibilidade}.
        \chapter{PROCEDIMENTOS DE ACOMPANHAMENTO E AVALIAÇÃO DOS PROCESSOS DE ENSINO E APRENDIZAGEM}
\label{cap:procedimentos-de-acompanhamento}
% regis

Com o intuito de oferecer aos alunos de \nomedocurso formação de alto nível e conhecimento extenso e aprofundado nas diversas áreas que compõem essa formação, é realizado um acompanhamento detalhado do processo de ensino-aprendizagem (que contempla os aspectos de acessibilidade, como apresentado na seção \ref{sec:acessibilidade}) nos componentes curriculares e demais atividades do curso. Nesse sentido, através de atividades de avaliações escritas e práticas, os docentes podem acompanhar o amadurecimento dos alunos durante o curso, registrando os resultados individuais e gerais das turmas.

A avaliação do ensino-aprendizagem, dentro do contexto das disciplinas, caracteriza-se tanto formativa quanto somativa. As avaliações formativas são aplicadas para alterar ou aprimorar o que é visto nas aulas enquanto continuam em andamento \cite{russell2014avaliacao}. Assim, espera-se que a avaliação formativa ocorra ao longo do desenvolvimento da disciplina, com a finalidade de proporcionar informações úteis destinadas ao aprimoramento das ações executadas \cite{depresbiteris2017diversificar}. Já as avaliações somativas são úteis para avaliar os resultados do que foi ensinado \cite{russell2014avaliacao}. Espera-se, portanto, que esse tipo de avaliação determine o mérito, o valor final de um programa, com o objetivo de propiciar a tomada de decisões sobre sua continuidade ou não \cite{depresbiteris2017diversificar}. Em resumo, a avaliação formativa possibilita melhorias no processo de ensino e aprendizagem, e a somativa tem caráter mais final, mais conclusivo \cite{depresbiteris2017diversificar}.

Na maior parte das disciplinas do curso de \nomedocurso, a prática didática recorre a avaliações formativas, ou seja, ao longo do semestre letivo, os professores observam e analisam o comportamento e desempenho dos alunos. Ainda assim, muitas vezes é necessário recorrer a avaliações somativas, a fim de se ter uma visão mais sistemática da situação do aluno.

Seja de maneira formativa ou somativa, o processo avaliativo na totalidade exige um conjunto de atividades formais, sistemáticas, que levam o professor a ter condições de apresentar juízo de valor sobre determinado aspecto educacional de interesse \cite{depresbiteris2017diversificar}. Entretanto, esse juízo de valor está frequentemente carregado de certo grau de subjetividade, deixando o processo de análise e tomada de decisão em relação à aprendizagem ainda mais difícil. Para equilibrar esse fator, é interessante, portanto, que o julgamento final sobre o aluno seja emitido baseado em múltiplas situações e múltiplos instrumentos de avaliação \cite{depresbiteris2017diversificar}.

Dessa forma, vê-se o quanto são importantes a escolha e a aplicação de variados instrumentos de avaliação durante o processo de ensino-aprendizagem. Quando se fala em avaliação, costuma-se pensar em provas. Entretanto, é importante reconhecer que nem todas as decisões de avaliação exigem o uso de provas ou mensurações \cite{russell2014avaliacao}. Além disso, historicamente, a avaliação educacional vem sofrendo uma transformação radical com a mudança da cultura da prova (\textit{testing}) para a cultura da avaliação (\textit{assessment}), pois esta pressupõe uma discussão mais ampla, a partir de suas finalidades \cite{depresbiteris2017diversificar}.

Em geral, nas disciplinas do curso, as provas tradicionais são utilizadas apenas como parte do processo avaliativo. Elas não são o único instrumento, sendo, portanto, aplicadas em conjunto com outros formatos, entre os quais se destacam:
\begin{itemize}
    \item Exercícios: podem ser individuais ou em grupo; resolvidos em sala ou em casa; em forma de texto, de imagens ou de códigos; únicos ou em listas.
    \item Discussões e seminários: debates entre professores e alunos a partir de leituras recomendadas ou seminários preparados pelos alunos.
    \item Elaboração de peças e produtos específicos: os alunos são estimulados a desenvolver criativamente peças físicas (papel, argila) ou digitais dentro do contexto do que é aprendido em sala de aula.
    \item Autoavaliação: acontece quando os alunos atribuem notas a si ou ao grupo a que pertencem, contribuindo para a reflexão e crítica do que eles próprios desenvolveram ao longo da disciplina.
    \item Pesquisas: podem ser bibliográficas, de tecnologias, de campo, iconográficas. Nesses casos, os alunos buscam conhecer o que já existe.
    \item Projetos: podem ser projetos de concepção e desenvolvimento de soluções, envolvendo programação ou não; podem ser individuais ou em grupo; podem ser desenvolvidos e apresentados no final da disciplina (avaliação somativa), ou durante o semestre letivo, com o acompanhamento do professor (avaliação somativa).
    %\item Apresentação em bancas: nas disciplinas de Projeto integrado, os alunos apresentam seus projetos a uma banca avaliadora, normalmente formada pelos professores das disciplinas que contribuíram diretamente com o projeto em questão.
\end{itemize}

Vê-se, portanto, que a diversidade de metodologias, técnicas e instrumentos de avaliação contribui para uma visão mais confiável e abrangente do processo de ensino e aprendizagem, já que “a ideia de diversificar os instrumentos de avaliação tem respaldo na necessidade de que se analise a aprendizagem do aluno sob diferentes ângulos e dimensões” \cite[p. 95]{depresbiteris2017diversificar}. Além disso, a variedade de instrumentos e práticas avaliativas favorece o atendimento à diversidade das necessidades dos estudantes, pois, caso algum aluno, especialmente por questões de acessibilidade, não possa realizar uma determinada atividade avaliativa, o professor ainda assim terá alternativas para avaliá-lo.

Os procedimentos de avaliação aplicados no curso estão ainda alinhados a alguns princípios norteadores como “Respeito às diferenças e à diversidade humana”, “Desenvolvimento da capacidade crítica e da proatividade do educando” e “Integração entre teoria e prática” e a estratégias metodológicas como “Metodologias ativas” e “Acessibilidade metodológica”.

É importante esclarecer também que, independentemente do formato adotado pelos professores individualmente em suas componentes curriculares, a avaliação discente segue o Regimento Geral \cite{ufc_regimento_geral_2019} e a Resolução Nº 12/CEPE, de 19 de junho de 2008 \cite{ufc_resolucao_12_cepe_2008}, nos artigos que tratam das regras para aprovação e reprovação por nota ou por falta.

Em geral, são realizadas, no mínimo, uma avaliação no decorrer da primeira metade do semestre letivo e uma segunda avaliação no decorrer da segunda metade do semestre, para cálculo da nota final do aluno, em cada componente curricular, não se restringindo apenas a isso. Os docentes do curso são orientados a discutir os resultados das avaliações, pois esses também são momentos de aprendizado.

A avaliação do rendimento escolar por disciplina abrange a assiduidade e a eficiência, ambas eliminatórias. Com relação à assiduidade, será aprovado o aluno que frequentar 75\% (setenta e cinco por cento) ou mais da carga horária, no caso de disciplina, vedado o abono de faltas. Quando se tratar de componente do tipo atividade (Estágios e TCC), o aluno deverá frequentar 90\% (noventa por cento) ou mais da carga horária. No caso das Atividades Complementares e Atividades de Extensão, o aluno deve apresentar o comprovante adequado, de acordo com os manuais de regulamentação das respectivas atividades.

Na verificação da eficiência, será aprovado por média o aluno que, em cada disciplina, apresentar média aritmética das notas resultantes das avaliações progressivas, igual ou superior a sete. O aluno que apresentar média igual ou superior a quatro e inferior a sete será submetido à avaliação final. O aluno que apresentar a média inferior a quatro está reprovado. Na hipótese de o aluno necessitar da avaliação final, deverá obter uma nota superior ou igual a quatro, e a média dessa avaliação com a média das avaliações progressivas deve resultar em um valor superior ou igual a cinco para ser considerado aprovado. A verificação do rendimento na perspectiva do curso é realizada por meio do Trabalho de Conclusão de Curso I e II e Estágio Supervisionado, atividades obrigatórias para a conclusão do curso de \nomedocurso.

O estudante que contrair duas reprovações por falta, no mesmo componente, ou atingir um total de quatro reprovações por falta em componentes do curso, terá sua matrícula do semestre subsequente bloqueada. Para desbloqueá-la, o aluno precisa preparar um plano de estudos com a Coordenação, que considera horários de aula, estudo, descanso e lazer, além de planejamento de disciplinas a cursar nos semestres seguintes, e assinar um termo de compromisso, obrigando-se a cumprir com aquilo que foi planejado.

Mesmo antes de ser bloqueada, a Coordenação estimula o aluno com mais dificuldades a procurá-la para preparar esse plano de estudos, minimizando ao máximo os efeitos negativos das reprovações no percurso formativo do aluno. Outra ação da Coordenação, no mesmo sentido, é um acompanhamento mais atento durante o período de matrículas, antes do início de cada semestre, solicitando aos alunos que voltem a matricular-se nos componentes curriculares que tenham reprovado anteriormente, e demandando à Coordenação de programas acadêmicos, na medida da necessidade e para evitar represamentos, a abertura de vagas extraordinárias.

Um dos princípios básicos da avaliação da aprendizagem é a transparência, aos sujeitos avaliados, dos elementos passíveis de avaliação, bem como de seus mecanismos e instrumentos. No curso de \nomedocurso, assim como nos demais cursos do campus, essa transparência é estimulada através da publicação, nos primeiros dias de aula, dos planos de ensino das disciplinas. O plano é elaborado pelo professor, preferencialmente no início do semestre. Nele, além das informações básicas do componente curricular, como justificativa, objetivos, ementa e bibliografia, constam também informações específicas do andamento do componente no semestre correspondente, como a metodologia de ensino, as atividades discentes e as formas e cálculos de avaliação. Todos os planos de ensino são, obrigatoriamente, disponibilizados no SIGAA, sistema oficialmente utilizado na UFC.

Para os casos de extraordinário desempenho acadêmico dos discentes, o adiantamento de seus estudos poderá ser realizado mediante Resolução Nº 09/CEPE, de 1º de novembro de 2012 \cite{ufc_resolucao_09_cepe_2012}. Segundo esse documento, é possível conceder abreviação de estudos de componentes curriculares dos cursos de graduação, tendo o aluno de satisfazer todas as exigências preconizadas no texto do documento, bem como obter aprovação em processo avaliativo.
        \chapter{PROCESSOS DE AVALIAÇÃO INTERNA E EXTERNA DO CURSO}
\label{cap:processos-de-avaliacao-interna}

O planejamento e as ações relacionadas à gestão do curso são regularmente avaliados por meio de processos de avaliação internos e externos, cujos resultados subsidiam ajustes no planejamento. Os processos são todos tratados como ferramentas de Avaliação Institucional, mesmo quando externos ou realizados em outros cursos do campus, com os quais o curso de \nomedocurso compartilha parte de infraestrutura, corpo docente e técnico-administrativo.

A Avaliação Institucional tem como objetivo identificar o perfil e o significado da atuação da IES, através das suas atividades, cursos, programas, projetos e setores. Esse processo é norteado pelo SINAES (BRASIL, 2004b), que adota três macroprocedimentos visando valorar o mérito e a excelência de uma IES, a saber: a Autoavaliação Institucional, a Avaliação das Condições de Ensino dos Cursos de Graduação (ACE) e o ENADE.

Na UFC, a Comissão Própria de Avaliação (CPA) e as Comissões Setoriais de Avaliação (CSA) conduzem o Programa de Autoavaliação Institucional, realizando sua divulgação nas unidades acadêmicas e estimulando as coordenações de curso a realizar a discussão dos seus resultados entre os alunos e professores. No campus, tem como princípio ser um processo contínuo, viabilizado por práticas tanto de pesquisa quanto de gestão do conhecimento.

O Programa de Autoavaliação Institucional da UFC é operacionalizado através do Sistema Integrado de Gestão Acadêmica (SIGAA). A avaliação permite que os alunos expressem, semestralmente, suas opiniões sobre os trabalhos dos docentes, em quatro dimensões com diferentes pesos para o cálculo da nota do docente: planejamento pedagógico, didático e domínio do conteúdo (peso: 40\%); relacionamento e postura com os discentes (peso: 20\%); formas e usos da avaliação do aprendizado discente (peso: 20\%); e pontualidade e assiduidade às aulas (peso: 20\%).

Como parte deste processo, os discentes que cursam disciplinas com carga horária de extensão são convidados a avaliar a qualidade e o êxito dessas práticas. Essa avaliação discente permite analisar os resultados alcançados pelas atividades junto ao público participante, fornecendo subsídios diretos para o aperfeiçoamento das metodologias e da relação com a comunidade externa. Esses dados são utilizados pelo curso, por meio do Núcleo Docente Estruturante (NDE) e do Colegiado, para realizar uma análise contínua quanto as possibilidades, critérios, formas, pertinência, contribuição e resultados esperados das atividades de extensão, bem como sua contribuição para os objetivos do PPC, visando o aperfeiçoamento das características essenciais da extensão e garantindo a sua articulação com o ensino e a pesquisa.

Além da avaliação dos docentes, os discentes realizam autoavaliação e, anualmente, avaliam a infraestrutura e a coordenação do curso. Com relação à infraestrutura, são respondidas questões como se os ambientes de aprendizagem possuem tamanho adequado à quantidade de alunos da turma, se possuem adequada climatização, acústica, iluminação, mobiliários e equipamentos adequados ao ensino, além de laboratórios e acervo bibliográfico. Avalia-se também se os banheiros são limpos e adequados ao uso, e se os espaços comuns, as vias de acesso aos ambientes de aprendizagem e a biblioteca estão adaptados ao atendimento de alunos com deficiências.

Sobre a coordenação, os discentes avaliam se ela é acessível, se presta orientação e os auxilia quando necessário, se promove e divulga o PPC e estimula a participação dos alunos em encontros científicos e nos processos avaliativos do curso, se promove momentos de diálogo com os alunos sobre a formação acadêmica, currículo e mercado de trabalho, finalizando com uma avaliação geral sobre a satisfação com a coordenação do curso.

Quanto à autoavaliação realizada pelos discentes, eles respondem sobre o seu nível de assiduidade, pontualidade, envolvimento e esforço na disciplina, o nível em que os seus conhecimentos prévios contribuíram para o aprendizado e a ampliação dos conhecimentos, e sobre as competências e habilidades deles como resultado do que foi visto na disciplina.

Os docentes também avaliam seu próprio trabalho e os alunos das disciplinas que ministram. Respondem se os alunos foram assíduos e pontuais, se demonstraram motivação para o aprendizado, envolvimento com as atividades de ensino-aprendizado e responsabilidade na execução das atividades acadêmicas solicitadas, se tiveram postura adequada ao processo de ensino e aprendizado e se tinham as competências cognitivas adequadas para cursar a disciplina.

Além destas, o formulário de avaliação possui um campo aberto no qual os estudantes podem fazer comentários diretos, utilizando as próprias palavras, sobre o item que estão avaliando. Todo o processo avaliativo é feito de maneira anônima e nem docente nem coordenação são capazes de identificar o discente avaliador.
        \chapter{GESTÃO ACADÊMICA DO CURSO}
\label{cap:gestao-academica-do-curso}
\section{Coordenação do Curso}

A Coordenação de Curso de graduação é regida pelo Estatuto da UFC \cite{ufc2025estatuto} e pelo Regimento Geral da UFC \cite{ufc_regimento_geral_2019}. A Coordenação de Curso de graduação é exercida, no plano deliberativo e consultivo, pelo Colegiado de curso e, no plano executivo, pelo Coordenador de curso.

O coordenador de curso é um gestor pedagógico que deve ter o compromisso com a melhoria da qualidade do curso, atuando nas dimensões didáticas, pedagógicas, administrativas e políticas, por meio do exercício da liderança democrática, desenvolvendo ações propositivas e proativas. Na UFC, será um professor associado ou titular, ou que possua o título de doutor; na inexistência ou impossibilidade destes, um professor adjunto e, em último caso, um professor assistente, eleito em escrutínio secreto pelos integrantes do Colegiado do curso, entre os seus pares representantes de UC, para um mandato de três anos, permitida uma única recondução.

Concomitantemente com a eleição do coordenador de curso e segundo as mesmas normas, é realizada a eleição do vice-coordenador, para cumprir mandato de igual duração, a quem caberá substituir o coordenador durante suas faltas e impedimentos, bem como concluir o mandato do titular nos casos de renúncia ou afastamento definitivo.

Nas faltas e impedimentos simultâneos do coordenador e do vice-coordenador, a coordenação do curso será exercida pelo professor mais antigo, entre os seus pares representantes de UC e, no caso de empate, pelo mais idoso.

O coordenador de curso exerce o seu mandato em dedicação exclusiva ou em regime de tempo integral \cite[p. 30]{ufc2025estatuto}. A coordenação do curso de \nomedocurso é assessorada diretamente pela Secretaria Acadêmica do campus, com pelo menos um secretário dedicado ao curso.

A coordenação trabalha articuladamente com o colegiado do curso, o NDE e os discentes, no compromisso com a melhoria contínua da qualidade do curso, atuando nas dimensões didáticas, pedagógicas, administrativas e políticas, desenvolvendo ações propositivas e proativas e favorecendo a integração e a melhoria contínua das atividades realizadas no curso. As atribuições mais habituais da Coordenação do curso de \nomedocurso são:
\begin{itemize}
    \item Matrícula: durante o período de matrículas (regular e ajuste), o coordenador e/ou o vice observam, acompanham e orientam os alunos, tirando dúvidas sobre os componentes curriculares, auxiliando nas suas escolhas levando em consideração seus objetivos pessoais e procurando solucionar os problemas que porventura possam surgir, como a solicitação de vagas extras em componentes curriculares de outros cursos ou o reequilíbrio de número de vagas em disciplinas do curso.

    \item Atendimento aos discentes: a coordenação do curso mostra-se disponível para atender os alunos e orientá-los sobre questões de diferentes tipos, com relação e interferência na vida acadêmica dos discentes, como a composição das disciplinas a serem matriculadas no semestre e ao longo do curso, reclamações e sugestões sobre o andamento das aulas, a relação com professores, assuntos relacionados à estrutura do campus e orientação profissional.

    \item Atendimento aos docentes: os professores têm acesso facilitado à coordenação e partilham com ela demandas relacionadas às atividades acadêmicas que desenvolvem. Para promover as potencialidades do corpo docente, na medida do possível, a coordenação consulta os professores sobre suas preferências de disciplinas a cada início de semestre e comunica sobre eventos e congressos na área, que possam incentivar o desenvolvimento de pesquisa por parte dos docentes. A coordenação também auxilia no planejamento de atividades complementares dentro e fora do campus (como viagens de estudo, reserva de laboratórios específicos, etc.), com questões pedagógicas com discentes ou turmas, com assessoria e apoio em questões administrativas (como pedido de afastamento, progressão funcional, etc.), e no desenvolvimento de trabalho específico com professores que apresentem resultados com potencial de melhora na Avaliação Institucional, ferramenta que permite aos alunos expressarem, semestralmente, suas opiniões sobre a experiência de cursar cada disciplina em dimensões apropriadas para uma avaliação pedagógica, como apresentada no Capítulo \ref{cap:procedimentos-de-acompanhamento} deste documento.

    \item Recepção e orientação de ingressantes: cada nova turma de ingressantes deve ser recepcionada com muita atenção e zelo pela coordenação. Na primeira semana de aula, a coordenação deve apresentar detalhes sobre o curso (PPC, características, perfil do egresso, matriz curricular); recomendações e conselhos sobre a vida acadêmica; informações sobre serviços de apoio ao discente, política de bolsas e auxílios (Seção \ref{sec:bolsas-e-auxilios}), entre outros. Também nesse encontro, a coordenação começa a conhecer os ingressantes (cidades de origem, interesses, motivações e expectativas). Os alunos aproveitam a participação nesses encontros como atividade complementar.

    % \item Promoção de eventos: a Coordenação tem organizado algumas palestras de diversas áreas de interesse dos alunos do curso e do campus, por exemplo, dES – Dia da \nomedocurso, evento organizado pela coordenação do curso com o objetivo de promover discussões sobre o curso e áreas temáticas da \nomedocurso entre os alunos, egressos e professores. Também participa diretamente da promoção e organização de eventos tradicionais do campus como o WTISC, os Encontros Universitários, o InfoGirl e o Empreenday, tanto sugerindo e convidando palestrantes da área, como promovendo cursos e oficinas.

    \item Outras funções administrativas: definir pautas e convocar reuniões do colegiado do curso; tratar e encaminhar processos a secretarias e coordenadorias da UFC; auxiliar a Coordenação de Programas Acadêmicos do campus na definição de componentes curriculares ofertadas e alocação de professores, horários e salas; atribuir e acompanhar tarefas da Secretaria Acadêmica; participar do Conselho do campus (que funciona como órgão deliberativo, em todos os assuntos de sua competência, e como órgão consultivo de sua Diretoria) \cite{ufc_regimento_geral_2019}.
\end{itemize}

A atuação do coordenador tem como importante insumo os resultados da Avaliação Institucional, em que discentes se autoavaliam e avaliam a atuação docente, infraestrutura e coordenação; e docentes autoavaliam-se e avaliam suas condições de trabalho. Os dados coletados subsidiam a elaboração de um documento denominado Plano de Melhorias, que tem seu conteúdo debatido no Colegiado do curso e no NDE.

Nas atividades da coordenação há o exercício constante de promoção de uma liderança democrática, onde todos os entes envolvidos têm voz. Essa característica é reforçada nos encontros promovidos pela coordenação com os discentes, na apresentação dos resultados da Avaliação Institucional, na disponibilização de horários de atendimentos para professores e alunos e no esforço constante de criação de um ambiente saudável de trabalho e aprendizagem.

\section{Colegiado do Curso}

Na UFC, o Colegiado é regido pelas resoluções CEPE/UFC, N° 03 de 29 de janeiro de 2016 \cite{ufc2016resolucao03} e CEPE/UFC, n.º 07 de 08 de abril de 1994 \cite{ufc1994resolucao07}. O Colegiado é a instância máxima no plano deliberativo e consultivo do curso, onde são propostas, apreciadas e avaliadas as políticas e ações de gestão, e compõe, junto à coordenação, a esfera administrativa do curso.

O Colegiado do curso é formado por representação docente através das UC, com os representantes eleitos pelos pares, e por representação estudantil, também com representantes eleitos por pares, estes na proporção de 1/5 do total de docentes representantes de UC. No curso de \nomedocurso, o Colegiado é constituído pelos professores titulares representantes das UCs e seus respectivos suplentes, além de um aluno titular e seu suplente, que se reúnem mensalmente, considerando-se como pré-agendadas a realização de reuniões ordinárias na última semana do mês. Caso seja necessário, reuniões extraordinárias podem ser marcadas. %Mais detalhes sobre a composição do Colegiado do curso podem ser encontrados em \url{https://es.quixada.ufc.br/coordenacao/nde-e-colegiado/}.

Todas as deliberações são registradas em ata que, juntamente com os demais documentos de trabalho do Colegiado, ficam disponíveis em arquivos online. Caso alguma deliberação necessite de aprovação superior, a Coordenação conduz a pauta, após discutida em Colegiado, para a próxima reunião do Conselho do campus, onde serão realizados os encaminhamentos necessários.

Além das reuniões mensais, o Colegiado do curso de \nomedocurso amplia suas atividades de deliberação empregando recursos colaborativos online. A utilização de grupo de e-mails e pasta compartilhada agiliza as discussões e permite um aprofundamento em questões mais sensíveis que, sem a utilização desses recursos, não teriam possibilidade de ser debatidas, na mesma profundidade, durante as reuniões.

Uma das importantes ações do Colegiado é deliberar a respeito da demanda por componentes curriculares do curso antes do início de cada semestre. Depois de definida, a Coordenação do curso cadastra as disciplinas correspondentes no sistema ``Oferta Acadêmica'', criado pelo NPI, para auxiliar neste processo. O sistema tem funcionalidades como: importação das disciplinas ofertadas em semestres anteriores; solicitação de vagas em turmas compartilhadas entre os cursos; detalhamento das necessidades de cada disciplina, como aulas em laboratório, turmas divididas, indicação de professores, turnos preferenciais. Concluída esta etapa, a Coordenadoria de Programas Acadêmicos do campus (setor que tem a responsabilidade de acompanhar a programação acadêmica da unidade) trabalha na definição de quais e como as disciplinas serão ofertadas. Para auxiliar nesse processo, um segundo sistema (Sistema de Alocação), também desenvolvido internamente, faz a alocação automática de componentes curriculares, docentes e horários, a partir daquilo que foi demandado por cada Colegiado de curso e das regras específicas cadastradas.

Os principais insumos para a atuação do Colegiado são: os resultados das Avaliações Institucionais realizadas semestralmente (informações detalhadas no Capítulo \ref{cap:processos-de-avaliacao-interna}); demandas de alunos ou docentes encaminhadas diretamente à Coordenação do curso ou a algum dos membros do colegiado; demandas oriundas da Direção ou Conselho do campus. Durante as reuniões, pode ocorrer a participação de convidados (por exemplo, grupos de estudantes) para temas específicos. A análise desses insumos leva à elaboração do relatório de gestão de melhorias (Plano de Melhorias), que pauta a atuação da Coordenação.

\section{Núcleo Docente Estruturante - NDE}
\label{sec:NDE}

Na UFC, o Núcleo Docente Estruturante (NDE) é regido pelas resoluções CEPE/UFC nº 10/2012 (UFC, 2012b) e MEC/CONAES nº 1/2010 (BRASIL, 2010a). O NDE constitui segmento da estrutura de gestão acadêmica em cada curso de graduação, com atribuições consultivas, propositivas e de assessoria sobre matéria de natureza acadêmica e pedagógica, corresponsável pela elaboração, implementação, acompanhamento, atualização e consolidação do Projeto Pedagógico do Curso.

Com caráter de instância autônoma, colegiada e interdisciplinar, vinculada à coordenação de curso, o NDE é composto pelo coordenador do curso e, no mínimo, 5 outros docentes que possuam, preferencialmente, o título de doutor, que atuam no desenvolvimento do curso e exercem liderança acadêmica, percebida na produção de conhecimentos na área, no desenvolvimento do ensino e em outras dimensões entendidas como importantes pela instituição. A escolha dos representantes docentes é feita pelo Colegiado de curso para um mandato de três anos, com possibilidade de uma recondução. A renovação dos membros dá-se pela finalização do mandato ou por necessidade individual, de modo que parte deles permaneça, como modo de preservar o espírito do curso. %Mais detalhes sobre a composição do NDE podem ser encontrados em \url{https://es.quixada.ufc.br/coordenacao/nde-e-colegiado/}.

São atribuições do NDE:
\begin{itemize}
    \item Avaliar, periodicamente, pelo menos a cada três anos no período do ciclo avaliativo do Sistema Nacional de Avaliação da Educação Superior (SINAES) e, sempre que necessário, elaborar propostas de atualização para o PPC e encaminhá-las para apreciação e aprovação do Colegiado do curso;

    \item Fazer o acompanhamento curricular do curso, tendo em vista o cumprimento da missão e dos objetivos definidos em seu Projeto Pedagógico;
    
    \item Zelar pela integração curricular interdisciplinar entre as diferentes atividades de ensino constantes no currículo;
    
    \item Contribuir para a consolidação do perfil profissional do egresso do curso;

    \item Indicar formas de incentivo ao desenvolvimento de linhas de pesquisa e extensão, oriundas de necessidades da graduação, de exigências do mundo do trabalho e afinadas com as políticas públicas relativas à área de conhecimento do curso;
    
    \item Zelar pelo cumprimento das Diretrizes Curriculares Nacionais para os cursos de graduação;

    \item Sugerir e fomentar ações voltadas para a formação e o desenvolvimento dos docentes vinculados ao curso.
\end{itemize}

No curso de \nomedocurso, o NDE realiza todas essas atividades em parceria com a Coordenação e o Colegiado do curso, sempre com a preocupação de confirmar que o trabalho desenvolvido está alinhado com o perfil do egresso e com as novas demandas do mundo do trabalho.

O NDE do curso de \nomedocurso encontra-se constituído e atuante, reunindo-se ordinariamente uma vez ao semestre, com suas deliberações registradas em ata e alinhado com a regulamentação da UFC.

% \section{Integração com as redes públicas de ensino/Integração do curso com o sistema local e regional de saúde-SUS}

\section{Apoio ao discente}
Fortalecer o vínculo institucional do estudante pelas condições de acesso, permanência, melhoria contínua e qualidade de vida é a missão da Pró-reitoria de Assuntos Estudantis (PRAE) da UFC. Suas atribuições são ampliar as condições de permanência dos estudantes em situação de vulnerabilidade socioeconômica na UFC, viabilizar a igualdade de oportunidades entre os estudantes, contribuir para a melhoria do desempenho acadêmico individual e agir, preventivamente, nas situações de retenção e evasão decorrentes da insuficiência de condições financeiras.

No Campus da UFC em Quixadá, o Núcleo de Atendimento Social (NAS) implementa e acompanha os projetos e ações promovidos pela PRAE.

\subsection{NÚCLEO DE ATENDIMENTO SOCIAL (NAS)}
\label{sec:NAS}

O NAS foi previsto no Planejamento Estratégico realizado no Campus da UFC em Quixadá em 2013. Dentre as atividades e ações de assistência estudantil, foi prevista a criação de um setor de apoio à saúde e bem-estar dos discentes para promoção de serviços de assistência psicopedagógica dentro do campus.

Grande parte dos alunos do campus é de fora da cidade de Quixadá e passa por diversos desafios no decorrer do seu percurso acadêmico. Para auxiliá-los, o NAS promove a implementação de políticas, programas e ações de acompanhamento e orientação, além de trabalharem em conjunto com as coordenações dos cursos, no desenvolvimento de estratégias de acessibilidade metodológica e instrumental. O núcleo conta com uma equipe interdisciplinar, formada por um psicólogo, um assistente social e um nutricionista, pautada no diálogo e preparada para elaborar, implementar e coordenar projetos que promovam a articulação entre as políticas de ensino superior à assistência estudantil, integrando processos psicossociais e educacionais, identificando e superando desafios; além de viabilizar ações, serviços e programas que previnam a evasão estudantil.

O núcleo também gerencia os processos de seleção e acompanhamento de Bolsa de Iniciação Acadêmica, Auxílio Moradia, Renovação do Auxílio Moradia, Auxílio Emergencial, Auxílio Creche, Isenção do RU. Dentre as atividades desenvolvidas nesses processos, estão: a formulação de edital, a definição de cronogramas, a formação de comissão, o recebimento de documentos, a realização de entrevistas, a análise documental, a divulgação dos resultados (apoiada pelo Núcleo de Comunicação do campus), o controle das listas de frequência dos bolsistas e o acompanhamento de seu desempenho acadêmico. %Em 2017, cerca de 500 alunos do campus participaram de uma ou mais dessas atividades.

Uma das atividades diretamente relacionadas com o acolhimento e a permanência de discentes é a palestra de recepção aos ingressantes, com apresentação dos serviços e pessoal do núcleo aos  alunos que chegam ao campus no início do ano letivo. %300

Com esse mesmo objetivo, a coordenação do Projeto de Orientação Acadêmica (POA) é outra atividade importantíssima promovida pelo NAS. O projeto propõe o acompanhamento dos alunos, principalmente, dos anos iniciais, com o objetivo de aprimorar sua visão sobre o ambiente universitário e a vida acadêmica, dedicando-se ao aproveitamento nas disciplinas e em atividades que complementam sua formação. %Foi iniciado também em 2017.2, o processo de Orientação Acadêmica em Grupo, baseado na orientação acadêmica individual, mas tentando atender um maior número de alunos. Em 2021, foram realizados onze encontros, com temas relacionados ao âmbito acadêmico e à assistência estudantil. %No último ano, cerca de 43 alunos foram acompanhados por um Orientador Acadêmico no campus.  % através da tutoria de um servidor experiente (na função de orientador acadêmico)

No início do ano letivo, é aplicado o Questionário sobre o Perfil do Aluno, cujo objetivo é conhecer mais profundamente os discentes recém-ingressos, abordando questões como características socioeconômicas e demográficas, perfil do estudante, motivos da escolha do curso, expectativas e receios (SOUZA, 2016) (SOUZA, 2017).

A Jornada de Iniciação Acadêmica (JOIA) é um exemplo de evento realizado pelo NAS. Anualmente, e especialmente para os alunos participantes do processo de seleção da Bolsa de Iniciação Acadêmica, o núcleo apresenta as características da bolsa e promove ações de interação entre os discentes, com o objetivo de identificar possíveis relações entre seus perfis e os projetos desenvolvidos no campus aos quais eles podem vincular-se. %Em 2021, cerca de 108 alunos participaram do JOIA.

%Em abril de 2018, foi realizada uma mesa redonda organizada pelo NAS e pelo Núcleo de Cultura e Arte do campus com o tema “Uma discussão sobre o machismo pela perspectiva de gênero”, que contou com a participação de convidados da UFC, da Universidade da Integração Internacional da Lusofonia Afro-Brasileira (UNILAB) e da FECLESC. Destacando a importância de falar sobre o machismo de maneira geral, mas também trazendo para a realidade de mulheres na área de TIC, assim como no âmbito universitário, o evento procurou esclarecer questões sobre machismo, feminismo e dar luz a situações que mulheres passam e que acabam não percebendo, dado o enraizamento do machismo na cultura.

% Existe a previsão de ampliação das atividades do NAS, com o reforço à saúde dos discentes a partir da oferta de atendimento odontológico no campus.

A seguir, são apresentados os serviços que compõem o NAS e as ações especificamente por ele promovidas.

\subsubsection{Serviço de Psicologia}

O Serviço de Psicologia do Campus da UFC em Quixadá visa à promoção de um ambiente educativo harmônico e produtivo e tem por finalidade apoiar, orientar e acompanhar o aluno nos diversos processos de adaptação frente à realidade do contexto universitário. Com esse objetivo, o serviço procura analisar e corrigir os equívocos institucionais que comprometem a qualidade do processo de ensino-aprendizagem e as relações intersubjetivas estabelecidas entre os diversos atores que compõem a vida acadêmica, contribuindo para o bem-estar e para a qualidade de vida dos alunos do campus.

O serviço oferece aconselhamento e orientação psicológica aos alunos cuja dificuldade esteja centrada nos fatores psíquicos ligados às questões acadêmicas, focando nos aspectos cognitivos, sociais e afetivos que geram resistência e dificultam o processo de aprendizagem, o desempenho acadêmico e o bem-estar estudantil.

\subsubsection{Serviço de Nutrição}

A alimentação saudável é essencial para a manutenção da saúde e deve estar baseada em práticas alimentares que tenham significado social e cultural, além de ser acessível do ponto de vista físico e financeiro. O Serviço de Nutrição tem como objetivo auxiliar os alunos a adotar hábitos mais saudáveis a partir de duas ações principais: a orientação nutricional individualizada e a coordenação do Refeitório Universitário (RU) do campus.

O RU oferece refeições balanceadas e de qualidade, a um valor monetário simbólico, subsidiado pela universidade, aos alunos (R\$1,10 na recarga via GRU e R\$3,00 no guichê de vendas), e constitui-se como espaço de convivência e integração de estudantes, docentes e servidores técnico-administrativos.

% Durante o período da pandemia, as atividades do RU foram suspensas. Em 2019 (último ano antes da pandemia), foram servidas no campus, aproximadamente, 152.919 refeições, sendo 95.172 almoços e 57.747 jantares. No segundo semestre, teve início o processo de informatização do sistema de tickets, permitindo a consulta, em tempo real, do número de refeições em toda a instituição. O dia 23/08/2017 ficou marcado por ter sido a primeira vez que um tíquete eletrônico foi utilizado no campus, durante o jantar.

O Serviço de Nutrição avalia a qualidade do alimento servido no RU, essencialmente, de duas formas. Primeiro, a partir da promoção diária de uma pesquisa de satisfação sobre as refeições, com o objetivo de recolher a opinião da comunidade acadêmica. Após a refeição, cada comensal coloca em uma urna, conforme a sua percepção, uma avaliação que varia em três níveis, separados por tipo de proteína consumida (carne vermelha, carne branca, vegetariana). O segundo método é feito por meio da pesagem dos alimentos recebidos e do lixo descartado, conhecida como relação resto-ingestão. O estudo do resultado desses métodos permite ao Serviço de Nutrição acompanhar a qualidade do alimento servido e dá informações preciosas sobre a necessidade de alterações no cardápio do RU.

%No ano de 2021, foram atendidos, de forma remota, 47 alunos nesse serviço. 
Para além das atividades de gestão do RU, o serviço oferece atendimento nutricional à comunidade acadêmica. Após o atendimento pelo nutricionista, uma nova consulta é marcada para a entrega do plano nutricional individualizado proposto. %Além dessas, outras consultas informais, não contabilizadas, foram realizadas para fins de acompanhamento, esclarecimento de dúvidas, reformulação do plano proposto, etc.

%Complementando o atendimento individualizado, desde 2018, o Serviço de Nutrição promove o grupo “Mitos e verdades sobre alimentação”. Em cada encontro é debatido um tema específico da área, com uma breve exposição seguida de debate. Além disso, em 2020, foi criado um canal no Telegram, denominado “Shots de Nutrição”, com o objetivo de ajudar a repensar a alimentação no dia a dia, com o envio de informações gerais sobre alimentação saudável.

\subsubsection{Serviço Social}

Na área de Serviço Social, o assistente social realiza suas intervenções no atendimento à população e/ou na formulação e execução de políticas públicas que possibilitam o acesso aos direitos sociais, com base em uma formação crítica, que o capacita para realizar a análise da realidade e intervir nas várias questões apresentadas.

O Serviço Social do Campus da UFC em Quixadá volta-se para o atendimento das demandas dos discentes, com base na Política Nacional de Assistência Estudantil (PNAES, Decreto n.º 7.234/2010) \cite{decreto_7234_2010}. Além do trabalho direto com os auxílios oferecidos pela PRAE, o Serviço Social realiza orientações aos alunos e encaminhamentos para a rede socioassistencial do município de Quixadá e adjacências. Realiza ainda pesquisas e ações educativas relativas às mais variadas expressões da questão social, tais como violência, vulnerabilidade socioeconômica, dentre outros. Os atendimentos para orientação são realizados através do agendamento, realizado na sala do Serviço Social ou por e-mail.

\subsection{APOIO PEDAGÓGICO E ACADÊMICO}

Todos os docentes do curso são responsáveis pelo acompanhamento e apoio pedagógico, de maneira sistemática, com horários de atendimento aos discentes fora de sala de aula. Além deles, a coordenação tem papel importante nessa área, com o acompanhamento individual de matrícula, orientação sobre carga-horária adequada ao discente, planejamento do fluxo curricular para alunos com reprovações e o planejamento e adequação da oferta de componentes curriculares, como a inclusão de turmas extras na medida da necessidade, visando minimizar o represamento curricular.

Complementando o trabalho desenvolvido por docentes e coordenação, é importante ressaltar duas outras ações desenvolvidas no campus, alinhadas às políticas institucionais descritas no PDI da UFC: o Programa de Orientação Acadêmica (POA) e a mobilidade acadêmica.

\subsubsection{Programa de Orientação Acadêmica (POA)}
\label{sec:POA}

Ao longo do percurso formativo, inúmeros aspectos podem interferir no processo de aprendizagem, dentre eles a adaptação do aluno ao contexto universitário. O ingresso na vida acadêmica pode constituir-se como momento de crise na vida do sujeito, visto que a transição do Ensino Médio para o Ensino Superior implica no aumento de responsabilidades e na necessidade de desenvolver autonomia. Trata-se de uma fase geradora de amadurecimento e, ao mesmo tempo, desencadeadora de sentimentos de vulnerabilidade e desamparo. Soma-se a isso o fato de a maior parte dos alunos do campus vir de outras cidades e terem, ainda muito novos e inexperientes, que sair da casa dos pais, para morar em Quixadá, o que os leva a assumir responsabilidades da “vida adulta” que até então não conheciam.

Diante desse cenário, foi criado o Programa de Orientação Acadêmica (POA), desenvolvido como uma política inovadora do Campus da UFC em Quixadá, que busca promover a integração dos alunos à vida universitária, orientando-os quanto às suas atividades acadêmicas, prioritariamente nos dois anos iniciais do curso, contribuindo, dessa forma, para o processo de socialização e ambientação dos alunos ao campus.

%São duas modalidades de acompanhamento. A individual, onde cada aluno participante do programa tem um orientador específico (docente ou servidor técnico-administrativo), com quem mantém encontros regulares. O planejamento é realizado de forma singular, pensado a partir da realidade de cada aluno.

As atividades ocorrem por meio de orientação grupal, nas quais, periodicamente, são realizados encontros para trabalhar temáticas relacionadas à trajetória acadêmica (exemplos de temáticas: “Mercado de Trabalho x Área Acadêmica”, “Encontros Universitários: Por que participar?”, “Desmistificando a Universidade”, dentre outras).

%No POA, articulam-se os membros da comunidade acadêmica em diferentes níveis. São atribuições dos orientadores acadêmicos: pensar, junto ao aluno, considerando a programação acadêmica do seu curso, um fluxo curricular compatível com seus interesses e possibilidades de desempenho acadêmico; orientar a tomada de decisões relativas à matrícula; apresentar aos alunos o projeto pedagógico do curso de graduação e a estrutura universitária; encontrar-se pelo menos 2 vezes por semestre com seus alunos; e entregar à secretaria acadêmica, ao final de cada semestre letivo, relatório das atividades realizadas por cada aluno.

%Os alunos que participam do programa têm o compromisso de: manter contato com o orientador para o agendamento dos encontros; participar ativamente da construção do plano de estudos e de outras atividades propostas pelo orientador, atuando como protagonista no processo; agendar novos encontros com o orientador sempre que julgar necessário.

A comunidade acadêmica compromete-se a: acolher os estudantes no contexto universitário, viabilizando a sua integração; colaborar para a promoção de estratégias dialógicas de ensino-aprendizagem; favorecer processos comunicacionais envolvendo servidores e discentes; desenvolver a autonomia e o protagonismo dos estudantes na busca de soluções para os desafios do cotidiano universitário; e sanar os fatores de retenção, desistência e abandono, promovendo ações que identifiquem e minimizem os problemas no âmbito de cada curso.

\subsubsection{Mobilidade Acadêmica}

A mobilidade acadêmica é o processo que possibilita ao discente matriculado em uma IES estudar em outra e, após a conclusão dos estudos, obter um comprovante de estudos e, possivelmente, o aproveitamento de disciplinas em sua instituição de origem.
A mobilidade acadêmica envolve a existência de condições apropriadas, que contribuem com a formação e o aperfeiçoamento dos quadros docente e discente, objetivando a aquisição de novas experiências e a interação com outras culturas.

Os discentes do campus têm acesso a duas modalidades de mobilidade acadêmica, oferecidas pela PROGRAD e pela Pró-Reitoria de Relações Internacionais (PROINTER).

O Programa Andifes de Mobilidade Acadêmica alcança somente alunos regularmente matriculados em cursos de graduação de universidades federais, que tenham concluído pelo menos 20\% (vinte por cento) da carga horária de integralização do curso de origem e com no máximo duas reprovações acumuladas nos dois períodos letivos que antecedem o pedido de mobilidade. O estudante da UFC pode solicitar a mobilidade acadêmica a qualquer tempo, mas deve buscar informações junto à IFES de seu interesse sobre seus prazos e procedimentos.

É possível também a mobilidade acadêmica entre a UFC e instituições no exterior, em programas promovidos pela PROINTER, órgão que coordena as relações da universidade com instituições estrangeiras de educação, ciência e cultura, bem como oferece o suporte necessário à execução de convênios e acordos internacionais através das atividades desenvolvidas pelas unidades que lhe são subordinadas (no campus, a Coordenadoria de Assuntos Internacionais).

Os pedidos de inscrição dos alunos que desejem participar de programas de mobilidade acadêmica são realizados mediante encaminhamento do coordenador do curso à PROINTER, juntamente com o plano de estudos elaborado pelo aluno, contendo as disciplinas que cursará na IES desejada. Cabe ao coordenador do curso analisar as solicitações de afastamento temporário, bem como os programas das disciplinas a serem cursadas, de modo a permitir, inequivocamente, a posterior e obrigatória concessão de equivalência e consequente dispensa. O coordenador emitirá parecer conclusivo sobre as solicitações e informará a PROINTER para que esta providencie junto a IES pretendida a efetivação do Intercâmbio.

A UFC é conveniada a diferentes programas de mobilidade internacional oferecidos por diversos países, como os Programas BRAFITEC e Duplo Diploma de Graduação em engenharia (com a França), e UNIBRAL e PROBRAL (com a Alemanha), além do programa Erasmus Mundus que já beneficiou vários estudantes da UFC através de projetos coordenados pela Universidade Técnica de Munique (Alemanha), pela Universidade de Santiago de Compostela (Espanha) e pela Universidade do Porto (Portugal).

A UFC também participou do PROGRAMA ALFA, programa de cooperação entre IES da União Europeia e da América Latina, e PROGRAMA ALBAN, programa de cooperação entre União Europeia e países latino-americanos, destinado a estudantes e profissionais latino-americanos e futuros acadêmicos.

\subsection{POLÍTICA DE BOLSAS E AUXÍLIOS}
\label{sec:bolsas-e-auxilios}

%A UFC disponibiliza diversos programas de auxílios financeiros para permanência e integração acadêmica do discente, como o Programa de Bolsas de Auxílio Moradia e de Iniciação Acadêmica, que constituem auxílios para alunos socialmente vulneráveis, além da oferta de bolsas de mérito acadêmico, como as dos programas de Iniciação à Docência e de Iniciação Científica. Além disso, o campus também conta com a Coordenadoria de Estágios, responsável por apoiar e acompanhar o planejamento dos estágios curriculares obrigatórios e não-obrigatórios, como apresentado no Capítulo 10. A seguir, são apresentados os programas de bolsas e auxílios que constituem a política institucional nesse âmbito.

A UFC conta com diversas políticas e programas para apoio ao estudante, visando a permanência e a integração acadêmica dos discentes. Dentre os programas existentes, alguns oferecem auxílio financeiro para alunos em vulnerabilidade socioeconômica, como o Programa de Auxílio Moradia e o Programa de Bolsas de Iniciação Acadêmica (BIA); enquanto outros oferecem bolsas por mérito acadêmico, como os Programas de Iniciação à Docência e de Iniciação Científica.  O Campus da UFC em Quixadá também conta com a Coordenadoria de Estágios, obrigatórios e não obrigatórios, como apresentado no Capítulo \ref{cap:estagio-curricular-supervisionado}. Os programas que constituem a política institucional de bolsas e auxílios são apresentados a seguir.

\subsubsection{Programa Ajuda de Custo}

%Concede apoio aos estudantes dos cursos de graduação que desejam apresentar trabalhos em eventos de naturezas diversas. Apoia o Diretório Central dos Estudantes (DCE), os Centros Acadêmicos (CA) e as Associações Atléticas na participação em eventos, com representação de delegados e equipes de modalidades esportivas e na promoção de eventos acadêmicos, políticos, culturais e esportivos.
Visa complementar as despesas referentes ao deslocamento de estudantes, concedendo apoio financeiro aos estudantes dos cursos de graduação que desejam participar de eventos de naturezas diversas, representando a Universidade. Também apoia o Diretório Central dos Estudantes (DCE), os Centros Acadêmicos (CA) e as Associações Atléticas na participação em eventos, com representação de delegados e equipes de modalidades esportivas e na promoção de eventos acadêmicos, políticos, culturais e esportivos.

\subsubsection{Programa Auxílio Moradia}


%Tem por objetivo viabilizar a permanência de estudantes matriculados nos cursos de graduação, localizados fora dos municípios de residência e que estejam em comprovada situação de vulnerabilidade socioeconômica assegurando-lhes auxílio institucional para complementação de despesas com moradia e alimentação durante todo o período do curso ou enquanto persistir a situação. A vinculação dos estudantes ao programa não os impede de receber, por mérito, qualquer uma das bolsas dos diversos programas da UFC, de agências de fomento ou de empresas.
Tem por objetivo viabilizar a permanência de estudantes matriculados nos cursos de graduação, localizados fora dos municípios de residência, e que estejam em comprovada situação de vulnerabilidade socioeconômica. Assegura a esses estudantes um auxílio institucional para complementação de despesas com moradia e alimentação durante todo o período do curso ou enquanto persistir a situação de vulnerabilidade. A vinculação dos estudantes ao programa não os impede de receber, por mérito, qualquer uma das bolsas dos diversos programas da UFC, de agências de fomento ou de empresas.

\subsubsection{Auxílio Emergencial}

É um auxílio em dinheiro destinado aos estudantes dos cursos presenciais de graduação, que apresentem vulnerabilidade socioeconômica comprovada e que atendam a um dos requisitos que estejam impedindo que o aluno frequente normalmente as aulas. São atendidos pelo auxílio estudantes que moram em cidades diferentes do Campus  em que está matriculado, com dificuldades financeiras para realizar o  deslocamento diário;  estudantes com necessidades médicas (inclusive de saúde mental) sem os quais ficará difícil frequentar normalmente as aulas, mediante apresentação de laudo médico; estudantes que apresentem dificuldades para aquisição de material acadêmico, mediante solicitação formal do professor da disciplina e orçamento, que sejam essenciais para cursar a matéria em questão, exceto para pagamento de fotocópia e compra de livros; estudantes que estejam em situação agravante de vulnerabilidade ou risco social que não conseguiram acessar nenhum outro auxílio, bolsa ou estágio, avaliados por meio de estudo social.

\subsubsection{Auxílio Creche}

O Auxílio Creche foi regulamentado por meio do anexo \textit{ad referendum} XXI – da resolução n.º 08/2013 do CEPE. Podem solicitar este benefício estudantes regularmente matriculados em cursos de graduação da UFC, com a carga horária mínima exigida em edital, que atendam aos critérios do Programa Nacional de Assistência Estudantil (PNAES) e tenham filhos de 4 meses a 6 anos incompletos, sob sua guarda e convivendo no mesmo domicílio. O auxílio creche tem por objetivo apoiar nas despesas com o cuidado da criança, permitindo que o estudante frequente as aulas regularmente, mantenha bom desempenho acadêmico e conclua o curso no tempo previsto.



\subsubsection{Programa de Promoção da Cultura Artística}

%Mais conhecido como Bolsa Arte, foi instituído pela Resolução nº 08 do CEPE – Conselho de Ensino, Pesquisa e Extensão, em sua reunião de 26 de abril de 2013. O programa tem como objetivo principal oferecer aos estudantes, servidores docentes e técnico-administrativos da UFC condições para produção, realização e fruição de bens artístico-culturais e tem duração de um ano.
%O programa é gerido pela Secretaria de Cultura Artística (Secult-Arte/UFC), que tem por objetivo trabalhar pela articulação das iniciativas relacionadas às artes na instituição, incentivando e apoiando ações e projetos. A partir do apoio dispensado em diferentes ações, visa fortalecer a cultura artística, compreendida como dimensão inalienável da vida universitária, buscando criar estratégias para o incremento da produção estética nas diversas linguagens das artes, e estimulando a reflexão crítica sobre esta mesma produção.

Considerando a necessidade de fazer com que a cultura artística esteja presente nos processos de formação universitária desenvolvidos na UFC, o CEPE instituiu  em 26 de abril de 2013 o Programa de Promoção da Cultura Artística, também conhecido como Bolsa Arte. O programa tem como objetivo principal oferecer aos estudantes, servidores docentes e técnico-administrativos da UFC condições para a produção, realização e fruição de bens artístico-culturais e tem duração de um ano.

O programa é gerido pela Secretaria de Cultura Artística (Secult-Arte/UFC), que trabalha pela articulação das iniciativas relacionadas às artes na instituição, incentivando e apoiando ações e projetos. A partir do apoio dispensado em diferentes ações, visa fortalecer a cultura artística, compreendida como dimensão inalienável da vida universitária, buscando criar estratégias para o incremento da produção estética nas diversas linguagens das artes, e estimulando a reflexão crítica sobre esta mesma produção.

\subsubsection{Programa de Educação Tutorial (PET)}

O campus conta com dois grupos PET, a saber: PET-TI (Conexão dos Saberes, Tecnologia da Informação) e PET-SI (Sistemas de Informação). Os grupos desenvolvem ações de pesquisa, ensino e extensão, contando com o apoio de um tutor, 
com o objetivo de fortalecer os vínculos entre a instituição e a comunidade de Quixadá. O PET-TI é um grupo multidisciplinar, o qual pode ser constituído por discentes de todos os cursos do campus. O PET-SI é um grupo exclusivo para alunos do curso de Sistemas da Informação. No entanto, também desenvolve atividades em colaboração com o PET-TI.  % e, em 2022, de uma equipe de oito alunos, dois bolsistas são do curso de \nomedocurso.

A equipe de bolsistas é renovada à medida que os membros mais antigos terminam o curso. O processo de seleção é realizado por uma comissão composta por membros atuais do PET, professores e o tutor, e os bolsistas, quando iniciam suas atividades, devem ter disponibilidade de participar por pelo menos dois anos do programa.

Os grupos PETs são responsáveis pela promoção de um número importante de atividades no campus, como: a recepção e orientação dos alunos ingressantes, oferecendo informação acerca do campus, do curso e da vida acadêmica; o InfoGirl, realizado anualmente, busca atrair mulheres para a área de TI, apresentando a área para meninas do ensino médio de escolas públicas por meio de palestras, oficinas, \textit{workshops} e roda de conversa com profissionais e alunas do campus; o ensino de programação nas escolas, com minicursos ofertados em escolas públicas de Ensino Fundamental (Anos Finais) e Médio; o grupo de preparação para a OBI (Olimpíada Brasileira de Informática) e Maratona de Programação, eventos promovidos pela Sociedade Brasileira de Computação, para a resolução de desafios de lógica e programação; o ``\textit{Workshop} de Tecnologia da Informação do Sertão Central (WTISC)'', realizado anualmente, promove palestras, minicursos, mesas redondas e \textit{hackathons} em diversas  áreas de Tecnologia da Informação de forma a fortalecer e motivar a formação técnico-profissional dos estudantes; e o PET-Recebe, projeto colaborativo entre os PET-TI e PET-SI, que organiza visitas guiadas de alunos de escolas de ensino médio ao Campus.

%o Ecopet, que tem o objetivo de conscientizar os alunos e servidores sobre a importância da preservação ambiental, promovendo práticas sustentáveis dentro e fora da universidade por meio de desafios, apresentações e competições que envolvam alunos e servidores; os seminários de pesquisa, visando acompanhar as atividades individuais e coletivas do grupo, são realizados seminários de pesquisas semanais; a edição regional do “Festival latino-americano de instalação de software livre – FLISoL”, evento anual cujo objetivo é promover o uso de Software Livre, mostrando ao público em geral sua filosofia, abrangência, avanços e desenvolvimento; o “Dojo de Programação” uma série de reuniões entre desenvolvedores com diferentes níveis de experiência, com a finalidade de aprimorar técnicas e metodologias de programação; e o 

%\subsubsection{Programa Institucional de Bolsas de Administração (PIBAD)}

%O PIBAD tem como finalidade promover a inserção dos estudantes nas unidades administrativas e acadêmicas da universidade, por meio da interação dos conhecimentos inerentes a rotinas administrativas necessárias à gestão destas unidades, nos moldes estabelecidos em seus editais. O programa é gerenciado pela Pró-reitoria de Pesquisa e Pós-Graduação (PRPPG).

\subsubsection{Programa de Bolsas de Extensão Universitária}

%O programa destina bolsas ao estudante de graduação vinculado a uma ação de extensão, orientado e acompanhado por um professor ou servidor técnico-administrativo de nível superior vinculado ao quadro da UFC, e tem como objetivos: apoiar, por meio da concessão de bolsas, alunos regularmente matriculados em cursos de graduação da UFC, proporcionando o desenvolvimento de ações de extensão, com vistas à formação cidadã e à transformação social; viabilizar a participação de discentes no processo de interação entre a universidade e outros setores da sociedade através de atividades acadêmicas que contribuam para a sua formação acadêmica, profissional e para o exercício da cidadania; incentivar os processos educativos, culturais, científicos e tecnológicos, como forma de aprendizagem da atividade extensionista, articulados com o ensino e a pesquisa de forma indissociável e que viabilizem a relação transformadora entre a universidade e outros setores da sociedade, contribuindo para a inclusão social; fomentar o interesse em extensão universitária e incentivar novos talentos potenciais entre estudantes de graduação, assim como contribuir para a formação e a qualificação de cidadãos socialmente comprometidos.

O programa destina bolsas como auxílio financeiro ao estudante de graduação vinculado a um Programa ou a um Projeto de Extensão, orientado e acompanhado por um professor ou servidor técnico-administrativo de nível superior vinculado ao quadro da UFC. O programa tem como objetivos: apoiar, por meio da concessão de bolsas, alunos regularmente matriculados em cursos de graduação da UFC, proporcionando o desenvolvimento de ações de extensão, com vistas à formação cidadã e à transformação social; 
possibilitar a implementação do processo de curricularização da
extensão nos cursos de graduação da UFC no âmbito da Unidade Curricular
Especial de Extensão;
viabilizar a participação de discentes no processo de interação entre a universidade e outros setores da sociedade através de atividades acadêmicas que contribuam para a sua formação acadêmica, profissional e para o exercício da cidadania; 
incentivar os processos educativos, culturais, científicos e tecnológicos, como forma de aprendizagem da atividade extensionista, articulados com o ensino e a pesquisa de forma indissociável e que viabilizem a relação transformadora entre a universidade e outros setores da sociedade, contribuindo para a inclusão social; 
fomentar o interesse em extensão universitária e incentivar novos talentos potenciais entre estudantes de graduação, assim como contribuir para a formação e a qualificação de cidadãos socialmente comprometidos.

A extensão no campus guarda consonância com as orientações da PREx, com o esforço de integração entre ensino, pesquisa e extensão, no sentido de levar conhecimento prático à comunidade. Em 2025, um total de 32 servidores docentes e 12 servidores técnico-administrativos do campus esteve envolvido em 29 ações de extensão, com a participação de 46 alunos.

%\subsubsection{Programa de Desenvolvimento Institucional em Tecnologia da Informação}


\subsubsection{Programa de Iniciação à Docência (PID)}

O Programa de Iniciação à Docência (PID) é desenvolvido em duas modalidades, monitoria remunerada e monitoria voluntária. Na primeira, o monitor recebe uma bolsa-auxílio para desempenhar as funções e, por isso, não deve participar de qualquer outra atividade remunerada, seja pública ou privada. Na segunda, o monitor desempenha as atividades de maneira voluntária, sem o recebimento do auxílio.

A carga horária da monitoria é de 12 horas semanais e deve ser cumprida sem afetar as demais atividades acadêmicas do estudante. A duração é de 10 meses, mas a monitoria pode ser renovada uma vez, por igual período, caso o bolsista seja novamente aprovado em processo seletivo. A função de monitor não constitui cargo ou emprego, nem representa vínculo empregatício de qualquer natureza com a universidade, e é uma importante estratégia para o nivelamento dos discentes com mais dificuldades, que têm a oportunidade de reforçar seus estudos com colegas mais experientes.

Algumas das atividades dos bolsistas de PID são: elaborar, juntamente com o professor-orientador, o plano de trabalho da monitoria; participar das tarefas didáticas, inclusive na programação de aulas e em trabalhos escolares; auxiliar o professor-orientador na realização de trabalhos práticos e experimentais, na preparação de material didático e em atividades de classe e/ou laboratório; contribuir, juntamente com o professor-orientador, para a avaliação do andamento da disciplina ou da área. Em 2025, o Campus da UFC em Quixadá teve 18 projetos aprovados, com 25 vagas de bolsas remuneradas e 20 vagas para voluntários.

%Em 2021, 25 bolsistas de PID eram alunos do Campus da UFC em Quixadá.

\subsubsection{Programa Institucional de Bolsas de Iniciação Científica (PIBIC)}

O Programa Institucional de Bolsas de Iniciação Científica (PIBIC) é um programa voltado para o desenvolvimento do pensamento científico e iniciação à pesquisa de estudantes de graduação do ensino superior. Atualmente, o PIBIC é o principal programa de iniciação científica da universidade, resultado de um convênio entre a UFC, o CNPq e a Fundação Cearense de Apoio ao Desenvolvimento Científico e Tecnológico (Funcap).


%Em 2022, o PIBIC contou com 601 cotas de bolsas do CNPq. Essas bolsas são destinadas às várias unidades da graduação. Cada orientador poderá ter, no máximo, dois bolsistas. Pesquisadores que já possuem bolsa BPI da Funcap podem ter, no máximo, uma bolsa. As bolsas BPI são uma espécie de bolsa de produtividade em pesquisa da agência estadual.

Podem concorrer às quotas de bolsa, além de docentes, servidores técnico-administrativos da universidade portadores do título de doutor, em regime de dedicação exclusiva ou 40 horas. São usados critérios de produtividade, produção científica, tecnológica e artística no julgamento dos projetos.

Nos últimos anos, os editais PIBIC reservaram 20\% das cotas de bolsa da UFC para estudantes dos campi no interior do estado. Em 2024, o campus da UFC em Quixadá contou com 11 bolsas PIBIC de diferentes projetos. 

\subsubsection{ Programa Institucional de Bolsas de Iniciação em Desenvolvimento Tecnológico e Inovação (PIBITI)}

O Programa Institucional de Bolsas de Iniciação em Desenvolvimento Tecnológico e Inovação (PIBITI) tem por objetivo principal estimular os jovens do ensino superior nas atividades, metodologias, conhecimentos e práticas próprias ao desenvolvimento tecnológico e processos de inovação. 

Outros objetivos do PIBITI são: i) Contribuir para a formação e inserção de estudantes em atividades de pesquisa, desenvolvimento tecnológico e inovação; ii) Contribuir para a formação de recursos humanos que se dedicarão ao fortalecimento da capacidade inovadora das empresas no país; e iii) Contribuir para a formação do cidadão pleno, com condições de participar de forma criativa e empreendedora na sua comunidade.


%4 - 2022
%4 - 2021
% 2023 - 4 e 4 voluntárias
% 2024 - 1 e 2 volntária

Os responsáveis pela gestão do programa PIBITI são a Coordenadoria de Pesquisa da Pró-Reitoria de Pesquisa e Pós-Graduação (PRPPG) da UFC, com a ajuda do Comitê Interno. Periodicamente, o programa é avaliado pelo Comitê Externo composto por pesquisadores destacados de outras Instituições de Ensino Superior.

Em 2021 e 2022, o Campus Quixadá contou com 4 bolsas PIBITI. Em 2023, foram 4 bolsas remuneradas e 4 voluntários. Em 2024, houve um bolsista remunerado e dois voluntários.

\subsubsection{Programa de Iniciação Acadêmica}

O Programa de Iniciação Acadêmica tem por objetivo  propiciar aos estudantes de cursos de graduação presenciais da UFC – em situação de vulnerabilidade socioeconômica comprovada, especialmente os de semestres iniciais – condições financeiras para sua permanência e desempenho acadêmico satisfatório por meio  da concessão de uma bolsa e integração ao ambiente universitário. O bolsista deve atuar, em caráter de iniciação acadêmica, nas diversas unidades da instituição. Além do auxílio financeiro, o programa também proporciona ao estudante um espaço para conhecer e participar de ações de iniciação à universidade, ensino, pesquisa e extensão, artes, gestão, dentre outros.

O campus Quixadá dispõe anualmente de 80 vagas para Bolsas de Iniciação Acadêmica. Ao longo do ano, é comum que alguns bolsistas se desliguem do programa para assumir outras oportunidades. Nesses casos, novas chamadas são realizadas a partir da lista de espera, o que pode resultar em um número total de beneficiados superior ao número inicial de vagas.


%O Programa Bolsa de Iniciação Acadêmica tem por objetivo propiciar aos estudantes de cursos de graduação presenciais da UFC – em situação de vulnerabilidade socioeconômica comprovada, especialmente os de semestres iniciais – condições financeiras para sua permanência e desempenho acadêmico satisfatório, mediante atuação, em caráter de iniciação acadêmica, nas diversas unidades da instituição. 

%Site O Programa de Iniciação Acadêmica é uma forma de oferecer um espaço para que o estudante possa conhecer e participar de ações de iniciação à universidade, ensino, pesquisa e extensão, artes, gestão, dentre outros; além de ter a oferta direta de uma bolsa

%Em 2021, um total de 88 bolsas foram disponibilizadas para os alunos do campus e 11 (treze) alunos de \nomedocurso participaram da bolsa este ano em projetos.

\subsubsection{Programa de Incentivo ao Desporto}

O Programa Bolsa de Incentivo ao Desporto incentiva os estudantes a incrementarem o seu desempenho desportivo e acadêmico, mediante atuação em atividades relativas à gestão de atléticas, à assessoria desportiva e ao rendimento desportivo. Cada Associação Atlética possui pelo menos um bolsista para organizar eventos e representá-la ante ao Desporto Universitário. Com isso, a Bolsa de Incentivo ao Desporto é organizada nas seguintes modalidades: Gestão de Atléticas; Assessoria Desportiva; e Rendimento Esportivo.

\subsubsection{Programa de Acolhimento e Incentivo à Permanência (PAIP)}

O Programa de Acolhimento e Incentivo à Permanência (PAIP) tem como objetivo reduzir a evasão nos cursos de graduação da UFC por meio da oferta de bolsas a estudantes que participam de projetos voltados à articulação, acompanhamento e avaliação das ações acadêmicas. Desenvolvidos em diversas áreas e unidades da universidade, esses projetos auxiliam na adaptação dos alunos nos semestres iniciais, fortalecem a qualidade do ensino e da aprendizagem e contribuem para a autoavaliação dos cursos.

Há duas modalidades de bolsa PAIP: a Bolsa de Apoio a Projetos de Graduação, que visa promover a articulação, o acompanhamento e a avaliação de ações acadêmicas,  ampliar as iniciativas de atividades alternativas e inovadoras de ensino que contribuam para a redução de fatores determinantes para a reprovação e evasão dos estudantes; e a Bolsa de Apoio à Gestão Acadêmica da Prograd, que visa propiciar oportunidade ao estudante de graduação para o aprimoramento de sua formação, junto à gestão universitária, permitindo-lhe uma ampliação de seus conhecimentos acerca da estrutura e da dinâmica acadêmica. 

%A bolsa tem dois objetivos principais: de propiciar oportunidade ao estudante de graduação para o aprimoramento de sua formação, junto à gestão universitária, permitindo-lhe uma ampliação de seus conhecimentos acerca da estrutura e da dinâmica acadêmica; criar espaço de participação discente nas ações desenvolvidas pelo Gabinete da PROGRAD e pelas coordenadorias que compõem da pró-reitoria. 

A bolsa tem vigência de no máximo dez meses, relativos ao período de março a dezembro de cada ano, com carga horária de 12 horas semanais, nos turnos da manhã e/ou da tarde, de acordo com as indicações dos projetos. Os bolsistas PAIP devem participar do planejamento e da execução das ações da coordenadoria ou projeto ao qual estão vinculados; redigir relatórios parciais e final das atividades desenvolvidas; e apresentar trabalho nos Encontros Universitários da UFC. Em 2025, o Campus de Quixadá contou com 8 vagas remuneradas do Programa PAIP. 

%Em 2021, um total de 13 alunos do campus foram bolsistas desse programa.


%São atividades dos bolsistas PAIP: organizar os dias e turnos de atividades (manhã, tarde ou noite), considerando a disponibilidade de horário do bolsista e, também, as necessidades da coordenadoria ou projeto ao qual está vinculado; participar do planejamento e da execução das ações da coordenadoria ou projeto ao qual está vinculado, bem como redigir relatórios parciais e final das atividades desenvolvidas; e apresentar trabalho nos Encontros Universitários da UFC. Em 2021, um total de 13 alunos do campus foram bolsistas desse programa.

\subsection{ASSISTÊNCIA EM ACESSIBILIDADE} 
\label{sec:AssistenciaEmAcessibilidade}

%Desde agosto de 2010, a UFC conta com um setor exclusivo para elaborar ações rumo à inclusão de pessoas com deficiência, a Secretaria de Acessibilidade UFC Inclui, que, como o nome sugere, busca integrar pessoas cegas, surdas, cadeirantes e com outras limitações de mobilidade no dia a dia da instituição.

Desde agosto de 2010, a UFC conta com um setor exclusivo para elaborar ações rumo à inclusão de pessoas com deficiência, a Secretaria de Acessibilidade UFC-INCLUI. A UFC-INCLUI busca integrar pessoas cegas, surdas, cadeirantes e com outras limitações de mobilidade no dia a dia da instituição.


Com três eixos de atuação – tecnológico, atitudinal e pedagógico – a secretaria trabalha na formulação de uma política central de acessibilidade na UFC, agindo para que esta seja respeitada e implementada nos diversos espaços da universidade.

%Não se trata de um órgão executor – embora ofereça serviços como digitalização de textos, ledores, revisão de projetos arquitetônicos, entre outros – e, sim, de um núcleo de fomentação e acompanhamento de ações intersetoriais.



%Livia: parei aqui


%Não é objetivo da secretaria absorver todas as ações referentes à inclusão, uma vez que a tarefa de acolher pessoas com deficiência diz respeito a toda a sociedade, cabendo ao órgão disseminar a cultura inclusiva e despertar na comunidade universitária o compromisso com o respeito aos direitos desse público. É por isso que a secretaria trabalha na descentralização das iniciativas de acessibilidade, oferecendo suporte e orientação a professores, coordenadores, chefes de departamento, servidores técnico-administrativos e estudantes interessados em fazer sua parte desse desafio.

Uma vez que a tarefa de acolher pessoas com deficiência diz respeito a toda a sociedade, não é objetivo da secretaria absorver todas as ações referentes à inclusão, cabendo ao órgão disseminar a cultura inclusiva e despertar, na comunidade universitária, o compromisso com o respeito aos direitos desse público. A secretaria atua, portanto, na descentralização das iniciativas de acessibilidade, oferecendo suporte e orientação a professores, coordenadores, chefes de departamento, servidores técnico-administrativos e estudantes interessados em fazer sua parte nesse desafio. Além disso, embora ofereça serviços como digitalização de textos, ledores, revisão de projetos arquitetônicos, entre outros, é importante ressaltar que a UFC-INCLUI não se trata de um órgão executor, mas sim de um núcleo de fomentação e acompanhamento de ações intersetoriais.

São atribuições da secretaria: elaborar e gerenciar ações de acessibilidade; oferecer suporte às unidades acadêmicas para a efetivação da acessibilidade na UFC; estimular a inserção de conteúdos sobre acessibilidade nos projetos pedagógicos de cursos de graduação, contribuindo para a formação de profissionais sensíveis ao tema; identificar e acompanhar os alunos com deficiência na UFC; identificar metodologias de ensino que representam barreiras para os alunos com deficiência e propor estratégias alternativas; estimular o desenvolvimento de uma cultura inclusiva na universidade; oferecer serviços de apoio a esse público, como digitalização e leitura de textos acadêmicos, cursos de Língua Brasileira de Sinais (Libras), revisão de processos arquitetônicos com base em critérios de acessibilidade, entre outras ações; promover a formação de recursos humanos em gestão de políticas relacionadas às pessoas com deficiência, qualificando-os para um atendimento adequado; promover eventos para informar e sensibilizar a comunidade universitária; estimular o desenvolvimento de pesquisas de Avaliação Pós-Ocupação nos prédios da UFC; estimular a acessibilidade em ambientes virtuais e nos produtos e eventos de comunicação e marketing; e oferecer orientação e apoio pedagógico a coordenadores e professores, estabelecendo um canal de comunicação entre estes e os estudantes com deficiência.

%O trabalho da secretaria é desenvolvido em três eixos. O eixo atitudinal relaciona-se à ideia de que a inclusão é uma questão de atitude e de sensibilidade. É preciso ajudar a comunidade acadêmica a enfrentar o preconceito e incentivar mudanças de atitude, visando à remoção de barreiras que impedem a acessibilidade. O eixo tecnológico tem a ver com o incentivo de pesquisas e ações em tecnologias assistivas, para o desenvolvimento de equipamentos, serviços e estratégias que permitam o acesso ao conhecimento com autonomia. E o eixo pedagógico concentra-se na ideia de que não basta fazer com que o estudante com deficiência ingresse na Universidade – é preciso oferecer condições para que ele tenha a mesma formação que os colegas. Por isso, a secretaria promove ações que facilitem o ensino-aprendizagem, com alternativas de avaliação.

O trabalho da secretaria é desenvolvido em quatro eixos. O eixo atitudinal relaciona-se à ideia de que a inclusão é uma questão de atitude e de sensibilidade. Assim, a secretaria atua ajudando a comunidade acadêmica a enfrentar o preconceito e a incentivar mudanças de atitude, visando à remoção de barreiras que impedem a acessibilidade. No eixo tecnológico, o objetivo é incentivar pesquisas e ações em tecnologias assistivas, para o desenvolvimento de equipamentos, serviços e estratégias que permitam o acesso ao conhecimento com autonomia. O eixo pedagógico concentra-se na ideia de que não basta fazer com que o estudante com deficiência ingresse na Universidade – é preciso oferecer condições para que ele tenha a mesma formação que os colegas. Por isso, a secretaria promove ações que facilitem o ensino-aprendizagem, com alternativas de avaliação. Por fim, o eixo comunicacional busca oferecer recursos, atividades e bens culturais que promovam independência e autonomia aos indivíduos que necessitam de serviços específicos para acessar o conteúdo proposto, tais como audiodescrição, legendas, janela de Libras, impressões em braille, dublagem, etc.

\subsection{Política de Cotas para Pessoas com Deficiência}
Conforme a Lei n.º 12.711, sancionada em agosto de 2012 (BRASIL, 2012a), a UFC reserva 50\% de suas vagas para alunos que tenham cursado integralmente o Ensino Médio público, em cursos regulares ou da educação de jovens e adultos. Os demais 50\% das vagas permanecem para ampla concorrência. Destas vagas reservadas para a escola pública, metade é destinada a estudantes com renda mensal familiar de até um salário mínimo e meio. O preenchimento das vagas leva em conta ainda critérios de cor ou raça, seguindo dados estatísticos do Instituto Brasileiro de Geografia e Estatística (IBGE).

A partir do SiSU 2018, conforme a Lei n.º 13.409, sancionada em dezembro de 2016 (BRASIL, 2016a), o preenchimento das vagas começou a levar em consideração também uma reserva, em cada modalidade de cota, para pessoas com deficiência, no mínimo igual à proporção na população do estado do Ceará, de acordo com o IBGE.

Somente em 2018, na chamada regular do SiSU, a UFC adotou cotas para pessoas com deficiência. Em 2022, a Universidade recebeu 69 candidatos com deficiência, via SiSU, que ingressaram na instituição. Desse total, 08 vieram para o Campus Quixadá. A importância da adoção de cotas específicas para pessoas com deficiência no processo de democratização do ensino superior é um avanço importante para a universidade, que assume seu papel de atendimento à comunidade e amplia a inclusão de modo mais efetivo, com todas as cotas existentes hoje.


\subsection{Infraestrutura para facilitação da acessibilidade }
Para fazer frente ao esperado aumento do número de pessoas com deficiência na universidade, a secretaria UFC-INCLUI está se preparando para ampliar o atendimento a esse público, mediante a aquisição de novos equipamentos, contratação de intérpretes e ampliação da interlocução com as unidades acadêmicas, a fim de garantir formação de qualidade a todos.

O campus prevê em sua infraestrutura a facilitação da acessibilidade a pessoas com dificuldades de locomoção ou visão, contando com plataformas elevatórias, portas largas de acesso às salas e laboratórios, banheiros com cabine específica para deficientes e plaquetas com inscrições em braille, além de rampas de acesso ou mesmo ausência de degraus desde o estacionamento até todos os ambientes térreos. Além disso, trabalha em estreita ligação com a secretaria no atendimento aos discentes com deficiência.

\section{Gestão do curso com base nos processos de avaliação interna e externa}

A coordenação tem um papel fundamental no processo de avaliação, ao analisar os dados dos relatórios da Autoavaliação Institucional e promover a participação massiva dos discentes, para a promoção da melhoria dos cursos. A coordenação do curso de \nomedocurso busca assim promover uma série de ações para conscientizar os alunos da importância da participação na Autoavaliação Institucional. %É realizada a divulgação comum em sala de aula para que os alunos participem da Autoavaliação Institucional.

Idealmente, a Autoavaliação Institucional busca a participação responsável e efetiva da maioria dos seus agentes, egressos, discentes, docentes e servidores técnico-administrativos. O objetivo é construir uma cultura interna favorável à autoavaliação, que possibilita maior conscientização acerca da missão, bem como das finalidades acadêmica e social da UFC, consolidando assim a noção de que a Autoavaliação Institucional é importante via para a reflexão coletiva e, por conseguinte, para o planejamento institucional participativo.

Além da campanha de sensibilização, esse bom resultado é creditado também à política de divulgação dos resultados do processo. Semestralmente, a coordenação, após a sistematização dos dados, reúne-se com o corpo discente e os apresenta num evento chamado “Seminário de Autoavaliação Institucional”, onde são respondidas questões sobre a importância da avaliação, como os discentes podem participar e as ações a serem tomadas a partir das informações recolhidas.

Durante o seminário, a coordenação do curso de \nomedocurso apresenta a comparação entre os resultados das avaliações dos discentes com o resultado do campus e os da UFC. Esses seminários também são oportunidades para discutir com os alunos questões gerais sobre o curso e sobre o PPC, contando inclusive como atividade complementar na categoria Vivências em Gestão, pois se considera que esse momento aproxima os alunos do acompanhamento e dos processos de tomada de decisão do curso.

A partir da análise dos dados, são estabelecidas metas de resultados positivos e negativos para cada quesito avaliado, e tomadas atitudes em caso de má avaliação. Por exemplo, se um docente tiver uma avaliação negativa na dimensão ``Planejamento pedagógico, didático e domínio do conteúdo'', isso implica que a coordenação analise a situação sob dois cenários: primeiro, quando o professor se autoavalia com ótimo desempenho, tem-se uma situação que exige intervenção imediata da coordenação; já quando o professor se autoavalia reconhecendo suas dificuldades, implica em outro tipo de intervenção por parte da coordenação. A partir da distinção dos diferentes cenários, o coordenador tem informações objetivas que permitem melhor gerenciar as potencialidades do corpo docente e favorecer a melhoria contínua de suas atividades.

A análise conjunta dos diferentes dados da Autoavaliação Institucional proporciona subsídios importantíssimos para a gestão efetiva do curso. A partir desses resultados e das reuniões de apresentação, é elaborado um Plano de Melhorias, acompanhado pela CPA, com ações de melhoria e cronograma de trabalho aprovados pelo Colegiado do curso e pelo NDE.

Na área administrativa, há um esforço contínuo de mapeamento de processos, que naturalmente induz à melhoria das rotinas do curso. Regularmente, também são realizadas pesquisas socioeconômicas com discentes, que apontam as principais questões acadêmicas e sociais que interferem no desempenho ou permanência do aluno no curso.

\subsection{Acompanhamento de egressos do curso}

Para o curso de \nomedocurso, é importante manter o vínculo com os egressos, pois se trata de uma maneira importante de avaliar aspectos diversos do curso. O egresso, estando em plena atuação profissional ou não, tem o potencial de trazer uma percepção atualizada das exigências da sociedade e do mercado, mostrando-se como um elemento relevante para as atividades do curso.

Em 2009, foi criada a Associação dos ex-alunos da UFC (ASSOEX) com o objetivo de congregar aqueles que já passaram pelas salas de aula da instituição, tendo em vista a manutenção de sua proximidade com a universidade, e de criar mecanismos que promovam a sua plena integração à vida acadêmica, política e cultural da instituição. Também é meta da associação despertar nos ex-alunos o interesse pela promoção sociocultural da UFC, garantindo o acesso deles às instalações acadêmicas, esportivas e culturais em iguais condições de tratamento dos atuais alunos e professores.

Além da possibilidade de fazer parte da ASSOEX, listamos abaixo algumas ações planejadas para manter o vínculo com os egressos de \nomedocurso:

\begin{itemize}
    \item Manutenção de cadastro atualizado dos egressos: a cada semestre, os dados de contato dos alunos formados serão coletados, registrados e mantidos pela secretaria acadêmica e Coordenação do curso.
    
    \item Manutenção de uma conta institucional do curso em redes sociais profissionais para o acompanhamento dos egressos.
    
    \item Promoção de eventos com participação de egressos: o Campus da UFC em Quixadá frequentemente promove eventos com a participação de egressos em palestras, minicursos e mesas redondas. Alguns desses eventos são o ExpoSE, WTISC, Flisol, os Encontros Universitários e o InfoGirl.
    
    \item Pesquisa com egressos: desde 2016, tem-se realizado anualmente uma pesquisa de levantamento com ingressantes dos cursos do Campus da UFC em Quixadá. Pretende-se comparar os dados obtidos com os ingressantes anualmente, com os dados que serão coletados quando estes alunos estiverem próximos da formatura, a fim de entender melhor o percurso por eles traçado e compreender suas expectativas profissionais e acadêmicas ao final do curso. Ainda em 2016, a pesquisa de mestrado de uma servidora do Campus da UFC em Quixadá focou na inserção dos egressos de campi do interior no mercado regional, trazendo importantes informações a respeito da realidade desses alunos (NUNES, 2016).
\end{itemize}
        \chapter{INFRAESTRUTURA DO CURSO}
\label{cap:infraestrutura-do-curso}

%INFRAESTRUTURA
\def\ufcNumEquipEletro{266\xspace}
\def\ufcNumNotebooks{28\xspace}
\def\ufcNumProjetores{33\xspace}
\def\ufcNumDesktops{213\xspace}
\def\ufcNumServidores{6\xspace}
\def\ufcNumVoip{16\xspace}
\def\ufcNumNobreaks{3\xspace}
\def\ufcNumLab{sete\xspace}
\def\ufcNumGabinetes{35\xspace}
% RECURSOS HUMANOS
\def\ufcNumDocentes{59\xspace}
\def\ufcNumDocSubstitutos{1\xspace}
\def\ufcNumDocQuarentaHoras{1\xspace}
\def\ufcNumDocDE{57\xspace}
\def\ufcDocMestres{14\xspace}
\def\ufcDocDoutores{45\xspace}
\def\ufcDocDoutorandos{7\xspace}
\def\ufcDocAfasDoutorado{1\xspace}
\def\ufcNumSTA{29\xspace}
\def\ufcNumSTAMedio{13\xspace}
\def\ufcNumSTASuperior{16\xspace}
\def\ufcSTAEscMedio{2\xspace}
\def\ufcSTAEscGrad{10\xspace}
\def\ufcSTAEscEspec{12\xspace}
\def\ufcSTAEscMestre{5\xspace}
%BIBLIOTECA
\def\ufcBCQAnoReferencia{2025\xspace}
\def\ufcBCQTitulos{1.425\xspace}
\def\ufcBCQExemplares{9.524\xspace}
\def\ufcBCQDocsRepositorio{74 mil\xspace}


% Para o devido aprendizado do corpo discente e de modo a permitir espaços e situações de experimentação, os docentes e discentes dispõem de toda a infraestrutura do campus da UFC de Quixadá, conforme será visto a seguir.

O Campus da UFC em Quixadá conta com uma estrutura física que contempla um estacionamento com vagas para carros de passeio (algumas reservadas para pessoas com deficiência), vagas para motocicletas, vagas para ônibus e bicicletário; espaço de convivência que acomoda a cantina e salas de centro e diretório acadêmico; cinco blocos, sendo quatro blocos didáticos medindo 1.400 m$^2$, cada um, e um bloco administrativo medindo 1.500 m$^2$. Todos os blocos contam com dois pavimentos, térreo e superior. Portanto, o Campus da UFC em Quixadá conta hoje com a infraestrutura arquitetônica descrita abaixo:%, excluindo-se os blocos ainda em construção:

\begin{itemize}
    \item Bloco I (1400 m$^2$): salas de aula, salas de projetos (PET), salas para serviços (secretaria acadêmica, atendimento nutricional, videoconferência), banheiros, gabinetes para professores, laboratórios, copa, plataforma elevatória.

    \item Bloco II (1400 m$^2$): salas de aula, salas de projetos (empresa júnior), salas de estudo em grupo, biblioteca, salão multiuso, banheiros, gabinetes para professores, laboratórios, copa, plataforma elevatória.

    \item Bloco III (1400 m$^2$): salas de aula, salas de desenho, ateliê, salas de estudo em grupo, laboratórios, banheiros, gabinetes para professores, copa, plataforma elevatória.

    \item Bloco IV (1400 m$^2$): salas de aula, salas de projetos (NPI), salas de estudo em grupo, banheiros, gabinetes para professores, laboratórios, copa, plataforma elevatória.

    \item Bloco Administrativo (1500 m$^2$): refeitório universitário, salas para serviços acadêmicos e administrativos (venda de tickets, direção, secretaria da direção, prefeitura, coordenação de programas acadêmicos, núcleo de TIC), salas das coordenações dos cursos de graduação, sala de reuniões, sala do núcleo de inovação, laboratório de redes de alta velocidade, salas de atendimento do NAS (nutrição, serviço social, serviço de psicologia), copa, banheiros.

    \item Área de Convivência (500 m$^2$): cantina, salas de centros e diretórios acadêmicos, área comum coberta, banheiros.
\end{itemize}

O campus conta com \ufcNumEquipEletro equipamentos eletrônicos para atender a demanda interna, sendo estes: \ufcNumNotebooks notebooks; \ufcNumProjetores projetores; \ufcNumDesktops desktops; \ufcNumServidores máquinas servidoras para telemática; \ufcNumVoip aparelhos telefônicos VOIPs; \ufcNumNobreaks nobreaks. Os equipamentos de TI estão distribuídos em salas administrativas e nos laboratórios apresentados adiante.

As seções a seguir descrevem os vários ambientes e condições estruturais que contribuem diretamente para a formação do aluno de \nomedocurso.

\section{SALAS DE AULA}

O campus disponibiliza quatorze salas de aula, distribuídas nos quatro blocos didáticos, onde ocorrem prioritariamente as aulas teóricas. Algumas salas comportam confortavelmente até 60 alunos e outras 30, de maneira a atender às necessidades institucionais e do curso de \nomedocurso. As salas de aula passam por limpeza diária e manutenção periódica, além da disponibilidade de recursos tecnológicos adequados às atividades desenvolvidas: a maior parte dispõe de projetores digitais instalados e duas delas possuem sistema de som integrado. As cadeiras podem ser dispostas em diferentes configurações, oportunizando distintas situações de ensino-aprendizagem. Outras atividades didáticas podem acontecer na sala multiuso, sala de videoconferência, além de nos laboratórios, especificados a seguir.

\section{LABORATÓRIOS DIDÁTICOS}

O Campus da UFC em Quixadá oferece estrutura de laboratórios compreendendo principalmente o uso de ferramentas computacionais. No total, há \ufcNumLab laboratórios de computadores para uso didático, todos com projetores digitais. Um laboratório é utilizado como estrutura permanente de apoio discente, no qual os alunos podem realizar estudos, pesquisas e desenvolvimento de trabalhos em horários extraclasse, não havendo aulas alocadas neste espaço. Os demais laboratórios são destinados às aulas e possuem capacidade variando entre 25 e 30 computadores, podendo assim comportar turmas de até 60 alunos, com dois alunos por computador. Todos os laboratórios possuem acesso à rede Wifi. 

Os computadores são disponibilizados com acesso à internet, além de ferramentas de criação e execução apropriadas para os ambientes de aprendizado, contando com softwares open source condizentes com as práticas efetivadas no mercado de trabalho e destinados a procedimentos como edição de imagem e som, edições vetoriais, diagramação, modelagem tridimensional, arquitetura da informação, prototipagem de telas, ambientes de desenvolvimento e programação, além de pacotes de escritório. As estações de trabalho nos laboratórios são configuradas em dual boot, executando Windows e Linux como sistemas operacionais, de modo que alunos e professores possam lidar com ambientes de acordo com as necessidades de aprendizado.

A manutenção dos laboratórios é feita em dois sentidos: estruturalmente, o controle é realizado por meio de levantamento patrimonial da Universidade, identificando o estado de conservação e local onde se encontra cada unidade registrada. Problemas relacionados ao desgaste físico de materiais são solucionados pontualmente, sob demanda. Já em relação à manutenção do software, anualmente os computadores são formatados e os programas necessários são instalados de acordo com a requisição dos docentes, de modo a manter as máquinas rápidas e atualizadas com as versões mais recentes de sistemas operacionais e programas requeridos.

\section{ESTRUTURAS DE HOSPEDAGEM E COMPUTAÇÃO EM NUVEM}

O campus conta com serviços de hospedagem de subdomínio, para o caso de estudos e projetos que necessitem de páginas web como forma de divulgação ou experimentação. Professores e alunos contam com apoio do Núcleo de Tecnologia da Informação e Comunicação (NTIC) do campus para a configuração dos recursos necessários a essas estruturas e acesso simplificado dentro e fora do campus.

O campus de Quixadá está credenciado no programa AWS Academy, parceria com a Amazon Web Services. Os professores cadastrados podem criar laboratórios virtuais nos quais os alunos têm acesso aos principais serviços da nuvem AWS. Professores e alunos contam com apoio do Núcleo de Tecnologia da Informação e Comunicação (NTIC) do campus para a configuração dos recursos necessários a essas estruturas e acesso simplificado dentro e fora do campus.

\section{ACESSO À INTERNET}

Todos os espaços do Campus da UFC em Quixadá contam com acesso à Internet sem fio (Wifi), estando o campus conectado por um link de 200 Mbps ao Cinturão Digital do Ceará (CDC) e, através deste, com a Rede Nacional de Pesquisa e à Internet. Isso significa que há redes disponíveis nas salas de aula, nos laboratórios, na biblioteca e em demais partes do campus, com acesso livre e descomplicado.

A comunidade acadêmica pode se conectar por meio da rede Eduroam, uma rede internacional segura desenvolvida para a comunidade de ensino e pesquisa. Com ela, estudantes, docentes e técnico-administrativos podem se conectar à internet utilizando suas credenciais institucionais, tanto no campus da UFC quanto em qualquer outra instituição participante da iniciativa, no Brasil e no exterior.

\section{ESPAÇOS DE TRABALHO}

Os docentes do campus desfrutam de vários ambientes para desenvolverem seu trabalho. Além dos espaços didáticos (salas de aula e laboratórios), o campus disponibiliza para os docentes uma sala climatizada, na qual eles podem executar ações acadêmicas extraclasse (planejamento, acompanhamento, orientação, pesquisa). Ao todo, o campus possui \ufcNumGabinetes gabinetes que comportam até dois professores cada. Nos gabinetes, o espaço individual de trabalho de cada docente conta com mesa, cadeira para si e para alunos em atendimento, armário com chave, garantindo privacidade para uso dos recursos, para o atendimento a discentes e orientandos, e para a guarda de material e equipamentos pessoais, com segurança.

Para cada um dos cursos de graduação do campus, há uma sala da coordenação ampla e climatizada, onde trabalham o Coordenador e o Vice-coordenador. Nas coordenações, dispõe-se de mesas, cadeiras para os docentes e para os discentes, armários com chave e telefones VOIP. O ambiente da coordenação possibilita o atendimento individual e em pequenos grupos.

Além das salas de trabalho, os coordenadores e docentes podem usufruir de outros ambientes de apoio às ações acadêmicas como sala de seminários, sala de reunião, sala de videoconferência e sala multiuso. Eles também têm acesso a equipamentos como projetores e notebooks.

\section{BIBLIOTECA}
\label{sec:biblioteca}

A Biblioteca Universitária (BU) é um órgão subordinado à Reitoria, a qual compete prover a UFC de um sistema central de informação, de forma a proporcionar serviços biblioteconômicos e documentais eficientes que possibilitem o desenvolvimento das atividades de ensino, pesquisa e extensão. Este é o órgão responsável por gerenciar todas as bibliotecas da UFC, mantendo acervos especializados tombados e catalogados que visam atender a demanda da comunidade acadêmica em geral. Seus serviços são direcionados ao atendimento de alunos, docentes, servidores técnico-administrativos e pesquisadores. Em 2025, o Sistema de Bibliotecas da UFC conta com 29 bibliotecas distribuídas entre seus campi em Fortaleza e no interior do estado.

A Biblioteca Universitária tem em sua estrutura vários setores (diretorias, secretarias, divisões e seções), com atribuições diversas. Em relação ao registro, catalogação e manutenção do acervo, são executadas ações como: a) classificar o material bibliográfico e documental; b) catalogar, de acordo com as normas vigentes, os conteúdos informacionais do Sistema de Bibliotecas da UFC; c) desenvolver ações educativas voltadas para a sensibilização dos usuários em relação à necessidade de preservação do acervo; d) executar as ações de preservação, conservação preventiva, reparadora e restauração dos acervos documentais e bibliográficos.

Além de cuidar do acervo bibliográfico em si, a BU também auxilia no processo de acompanhamento das bibliografias básicas e complementares dos componentes curriculares dos cursos de graduação da UFC. Neste sentido, são atribuições de setores da BU, por exemplo:
\begin{enumerate}
    \item organizar, manter e atualizar na biblioteca os arquivos dos planos de ensino dos cursos de graduação; e b) analisar as bibliografias do plano de ensino conforme os critérios do Instrumento de Avaliação de Cursos de Graduação Presencial e a Distância do MEC. Este é um dos esforços da biblioteca para, juntamente com a Coordenação, o Colegiado e as UC do curso, manter o acervo das bibliografias básica e complementar adequado e atualizado em relação aos componentes curriculares e aos conteúdos previstos no PPC do curso.
    \item atualização do material bibliográfico e documental através de compras. Para isso, ações como as listadas abaixo são executadas: a) receber e preparar a solicitação de compra de material bibliográfico e documental, mediante indicação das sugestões do corpo docente, discente e técnico-administrativo em educação; b) elaborar o processo de licitação do material bibliográfico; c) acompanhar o processo de compra e recebimento de material bibliográfico.
    \item atividades de gestão que auxiliem os processos de tomada de decisão como: a) coordenar e controlar os relatórios patrimoniais do material bibliográfico do Sistema de Bibliotecas da UFC; b) coordenar o inventário anual do material bibliográfico e documental;
    \item coletar e analisar os dados gerados a partir dos relatórios automatizados, mantendo estatísticas que subsidiem estudos na área.    
\end{enumerate}

Além das atribuições próprias da BU, a administração de cada Biblioteca Setorial, como é o caso da Biblioteca do Campus de Quixadá (BCQ), tem em sua estrutura a Diretoria Setorial, a Seção de Representação Descritiva e Temática da Informação, a Seção de Atendimento ao Usuário, a Seção de Preservação, Conservação e Restauração do Acervo, a Seção de Coleções Especiais e a Seção de Atendimento às Pessoas com Deficiência. Algumas funções básicas de gestão do acervo desempenhadas por algumas das Seções mencionadas são:
\begin{itemize}
    \item Catalogar e classificar, de acordo com as normas vigentes, todo o conteúdo informacional pertinente a área de atuação da biblioteca;
    \item Realizar estudo da bibliografia adotada nos cursos atendidos pela biblioteca setorial e propor a aquisição do material bibliográfico que preencha os requisitos necessários ao pleno desenvolvimento das disciplinas ofertadas;
    \item Acompanhar os relatórios mensais de aquisição de material bibliográfico;
    \item Manter atualizado os repositórios locais, nacionais e internacionais;
    \item Supervisionar os serviços de atendimento ao usuário no que diz respeito à circulação de material bibliográfico;
    \item Orientar os usuários na busca de informações e no uso dos acervos existentes nas Bibliotecas do sistema e fora delas, auxiliando-os em suas necessidades de estudo e pesquisa;
    \item Localizar e fornecer documentos e informações solicitadas;
    \item Desenvolver ações educativas voltadas para a sensibilização dos usuários em relação à necessidade de preservação do acervo;
    \item Zelar pela conservação e funcionamento de equipamentos, máquinas e aparelhos da Seção.
\end{itemize}

%Em relação ao processo de aquisição de material bibliográfico, o campus dispõe de um sistema específico (Gestão de Aquisição de Livros – GAL), para auxiliar no processo de compras. Neste sistema, todas as disciplinas e toda a bibliografia de todos os cursos são cadastrados. Registra-se também os títulos que se deseja para a compra. Baseado nesses dados, o sistema compara a quantidade de exemplares existente no acervo com a quantidade necessária, sugerindo assim a quantidade de exemplares a ser adquirida. O GAL é mais um sistema inovador, desenvolvido no campus, para auxiliar os processos administrativos e acadêmicos dos cursos que aqui estão.

Além destas funções administrativas, a Biblioteca Universitária realiza, periodicamente ou sob demanda, cursos e treinamentos de capacitação no uso otimizado dos recursos informacionais para alunos, professores e técnico-administrativos\footnote{\url{http://www.biblioteca.ufc.br/servicos-e-produtos/cursos-e-treinamentos/}}. Alguns treinamentos específicos fazem parte do calendário oficial da BU: Normalização de Trabalhos Acadêmicos, Referências e Citações, Treinamento em Bases de Dados. Além destes, a Biblioteca promove o projeto “Descobrindo a Biblioteca”, que acontece sempre no início de cada semestre letivo e visa apresentar aos novos alunos os serviços ofertados pelo Sistema de Bibliotecas da UFC. São abordados conteúdos como: regulamento, guia de serviços, acervo, catálogo on-line, meios de acesso do usuário, livros eletrônicos, eventos da Biblioteca Universitária, dentre outros. A iniciativa faz parte da programação de recepção dos recém-ingressos, que tradicionalmente inclui cursos e palestras em suas unidades acadêmicas.

Especificamente, em Quixadá, além de todas as iniciativas listadas acima, desde 2013, a Biblioteca do campus promove a “Maratona do Conhecimento do Sertão Central”, que atualmente é ocorre durante os Encontros Universitários da UFC. O objetivo do evento é difundir ferramentas e fontes de informação acadêmica, promovendo assim o desenvolvimento de habilidades relacionadas à busca, acesso e utilização de informações para construção do conhecimento.

O curso de \nomedocurso dispõe de acervo atualizado, contemplando títulos adotados como bibliografia básica e complementar devidamente indicados nas ementas das disciplinas. Além dos livros físicos, a biblioteca do campus oferece acesso a plataformas online nas quais se encontram livros eletrônicos e artigos de periódicos voltados às principais temáticas abordadas em sala de aula.

Em \ufcBCQAnoReferencia, a BCQ conta com um acervo de \ufcBCQExemplares exemplares físicos, correspondentes a \ufcBCQTitulos títulos, principalmente nas seguintes áreas do conhecimento: Ciências Exatas e da Terra, Engenharias, Ciências Sociais Aplicadas, Ciências Humanas, Linguística, Letras e Artes. Uma descrição mais detalhada das estruturas de conteúdo eletrônico e físico é trazida nas seções a seguir.

Em relação à infraestrutura física, a BCQ conta com salão de estudo climatizado e computadores disponíveis para consulta ao catálogo eletrônico. Além disso, a biblioteca é responsável por gerenciar a ocupação das salas de estudo em grupo disponíveis aos discentes. Vê-se, portanto, que os eventos e treinamentos promovidos, as ações de conscientização e a infraestrutura de acesso aos títulos bibliográficos e de acomodações físicas são exemplos de soluções de apoio à leitura, estudo e aprendizagem disponibilizados pela biblioteca aos discentes e demais membros da comunidade acadêmica.

\subsection{ACERVOS DIGITAIS}

Todo o acervo físico do Sistema de Bibliotecas da UFC está catalogado e disponível digitalmente para a comunidade acadêmica. Através do sistema Pergamum\footnote{\url{https://pergamum.ufc.br/pergamum/biblioteca/index.php}}, os usuários acessam o catálogo online do Sistema de Bibliotecas da UFC, ou seja, os registros de todo o acervo das bibliotecas, inclusive documentos eletrônicos em texto completo, tais como: livros, teses e dissertações, monografias, periódicos, artigos, obras raras e CDs/DVDs, dentre outros.

Além do acervo físico, os usuários das bibliotecas da UFC têm acesso a um enorme conjunto de material digital, conforme descrito a seguir.

\subsection{PORTAL DE ACESSO A CONTEÚDO CIENTÍFICO DIGITAL}

A UFC oferece materiais de estudo por meio de acesso à plataforma Minha Biblioteca\footnote{\url{https://dliportal.zbra.com.br/Login.aspx?key=UFC}} e à coleção da editora Springer\footnote{\url{http://ufc.dotlib.com.br/}}. O acesso interno é feito de qualquer dos campi da UFC. Para acesso remoto, alunos e professores podem realizar configurações de Proxy, por meio de CPF e senha do sistema acadêmico SIGAA/SIGPRH. Assim, todos os alunos matriculados possuem acesso eletrônico a livros e artigos disponibilizados nessas plataformas que contam com mais de 14 mil títulos em texto completo, que podem ser pesquisados pelos usuários no ambiente da UFC ou de qualquer ponto de Internet, sem limites de utilização, visualização ou restrições de usuários simultâneos.

\subsection{PERIÓDICOS}

A biblioteca conta com acesso online às 27 revistas de publicação própria da UFC, por meio do Portal de Periódicos da UFC\footnote{\url{http://periodicos.ufc.br/}}. Além disso, professores e alunos contam com acesso ao Portal Periódicos da CAPES\footnote{\url{https://www.periodicos.capes.gov.br/}}, que disponibiliza documentos periódicos (internacionais e nacionais), livros eletrônicos, bases de dados contendo artigos, referências e resumos de trabalhos acadêmicos e científicos, normas técnicas, patentes, teses e dissertações e outros tipos de materiais, em todas as áreas do conhecimento, totalizando mais de 37 mil títulos. O acesso interno é liberado automaticamente por meio de faixa de IPs da UFC, já o acesso remoto é liberado por meio da Comunidade Acadêmica Federada (CAFe).

\subsection{REPOSITÓRIO INSTITUCIONAL}

Finalmente, há o Repositório Institucional (RI)\footnote{\url{https://repositorio.ufc.br/}} da UFC, de acesso aberto via Internet, que tem como propósito reunir, armazenar, organizar, recuperar, preservar e disseminar a produção científica e intelectual da comunidade universitária (docentes, pesquisadores, técnicos e alunos de pós-graduação stricto sensu) pertencente à UFC. Em \ufcBCQAnoReferencia, este repositório conta com mais de \ufcBCQDocsRepositorio documentos, dentre estes, artigos de periódicos, dissertações, teses, capítulos de livros, artigos publicados em eventos, além das monografias e trabalhos de conclusão dos cursos de graduação da UFC. Dessa forma, o corpo discente tem acesso imediato a produções de alunos situados em outros campi da UFC.

\subsection{OUTROS RECURSOS DIGITAIS}

Além das fontes de bibliografia digital destacadas acima, a Biblioteca Universitária disponibiliza ainda os seguintes recursos:
\begin{itemize}
    \item Portal de Livros Eletrônicos UFC\footnote{\url{http://livros.ufc.br/ojs/}}: acesso aos Livros da Coleção de Estudos da Pós-Graduação da Universidade. Acesso livre.
    \item Coleção de Normas ABNT\footnote{\url{https://www.gedweb.com.br/UFC/}}: acesso a mais de 9 mil Normas Técnicas Brasileiras e Normas Técnicas do Mercosul em texto completo. Acesso remoto via proxy.
\end{itemize}

\subsection{BIBLIOTECA ACESSÍVEL}

A fim de proporcionar ambientes de estudo adequados e um maior acesso à informação aos usuários com deficiência, o Sistema de Bibliotecas da UFC oferece um atendimento pautado na prestação de serviços especializados, na aquisição de equipamentos e tecnologias assistivas desenvolvidas especialmente para esses usuários e na inserção da acessibilidade arquitetônica em suas edificações.

Neste contexto, os alunos com necessidades especiais podem contar com os seguintes especializados, oferecidos em parceria com a Secretaria de Acessibilidade da UFC:
\begin{itemize}
    \item Digitalização e conversão de materiais bibliográficos em formatos acessíveis: neste serviço a bibliografia solicitada pelo professor passa pelo processo de digitalização e/ou edição e é convertida em arquivo digital acessível para posteriormente ser disponibilizada no catálogo da Biblioteca Universitária com acesso restrito aos usuários com deficiência visual.
    \item Orientação à pesquisa bibliográfica para usuários com deficiência visual: serviço realizado mediante treinamentos de uso das bases de dados on-line do Portal da CAPES, Biblioteca Digital de Teses e Dissertações (BDTD) e livros eletrônicos com o auxílio de softwares leitores de tela.
    \item Levantamento bibliográfico para usuários com deficiência visual: o serviço consiste na pesquisa bibliográfica demandada pelos usuários com deficiência visual, transformada posteriormente em arquivo digital acessível para seu uso exclusivo.
\end{itemize}

Em relação a recursos e tecnologias assistivas, a Biblioteca Universitária disponibiliza: a) nos terminais de consulta online, programas leitores de telas; b) recurso de transcrição de textos em Braille, com impressoras Braille à disposição dos alunos (em Fortaleza); c) recurso de acessibilidade linguística, através do projeto “Descobrindo a Biblioteca em Libras”, com a apresentação dos serviços oferecidos pelas bibliotecas para a comunidade acadêmica por uma intérprete da Língua Brasileira de Sinais (Libras).

Especificamente na BCQ, os alunos com deficiência visual têm acesso a um equipamento para leitura autônoma, instalado na biblioteca. Em relação à acessibilidade arquitetônica, a biblioteca tem portas de acesso largas, balcões e mesas de estudo rebaixados, espaço de circulação amplo entre prateleiras, placas com inscrições em Braille e piso tátil.

\section{ACESSIBILIDADE FÍSICA}

O campus possui em sua infraestrutura a facilitação da acessibilidade a pessoas com dificuldades de locomoção ou de visão, contando com plataformas elevatórias, portas largas de acesso às salas e laboratórios, banheiros com cabines específicas para deficientes e plaquetas com inscrições em braille, além de rampas de acesso ou mesmo ausência de degraus desde o estacionamento até todos os ambientes térreos.

\section{RECURSOS HUMANOS}

\def\ufc{\xspace}
\def\ufc{\xspace}

O corpo docente do Campus da UFC em Quixadá, em \ufcAnoReferencia, é formado por \ufcNumDocentes professores, sendo \ufcNumDocSubstitutos substituto, \ufcNumDocQuarentaHoras efetivo em regime de 40h e \ufcNumDocDE efetivos em regime de dedicação exclusiva. Dos efetivos, \ufcDocDoutores são doutores e \ufcDocMestres são mestres. Dos \ufcDocMestres docentes mestres, \ufcDocDoutorandos são doutorandos.\footnote{A lista completa de docentes pode ser encontrada em \url{https://www.quixada.ufc.br/docente/}.}

Em relação ao estímulo à titulação dos docentes, foi aprovada a Resolução nº 01/Conselho do Campus de Quixadá, de 15 de maio de 2014, para disciplinar os afastamentos de docentes para mestrado, doutorado e pós-doutorado. Salienta-se que foi desenvolvido no campus o Sistema de Afastamento (SIAF), criado para gerenciar o afastamento dos docentes. Neste sistema, os docentes fazem uma solicitação de reserva de afastamento e podem acompanhar, em tempo real, sua posição em um ranking construído a partir das regras de afastamento definidas na Resolução. Em \ufcAnoReferencia, há \ufcDocAfasDoutorado professor afastado para o doutorado.

Quanto ao quadro de servidores técnico-administrativos, o campus dispõe, em \ufcAnoReferencia, de \ufcNumSTA servidores. Destes, \ufcNumSTAMedio são cargos de nível médio e \ufcNumSTASuperior de nível superior. Em relação à escolaridade, \ufcSTAEscMedio têm o nível médio, \ufcSTAEscGrad possuem graduação completa, \ufcSTAEscEspec são especialistas e \ufcSTAEscMestre são mestres, demonstrando a qualidade do corpo técnico-administrativo do campus.\footnote{A lista completa dos servidores técnicos-administrativos pode ser encontrada em \url{https://www.quixada.ufc.br/sta/}}

\section{COMITÊ DE ÉTICA EM PESQUISA}

O Comitê de Ética em Pesquisa da UFC (CEP/UFC/PROPESQ), vinculado à Pró-reitoria de Pesquisa e Pós-Graduação, foi instituído em 20 de outubro de 2005. É credenciado junto à Comissão Nacional de Ética em Pesquisa (CONEP) do Ministério da Saúde e constitui um colegiado interdisciplinar, independente e normativo, com ``munus público'', sem fim lucrativo, de caráter consultivo, deliberativo e educativo, criado para defender os interesses dos participantes da pesquisa em sua integridade e dignidade, que obedece aos princípios da Bioética: autonomia, não maleficência, beneficência e justiça, e visa contribuir no desenvolvimento da pesquisa dentro de padrões éticos (Normas e Diretrizes Regulamentadoras da Pesquisa Envolvendo Seres Humanos – Resolução CNS 466/12, II.4). O CEP é responsável pela avaliação e acompanhamento dos aspectos éticos de todas as pesquisas envolvendo seres humanos. No campus, a Coordenadoria de Pesquisa, além de incentivar, apoiar e acompanhar o desenvolvimento das atividades de pesquisa, apoia os contatos entre a comunidade acadêmica e o CEP/UFC/PROPESQ.
	
	%Elementos pós-textuais	
	% \bibliography{3-pos-textuais/referencias}
%	\imprimirglossario	
	\imprimirapendices
		% Adicione aqui os apendices do seu trabalho
		\apendice{Ementário e bibliografias}
\label{ap:ementario_bibliografia}

\newcommand{\DetalhesDisciplina}[4]{
\begin{table}[h!]
    \centering
    \small
    \renewcommand{\arraystretch}{1.5}
    \begin{tabular}{|>{\columncolor{gray!15}}p{2.8cm}|p{11.2cm}|}
        \hline        
        \textbf{Componente} & \textbf{#1} \\
        \hline        
        \textbf{Ementa} & #2 \\
        \hline        
        \textbf{Bibliografia \newline Básica} & #3 \\
        \hline        
        \textbf{Bibliografia Complementar} & #4 \\
        \hline
    \end{tabular}
\end{table}
\clearpage
}

\DetalhesDisciplina{
    <03> - Agentes Inteligentes
}{% EMENTA
    Conceitos fundamentais de agentes inteligentes. Arquitetura, Racionalidade. Natureza dos ambientes: totalmente observável, parcialmente observável, determinístico, estocástico, episódico, sequencial, estático, dinâmico, discreto, contínuo, único agente, multiagentes. Programas de agentes: agentes reflexivos simples, agentes baseados em modelos, agentes baseados em objetivos, agentes baseados em utilidade e agentes que aprendem. Introdução a sistemas multiagentes.
}{% BIBLIOGRAFIA BÁSICA
    RUSSELL, Stuart J.; NORVIG, Peter. Inteligência Artificial: Uma Abordagem Moderna. 4. ed. Rio de Janeiro: GEN LTC, 2022. E-book. p.Capa. ISBN 9788595159495. 
\newline \newline 
WOOLDRIDGE, M. Introduction to Multiagent Systems. Wiley, 2009. ISBN: 9780470519462
\newline \newline 
COPPIN, B Inteligência artificial. LTC, 2010. ISBN: 9788521617297
}{% BIBLIOGRAFIA COMPLEMENTAR
    SILVA, Fabrício M.; LENZ, Maikon L.; FREITAS, Pedro H C.; et al. Inteligência artificial. Porto Alegre: SAGAH, 2018. E-book. p.1. ISBN 9788595029392. Disponível em: https://app.minhabiblioteca.com.br/reader/books/9788595029392/. Acesso em: 22 mai. 2025.
\newline \newline 
BRAKTO, I.Prolog Programming for Artificial Intelligence. 4 ed. Addison Wesley, 2011. ISBN 0321417461
\newline \newline 
SHOHAM, Y. Multiagent systems: algorithms, game theoretic. Cambridge University, 2009. ISBN: 9780521899437
\newline \newline 
KAUFMAN, Dora. Desmistificando a inteligência artificial. São Paulo: Autêntica Editora, 2022. E-book. p.1. ISBN 9786559281596. Disponível em: https://app.minhabiblioteca.com.br/reader/books/9786559281596/. Acesso em: 30 mai. 2025.
\newline \newline 
LIMA, Isaías. Inteligência Artificial. Rio de Janeiro: GEN LTC, 2014. E-book. p.Capa. ISBN 9788595152724. Disponível em: https://app.minhabiblioteca.com.br/reader/books/9788595152724/. Acesso em: 22 mai. 2025.
}


\DetalhesDisciplina{
    QXD0116 - Álgebra Linear
}{% EMENTA
    Matrizes. Sistemas de equações lineares. Vetores. Produto interno. Dependência e independência linear. Espaços vetoriais. Bases de espaços vetoriais. Transformações lineares. Autovalores e autovetores. Diagonalização de operadores.
}{% BIBLIOGRAFIA BÁSICA
    VUJICIC, Milan; SANDERSON, Jeffrey SPRINGERLINK (ONLINE SERVICE). Linear Algebra Thoroughly
Explained. Springer e-books Berlin, Heidelberg: Springer-Verlag Berlin Heidelberg, 2008. ISBN
9783540746393
\newline \newline 
CABRAL, Marco Aurélio Palumbo; GOLDFELD, Paulo. Curso de álgebra linear: fundamentos e aplicações. 3.
ed. Rio de Janeiro: UFRJ, 2012. 257 p. Disponível em: 94 [recurso eletrônico]
\newline \newline 
LIPSCHUTZ, Seymour; LIPSON, Marc; Álgebra Linear; Coleção Schaum; Editora Bookman; 2011.E-book.
Disponível em: https://integrada.minhabiblioteca.com.br/\#/books/9788540700413. Acesso em: 12 de set.
2022
}{% BIBLIOGRAFIA COMPLEMENTAR
    BOLDRINI, José Luiz; Algebra Linear; 3. Edição; Editora Harbra
\newline \newline 
STEINBRUCH, Alfredo; Introduçao à Algebra Linear; Makron Books Editora; 1990
\newline \newline 
IEZZI, GELSON et al; Fundamentos de matemática elementar v.4: Sequências, matrizes, determinantes e
sistemas, 8 ed., 2004, Saraiva
\newline \newline 
ROMAN, Steven SPRINGERLINK (ONLINE SERVICE). Advanced Linear Algebra. Springer e-books Third
Edition. New York, NY: Springer Science+Business Media, LLC, 2008. (Graduate Texts in Mathematics, 135)
ISBN 9780387728315
\newline \newline 
ESPINOSA, Isabel Cristina de Oliveira Navarro; BISCOLLA, Laura Maria da Cunha Canto Oliva; BARBIERI
FILHO, Plinio. Álgebra linear para computação. Rio de Janeiro, RJ: LTC; 2007. x, 286p. (Fundamentos de
informática). ISBN 9788521615521 (broch.)
\newline \newline 
ANTON, Howard; RORRES, Chris. Álgebra Linear com Aplicações. São Paulo: Bookman, 2012. E-book.
Disponível em: https://integrada.minhabiblioteca.com.br/\#/books/9788540701700. Acesso em: 3 de out.
2022
}


\DetalhesDisciplina{
    <21> - Análise Exploratória de Dados
}{% EMENTA
    Tipos e organização de dados. Coleta e tratamento de Dados. Representação gráfica e tabular de dados qualitativos. Medidas descritivas de dados quantitativos. Representação gráfica de dados quantitativos. Associação entre variáveis qualitativas.Associação entre uma variável quantitativa e qualitativa. Associações entre variáveis quantitativas.Ferramentas e linguagens para análise exploratória de dados.

}{% BIBLIOGRAFIA BÁSICA
    LARSON, Ron; FARBER, Elizabeth. Estatística aplicada. 4. ed. São Paulo: Pearson Prentice Hall, 2010.
xiv,637 p. ISBN 978-85-7605-372-9
\newline \newline 
TRIOLA, Mario F.. Introdução à estatística. 10. ed. Rio de Janeiro: LTC, c2008. xxvi, 696 p. ISBN 978-85-
216-1586-6.
\newline \newline 
TRIOLA, Mario F.; FLORES, Vera Regina Lima de Farias e. Introdução à estatística: atualização da tecnologia. 11. ed. Rio de Janeiro: LTC, c2013. xxviii ; 707 p. ISBN 97788521622060
}{% BIBLIOGRAFIA COMPLEMENTAR
    BEIGHLEY, Lynn. Use a cabeça!: SQL. Rio de Janeiro: Alta Books, 2008. xxxiv, 454 p. (Use a cabeça).
ISBN 9788576082101.
\newline \newline 
BARRY, Paul. Use a cabeça: Python. Rio de Janeiro: Alta Books, 2012. xxxiv, 457 p. (Use a cabeça). ISBN 978-85-7608-743-4.
\newline \newline 
IEZZI, Gelson; HAZZAN, Samuel; DEGENSZAJN, David Mauro. Fundamentos de matemática elementar 11 : matemática comercial, matemática financeira, estatística descritiva . São Paulo, SP: Atual, 2004. 232 p.
ISBN 8535704620 (broch.)
\newline \newline 
FERREIRA, Rafael G C.; MIRANDA, Leandro B. A de; PINTO, Rafael A.; et al. Preparação e Análise Exploratória de Dados. Porto Alegre: SAGAH, 2021. E-book. p.Capa. ISBN 9786556902890. Disponível em: https://app.minhabiblioteca.com.br/reader/books/9786556902890/ . Acesso em: 03 jun. 2025.
\newline \newline 
CARVALHO, André C. P. L. F de; MENEZES, Angelo G.; BONIDIA, Robson P. Ciência de Dados - Fundamentos e Aplicações. Rio de Janeiro: LTC, 2024. E-book. p.Capa. ISBN 9788521638766. Disponível em: https://app.minhabiblioteca.com.br/reader/books/9788521638766/ . Acesso em: 30 mai. 2025.
\newline \newline 
BARBETTA, Pedro A.; REIS, Marcelo M.; BORNIA, Antonio C. Estatística para Cursos de Engenharia, Computação e Ciência de Dados. 4. ed. Rio de Janeiro: LTC, 2024. E-book. p.29. ISBN 9788521638827. Disponível em: https://app.minhabiblioteca.com.br/reader/books/9788521638827/ . Acesso em: 03 jun. 2025.
}


\DetalhesDisciplina{
    <12> - Aprendizado por Reforço
}{% EMENTA
    Fundamentos do aprendizado por reforço. Conceitos básicos: agente, ambiente, estado, ação, recompensa e política. Processos de Decisão de Markov (MDPs) e funções de valor e política. Programação dinâmica através da avaliação de política e iteração de política. Métodos de Monte Carlo. Temporal Difference Learning. Algoritmos Q-Learning e SARSA. Aproximação de funções de valor. Aprendizado por reforço profundo (Deep RL) com Deep Q-Networks (DQN) e variações. Métodos de Policy Gradient com estimação de valor. Actor-Critic Methods (A2C, A3C e PPO). Exploração versus explotação. Aprendizado por imitação. Aprendizado por reforço multi-agente. Aplicações em jogos, robótica, sistemas autônomos e otimização de estratégias.
}{% BIBLIOGRAFIA BÁSICA
    RUSSELL, Stuart J.; NORVIG, Peter. Inteligência Artificial: Uma Abordagem Moderna. 4. ed. Rio de Janeiro: GEN LTC, 2022. E-book. p.Capa. ISBN 9788595159495. Disponível em: https://app.minhabiblioteca.com.br/reader/books/9788595159495/. 
\newline \newline 
SEJNOWSKI, Terrence. A Revoluçao do Aprendizado Profundo. Rio de Janeiro: Editora Alta Books, 2020. E-book. p.176. ISBN 9788550814353. Disponível em: https://app.minhabiblioteca.com.br/reader/books/9788550814353/.
\newline \newline 
WANG, Xu; WANG, Sen; LIANG, Xingxing; ZHAO, Dawei; HUANG, Jincai; XU, Xin; DAI, Bin; MIAO, Qiguang. Deep reinforcement learning: A survey. IEEE Transactions on Neural Networks and Learning Systems, v. 35, n. 4, p. 5064–5078, 2024. DOI: 10.1109/TNNLS.2022.3207346. Disponível em: https://ieeexplore.ieee.org/document/9904958.

}{% BIBLIOGRAFIA COMPLEMENTAR
    HARIOM, Tatsat,; SAHIL, Puri,; BRAD, Lookabaugh,. Blueprints de aprendizado de máquina e ciência de dados para finanças: desenvolvendo desde estratégias de trades até robôs Advisors com Python. Rio de Janeiro: Editora Alta Books, 2024. E-book. p.i. ISBN 9788550821726. Disponível em: https://app.minhabiblioteca.com.br/reader/books/9788550821726/.
\newline \newline 
NAEEM, Muddasar; RIZVI, Syed Tahir Hussain; CORONATO, Antonio. A gentle introduction to reinforcement learning and its application in different fields. IEEE Access, New York, v. 8, p. 209 320–209 344, 17 nov. 2020. DOI: 10.1109/ACCESS.2020.3038605. Disponível em: https://ieeexplore.ieee.org/document/9261348. 
\newline \newline 
HENDERSON, P.; ISLAM, R.; BACHMAN, P.; PINEAU, J.; PRECUP, D.; MEGER, D. Deep reinforcement learning that matters. Proceedings of the AAAI Conference on Artificial Intelligence, v. 32, n. 1, 29 abr. 2018. DOI: 10.1609/aaai.v32i1.11694. Disponível em: https://ojs.aaai.org/index.php/AAAI/article/view/11694.
}


\DetalhesDisciplina{
    <09> - Aprendizado Profundo
}{% EMENTA
    Redes neurais artificiais. Aprendizado profundo: perceptron, redes multicamadas (MLPs) e arquiteturas feedforward. Backpropagation e técnicas de otimização. Funções de ativação. Técnicas de regularização (dropout, batch normalization e data augmentation). Redes neurais convolucionais (CNNs), redes neurais recorrentes (RNNs, LSTM, GRU). Arquiteturas avançadas como ResNet, autoencoders, redes generativas (GANs, VAEs) e transformadores com mecanismos de atenção. Transfer learning, fine-tuning e uso de redes pré-treinadas. Aplicações em visão computacional, processamento de linguagem natural (NLP) e outras áreas do aprendizado profundo.
}{% BIBLIOGRAFIA BÁSICA
    FACELI, Katti; LORENA, Ana C.; GAMA, João; AL, et. Inteligência Artificial - Uma Abordagem de Aprendizado de Máquina. 2. ed. Rio de Janeiro: LTC, 2021. E-book. p.Capa. ISBN 9788521637509. Disponível em: https://app.minhabiblioteca.com.br/reader/books/9788521637509/.
\newline \newline 
RUSSELL, Stuart J.; NORVIG, Peter. Inteligência Artificial: Uma Abordagem Moderna. 4. ed. Rio de Janeiro: GEN LTC, 2022. E-book. p.Capa. ISBN 9788595159495. Disponível em: https://app.minhabiblioteca.com.br/reader/books/9788595159495/.
\newline \newline 
SEJNOWSKI, Terrence. A Revoluçao do Aprendizado Profundo. Rio de Janeiro: Editora Alta Books, 2020. E-book. p.157. ISBN 9788550814353. Disponível em: https://app.minhabiblioteca.com.br/reader/books/9788550814353/. 
}{% BIBLIOGRAFIA COMPLEMENTAR
    HARIOM, Tatsat,; SAHIL, Puri,; BRAD, Lookabaugh,. Blueprints de aprendizado de máquina e ciência de dados para finanças: desenvolvendo desde estratégias de trades até robôs Advisors com Python. Rio de Janeiro: Editora Alta Books, 2024. E-book. p.i. ISBN 9788550821726. Disponível em: https://app.minhabiblioteca.com.br/reader/books/9788550821726/.
\newline \newline 
LECUN, Yann; BENGIO, Yoshua; HINTON, Geoffrey. Deep learning. Nature, Londres, v. 521, n. 7553, p. 436–444, 1 maio 2015. DOI: 10.1038/nature14539. Disponível em: https://doi.org/10.1038/nature14539. 
\newline \newline 
JANIESCH, Christian; ZSCHECH, Patrick; HEINRICH, Kai. Machine learning and deep learning. Electronic Markets, v. 31, n. 3, p. 685–695, 1 set. 2021. DOI: 10.1007/s12525-021-00475-2. Disponível em: https://doi.org/10.1007/s12525-021-00475-2. 
}


\DetalhesDisciplina{
    <06> - Aprendizado Supervisionado
}{% EMENTA
    Princípios teóricos e práticos do aprendizado supervisionado. Classificação e Regressão. Modelos lineares e não lineares, k-NN, regressão linear, regressão polinomial, regressão logística, regularização (Ridge, Lasso, Elastic Net), SVM, Naive Bayes, árvores de decisão, métodos de ensemble (Random Forest, Gradient Boosting), e redes neurais. Pipeline de modelagem: pré-processamento, particionamento de dados, Métricas de avaliação e validação cruzada, e ajuste de hiperparâmetros. Tratamento de dados desbalanceados e overfitting. Seleção e engenharia de features. Interpretabilidade de modelos. Aplicações práticas de aprendizado supervisionado.
}{% BIBLIOGRAFIA BÁSICA
    FACELI, Katti; LORENA, Ana C.; GAMA, João; AL, et. Inteligência Artificial - Uma Abordagem de Aprendizado de Máquina. 2. ed. Rio de Janeiro: LTC, 2021. E-book. p.Capa. ISBN 9788521637509. Disponível em: https://app.minhabiblioteca.com.br/reader/books/9788521637509/.
\newline \newline 
COPPIN, Ben. Inteligência Artificial. São Paulo: LTC, 2010. E-book. Disponível em: https://integrada.minhabiblioteca.com.br/\#/books/978-85-216-2936-8.
\newline \newline 
NORVIG, Peter. Inteligência Artificial. São Paulo: GEN LTC, 2013. E-book. Disponível em: https://integrada.minhabiblioteca.com.br/\#/books/9788595156104. 
}{% BIBLIOGRAFIA COMPLEMENTAR
    HARIOM, Tatsat,; SAHIL, Puri,; BRAD, Lookabaugh,. Blueprints de aprendizado de máquina e ciência de dados para finanças: desenvolvendo desde estratégias de trades até robôs Advisors com Python. Rio de Janeiro: Editora Alta Books, 2024. E-book. p.i. ISBN 9788550821726. Disponível em: https://app.minhabiblioteca.com.br/reader/books/9788550821726/.
\newline \newline 
HAYKIN, Simon. Redes neurais princípios e prática.. São Paulo: Bookman, 2001. E-book. Disponível em: https://integrada.minhabiblioteca.com.br/\#/books/9788577800865.
\newline \newline 
LIMA, Isaías. Inteligência Artificial. Rio de Janeiro: GEN LTC, 2014. E-book. p.Capa. ISBN 9788595152724. Disponível em: https://app.minhabiblioteca.com.br/reader/books/9788595152724/. 
}


\DetalhesDisciplina{
    QXD0005 - Arquitetura de Computadores
}{% EMENTA
    Sistemas numéricos. Aritmética binária: ponto fixo e ponto flutuante.
Organização de computadores: memórias, unidade central de processamento,
unidades de entrada e unidades de saída. Linguagens de montagem. Modos
de endereçamento, conjunto de instruções. Mecanismos de interrupção e de
exceção. Barramento, comunicações, interfaces e periféricos. Organização de
memória. Memória auxiliar. Arquiteturas RISC e CISC. Pipeline. Paralelismo de
baixa granularidade. Multiprocessadores. Processadoressuperescalares 
Multicomputadores. Arquiteturas e superpipeline paralelas e não convencionais.
}{% BIBLIOGRAFIA BÁSICA
    STALLINGS,
William; VIEIRA, Daniel. Arquitetura e Organização de
Computadores. São Paulo: Pearson Prentice Hall, 2010. xiv, 624 p. ISBN
9788576055648
\newline \newline 
TANENBAUM, A. Organização Estruturada de Computadores. Pearson
Prentice Hall. São Paulo. 2007. 5ª Ed ISBN 8576050676
\newline \newline 
ENGLANDER, Irv. A Arquitetura de Hardware Computacional, Software de
Sistema e Computação e Comunicação em Rede. 4 Edição. 2011. Editora
LTC. ISBN: 9788521617914
\newline \newline 
MURDOCCA, Miles J.; HEURING, Vincent P., Introdução à arquitetura de
computadores. Rio de Janeiro: ​Campus,​ 2000. 512 p. ISBN 8535206841
}{% BIBLIOGRAFIA COMPLEMENTAR
    DELGADO, José; RIBEIRO, Carlos. Arquitetura de computadores. 4. ed. rev.
atual. FCA, 2010. ISBN 9789727226665
\newline \newline 
HENNESSY, John L.; PATTERSON, David A. Arquitetura de computadores:
uma abordagem quantitativa. 4. ed. Rio de Janeiro: Elsevier, 2008.
\newline \newline 
WEBER, Raul Fernando. Fundamentos de arquitetura de computadores. 3. ed.
Porto Alegre, RS: Sagra Luzzatto, 2008.
\newline \newline 
FEDELI, Ricardo Daniel; POLLONI, Enrico Giulio Franco; PERES, Fernando
Eduardo. Introdução à ciência da computação. 2. ed. atual. Cengage Learning,
2010. 250 p. ISBN 139788522108459.
\newline \newline 
MONTEIRO, Mario. A. Introdução à Organização de Computadores. 5. ed. Rio
de Janeiro, RJ: LTC, 2007. ISBN 9788521615439.
\newline \newline 
MACHADO, Francis B; MAIA, Luiz Paulo. Arquitetura de sistemas operacionais. 4. ed. Rio de Janeiro: LTC, 2007. 308 p.
}


\DetalhesDisciplina{
    <07> - Busca e Otimização
}{% EMENTA
    Fundamentos da resolução de problemas por meio de busca. Formulação de problemas: espaço de estados, funções de custo, ações aplicáveis e testes de objetivo. Algoritmos de busca não informada: Busca em largura, Busca em profundidade, Busca de custo uniforme e Busca em profundidade iterativa. Princípios de Busca Heurística. Algoritmos de busca informada (heurística): Busca gulosa, A* (A estrela). Busca Local: Subida de encosta (Hill Climbing), Têmpera Simulada (Simulated Annealing), Busca em Feixe Local, Algoritmos genéticos. Problemas de Satisfação de restrições (CSPs).
}{% BIBLIOGRAFIA BÁSICA
    RUSSELL, Stuart J.; NORVIG, Peter. Inteligência Artificial: Uma Abordagem Moderna. 4. ed. Rio de Janeiro: GEN LTC, 2022. E-book. p.Capa. ISBN 9788595159495. 
\newline \newline 
CORMEN, Thomas H.; LEISERSON, Charles E.; Ronald L. Rivest; et al. Algoritmos. 4. ed. Rio de Janeiro: GEN LTC, 2024. E-book. p.Capa. ISBN 9788595159914. 
\newline \newline 
COPPIN, B Inteligência artificial. LTC, 2010. ISBN: 9788521617297
}{% BIBLIOGRAFIA COMPLEMENTAR
    KLEINBERG, Jon; TARDOS, Éva. Algorithm design. Boston, Massachusetts: Pearson/Addison Wesley, c2006. 838 p. ISBN 0321295358.
\newline \newline 
TOSCANI, Laira V.; VELOSO, Paulo A. S. Complexidade de algoritmos: análise, projeto e métodos. 3. ed. Porto Alegre: Sagra Luzzato, 2012. 262 p. (Serie Livros Didáticos Informática UFRGS ; 13). ISBN 9788540701380 (broch.).
\newline \newline 
ZIVIANI, Nivio; BOTELHO, Fabiano Cupertino. Projeto de algoritmos: com implementações em java e C++. São Paulo, SP: Thomson Learning, 2007. vii, 620 p. ISBN 8522105251.
\newline \newline 
DASGUPTA, Sanjoy; PAPADIMITRIOU, Christos H.; VAZIRANI, Umesh. Algoritmos . São Paulo: McGraw-Hill, c2009. xiv, 320 p. ISBN 9788577260324.
\newline \newline 
HILLIER, FREDERICK S.; LIEBERMAN, GERALD J. Introdução a Pesquisa Operacional. MCGRAW HILL. 9a edição. (ISBN: 8580551188)
}


\DetalhesDisciplina{
    QXD0006 - Cálculo Diferencial e Integral I
}{% EMENTA
    Derivada; Aplicações; Integrais.
}{% BIBLIOGRAFIA BÁSICA
    IEZZI, Gelson; MURAKAMI, Carlos; MACHADO, Nílson José. Fundamentos de matemática elementar: 8 :limites, derivadas, noções de integral . 6. ed. São Paulo, SP: Atual, 2005. 263 p. ISBN 8535705473 (broch.)
\newline \newline 
LEITHOLD, Louis. O Cálculo com geometria analítica. 3. ed. São Paulo: Harbra, c1994. 2 v. ISBN 8529400941v.1 (broch.)
\newline \newline 
GUIDORIZZI, Hamilton Luiz. Um Curso de Cálculo - Vol. 1, 6ª edição. São Paulo: LTC, 2018. E-book. Disponívelem: https://integrada.minhabiblioteca.com.br/\#/books/9788521635574. Acesso em: 3 de out. 2022
}{% BIBLIOGRAFIA COMPLEMENTAR
    DEMANA, Franklin D. et al. Pré-cálculo. São Paulo, SP: Addison-Wesley, 2009. xv, 380 p. ISBN 9788588639379
\newline \newline 
IEZZI, G. Fundamentos de matemática elementar v.10: Geometria Espacial. 6 ed. Atual, 2005
\newline \newline 
IEZZI, G. Fundamentos de matemática elementar v.7: Geometria Analítica.7ed. Atual, 2005
\newline \newline 
GILAT, A.; SUBRAMANIAM, V. Numerical Methods With Matlab. E-Wiley,2007. ISBN: 9780471734406
\newline \newline 
DEMANA, F. D. Precalculus graphical, numerical, algebraic: media update. 7 ed. Addison-Wesley, 2010. ISBN:9780321356932
\newline \newline 
THOMAS, George Brinton; WEIR, Maurice D.; HASS, Joel. Cálculo. 12. ed. São Paulo, SP: Pearson Education doBrasil, 2012. 2 v. ISBN 9788581430867 v. 1 (broch.)
}


\DetalhesDisciplina{
    QXD0134 - Cálculo Diferencial e Integral II
}{% EMENTA
    Derivadas parciais. Funções transcendentes. Funções hiperbólicas. Noções
de coordenadas polares. As técnicas de integração. Integrais impróprias.
Integrais múltiplas. Séries.
}{% BIBLIOGRAFIA BÁSICA
    THOMAS, George Brinton; WEIR, Maurice D.; HASS, Joel. Cálculo.
12. ed. São Paulo, SP: Pearson Education do Brasil, 2012. 2 v. ISBN
9788581430874 v. 2 (broch.)
\newline \newline 
LEITHOLD, Louis. O Cálculo com geometria analítica. 3. ed. São Paulo,
SP: Harbra, c1994. 2 v. 426 p. ISBN 8529402065 v.2 (broch.).
\newline \newline 
GUIDORIZZI, Hamilton Luiz. Um curso de cálculo. 5. ed. Rio de Janeiro:
Livros Técnicos e Científicos, 2002. 4 v. ISBN 9788521612803 v. 2(broch.).
}{% BIBLIOGRAFIA COMPLEMENTAR
    THOMAS, George Brinton; WEIR, Maurice D.; HASS, Joel. Cálculo.
12. ed. São Paulo, SP: Pearson Education do Brasil, 2012. 2 v. ISBN
9788581430867 v. 1 (broch.)
\newline \newline 
GUIDORIZZI, Hamilton Luiz. Um curso de cálculo. 5. ed. Rio de Janeiro:
Livros Técnicos e Científicos, c2002. v.3
\newline \newline 
BORTOLOSSI, Humberto. Cálculo diferencial a várias variáveis: uma
introdução à teoria de otimização . 5. ed. Rio de Janeiro, RJ: Ed. PUCRio; São Paulo, SP: Loyola, 2011
\newline \newline 
GUIDORIZZI, Hamilton Luiz. Um curso de cálculo. 5. ed. Rio de Janeiro:
Livros Técnicos e Científicos, 2002. 4 v. ISBN 9788521612599 v. 1(broch.)
\newline \newline 
LEITHOLD, L. O cálculo com geometria analítica v.1, 3ed. Harbra, 1994
}


\DetalhesDisciplina{
    QXD0135 - Cálculo Diferencial e Integral III
}{% EMENTA
    Estudo de vetores, retas e planos no R3. Funções de várias variáveis.
Diferenciação. Integrais múltiplas. Campos vetoriais. Jacobiano. Integrais
sobre curvas e superfícies. Teoremas de Green, Gauss e Stokes.
}{% BIBLIOGRAFIA BÁSICA
    THOMAS, George Brinton; WEIR, Maurice D.; HASS, Joel. Cálculo.
12. ed. São Paulo, SP: Pearson Education do Brasil, 2012. 2 v. ISBN
9788581430874 v. 2 (broch.)
\newline \newline 
LEITHOLD, Louis. O Cálculo com geometria analítica. 3. ed. São Paulo,
SP: Harbra, c1994. 2 v. 426 p. ISBN 8529402065 v.2 (broch.).
\newline \newline 
GUIDORIZZI, Hamilton Luiz. Um curso de cálculo. 5. ed. Rio de Janeiro:
Livros Técnicos e Científicos, 2002. 4 v. ISBN 9788521612575 v. 3 (broch.)
}{% BIBLIOGRAFIA COMPLEMENTAR
    GUIDORIZZI, Hamilton Luiz. Um curso de cálculo. 5. ed. Rio de Janeiro:
Livros Técnicos e Científicos, 2002. 4 v. ISBN 9788521612803 v. 2(broch.)
\newline \newline 
THOMAS, George Brinton; WEIR, Maurice D.; HASS, Joel. Cálculo.
12. ed. São Paulo, SP: Pearson Education do Brasil, 2012. 2 v. ISBN 9788581430867 v. 1 (broch.)
\newline \newline 
BOYCE, William E.; DIPRIMA, Richard C. Equações diferenciais elemen-
tares e problemas de valores de contorno. 8. ed. Rio de Janeiro, RJ: Livros
\newline \newline 
Técnicos e Científicos, c2006. xvi, 434 p. ISBN 8521614993 (broch.).
BORTOLOSSI, Humberto. Cálculo diferencial a várias variáveis: uma
introdução à teoria de otimização . 5. ed. Rio de Janeiro, RJ: Ed. PUC-
Rio; São Paulo, SP: Loyola, 2011
\newline \newline 
LEITHOLD, L. O cálculo com geometria analítica v.1, 3ed. Harbra, 1994
}


\DetalhesDisciplina{
    QXD0025 - Compiladores
}{% EMENTA
    Introdução à Compiladores, Análise Léxica, Análise Sintática e Abstração de Sintaxe, Análise Semântica, Tabela de Símbolos e Análise de Escopo;Seleção de Instruções;Análise de Longevidade;Seleção de Registradores;Emissão de Código, Tópicos Especiais em Compiladores.
}{% BIBLIOGRAFIA BÁSICA
    AHO, Alfred V.; LAM, Monica S.; SETHI, Ravi; ULLMAN, Jeffrey D.,
Compiladores: princípios, técnicas e ferramentas, Pearson, 2a Edição, 2007.
\newline \newline 
LOUDEN, Kenneth. Compiladores: princípios e práticas. Thomson Pioneira, 2004.
\newline \newline 
HOPCROFT, John E. Introdução à teoria de autômatos, linguagens e
computação. Rio de Janeiro: Elsevier, 2002. 560 p
}{% BIBLIOGRAFIA COMPLEMENTAR
    DELAMARO, M. E. Como Construir Um Compilador Utilizando Ferramentas Java, 1 ed, 2004 Novatec
\newline \newline 
MENEZES, Paulo Blauth. Linguagens formais e autômatos. 6. ed. Porto
\newline \newline 
Alegre: Bookman, 2011. 215 p. (Livros didáticos ; n.3 Série Livros Didáticos ; 3) ISBN: 9788577807659
\newline \newline 
SEBESTA, Robert W. Conceitos de linguagens de programação. 5. ed.
Porto Alegre: Bookman, 2003. ISBN 8536301716.
FORBELLONE, André Luiz Villar; EBERSPÄCHER, Henri Frederico.
\newline \newline 
Lógica de programação: a construção de algoritmos e estruturas de dados . 3. ed. São Paulo: Makron, 2005. xii, 218 p.
\newline \newline 
SIPSER, Michael. Introdução à teoria da computação. São Paulo: Cengage Learning, c2007.
}


\DetalhesDisciplina{
    QXD0079 - Computação em Nuvem
}{% EMENTA
    Introdução a Computação em Nuvem, Princípios da Computação em Nuvem,
Arquitetura da Computação em Nuvem, Modelos de Serviço: Infraestrutura
como um Serviço, Plataforma como um Serviço e Software como um Serviço,
Gerenciamento de Dados em Nuvem, Middlewares para a Computação em
Nuvem, Gerenciamento e Monitoramento da Nuvem, Migração de Aplicações
para Nuvem.
}{% BIBLIOGRAFIA BÁSICA
    COULOURIS, George F.; DOLLIMORE, Jean; KINDBERG, Tim. Sistemas
distribuídos: conceitos e projeto. 5 ed. Bookman, 2013. 1048 p. ISBN
9788582600535.
\newline \newline 
OZSU, M.Tamer. Principles of distributed database systems. 3nd. ed. New
York: Springer, 2011. ISBN 9781441988331
\newline \newline 
VERAS, M. Arquitetura de Nuvem - Amazon Web Services (AWS). 1a ed.
Brasport, 2013. ISBN 9788574525686 .
}{% BIBLIOGRAFIA COMPLEMENTAR
    JENNINGS, Roger. Cloud computing with the Windows Azure Platform.
Indianapolis, Indiana: Wiley Pub., 2009.
\newline \newline 
WHITE, Tom. Hadoop: the definitive guide. California: O´Reilly, 2009.
\newline \newline 
TANENBAUM, Andrew S.; STEEN, Van Maarten; MARQUES, Arlete
Simille. Sistemas distribuídos: princípios e paradigmas. 2. ed. São Paulo, SP:
Prentice Hall, 2007.
\newline \newline 
JOSUTTIS, Nicolai M. SOA na prática: a arte da modelagem de sistemas
distribuídos. Rio de Janeiro, RJ: Alta Books, 2008.
\newline \newline 
HAY, Chris; PRINCE, Brian H. Azure in action. Stamford, Ct: Manning, 2011.
\newline \newline 
TAURION, Cezar. Cloud Computing: computação em nuvem, transformando o
mundo da Tecnologia da Informação. Rio de Janeiro, RJ: Brasport, 2009.
}


\DetalhesDisciplina{
    QXD0183 - Computação Paralela
}{% EMENTA
    Conceitos de computação paralela. Modelos de computação paralela. Algoritmos paralelos. Ambientes de programação paralela.
}{% BIBLIOGRAFIA BÁSICA
    STALLINGS, William. Arquitetura e organização de computadores. 8.
ed. São Paulo: Pearson Prentice Hall, c2010.
\newline \newline 
TANENBAUM, Andrew S.; BOS, Herbert. Sistemas operacionais modernos. 4. ed. São Paulo: Pearson, 2015.
\newline \newline 
RAUBER, Thomas and Rünger Gudula; Parallel Programming: for Multicore and Cluster Systems; First Edition; Editora Springer; ISBN-10: 364204817X ISBN-13: 978-3642048173
}{% BIBLIOGRAFIA COMPLEMENTAR
    VIANA, Gerardo Valdisio Rodrigues. Meta-heuristicas e programacao paralela em otimizacao combinatoria. Fortaleza: Edições UFC, 1998.
\newline \newline 
NEVES, Julio Cezar. Programação SHELL LINUX. 9. ed. Rio de Janeiro, RJ: Brasport, 2013.
\newline \newline 
TANENBAUM, Andrew S.; STEEN, Van Maarten. Sistemas distribuídos:
princípios e paradigmas. 2. ed.São Paulo, SP: Pearson Education, 2007.
\newline \newline 
SILBERSCHATZ, Abraham; GALVIN, Peter Baer; GAGNE, Greg. Fundamentos de Sistemas Operacionais. São Paulo: LTC, 2015. Ebook. Disponível em: https://integrada.minhabiblioteca.com.br/\#/books/978-85-216-3001-2. Acesso em: 3 de out. 2022.
\newline \newline 
PATTERSON, David A.; HENNESSY, John L.. Arquitetura de computadores: uma abordagem quantitativa. Rio de Janeiro: Elsevier, c2014.
}


\DetalhesDisciplina{
    QXD0153 - Desafios de Programação
}{% EMENTA
    Estruturas de Dados Avançadas. Busca por Padrões. Combinatória. Teoria dos
Números. Backtracking. Algoritmos em Grafos. Programação Dinâmica.
Geometria Computacional.
}{% BIBLIOGRAFIA BÁSICA
    CORMEN,	Thomas	H.	Algoritmos:	teoria	e	prática.	Rio	de	Janeiro:	Elsevier,	2002.	xvii	,	916	p.	ISBN	8535209263	(broch.).	
\newline \newline 
DASGUPTA,	Sanjoy;	PAPADIMITRIOU,	Christos	H.;	VAZIRANI,	Umesh.	
Algoritmos.	São	Paulo:	McGraw-Hill,	c2009.	xiv,	320	p.	ISBN	
9788577260324	(broch.).	
\newline \newline 
KLEINBERG,	Jon;	TARDOS,	Éva.	Algorithm	design.	Boston,	Massachusetts: Pearson/Addison Wesley, c2006. 838 p. ISBN 0321295358 (enc.).
}{% BIBLIOGRAFIA COMPLEMENTAR
    ZIVIANI, Nivio; BOTELHO, Fabiano Cupertino. Projeto de algoritmos: com implementações em java e C++. São Paulo, SP: Thomson Learning, 2007. vii, 620 p. ISBN 8522105251 (broch.).
\newline \newline 
ROSEN, Kenneth H. Matemática discreta e suas aplicações. 6. ed. São
Paulo: McGraw-Hill, c2009. xxi, 982 p. ISBN 9788577260362 (broch.).
\newline \newline 
HALIM, S.; HALIM, F.; Competitive Programming. 1 ed. Ebook. Disponível
em <http://www.comp.nus.edu.sg/~stevenha/myteaching/competitive\_prog
ramming/cp1.pdf>. Acesso
em: 18 de jan de 2016.
\newline \newline 
BOAVENTURA NETTO, Paulo Oswaldo. Grafos: teoria, modelos, algoritmos. 5. ed. rev. e ampl. São Paulo: Edgard Blücher, c2012. xiii, 310
p. ISBN 9788521206804
\newline \newline 
SZWARCFITER, Jayme Luiz; MARKENZON, Lilian. Estruturas de dados e
seus algoritmos. 2. ed. rev. Rio de Janeiro: Livros Técnicos e Científicos,
c1994. 320 p.
\newline \newline 
TOSCANI, Laira V.; VELOSO, Paulo A. S. Complexidade de algoritmos:
análise, projeto e métodos. 3. ed. Porto Alegre: Sagra Luzzato, 2012. 262 p.
(Serie Livros Didáticos Informática UFRGS ; 13). ISBN 9788540701380
}


\DetalhesDisciplina{
    QXD0276 - Desenvolvimento de Software para Dispositivos Móveis
}{% EMENTA
    Visão geral sobre dispositivos móveis: Comparação entre dispositivos de
sensoriamento, celulares, tablets e computadores convencionais; Visão geral sobre
as plataformas de desenvolvimento mais utilizadas, como Android SDK, Iphone
SDK e Windows Mobile. Requisitos e desafios para computação móvel.
Arquitetura de Software Móvel. Comunicação para Software móvel. Middleware
e frameworks para Computação Móvel. Sensibilidade ao contexto e adaptação.
Plataforma Android. Activities e Intents. Interfaces e Layouts. Services.
Localização e Mapas. Sensores disponíveis.
}{% BIBLIOGRAFIA BÁSICA
    SALMRE, Ivo. Writing mobile code: essential software engineering for building
mobile applications. New Jersey: Addison-Wesley, 2005. xviii, 771p. ISBN
9780321269317 (broch.).
\newline \newline 
OLIVEIRA, Diego Bittencourt de; SILVA, Fabrício Machado da; PASSOS,
Ubiratan R C.; et al. Desenvolvimento para dispositivos móveis. [Digite o Local
da Editora]: Grupo A, 2019. E-book. ISBN 9788595029408. Disponível em:
https://app.minhabiblioteca.com.br/\#/books/9788595029408/. Acesso em: 27
mar. 2023.
\newline \newline 
DEITEL, Paul; DEITEL, Harvey; WALD, Alexander. Android 6 para
programadores. [Digite o Local da Editora]: Grupo A, 2016. E-book. ISBN
9788582604120. Disponível em:
https://app.minhabiblioteca.com.br/\#/books/9788582604120/. Acesso em: 27
mar. 2023.
}{% BIBLIOGRAFIA COMPLEMENTAR
    POSLAD, Stefan. Ubiquitous Computing: Smart Devices, Environments and
Interactions. 1 ed, Wiley Publishing, 2009. ISBN13: 9780470035603.
\newline \newline 
FREDERICK, Gail Rahn; LAL, Rajesh. Dominando o desenvolvimento web
para smartphone:construindo aplicativos baseados em JavaScript, CSS, HTML e
Ajax para iPhone, Android, Palm Pre, BlackBerry, Windows Mobile e Nokia
S60. 344 p. Rio de Janeiro: Alta Books, 2011.
\newline \newline 
LECHETA, Ricardo R. Google android: aprenda a criar aplicações para
dispositivos móveis com o Android SDK. 2. ed. rev. ampl. São Paulo, SP:
Novatec, 2010. 608 p. ISBN 9788575222447.
\newline \newline 
SIMAS, Victor L.; BORGES, Olimar T.; COUTO, Júlia M C.; et al.
Desenvolvimento para dispositivos móveis - Volume 2. [Digite o Local da
Editora]: Grupo A, 2019. E-book. ISBN 9788595029774. Disponível em: https://app.minhabiblioteca.com.br/\#/books/9788595029774/. Acesso em: 27
mar. 2023.
\newline \newline 
MORAIS, Myllena Silva de F.; MARTINS, Rafael L.; SANTOS, Marcelo da
Silva dos; et al. Fundamentos de desenvolvimento mobile. Grupo A, 2022. Ebook. ISBN 9786556903057. Disponível em:
https://app.minhabiblioteca.com.br/\#/books/9786556903057/. Acesso em: 27 mar. 2023.
}


\DetalhesDisciplina{
    QXD0253 - Desenvolvimento de Software para Web
}{% EMENTA
    A Internet e protocolos de comunicação; HTML; CSS; Acessibilidade, design responsivo e boas práticas na programação web; JavaScript; Aplicação de JavaScript no contexto do desenvolvimento web no navegador; Chamadas assíncronas com JavaScript.
}{% BIBLIOGRAFIA BÁSICA
    LUCKOW, Décio Heinzelmann; MELO, Alexandre Altair de. Programação Java para a Web. São Paulo, SP: Novatec, 2010. 638 p. ISBN 9788575222386.
\newline \newline 
GEARY, David; HORSTMANN, Cay. Core JavaServer Faces. 3. ed. Rio de Janeiro, RJ: Alta Books, 2012: ISBN: 9788576086420.
\newline \newline 
SILVA, Maurício Samy. HTML 5: a linguagem de marcação que revolucionou a web. 2. ed. São Paulo: Novatec, 2014. 335 p. ISBN 9788575224038.
\newline \newline 
FLANAGAN, David. JAVASCRIPT – O Guia Definitivo. Bookman, 6ª ED./2012, 856583719x/9788565837194.
}{% BIBLIOGRAFIA COMPLEMENTAR
    KURNIAWAN, Budi. Java para a Web com Servlets, JSP e EJB: Budi Kurniawan; tradução Savannah Hartmann; revisão técnica Alfredo Dias da Cunha Júnior. Rio de Janeiro, RJ: Ciência Moderna, 2002. xxiv, 807 p. ISBN 8573932104 (broch.).
\newline \newline 
FREEMAN, Elisabeth; FREEMAN, Eric. Use a cabeça!: HTML com CCS e XHTML. 2. ed. Rio de Janeiro, RJ: Alta Books, 2008. xxxi, 580 p. ISBN 9788576082187 (broch.).
\newline \newline 
URUBATAN, Rodrigo. Ruby on rails: desenvolvimento fácil e rápido de aplicações Web. São Paulo, SP: Novatec, 2009. 285 p. ISBN 9788575221846 (broch.).
\newline \newline 
GONÇALVES, Edson. Desenvolvendo aplicações Web com NetBeans IDE 6. Rio de Janeiro: Ciência Moderna, 2008. 581 p. : CD-ROM ISBN 97885739366742.
\newline \newline 
BASHAM, Bryan. Use a cabeça!: Servlets \& JSP. 2. ed. Rio de Janeiro, RJ: Alta Books, 2008. ISBN 9788576082941.
}


\DetalhesDisciplina{
     - Desenvolvimento de Software para Web 2
}{% EMENTA
    Arquitetura web. HTTP API. Arquitetura REST. Streaming API. Aplicações e Servidores Web. Servidores de Aplicação. Frameworks de front-end e back-end para web.
}{% BIBLIOGRAFIA BÁSICA
    ALMEIDA, Flávio. Mean: Full stack JavaScript para aplicações web com MongoDB, Express,
Angular e Node. São Paulo, SP: Casa do Código, [2015]. xxiii, 361 p. (Caelum). ISBN
9788555190469 (broch.)
\newline \newline 
OLIVEIRA, Cláudio Luís V.; ZANETTI, Humberto Augusto P. Node.js: programe de forma rápida e prática. Rio de Janeiro: Expressa, 2021. E-book. p.15. ISBN 9786558110217. Disponível em: https://app.minhabiblioteca.com.br/reader/books/9786558110217/. Acesso em: 06 nov. 2024.
\newline \newline 
JIN, Brenda; SAHNI, Saurabh; SHEVAT, Amir. Designing Web APIs: Building APIs That Developers Love. Editora O’Reilly Media, Inc., 2018.
}{% BIBLIOGRAFIA COMPLEMENTAR
    URUBATAN, Rodrigo. Ruby on rails: desenvolvimento fácil e rápido de aplicações Web. São Paulo, SP: Novatec, 2009. 285 p. ISBN 9788575221846 (broch.).
\newline \newline 
DUCKETT, Jon. PHP\&MYSQL: desenvolvimento web no lado do servidor. Rio de Janeiro: Editora Alta Books, 2024. E-book. p.173. ISBN 9786555205930. Disponível em: https://app.minhabiblioteca.com.br/reader/books/9786555205930/. Acesso em: 06 nov. 2024.
\newline \newline 
FLANAGAN, David. JAVASCRIPT – O Guia Definitivo. Bookman, 6ª ED./2012, 856583719x/9788565837194.
\newline \newline 
SMITH, Ben. JSON básico: conheça o formato de dados preferido da web. São Paulo: Novatec, 2015. 400 p. ISBN 9788575224366 (broch.).
\newline \newline 
FERRAZ, Reinaldo. Acessibilidade na web: boas práticas para construir sites e aplicações acessíveis. São Paulo: Casa do Código, 2020. 246 p. ISBN 978-65-86110-10-4.
\newline \newline 
ALVES, William P. Java para Web - Desenvolvimento de Aplicações. Rio de Janeiro: Érica, 2015. E-book. p.1. ISBN 9788536519357. Disponível em: https://app.minhabiblioteca.com.br/reader/books/9788536519357/. Acesso em: 16 jan. 2025.
}


\DetalhesDisciplina{
    QXD0232 - Educação Ambiental
}{% EMENTA
    Educação ambiental, conceitos e metodologias. Histórico da educação
ambiental (EA). Conferências e marcos legais da EA. Desenvolvimento
Sustentável. Perspectivas filosóficas do Desenvolvimento Sustentável.
Transdisciplinaridade e Educação Ambiental. A Práxis em Educação
Ambiental.
}{% BIBLIOGRAFIA BÁSICA
    REIGOTA, Marcos. O que é educação ambiental. 2.ed. São Paulo, SP:
Brasiliense, 2012. 107p. (Coleção Primeiros Passos; 292). ISBN
9788511001228 (broch.).
\newline \newline 
ESMERALDO, Gema Galgani Silveira Leite. Educação, v. 5. Fortaleza:
Fundação Demócrito Rocha/Assembléia Legislativa do Estado do Ceará/
Sistema de Transmissão Nordeste S.A, 2015. 52 p. (Convivência com o
semiárido). ISBN 9788575297063 (enc.).
\newline \newline 
OLIVEIRA NETO, João Martins de. Gestão, v. 6. Fortaleza: Fundação
Demócrito Rocha/ Assembléia Legislativa do Estado do Ceará/ Sistema de Transmissão Nordeste S.A, 2015. 50 p. + 1 DVD (Convivência com o semiárido). ISBN 9788575297070 (enc.).
}{% BIBLIOGRAFIA COMPLEMENTAR
    MANSUR, Ricardo. Governança de TI verde: o ouro verde da nova TI . Rio de
Janeiro: Ciência Moderna, 2011. 212 p. ISBN 9788539900459 (broch.).
\newline \newline 
SUPERINTENDÊNCIA ESTADUAL DO MEIO AMBIENTE. Programa de
Educação Ambiental do Ceará : PEACE . Fortaleza: SEMACE, 2009.
\newline \newline 
SUASSUNA, Joao. Água, v. 1. Fortaleza: Fundação Demócrito Rocha/
Assembléia Legislativa do Estado do Ceará/ 2015. 52 p. + 1 DVD
(Convivência com o semiárido). ISBN 9788575297025 (enc.).
\newline \newline 
MARTINS, Eduardo Sávio Passos Rodrigues; OLIVEIRA, Sônia Barreto
Perdigão de; CARVALHO, Margareth Silvia Benicio de Souza. Clima, v.3.
Fortaleza: Fundação Demócrito Rocha/ Assembléia Legislativa do Estado do
Ceará, Sistema de Transmissão Nordeste S.A, 2015.. 52 p. + 1 DVD
(Convivência com o semiárido). ISBN 9788575297049 (enc.).
\newline \newline 
FABRE, Nicolas Arnaud. Produção, v. 4. Fortaleza: Fundação Demócrito
Rocha/ Assembléia Legislativa do Estado do Ceará/ Sistema de Transmissão
Nordeste S.A, 2015.. 52 p. + 1 DVD (Convivência com o semiárido). ISBN
9788575297056 (enc.).
\newline \newline 
SILVA, José Borzacchiello da. Terra, v. 2. Fortaleza: Fundação Demócrito
Rocha/ Assembléia Legislativa do Estado do Ceará, Sistema de Transmissão
Nordeste S.A, 2015. 52 p. + 1 DVD (Convivência com o semiárido). ISBN
9788575297032 (enc.).
}


\DetalhesDisciplina{
    QXD0245 - Educação em Direitos Humanos
}{% EMENTA
    Direitos Humanos, democratização da sociedade, cultura e paz e cidadanias. O nascituro, a criança e o adolescente como sujeitos de direito: perspectiva histórica e legal. O ECA e a rede de proteção integral. Educação
em direitos humanos na escola: princípios orientadores e metodologias. O
direito à educação como direito humano potencializador de outros direitos. Movimentos, instituições e redes em defesa do direito à educação.
Igualdade e diversidade: direitos sexuais, diversidade religiosa e diversidade étnica. Os direitos humanos de crianças e de adolescentes nos meios
de comunicação e nas mídias digitais.
}{% BIBLIOGRAFIA BÁSICA
    MORAES, Alexandre D. Direitos Humanos Fundamentais. [Digite o Local
da Editora]: Grupo GEN, 2021. E-book. ISBN 9788597026825. Disponível
em: https://app.minhabiblioteca.com.br/\#/books/9788597026825.
Acesso em: 23 set. 2022.
\newline \newline 
NETO, Silvio B. Curso de Direitos Humanos. [Digite o Local da Editora]:
Grupo GEN, 2021. E-book. ISBN 9788597028249. Disponível em: https:
//app.minhabiblioteca.com.br/\#/books/9788597028249. Acesso em:
23 set. 2022.
\newline \newline 
SCARANO, Renan Costa V.; DORETO, Daniella T.; ZUFFO, Sílvia; et
al. Direitos humanos e diversidade. [Digite o Local da Editora]: Grupo
A, 2018. E-book. ISBN 9788595028012. Disponível em: https://app.
minhabiblioteca.com.br/\#/books/9788595028012. Acesso em: 23 set.
2022.
}{% BIBLIOGRAFIA COMPLEMENTAR
    BRASIL/SECRETARIA ESPECIAL DE DIREITOS HUMANOS. Estatuto da Criança e do Adolescente (Lei 8069/90). Brasília, 2008. Disponível
em: http://www.planalto.gov.br/ccivil\_03/LEIS/L8069.htm
\newline \newline 
CADERNO DE EDUCAÇÃO EM DIREITOS HUMANOS. Educação em
Direitos Humanos: Diretrizes Nacionais. Brasília, DF: Secretaria Nacional de Promoção e Defesa dos Direitos Humanos, 2013. 76p. Disponível em: https://www.gov.br/mdh/pt-br/navegue-por-temas/educacaoem-direitos-humanos/DiretrizesNacionaisEDH.pdf
\newline \newline 
COMITÊ NACIONAL DE EDUCAÇÃO EM DIREITOS HUMANOS/
SECRETARIA ESPECIAL DOS DIREITOS HUMANOS. Plano Nacional de Educação em Direitos Humanos. Brasília: MEC/MJ/UNESCO, 2009. Disponível em: http://portal.mec.gov.br/docman/2191plano-nacional-pdf/file
\newline \newline 
NOLETO, M.J. Abrindo espaços: educação e cultura para a paz. Brasília:UNESCO, 2004. Disponível em: https://unesdoc.unesco.org/ark:
/48223/pf0000131816
\newline \newline 
BRASIL. Lei Maria da Penha: um avanço no combate à violência contra a
mulher. Brasília, DF: Senado Federal / Secretaria Especial de Editoração
e Publicações, 2007. 26 p.
\newline \newline 
MONDAINI, Marco. Direitos Humanos. [Digite o Local da
Editora]: Grupo Almedina (Portugal), 2020. E-book. ISBN
9788562938368. Disponível em: https://app.minhabiblioteca.com.
br/\#/books/9788562938368. Acesso em: 23 set. 2022.
}


\DetalhesDisciplina{
    QXD0029 - Empreendedorismo
}{% EMENTA
    Conceito de empreendedorismo. A formação da personalidade. O processo
comportamental. Fatores de sucesso, o perfil do empreendedor.
Desenvolvimento de habilidades empreendedoras. Lições e práticas
internacionais. Empreendedorismo no Brasil. Importância das MPEs na
economia. Globalização dos mercados, dos negócios e das oportunidades.
Pesquisas Tecnológicas. Propriedade Intelectual. Transferência de
Tecnologia. Papel da inovação. Ambientes de pré-incubação e incubação de ideias. Incubadoras de empresas. Parques Tecnológicos. Capital de Risco.
Recursos de Fomento. Fontes de Financiamento. Fundos Setoriais.
Programas governamentais. Plano de Negócio. Ferramentas de Plano de
Negócios. Projetos.
}{% BIBLIOGRAFIA BÁSICA
    SALIM, César Simões. Construindo planos de negócios: todos os passos
necessários para planejar e desenvolver negócios de sucesso. 3. ed. rev. e
atual. Rio de Janeiro, RJ: Elsevier, 2005. xiv, 332 p. ISBN 9788535217360
\newline \newline 
DORNELAS, Jose Carlos Assis. Empreendedorismo: transformando ideias
em negócios. 3. ed. rev. atual. Rio de Janeiro: Elsevier, 2008. 232 p. ISBN
9788535232707
\newline \newline 
DOLABELA, Fernando. O segredo de Luísa: uma idéia, uma paixão e um
plano de negócios: como nasce o empreendedor e se cria uma empresa. Rio
de Janeiro: Sextante, 2008. 299 p. ISBN 9788575423387 (broch.).
}{% BIBLIOGRAFIA COMPLEMENTAR
    HISRICH, Roberto D. Empreendedorismo. 7. ed. Porto Alegre, RS: Bookman,
2009. 662 p. ISBN 9788577803460 (broch. ).
\newline \newline 
FARAH, Osvaldo Elias. Empreendedorismo estratégico. São Paulo: Cengage
Learning, 2008. 251 p. ISBN 9788522106080 (broch.).
\newline \newline 
FERRARI, Roberto. Empreendedorismo para computação: criando negócios
em tecnologia. Rio de Janeiro, RJ: Elsevier, 2010. 164 p. ISBN
9788535234176 (broch.).
\newline \newline 
CHIAVENATO, Idalberto. Empreendedorismo: dando asas ao espírito
empreendedor: empreendedorismo e viabilização de novas empresas, um
guia eficiente para iniciar e tocar seu próprio negócio. 3. ed. rev. e atual. São
Paulo, SP: Saraiva, 2008. 281 p. ISBN 9788502067448 (broch.).
\newline \newline 
FRIEDMAN, Thomas L. O mundo é plano: uma breve história do século XXI.
2. ed., rev. e atual. Rio de Janeiro: Objetiva, 2007. 557 p. ISBN
9788573028638.
}


\DetalhesDisciplina{
    QXD0019 - Engenharia de Software
}{% EMENTA
    Visão geral e introdutória dos princípios fundamentais e ético-profissionais
da Engenharia de Software. Introdução às atividades de engenharia de
requisitos; projeto de software; modelos de desenvolvimento; e gerenciamento (qualidade, estimativa de custo, configuração, etc.) na engenharia
de software.
}{% BIBLIOGRAFIA BÁSICA
    SOMMERVILLE, I. Engenharia de software. 9 ed. Addison Wesley, 2011.
ISBN: 9788579361081.
\newline \newline 
PRESSMAN, Roger; MAXIM, Bruce. Engenharia de Software. São
Paulo: AMGH, 2016. E-book. Disponível em: https://integrada.
minhabiblioteca.com.br/\#/books/9788580555349. Acesso em: 3 de
out. 2022.
\newline \newline 
LARMAN, Craig. Utilizando UML e padroes : uma introducao a analise
e ao projeto orientados a objetos. 5. ed. Porto Alegre: Bookman, 2007.
695 p. ISBN 856003152-9
}{% BIBLIOGRAFIA COMPLEMENTAR
    TELES, V. Extreme programming. Novatec. 2004. ISBN: 8575220470
\newline \newline 
MOLINARI, L. Gerência de configuração: técnicas e práticas no desenvolvimento do software. VISUAL BOOKS. ISBN: 8575022105
\newline \newline 
DELAMARO, Marcio. Introdução ao Teste de Software. São Paulo:
GEN LTC, 2016. E-book. Disponível em: https://integrada.
minhabiblioteca.com.br/\#/books/9788595155732. Acesso em: 3 de
out. 2022.
\newline \newline 
KERIEVSKY, J. Refatoração para padrões. Bookman, 2008.
ISBN:9788577802449
\newline \newline 
PEZZÉ, M.; YOUNG, M. Teste e análise de software: processos, princípios
e técnicas. Bookman, 2008. ISBN:9788577802623
\newline \newline 
PILONE, D.; MILES, R. Use a cabeça! desenvolvimento de software.
ALTA BOOKS, 2008.
}


\DetalhesDisciplina{
    QXD0179 - Estatística Multivariada
}{% EMENTA
    Introdução a análise multivariada. Análise dos componentes principais. Análise
de agrupamentos (clusters). Análise de discriminantes. Modelos de regressão.
}{% BIBLIOGRAFIA BÁSICA
    HAIR, Joseph F. et al. Análise multivariada de dados. 6. ed. Porto Alegre:
Bookman, 2009. 688 p. ISBN 9788577804023 (enc.).
\newline \newline 
CORRAR, Luiz J; PAULO, Edilson; DIAS FILHO, José Maria; (Cord.)
FUNDAÇÃO INSTITUTO DE PESQUISAS CONTÁBEIS, ATUARIAIS E
FINANCEIRAS. Análise multivariada: para os cursos de administração, ciências
contábeis e economia. São Paulo, SP: Atlas, 2007. xxiv, 541 p. ISBN
9788522447077 (broch.).
\newline \newline 
IZENMAN , Alan SPRINGERLINK (ONLINE SERVICE). Modern Multivariate
Statistical Techniques : Regression, Classification, and Manifold Learning .
Springer eBooks New York, NY: Springer-Verlag New York, 2008. (Springer
Texts in Statistics,) ISBN 9780387781891. Disponível em:
<http://dx.doi.org/10.1007/978-0-387-78189-1>. Acesso em : 21 set. 2010.
\newline \newline 
MINGOTI, Sueli Aparecida. Análise de dados através de métodos de estatística
multivariada: uma abordagem aplicada. Belo Horizonte: Editora UFMG, 2005.
295p. : (Didática.8) ISBN 857041451X (broch.)
}{% BIBLIOGRAFIA COMPLEMENTAR
    LARSON, Ron; FARBER, Betsy. Estatística aplicada. 4. ed. São Paulo, SP:
Pearson/ Prentice Hall, 2010. xiv,637 p. ISBN 9788576053729 (broch.).
OJA, H.; Multivariate Nonparametric Methods with R. Springer Ebooks New
York: Springer New York. Disponível em
<http://link.springer.com/book/10.1007/978-1-4419-0468-3>. Acesso em 19 de
janeiro de 2016.
\newline \newline 
BARBETTA, Pedro Alberto; REIS, Marcelo Menezes; BORNIA, Antonio Cezar.
Estatística para cursos de engenharia e informática. 3. ed. São Paulo, SP:
Atlas, 2010. 410 p.
\newline \newline 
WALPOLE, Ronald E. Probabilidade e estatística: para engenharia e ciências.
8. ed. São Paulo, SP: Pearson/ Prentice Hall, 2009. xiv, 491 p. ISBN
9788576051992 (broch.).
\newline \newline 
JOHNSON, Richard A.; WICHERN, Dean W. Applied Multivariate Statistical
Analysis. Phi Learning Private Limited, 2010. ISBN-10 8120345878 ISBN-13
9788120345874
\newline \newline 
FÁVERO, Luiz Paulo et al. Análise de dados: modelagem multivariada para
tomada de decisões. Rio de Janeiro, RJ: Elsevier, Campus, 2009. xx, 646 p.
ISBN 9788535230468
\newline \newline 
EVERITT, Brian; HOTHORN, Torsten. An introduction to applied multivariate
analysis with R. 978-3-642-13312-1. Springer Science \& Business Media, 2011. ISBN
Disponível em <http://dx.doi.org/10.1007/978-1-4419-9650-3>
}


\DetalhesDisciplina{
    QXD0010 - Estrutura de Dados
}{% EMENTA
    Noções de análise de algoritmos, Recursividade, Tipos Abstratos de Dados,
Algoritmos de Ordenação, Listas Sequenciais e Encadeadas, Pilhas, Filas,
Árvores.
}{% BIBLIOGRAFIA BÁSICA
    FEOFILOFF, Paulo. . Algoritmos em linguagem C. Rio de Janeiro: Elsevier,
2009. 208p. ISBN 9788535232493 (broch.).
\newline \newline 
DROZDEK, Adam. Estrutura de dados e algoritmos em C++. São Paulo:
Thomson, 2002. 579p ISBN 852210295 (broch.).
\newline \newline 
ZIVIANI, Nivio; BOTELHO, Fabiano Cupertino. Projeto de algoritmos: com
implementações em java e C++. São Paulo, SP: Thomson Learning, 2007. 620
p. ISBN 8522105251 (broch.).
}{% BIBLIOGRAFIA COMPLEMENTAR
    CORMEN, Thomas H. Algoritmos: teoria e prática. Rio de Janeiro: Elsevier,
2002. 916 p. ISBN 8535209263 (broch.)
\newline \newline 
SZWARCFITER, Jayme Luiz; MARKENZON, Lilian. Estruturas de dados e
seus algoritmos. 2. ed. rev. Rio de Janeiro: Livros Técnicos e Científicos,
c1994. 320 p. ISBN 8521610149 (broch.).
\newline \newline 
ASCENCIO, Ana Fernanda Gomes; CAMPOS, Edilene Aparecida Veneruchi
de. Fundamentos da programação de computadores: algoritmos, Pascal,
C/C++ e java. 2. ed. São Paulo, SP: Prentice Hall, 2007. viii, 434 p. ISBN
8576051480 (broch.).
\newline \newline 
AGUILAR, Luis. Fundamentos de programação: algoritmos, estrutura de dados
e objetos. São Paulo, SP: McGraw-Hill, 2008. 690 p. ISBN 9788586804960
(broch.).
\newline \newline 
CELES, Waldemar; CERQUEIRA, Renato; RANGEL, José Lucas. Introdução a
estrutura de dados: com técnicas de programação em C. Rio de Janeiro:
Elsevier, 2004. 294 p. (Campus, Sociedade Brasileira de Computação) ISBN
8535212280 (broch.).
}


\DetalhesDisciplina{
    QXD0115 - Estrutura de Dados Avançada
}{% EMENTA
    Balanceamento de árvores de busca. Filas de prioridade (​heaps)​ . Estruturas de
dados para conjuntos disjuntos (​union-find​). Grafos: representação e
caminhamento. Tabelas ​hash​ e tratamento de colisões.
}{% BIBLIOGRAFIA BÁSICA
    CORMEN, Thomas H. Algoritmos: teoria e prática. Rio de Janeiro: Elsevier,
2002. xvii , 916 p. ISBN: 8535209263
\newline \newline 
DROZDEK, Adam. Estrutura de dados e algoritmos em C++. São Paulo:
Thomson, 2002. 579p.
\newline \newline 
SZWARCFITER, Jayme Luiz; MARKENZON, Lilian. Estruturas de dados e
seus algoritmos. 2. ed. rev. Rio de Janeiro: Livros Técnicos e Científicos,
c1994. 320 p.
\newline \newline 
SEDGEWICK, Robert, WAYNE, Kevin. Algorithms (4th Edition). Addison-Wesley Professional; 4 edition, 2011. ISBN: 032157351X
}{% BIBLIOGRAFIA COMPLEMENTAR
    GOODRICH, Michael T.; TAMASSIA, Roberto. Data structures and algorithms
in Java. 5th ed. New York, NY: J. Wiley \& Sons, 2010. xxii, 714 p. ISBN
9780470383261 (enc.).
\newline \newline 
CELES, Waldemar; CERQUEIRA, Renato; RANGEL, José Lucas. Introdução a
estruturas de dados: com técnicas de programação em C. Rio de Janeiro, RJ:
Elsevier: Campus, 2004. xiv, 294 p. (Editora Campus). ISBN 8535212280
(broch.).
\newline \newline 
DASGUPTA, Sanjoy; PAPADIMITRIOU, Christos H.; VAZIRANI, Umesh.
Algoritmos. São Paulo: McGraw-Hill, c2009. xiv, 320 p. ISBN 9788577260324
(broch.).
\newline \newline 
ROSEN, Kenneth H. Matemática discreta e suas aplicações. 6. ed. São Paulo:
McGraw-Hill, c2009. xxi, 982 p. ISBN 9788577260362 (broch.).
\newline \newline 
MEHLHORN, Kurt; SANDERS, Peter SPRINGERLINK (ONLINE SERVICE).
Algorithms and Data Structures : The Basic Toolbox . Springer eBooks Berlin,
Heidelberg: Springer-Verlag Berlin Heidelberg, 2008. ISBN 9783540779780.
Disponível em : <http://dx.doi.org/10.1007/978-3-540-77978-0>. Acesso em :
21 set. 2010.
\newline \newline 
KARUMANCHI, Narasimha. Data Structures and Algorithms Made Easy.
Createspace Pub, 2011. ISBN 1468108867
}


\DetalhesDisciplina{
    QXD0508 - Ética e Legislação
}{% EMENTA
    Conceituação do direito moral e ética. Direito e garantias constitucionais. Direito do trabalho. Direito comercial. Direito do consumidor. Direito autoral. Direito civil. Ética profissional.
}{% BIBLIOGRAFIA BÁSICA
    BARGER, Robert N. Ética na computação: uma abordagem baseada em casos.
Rio de Janeiro, RJ: LTC, 2011. xiv, 226 p. ISBN 9788521617761 (broch.).
\newline \newline 
GONÇALVES, Victor Hugo P. Marco Civil da Internet Comentado. São Paulo:
Grupo GEN, 2016. E-book. ISBN 9788597009514. Disponíví el em:
https://app.minhabiblioteca.com.br/\#/books/9788597009514/. Acesso em:
06 out. 2022.
\newline \newline 
ALMEIDA, Guilherme Assis D.; CHRISTMANM, Martha O. Ética e Direito: Uma
Perspectiva Integrada, 3ª edição. São Paulo: Grupo GEN, 2009. E-book. ISBN
9788522467150. Disponíví el em:
https://app.minhabiblioteca.com.br/\#/books/9788522467150/. Acesso em:
05 out
}{% BIBLIOGRAFIA COMPLEMENTAR
    DINIZ, Debora, et al. (Org.). Ética em pesquisa: temas globais. Brasílíia, DF:
Letras Livres, Editora Universidade de Brasílíia, 2011. 404 p. ISBN
9788598070209. Disponíví el em:
<http://www.repositoriobib.ufc.br/00003e/00003e82.pdf>.
\newline \newline 
SÁ, Antônio Lopes D. Ética Profissional. São Paulo: Grupo GEN, 2019. E-book.
ISBN 9788597021653. Disponíví el em:
https://app.minhabiblioteca.com.br/\#/books/9788597021653/. Acesso em:
26 set. 2022.
\newline \newline 
LÉVY, Pierre. As tecnologias da inteligência: o futuro do pensamento na era da informática . 2. ed. Rio de Janeiro: Editora 34, 2010. 206 p. (Coleção TRANS). ISBN 9788585490157 (broch.)
\newline \newline 
FLEISHER, Soraya. SCHUCH, Patrice (Orgs.) Ética e regulamentação na
pesquisa antropológica. Brasílíia: LetrasLivres: Editora Universidade de
Brasílíia, 2010. ISBN 978-85-98070-24-7 ISBN 978-85-230-1246-5.
Disponíví el em: <http://www.repositoriobib.ufc.br/00003e/00003e7d.pdf>.
\newline \newline 
MASIERO, Paulo Cesar. Ética em computação. São Paulo, SP: EDUSP, 2008. 213
p. (Acadêmica ; 32). ISBN 9788531405754
}


\DetalhesDisciplina{
    <08> - Ética em Inteligência Artificial
}{% EMENTA
    Responsabilidade legal em algoritmos de IA (vieses, justiça). Privacidade, segurança e vigilância em sistemas inteligentes. Explicabilidade, auditabilidade em algoritmos de IA. Impactos da IA no mercado de trabalho e nas relações sociais. Regulação e governança da IA: Diretrizes éticas e iniciativas internacionais de governança (UNESCO, OCDE, UE). Lei Geral de Proteção de Dados Pessoais (LGPD). Marco legal brasileiro. Estudo de casos em cenários reais. 
}{% BIBLIOGRAFIA BÁSICA
    FEFERBAUM, Marina; SILVA, Alexandre Pacheco da; COELHO, Alexandre Z.; et al. Ética, Governança e Inteligência Artificial. São Paulo: Almedina, 2023. E-book. p.1. ISBN 9786556279145. Disponível em: https://app.minhabiblioteca.com.br/reader/books/9786556279145/. Acesso em: 22 mai. 2025.
\newline \newline 
BLACKMAN, Reid. Máquinas Éticas: seu guia conciso para uma IA totalmente imparcial, transparente e respeitosa. Rio de Janeiro: Editora Alta Books, 2024. E-book. p.1. ISBN 9788550822426. Disponível em: https://app.minhabiblioteca.com.br/reader/books/9788550822426/. Acesso em: 23 mai. 2025.
\newline \newline 
RUSSELL, Stuart J.; NORVIG, Peter. Inteligência Artificial: Uma Abordagem Moderna. 4. ed. Rio de Janeiro: GEN LTC, 2022. E-book. p.Capa. ISBN 9788595159495. 
}{% BIBLIOGRAFIA COMPLEMENTAR
    PINTO, Rodrigo Alexandre L.; NOGUEIRA, Jozelia. Inteligência Artificial e Desafios Jurídicos: Limites Éticos e Legais. São Paulo: Almedina, 2023. E-book. p.1. ISBN 9786556279268. Disponível em: https://app.minhabiblioteca.com.br/reader/books/9786556279268/. Acesso em: 23 mai. 2025.
\newline \newline 
UNESCO. Recomendação sobre a Ética da Inteligência Artificial. Paris: UNESCO, 2022. Disponível em: https://unesdoc.unesco.org/ark:/48223/pf0000381137\_por. Acesso em: 30 maio 2025.
\newline \newline 
STAHL, Bernd Carsten. Artificial Intelligence for a Better Future: An Ecosystem Perspective on the Ethics of AI and Emerging Digital Technologies. Cham: Springer, 2021. E-book. ISBN 978-3-030-69978-9. Disponível em: https://link.springer.com/book/10.1007/978-3-030-69978-9. Acesso em: 30 maio 2025.
\newline \newline 
BARGER, R. N. Ética na computação: UMA ABORDAGEM BASEADA EM CASOS. LTC, 2011. ISBN: 9788521617761
\newline \newline 
STAHL, Bernd Carsten; SCHROEDER, Doris; RODRIGUES, Rowena. Ethics of Artificial Intelligence: Case Studies and Options for Addressing Ethical Challenges. Cham: Springer, 2023. (SpringerBriefs in Research and Innovation Governance). Disponível em: https://doi.org/10.1007/978-3-031-17040-9. Acesso em: 30 maio 2025.
}

\DetalhesDisciplina{
    QXD0011 - Fundamentos de Banco de Dados
}{% EMENTA
    Visão geral do gerenciamento de banco de dados. Arquitetura de um Sistema Gerenciador de Banco de Dados. Modelagem e projeto de banco
de dados: Modelo Entidade-Relacionamento, Modelo Relacional e Projeto
de Bancos de Dados Relacionais. SQL. Projeto Avançado: Restrições de
Integridade e Normalização.
}{% BIBLIOGRAFIA BÁSICA
    SILBERSCHATZ, Abraham. Sistema de Banco de Dados. São Paulo:
GEN LTC, 2020. E-book. Disponível em: https://integrada.
minhabiblioteca.com.br/\#/books/9788595157552. Acesso em: 3 de
out. 2022.
\newline \newline 
ELMASRI, R.; NAVATHE, S. B. Sistemas de banco de dados. 6 ed.
Pearson/Addison-Wesley, 2011. ISBN: 9788579360855
\newline \newline 
HEUSER, Carlos Alberto. Projeto de banco de dados - V4 - UFRGS. São
Paulo: Bookman, 2011. E-book. Disponível em: https://integrada.
minhabiblioteca.com.br/\#/books/9788577804528. Acesso em: 3 de
out. 2022.
}{% BIBLIOGRAFIA COMPLEMENTAR
    RAMAKRISHNAN, Raghu; GEHRKE, Johannes. Sistemas de Gerenciamento de Bancos de Dados. São Paulo: AMGH, 2008. E-book. Disponível em: https://integrada.minhabiblioteca.com.br/\#/books/
\newline \newline 
9788563308771. Acesso em: 3 de out. 2022.
DATE, C.J. Introdução a Sistemas de Banco de Dados. 8 ed. Campus,
2004. ISBN. 9788535212730
\newline \newline 
OLIVEIRA, C.H. SQL: Curso prático. Novatec, 2002. ISBN:9788575220245
BEIGHLEY, Lynn. Use a cabeça! SQL. Alta Books, 2008. ISBN:
9788576022101
\newline \newline 
ALVES, William Pereira. Banco de Dados. São Paulo: Érica, 2014. Ebook. Disponível em: https://integrada.minhabiblioteca.com.br/\#/
books/9788536518961. Acesso em: 27 de set. 2022.
}


\DetalhesDisciplina{
    QXD0001 - Fundamentos de Programação
}{% EMENTA
    Algoritmos, Conceitos Fundamentais de Programação, Expressões, Controles
de Fluxo, Funções e Procedimentos, Ponteiros, Vetores e Matrizes, Cadeias
de Caracteres, Alocação Dinâmica, Tipos Estruturados e Arquivos.
}{% BIBLIOGRAFIA BÁSICA
    MEDINA, Marco; FERTIG, Cristina. Algoritmos e programação: teoria e prática
2ed. Novatec, 2004.ISBN: 9788575220733/857522073X 2. ed. São Paulo, SP:
Novatec, 2006. 384 p. ISBN 857522073X (broch.).
\newline \newline 
ASCENCIO, A. F. G.;CAMPOS, E. A. V. Fundamentos da programação de
computadores: algoritmos, Pascal, C/C++ e Java. 2 ed. Prentice Hall, 2007.
ISBN: 978576051480 3. ed. São Paulo: Pearson Education do Brasil, c2012.
x, 569 p. ISBN 9788564574168 (broch.).
\newline \newline 
CELES, W.; CERQUEIRA, R.; RANGEL, J. L. Introdução à estrutura de dados:
com técnica de programação em C. Elsevier, 2004. ISBN: 8535212280.
\newline \newline 
FORBELLONE, A.L.V.; EBERSPACHER, H.F. Lógica de programação: a
construção de algoritmos. 3 ed. Prentice Hall, 2005
}{% BIBLIOGRAFIA COMPLEMENTAR
    SCHILDT, Herbert. C completo e total. 3 ed., rev. atual. São Paulo:
Pearson/Makron Books, c1997. xx. 827 p + 1 CD-ROM ISBN 8534605955
(broch.).
\newline \newline 
DEITEL, Harvey. M.; C++ como programar. 5 ed. Prentice Hall, 2006. 5. ed.
São Paulo, SP: Pearson/Prentice Hall, 2006. xlii,1163 p. + cd-rom ISBN
8576050560 (broch.).
\newline \newline 
AGUILAR, Luis JOYANNES. Fundamentos de programação: algoritmos,
estruturas de dados e objetos. São Paulo: McGraw-Hill, c2008. xxix, 690 p.
ISBN 9788586804960 (broch.).
\newline \newline 
MONTGOMERY, Eduard. Programando em C: simples \& prático. Rio de
Janeiro, RJ: Alta Books, 2006. 157 p. ISBN 9788576081210 (bronch).
\newline \newline 
FEOFILOFF, PAULO. Algoritmos em linguagem C. Rio de Janeiro : Elsevier,
2009. ISBN : 9788535232493
}


\DetalhesDisciplina{
    QXD0254 - Gerência de Projetos
}{% EMENTA
    Conceitos, terminologia e contexto de gerência de projetos. Ciclo de vida de
produto e projeto. Interessados (stakeholders). Organização de empresas
(funcionais, matriciais e baseadas em projetos). Estratégias para seleção de
projetos. Processos de gerência de projetos. Gerência de escopo. Gerência de
tempo (definição de atividades, sequenciamento de atividades, estimativa de
recursos, estimativa de duração, desenvolvimento de cronograma e controle de
cronograma). Gerência de custos (estimativas, orçamento e controle). Gerência
de qualidade. Gerência de recursos humanos. Gerência de comunicação. Gerência
de riscos. Gerência de aquisições. Gerência de integração (desenvolver carta de
projeto, desenvolver escopo preliminar, desenvolver plano de gerência de projeto,
dirigir e gerenciar a execução de projetos, monitorar e controlar atividades de
projeto, controle de mudanças e fechamento do projeto).
}{% BIBLIOGRAFIA BÁSICA
    RUBIN, Kenneth S. Scrum essencial: um guia prático para o mais popular
processo ágil. [Digite o Local da Editora]: Editora Alta Books, 2017. E-book.
ISBN 9788550804118. Disponível em:
https://app.minhabiblioteca.com.br/\#/books/9788550804118/. Acesso em: 27
mar. 2023.
\newline \newline 
CARVALHO, Marly M. Fundamentos em Gestão de Projetos - Construindo
Competências para Gerenciar Projetos. [Digite o Local da Editora]: Grupo GEN,
2018. E-book. ISBN 9788597018950. Disponível em:
https://app.minhabiblioteca.com.br/\#/books/9788597018950/. Acesso em: 27
mar. 2023.
\newline \newline 
CAMARGO, Robson Alves de; RIBAS, Thomaz. Gestão ágil de projetos. [Digite
o Local da Editora]: Editora Saraiva, 2019. E-book. ISBN 9788553131891.
Disponível em: https://app.minhabiblioteca.com.br/\#/books/9788553131891/.
Acesso em: 27 mar. 2023.
}{% BIBLIOGRAFIA COMPLEMENTAR
    DINSMORE, Paul C.; CABANIS-BREWIN, Jeannette. AMA: manual de
gerenciamento de projetos. Rio de Janeiro, RJ: Brasport, 2009. 498p. ISBN
9788574523237 (broch.).
\newline \newline 
INSTITUTE, Project M. Um guia de conhecimento em gerenciamento de projetos (guia PMBOK®). [Digite o Local da Editora]: Editora Saraiva, 2014. E-book. ISBN 9788502223745. Disponível em:
https://app.minhabiblioteca.com.br/\#/books/9788502223745/. Acesso em: 27 mar. 2023.
\newline \newline 
PHILLIPS, Joseph. Gerência de projetos de tecnologia da informação: no
caminho certo, do início ao fim. Rio de Janeiro, RJ: Elsevier, 2003. 449 p. ISBN 9788535211832 (broch.).
\newline \newline 
VAZQUEZ, Carlos Eduardo; SIMÕES, Guilherme Siqueira.; ALBERT, Renato
Machado. Análise de pontos de função: medição, estimativas e gerenciamento de projetos de software . 10. ed.rev. e ampl. São Paulo, SP: Érica, 2013. 272p. ISBN 9788536504520 (broch.).
\newline \newline 
DINSMORE, Paul Campbell; BARBOSA, Adriane Monteiro Cavalieri. Como se tornar um profissional em gerenciamento de projetos: livro-base de 'Preparação para certificação PMP® - Project management professional'. 4. ed., rev. e ampl. Rio de Janeiro: Qualitimark, 2011. 383p. ISBN 9788573039788 (Broch.).
\newline \newline 
COHN, Mike; SILVA, Aldir José Coelho da. Desenvolvimento de software com scrum: aplicando métodos ágeis com sucesso . Porto Alegre: Bookman, 2011. 496 p. ISBN 9788577808076 (broch.).
}


\DetalhesDisciplina{
    QXD0035 - Inglês Instrumental I
}{% EMENTA
    Vocabulário ligado à Informática. Técnicas de Leitura e Compreensão.
Tradutores. Verbos e tempos verbais. Nomes e Pronomes.
}{% BIBLIOGRAFIA BÁSICA
    GALLO, L. R. Inglês instrumental para informática .Ícone Ed.2008.
\newline \newline 
THE OFFICIAL GUIDE TO the TOEFL test. 4th ed. New York: McGraw-Hill,
2012. 653 p. + 1 CD-ROM ISBN 9780071766586 (broch.).
\newline \newline 
MURPHY, Raymond. English grammar in use: a self-study reference and
practice book for intermediate students. 3rd. ed. Cambridge: Cambridge
University Press, 2007. 379 p. ISBN 052143680X.
\newline \newline 
SWAM, M. Practical English Usage. 3 ed. Oxford do Brasil, 2005. ISBN:
0194420981.
}{% BIBLIOGRAFIA COMPLEMENTAR
    COLLINS. COLLINS: dicionário escolar. Martins Fontes, 2009.
EFRAIM, T. Decision support and business. 9 ed. Prentice Hall, 2010.
\newline \newline 
LONGMAN gramática escolar da língua inglesa: gramática de referências com
exercícios e respostas . São Paulo, SP: Longman, 2004. 317 p. : ISBN
8587214470
\newline \newline 
MACMILLAN EDUCATION (EDITORA). MacMillan English Dictionary for
advanced learners of American English. Macmillan Education. ISBN
9780333966709
\newline \newline 
MARINOTTO, D. Reading on info tech: inglês para informática. 2 ed. Novatec,
2007.
}


\DetalhesDisciplina{
    QXD0036 - Inglês Instrumental II
}{% EMENTA
    Conversação, leitura de artigos e jornais da área. Escrita de trabalhos técnicos.
Apresentação de seminários. Noções avançadas de gramática e compreensão
de texto.
}{% BIBLIOGRAFIA BÁSICA
    GALLO, L. R. Inglês instrumental para informática .Ícone Ed.2008.
\newline \newline 
THE OFFICIAL GUIDE TO the TOEFL test. 4th ed. New York: McGraw-Hill,
2012. 653 p. + 1 CD-ROM ISBN 9780071766586 (broch.).
\newline \newline 
MARINOTTO, D. Reading on info tech: inglês para informática. 2 ed. Novatec,
2007.
}{% BIBLIOGRAFIA COMPLEMENTAR
    COLLINS. COLLINS: dicionário escolar. Martins Fontes, 2009.
\newline \newline 
EFRAIM, T. Decision support and business. 9 ed. Prentice Hall, 2010.
\newline \newline 
LONGMAN gramática escolar da língua inglesa: gramática de referências com
exercícios e respostas . São Paulo, SP: Longman, 2004. 317 p. : ISBN
8587214470
\newline \newline 
MACMILLAN English dictionary for advanced learners. 2 ed. Macmillan
Education, 2007. ISBN: 9780230025455
\newline \newline 
MURPHY, Raymond. English grammar in use: a self-study reference and
practice book for intermediate students. 3rd. ed. Cambridge: Cambridge
University Press, 2007. 379 p. ISBN 052143680X.
}


\DetalhesDisciplina{
    <22> - Inteligência Artificial Explicável
}{% EMENTA
    Necessidade de métodos de IA explicável e seus desafios. Interpretabilidade e explanabilidade. Modelos interpretávels vs. modelos caixa-preta. Explicabilidade pré-modelo. Técnicas de visualização de modelos. Métodos Agnósticos.
}{% BIBLIOGRAFIA BÁSICA
    RUSSELL, Stuart J.; NORVIG, Peter. Inteligência artificial. Rio de Janeiro: Elsevier, Campus, 2013. 988 p. ISBN 9788535237016 (broch.).
\newline \newline 
MOLNAR, Christoph. Interpretable machine learning: a guide for making black box models explainable. 3. ed. [S.l.: s.n.], 2025. Disponível em: https://christophm.github.io/interpretable-ml-book/. Acesso em: 23 maio 2025.
\newline \newline 
COPPIN, Ben. Inteligencia artificial. Rio de Janeiro, RJ: LTC, 2010. 636 p. ISBN 9788521617297.
}{% BIBLIOGRAFIA COMPLEMENTAR
    HOLZINGER, Andreas; GOEBEL, Randy; FONG, Ruth; MOON, Taesup; MÜLLER, Klaus-Robert; SAMEK, Wojciech (Orgs.). XXAI – Beyond Explainable AI: International Workshop, Held in Conjunction with ICML 2020, July 18, 2020, Vienna, Austria, Revised and Extended Papers. Open access. Cham: Springer, ©2022.
\newline \newline 
GENOVESI, Sergio; KAESLING, Katharina; ROBBINS, Scott (Orgs.). Recommender Systems: Legal and Ethical Issues. Open access. Cham: Springer, ©2023. Disponível em: https://link.springer.com/book/10.1007/978-3-031-28241-6. Acesso em: 6 jun. 2025.
\newline \newline 
HORA, Nina da. MyNews Explica Algoritmos. São Paulo: Edições 70, 2023. E-book. p.59. ISBN 9786554271943. Disponível em: https://app.minhabiblioteca.com.br/reader/books/9786554271943/. Acesso em: 06 jun. 2025.
\newline \newline 
BLACKMAN, Reid. Máquinas Éticas: seu guia conciso para uma IA totalmente imparcial, transparente e respeitosa. Rio de Janeiro: Editora Alta Books, 2024. E-book. p.1. ISBN 9788550822426. Disponível em: https://app.minhabiblioteca.com.br/reader/books/9788550822426/. Acesso em: 06 jun. 2025.
\newline \newline 
FEFERBAUM, Marina; SILVA, Alexandre Pacheco da; COELHO, Alexandre Z.; et al. Ética, Governança e Inteligência Artificial. São Paulo: Almedina, 2023. E-book. p.48. ISBN 9786556279145. Disponível em: https://app.minhabiblioteca.com.br/reader/books/9786556279145/. Acesso em: 06 jun. 2025.
}


\DetalhesDisciplina{
    <13> - Inteligência Artificial Generativa
}{% EMENTA
    Fundamentos, arquiteturas e aplicações da Inteligência Artificial Generativa. Redes adversariais generativas (GANs), Autoencoders Variacionais (VAEs) e modelos de difusão. Arquiteturas transformer e modelos de linguagem de grande escala (LLMs). Técnicas de prompt engineering e fine-tuning. Métricas de avaliação de modelos generativos. Implicações éticas e sociais da IA generativa. Criação e avaliação de modelos generativos. Aplicações em geração de imagens, texto, áudio e vídeo. 
}{% BIBLIOGRAFIA BÁSICA
    FEUERRIEGEL, Stefan; HARTMANN, Jonas; JANIESCH, Christian; et al. Generative AI. Business \& Information Systems Engineering, Heidelberg, v. 66, p. 111–126, fev. 2024. Disponível em: https://doi.org/10.1007/s12599-023-00834-7. 
\newline \newline 
MANDUCHI, Laura et al. On the Challenges and Opportunities in Generative AI. 2025. Pré-publicação (preprint) disponível em: arXiv:2403.00025 [cs.LG]. Disponível em: https://arxiv.org/abs/2403.00025.
\newline \newline 
RUSSELL, Stuart J.; NORVIG, Peter. Inteligência Artificial: Uma Abordagem Moderna. 4. ed. Rio de Janeiro: GEN LTC, 2022. E-book. p.Capa. ISBN 9788595159495. Disponível em: https://app.minhabiblioteca.com.br/reader/books/9788595159495/.
}{% BIBLIOGRAFIA COMPLEMENTAR
    HARIOM, Tatsat,; SAHIL, Puri,; BRAD, Lookabaugh,. Blueprints de aprendizado de máquina e ciência de dados para finanças: desenvolvendo desde estratégias de trades até robôs Advisors com Python. Rio de Janeiro: Editora Alta Books, 2024. E-book. p.i. ISBN 9788550821726. Disponível em: https://app.minhabiblioteca.com.br/reader/books/9788550821726/.
\newline \newline 
GUZDIAL, Matthew; SNODGRASS, Sam; SUMMERVILLE, Adam. Generative AI. In: Procedural Content Generation via Machine Learning: An Overview. Cham: Springer Nature Switzerland, 2025. p. 199–218. Disponível em: https://doi.org/10.1007/978-3-031-84756-1\_11.
\newline \newline 
FERRARA, Emilio. GenAI against humanity: nefarious applications of generative artificial intelligence and large language models. Journal of Computational Social Science, v. 7, n. 1, p. 549–569, 1 abr. 2024. DOI: 10.1007/s42001-024-00250-1. Disponível em: https://doi.org/10.1007/s42001-024-00250-1. 
}


\DetalhesDisciplina{
    QXD0256 - Interação Humano-Computador
}{% EMENTA
    Os conceitos de interação e interface humano-computador; Estilos e paradigmas de interação: interfaces gráficas, manipulação direta, ícones e
linguagens visuais. Teorias de IHC: Engenharia cognitiva e Engenharia
semiótica de sistemas interativos; Sistemas de Ajuda e Sistemas de Explicação; Design de Interação: modelagem de interfaces e concretização
do projeto de interface (prototipação de interfaces, ferramentas de apoio
à construção de interfaces); Avaliação de sistemas interativos: métodos
de inspeção, métodos empíricos, testes com usuários, aspectos éticos na
relação com os usuários; Acessibilidade: conceitos, Lei Nacional de Acessibilidade, recomendações W3C para um site acessível, ferramentas de apoio
ao design de sistemas acessíveis, avaliação de acessibilidade.
}{% BIBLIOGRAFIA BÁSICA
    BARBOSA, Simone D. J.; SILVA, Bruno Santana da. Interação humanocomputador. Rio de Janeiro: Elsevier, c2010. 384 p. (Série SBC, Sociedade
\newline \newline 
Brasileira de Computação). ISBN 9788535234183.
\newline \newline 
PREECE, J.; ROGERS,Y. Design de interação: além da interação homemcomputador. Porto Alegre: Bookman, 2005.ISBN: 9788536304946
\newline \newline 
BARRETO, Jeanine dos S.; JR., Paulo A P.; BARBOZA, Fabrício F M.; et al. Interface humano-computador. Grupo A, 2018. E-book. ISBN
9788595027374. Disponível em: https://app.minhabiblioteca.com.
br/\#/books/9788595027374. Acesso em: 28 set. 2022.
}{% BIBLIOGRAFIA COMPLEMENTAR
    BENYON, David. Interação humano-computador. 2. ed. São Paulo:
Pearson Prentice Hall, 2011. 442 p. ISBN 9788579361098
\newline \newline 
SANTA ROSA, José Guilherme; MORAES, Anamaria de. Design participativo: técnicas para inclusão de usuários no processo de ergodesign de
interfaces. Rio de Janeiro: Rio Book’s, 2012. 170 p. ISBN 9788561556167
\newline \newline 
CYBIS, Walter de Abreu; BETIOL, Adriana Holtz.; FAUST, Richard.
Ergonomia e usabilidade: conhecimentos, métodos e aplicações. 3. ed.
São Paulo: Novatec, 2015. 496 p. ISBN: 9788575224595
\newline \newline 
SOBRAL, Wilma S. DESIGN DE INTERFACES - INTRODUÇÃO. Editora Saraiva, 2019. E-book. ISBN 9788536532073. Disponível em: https:
//app.minhabiblioteca.com.br/\#/books/9788536532073. Acesso em:
28 set. 2022.
\newline \newline 
BATISTA, Claudia R.; ULBRICHT, Vania R.; FADEL, Luciane M. Design para acessibilidade e inclusão. Editora Blucher, 2017. E-book. ISBN
9788580393040. Disponível em: https://app.minhabiblioteca.com.
br/\#/books/9788580393040/. Acesso em: 28 set. 2022.
}

\DetalhesDisciplina{
    <02> - Introdução à Inteligência Artificial
}{% EMENTA
    Introdução aos conceitos básicos de Inteligência Artificial (IA): Fundamentos, História da IA, Estado da Arte, Riscos e Benefícios da IA. Visão geral sobre as principais áreas da IA: representação do conhecimento, planejamento automatizado, aprendizado de máquina, comunicação, percepção e atuação. Contextualização acerca da profissão, curso e interdisciplinaridade.
}{% BIBLIOGRAFIA BÁSICA
    RUSSELL, Stuart J.; NORVIG, Peter. Inteligência Artificial: Uma Abordagem Moderna. 4. ed. Rio de Janeiro: GEN LTC, 2022. E-book. p.Capa. ISBN 9788595159495. 
\newline \newline 
COPPIN, B Inteligência artificial. LTC, 2010. ISBN: 9788521617297
\newline \newline 
SHOHAM, Y. Multiagent systems: algorithms, game theoretic. Cambridge University, 2009. ISBN: 9780521899437
}{% BIBLIOGRAFIA COMPLEMENTAR
    KAUFMAN, Dora. Desmistificando a inteligência artificial. São Paulo: Autêntica Editora, 2022. E-book. p.1. ISBN 9786559281596. Disponível em: https://app.minhabiblioteca.com.br/reader/books/9786559281596/. Acesso em: 30 mai. 2025.
\newline \newline 
SANTAELLA, Lucia. A inteligência artificial é inteligente?. São Paulo: Edições 70, 2023. E-book. p.1. ISBN 9786554270588. Disponível em: https://app.minhabiblioteca.com.br/reader/books/9786554270588/. Acesso em: 22 mai. 2025.
\newline \newline 
LIMA, Isaías. Inteligência Artificial. Rio de Janeiro: GEN LTC, 2014. E-book. p.Capa. ISBN 9788595152724. Disponível em: https://app.minhabiblioteca.com.br/reader/books/9788595152724/. Acesso em: 22 mai. 2025.
\newline \newline 
SILVA, Fabrício M.; LENZ, Maikon L.; FREITAS, Pedro H C.; et al. Inteligência artificial. Porto Alegre: SAGAH, 2018. E-book. p.1. ISBN 9788595029392. Disponível em: 
https://app.minhabiblioteca.com.br/reader/books/9788595029392/pageid/0
\newline \newline 
BRACHMAN, R ; LEVESQUE, Hector. Knowledge representation and reasoning. MorganKaufmann, 2004. 
}


\DetalhesDisciplina{
    QXD0113 - Língua Brasileira de Sinais - LIBRAS
}{% EMENTA
    Fundamentos histórico culturais da Libras e suas relações com a educação
dos surdos. Parâmetros e traços linguísticos da Libras. Cultura e identidades surdas. Alfabeto datilológico. Expressões não-manuais. Uso do espaço.
Classificadores. Vocabulário da Libras em contextos diversos. Diálogos em
língua de sinais.
}{% BIBLIOGRAFIA BÁSICA
    QUADROS, Ronice M. de; KARNOPP, Lodenir B.. Língua de sinais brasileira.. São Paulo: ArtMed, 2003. E-book. Disponível em: https://integrada.minhabiblioteca.com.br/\#/books/9788536311746. Acesso
em: 3 de out. 2022.
\newline \newline 
CORRÊA, Ygor; CRUZ, Carina R. Língua Brasileira de Sinais e Tecnologias Digitais. [Digite o Local da Editora]: Grupo A, 2019. E-book.
ISBN 9788584291687. Disponível em: https://app.minhabiblioteca.
com.br/\#/books/9788584291687. Acesso em: 23 set. 2022
\newline \newline 
QUADROS, Ronice Müller de. Educação de surdos: a aquisição da linguagem.. São Paulo: ArtMed, 1997. E-book. Disponível em: https://
integrada.minhabiblioteca.com.br/\#/books/9788536316581. Acesso
em: 3 de out. 2022.
}{% BIBLIOGRAFIA COMPLEMENTAR
    GOES, Maria Cecília Rafael; SMOLKA, Ana Luiza B. A linguagem e o
outro no espaço escolar: Vygotsky e a construção do conhecimento. Campinas: Papirus, 2013.
\newline \newline 
GOLDFELD, Marcia. A Criança Surda: linguagem e cognição numa perspectiva sócio-interacionista. São Paulo: Plexus, 2002.
\newline \newline 
SOUZA, Margarida M. P. Voando com Gaivotas: um estudo das interações
na educação de surdos. Dissertação (Mestrado em Educação Brasileira).
Faculdade de Educação, UFC. 2008. 152 p. [online]
\newline \newline 
COSTA, Márcia Cunha Silva. Educação inclusiva e prática docente : tenho
um aluno surdo em minha sala. E agora? 2013. Dissertação (mestrado em
Educação) -  Universidade Federal do Ceará, Faculdade de Educação, Programa de Pós-Graduação em Educação Brasileira, Fortaleza, 2013. [online]
\newline \newline 
PINHEIRO, Kátia Lucy; LEITAO, Vanda Magalhaes. Práticas pedagógicas bilíngues para crianças do Instituto Cearense de Educação de Surdos.
2012. 164f. Dissertação (mestrado) -  Universidade Federal do Ceará, Programa de Pós-Graduação em Educação Brasileira , Fortaleza, 2012. [online]
}


\DetalhesDisciplina{
    QXD0016 - Linguagens de Programação
}{% EMENTA
    Conceitos básicos de LP: domínios de aplicação, influências no projeto,
paradigmas, métodos de implementação, critérios de avaliação, evolução das
linguagens. Análise léxica e sintática. Variáveis: identificadores, vinculações,
verificação de tipos, escopo. Tipos de dados. Expressões e a declaração de
atribuição. Abstração de processos: subprogramas. Abstração de dados e
orientação a objetos. Noções de programação funcional. Noções de
programação lógica.
}{% BIBLIOGRAFIA BÁSICA
    SEBESTA, R.W. Conceitos de linguagens de programção. 9 ed. Bookman,
2011. ISBN: 9788577807918
\newline \newline 
TUCKER, A. B.; NOONAN, R. Linguagens de programação: princípios e
paradigmas. 2 ed. McGrawHill, 2008. ISBN: 9788577260447
\newline \newline 
CHEN, Yinong; TSAI, Wei-Tek. Introduction to programming languages:
proggramming in C, C++,
Scheme, Prolog, C\#, and SOA. 2nd ed. xii, 383 p.
\newline \newline 
WATT, D.A. Programming language: concepts and paradigms. Prentice Hall,
1990.
}{% BIBLIOGRAFIA COMPLEMENTAR
    DEITEL, H. M. C++ como programar. 5 ed. Prentice Hall, 2006. ISBN:
8576050560
\newline \newline 
AHO, A. V.; SETHI, R.; ULLMAN, J. D. Compiladores: princípios, técnicas e
ferramentas. 2 ed. Pearson/Addison-Wesley, 2008. ISBN: 9788588639249
\newline \newline 
HOPCROFT, J. E.; ULLMAN, J. D.; MOTWANI, R. Introdução à teoria dos
autômatos: linguagens e computação. ​Campus,​ 2002. ISBN:9788535210729
\newline \newline 
BARWISE, J. Language, proof and logic. Seven Bridges, 2002. ISBN:
9781575863740
\newline \newline 
MENEZES, P.B. Linguagens formais e autômatos. 5 ed. Sagra Luzzato, 2008.
ISBN: 9788577807659
\newline \newline 
BARNES, D.J.; KOLLING, M. Programação orientada a objetos com Java: uma
introdução prática usando BLUEJ 4 ed. Prentice Hall Brasil, 2009. ISBN:
9788576051879
\newline \newline 
URUBATAN, R. Ruby on rails: desenvolvimento fácil e rápido. Novatec, 2009.
}


\DetalhesDisciplina{
    QXD0040 - Linguagens Formais e Autômatos
}{% EMENTA
    Introdução. Linguagens, gramáticas e expressões regulares, autômatos finitos. Linguagens e gramáticas livre-do contexto e autômatos de pilha.
Linguagens sensíveis ao contexto. Hierarquia de classes de linguagens. Tópicos especiais e aplicações das linguagens formais e autômatos.
}{% BIBLIOGRAFIA BÁSICA
    HOPCROFT, J. E.; ULLMAN, J. D.; MOTWANI, R., Introdução à teoria
dos autômatos: linguagens e computação, 1 ed, 2002, 10a tiragem Campus
\newline \newline 
MENEZES, P.B. Linguagens formais e autômatos. 5 ed. Sagra Luzzato, 2008. ISBN: 9788577807659
\newline \newline 
RAMOS, M. V.; NETO, J.J.; VEGA, I.S. Linguagens Formais: teoria,
modelagem e implementação. Bookman, 2009. ISBN: 9788577804535
}{% BIBLIOGRAFIA COMPLEMENTAR
    AHO, A. V.; SETHI, R.; ULLMAN, J. D. Compiladores: princípios,
técnicas e ferramentas. 2 ed. Pearson/Addison-Wesley, 2008. ISBN:
9788588639249
\newline \newline 
CARNIELLI, W.; EPSTEIN, R. L. Computabilidade, Funções Computáveis, Lógica e os Fundamentos da Matemática. UNESP, 2009. ISBN:
9788571398979
\newline \newline 
SEBESTA, R.W. Conceitos de linguagens de programação. 5 ed. Bookman, 2003. ISBN: 9788577807918
\newline \newline 
SIPSER, M. Introdução a teoria da computação. 2 ed. Thompson Learning, 2007. ISBN: 9788522104994
\newline \newline 
TUCKER, A. B.; NOONAN, R. Linguagens de programação: princípios e
paradigmas. 2 ed. McGrawHill, 2008. ISBN: 9788577260447
DIVERIO, Tiarajú Asmuz. Teoria da computação: máquinas universais e
computabilidade. 3. ed. Porto Alegre: Bookman, 2011. 288 p. (Livros
didáticos. n.5)
}


\DetalhesDisciplina{
    QXD0017 - Lógica para Computação
}{% EMENTA
    Lógica proposicional e de Primeira Ordem; Formalização de problemas;
Sistemas dedutivos: axiomático, natural e tableaux; Correção e Completude. Lógicas Temporais para a Validação de Sistemas.
}{% BIBLIOGRAFIA BÁSICA
    SOUZA, João Nunes de. Lógica para ciência da computação: fundamentos
de linguagem, semântica e sistemas de dedução. Rio de Janeiro: Elsevier,
2002. 309 p. ISBN 8535210938 (broch.).
\newline \newline 
SILVA, Flávio Soares Corrêa da; FINGER, Marcelo; MELO, Ana Cristina
Vieira de. Lógica para computação. São Paulo, SP: Thomson Learning,
2006. 234 p. ISBN 8522105170 (broch.).
\newline \newline 
BISPO, Carlos Alberto F.; CASTANHEIRA, Luiz B.; FILHO, Oswaldo
Melo S. Introdução à Lógica Matemática. [Digite o Local da Editora]:
Cengage Learning Brasil, 2017. E-book. ISBN 9788522115952. Disponível
em: https://app.minhabiblioteca.com.br/\#/books/9788522115952.
Acesso em: 07 out. 2022.
}{% BIBLIOGRAFIA COMPLEMENTAR
    ENDERTON, Herbert B. A mathematical introduction to logic. 2nd ed.
San Diego, California: Harcourt/Academic Press, c2001. 317 p. ; ISBN 0122384520 (enc.).
\newline \newline 
GERSTING, Judith L. Fundamentos matemáticos para a ciência da computação: um tratamento moderno de matemática discreta . 5. ed. Rio
de Janeiro: Livros Técnicos e Científicos, 2004. 597 p. ISBN 8521614225
(broch.).
\newline \newline 
CLARKE, E. M. Model checking. Cambridge: MIT Press, 1999. 314 p.
ISBN 9780262032704 (enc.).
\newline \newline 
HUTH, Michael; RYAN, Mark. . Lógica em ciência da computação: modelagem e argumentação sobre sistemas . 2. ed. Rio de Janeiro: LTC,
2008. 322 p. ISBN 9788521616108 (broch.).
\newline \newline 
ALENCAR FILHO, Edgard de. Iniciação à lógica matemática. São Paulo:
Nobel, [2002]. 203 p. ISBN 852130403X (broch).
}


\DetalhesDisciplina{
    QXD0056 - Matemática Básica
}{% EMENTA
    Lógica: conectivos lógicos, tabela verdade, fórmulas equivalentes. Conjuntos:
notação, operações, propriedades das operações, diagramas de Venn,
partição, cardinalidade, conjuntos das partes, produto cartesiano. Contagem:
princípios da multiplicação e da adição, princípio da exclusão, princípios das
casas dos pombos, permutações, combinações, teorema binominal, triângulo
de pascal. Relações: definições, terminologia, propriedades. Funções:
definições, terminologia, propriedades.
}{% BIBLIOGRAFIA BÁSICA
    GERSTING, Judith L. Fundamentos matemáticos para a ciência da
computação: um tratamento moderno de matemática discreta . 5. ed. Rio de
Janeiro: Livros Técnicos e Científicos, 2004. 597 p. ISBN 8521614225 (broch.).
\newline \newline 
SILVA, Sebastiao Medeiros da; SILVA, Elio Medeiros da; SILVA, Ermes
Medeiros da. Matemática básica para cursos superiores. São Paulo: Atlas,
2002. 227 p. ISBN 8522430357 (broch.).
\newline \newline 
IEZZI, Gelson; MURAKAMI, Carlos. Fundamentos de matemática elementar: 1:
conjuntos, funções. 410 p. 9.ed. ISBN: 9788535716801
}{% BIBLIOGRAFIA COMPLEMENTAR
    ROSEN, Kenneth H. Matemática discreta e suas aplicações. 6. ed. São Paulo,
SP: McGraw-Hill, 2009. xxi, 982 p. ISBN 9788577260362 (broch.).
\newline \newline 
ALENCAR FILHO, Edgard de. Iniciação a lógica matemática. São Paulo:
Nobel, [2002] 203p ISBN 852130403X (broch).
\newline \newline 
MENEZES, Paulo Blauth; UNIVERSIDADE FEDERAL DO RIO GRANDE DO
SUL. Matemática discreta para computação e informática. 3. ed. Porto Alegre,
RS: Bookman, 2010. 350 p (Livros didáticos. 16).
\newline \newline 
DEMANA, Franklin D. Pré-cálculo. São Paulo: Addison-Wesley, 2009. 380 p.
ISBN 9788588639379 (broch.).
\newline \newline 
MENEZES, Paulo Blauth; TOSCANI, Laira V.; GARCÍA LÓPEZ, Javier.
Aprendendo matemática discreta com exercícios. Porto Alegre, RS: Bookman,
2009. 356p. (Livros didáticos informática ufrgs; v. 19) ISBN 9788577804719
(broch.).
\newline \newline 
SCHEINERMAN, Edward R. Matemática discreta: uma introdução . São Paulo:
Cengage Learning, 2011. 573 p. ISBN 9788522107964 (broch.).
}


\DetalhesDisciplina{
    QXD0120 - Matemática Computacional
}{% EMENTA
    Noções de modelagem matemática de problemas. Noções de métodos
numéricos: erros de representação em ponto flutuante e perda de significância;
raízes de equações; integração numérica; interpolação; solução de sistemas
de equações lineares; minimização de funções. Programação linear: forma
padrão e dualidade; método simplex. Uso de pacotes computacionais.
}{% BIBLIOGRAFIA BÁSICA
    BARROSO, Leônidas Conceição et al. Cálculo numérico: (com aplicações). 2.
ed. São Paulo, SP: Harbra, c1987. 367 p. ISBN 8529400895 (broch.).
\newline \newline 
GOLDBARG, Marco Cesar. Otimização combinatória e programação linear:
modelos e algoritmos. Rio de Janeiro: Elsevier: ​Campus,​ 2005. xvi, 518 p. :
ISBN 9788535215205 (broch.)
\newline \newline 
HILLIER, FREDERICK S.; LIEBERMAN, GERALD J. Introdução a Pesquisa
Operacional. MCGRAW HILL. 9a edição. (ISBN: 8580551188)
}{% BIBLIOGRAFIA COMPLEMENTAR
    RUGGIERO, Marcia A. Gomes; LOPES, Vera Lucia da Rocha. Cálculo
numérico: aspectos teóricos e computacionais . 2. ed. Pearson, c1997. ISBN
8534602042.
\newline \newline 
PASSOS, Eduardo José Pedreira Franco dos. Programação linear como
instrumento da pesquisa operaciona: Eduardo José Pedreira Franco dos
Passos. São Paulo, SP: Atlas, 2008. xii, 451p. ISBN 9788522448395 (broch.).
\newline \newline 
GERSTING, Judith L. Fundamentos matemáticos para a ciência da
computação: um tratamento moderno de matemática discreta . 5. ed. Rio de
Janeiro: Livros Técnicos e Científicos, c2004. xiv, 597 p. ISBN: 8521614225.
\newline \newline 
CHENEY, Ward; KINCAID, David (David Ronald). Numerical mathematics and
computing. 3rd. ed. Pacific Grove, CA: Books/Cole, c1994. 578p. ISBN
0534201121
\newline \newline 
SPERANDIO, Décio; MENDES, João Teixeira; SILVA, Luiz Henry Monken e.
Cálculo numérico: características matemáticas e computacionais dos métodos
numéricos. São Paulo, SP: Prentice Hall, 2003. ix, 354 p. ISBN 8587918745
(broch.).
}


\DetalhesDisciplina{
    QXD0008 - Matemática Discreta
}{% EMENTA
    Técnicas de demonstração: exaustiva, direta, contraposição, absurdo, indução
(fraca e forte). Somatórios: notação, propriedades, séries aritméticas,
geométricas e harmônicas, algumas fórmulas de somatórios úteis. Teoria dos
números: divisibilidade, primos, teorema fundamental da aritmética, aritmética
modular, aplicações. Relações: fechos, ordem parcial e total, relações e
classes de equivalências. Grafos: terminologia, alguns grafos especiais,
isomorfismo, conectividade, árvores (definição e propriedades).
}{% BIBLIOGRAFIA BÁSICA
    GERSTING, Judith L. Fundamentos matemáticos para a ciência da
computação: um tratamento moderno de matemática discreta . 5. ed. Rio de
Janeiro: Livros Técnicos e Científicos, 2004. 597 p. ISBN 8521614225 (broch.).
\newline \newline 
MENEZES, Paulo Blauth; UNIVERSIDADE FEDERAL DO RIO GRANDE DO
SUL. Matemática discreta para computação e informática. 3. ed. Porto Alegre,
RS: Bookman, 2010. 350 p (Livros didáticos. 16).
\newline \newline 
ROSEN, Kenneth H. Matemática discreta e suas aplicações. 6. ed. São Paulo,
SP: McGraw-Hill, 2009. 982 p. ISBN 9788577260362 (broch.).
}{% BIBLIOGRAFIA COMPLEMENTAR
    ALENCAR FILHO, E. Iniciação à lógica matemática. 21. ed. São Paulo: Nobel,
2008. ISBN:9788521304036.
\newline \newline 
SILVA, Sebastiao Medeiros da; SILVA, Elio Medeiros da; SILVA, Ermes
Medeiros da. Matemática básica para cursos superiores. São Paulo: Atlas,
2002. 227 p. ISBN 8522430357 (broch.).
\newline \newline 
MENEZES, Paulo Blauth; TOSCANI, Laira V.; GARCÍA LÓPEZ, Javier.
Aprendendo matemática discreta com exercícios. Porto Alegre, RS: Bookman,
2009. 356p. (Livros didáticos informática ufrgs ; ; v. 19) ISBN 9788577804719
(broch.).
\newline \newline 
SCHEINERMAN, Edward R. Matemática discreta: uma introdução . São Paulo:
Cengage Learning, 2011. 573 p. ISBN 9788522107964 (broch.).
\newline \newline 
HUNTER, David J. Fundamentos da matemática discreta. Rio de Janeiro, RJ:
LTC, 2011. 235 p. ISBN 9788521618102 (broch.)
}


\DetalhesDisciplina{
    QXD0178 - Mineração de Dados
}{% EMENTA
    O que é ​data mining.​ Aplicações potenciais. O processo de Descoberta do
Conhecimento. ​Data mining, data warehouse e OLAP. Tarefas de mineração
de dados: classificação, agrupamento (​clustering​), regras de associação e
análise de desvios. Estudo de algoritmos para as principais tarefas de
mineração de dados. Avaliação dos resultados obtidos. Introdução às técnicas
de recuperação de informações e aplicações em mineração de textos e ​Web
mining.​
}{% BIBLIOGRAFIA BÁSICA
    TAN, Pang-Ning; Steinbach, Michael; Kumar, Vipin. Introdução ao Data Mining
- Mineração de Dados . Ciência Moderna, 2009. 900 p. ISBN-10 8573937610
ISBN-13 9788573937619.
\newline \newline 
PINHEIRO, Carlos A. R. Inteligência analítica: mineração de dados e
descoberta de conhecimento. Ciência Moderna, c2008. 397 p. ISBN
9788573937077 (broch.)
\newline \newline 
ELMASRI, Ramez; NAVATHE, Sham. Sistemas de banco de dados. 6. ed. -.
São Paulo, SP: Pearson Education do Brasil, 2011. xviii, 788 p. ISBN:
9788579360855
\newline \newline 
WITTEN, Ian H.; FRANK, Eibe; HALL, Mark A. Data Mining: Practical Machine
Learning Tools and Techniques. The Morgan Kaufmann Series in Data
Management Systems. Morgan Kaufmann, 3ª Ed., 2011. ISBN: 0123748569
}{% BIBLIOGRAFIA COMPLEMENTAR
    SILBERSCHATZ, Abraham; KORTH, Henry F.; Sudarshan, S. Sistema de
banco de dados. Rio de Janeiro: Elsevier: Campus, 2012. ISBN
9788535245356
\newline \newline 
RAMAKRISHNAN, Raghu; GEHRKE, Johannes. Sistemas de Gerenciamento
de banco de dados. São Paulo, SP: McGraw-Hill, 2008. xxvii, 884 p. ISBN:
9788577260270
\newline \newline 
DATE, C. J. Introdução a sistemas de banco de dados. 8. ed. Rio de Janeiro:
Elsevier, 2004. 865p ISBN: 8535212736
\newline \newline 
LESKOVEC, Jure; RAJARAMAN, Anand; ULLMAN, Jeffrey David. Mining of massive datasets [e-book]. 3. ed. Cambridge: Cambridge University Press, 2020. Disponível em: https://www.mmds.org/.
\newline \newline 
ZAKI, M. J.; MEIRA Jr., W. Data Mining and Analysis: Fundamental Concepts
and Algorithms. Cambridge 9780521766333. University Press, Disponível
May 2014. ISBN: em: http://www.dataminingbook.info/uploads/book.pdf
\newline \newline 
GRUS, Joel. Data Science from Scratch: First Principles with Python. O'Reilly
Media. 1st edition. 2015. ISBN 149190142X.
\newline \newline 
BISHOP, Christopher M. Pattern Recognition and Machine Learning. 1 ed.
Springer, 2007. ISBN: 0387310738
\newline \newline 
HAN, Jiawei; Kamber, Micheline; Pei, Jian. Data Mining: Concepts and
Techniques. The Morgan Kaufmann Series in Data Management Systems.
Morgan Kaufmann, 3ª Ed., 2011. ISBN: 0123814790
\newline \newline 
HASTIE, Trevor; TIBSHIRANI, Robert; Friedman, Jerome. The Elements of
Statistical Learning: Data Mining, Inference, and Prediction. Springer, 2ª Ed.,
2009. ISBN: 0387848576
}


\DetalhesDisciplina{
    <14> - MLOPs
}{% EMENTA
    Integração de práticas de engenharia de software, DevOps e aprendizado de máquina, abordando conceitos de Big Data e seus desafios. Ciclo de vida completo de projetos de aprendizado de máquina em produção, incluindo versionamento de dados e modelos. Reproducibilidade e governança de modelos. Arquiteturas e paradigmas para processamento distribuído. Técnicas de ingestão, armazenamento, processamento e visualização de dados em larga escala. Pipelines automatizados de CI/CD para aprendizado de máquina, automação. Ferramentas e ecossistemas para processamento de grandes volumes de dados. Containerização e orquestração. Computação em nuvem para Big Data. Ferramentas de deployment e plataformas de MLOps. Testes automatizados para modelos de aprendizado de máquina, monitoramento de modelos em produção.
}{% BIBLIOGRAFIA BÁSICA
    HUYEN, Chip. Projetando sistemas de Machine Learning: processo interativo para aplicações prontas para produção. Rio de Janeiro: Editora Alta Books, 2024. E-book. p.4. ISBN 9788550819648. Disponível em: https://app.minhabiblioteca.com.br/reader/books/9788550819648/.
\newline \newline 
SANTOS, Roger R.; BORDIN, Maycon V.; NUNES, Sergio E.; et al. Fundamentos de Big Data. Porto Alegre: SAGAH, 2021. E-book. p.Capa. ISBN 9786556901749. Disponível em: https://app.minhabiblioteca.com.br/reader/books/9786556901749/.
\newline \newline 
ERL, Thomas; MONROY, Eric B. Computação em Nuvem: Conceitos, Tecnologia, Segurança e Arquitetura. 2. ed. Porto Alegre: Bookman, 2024. E-book. p.i. ISBN 9788582606599. Disponível em: https://app.minhabiblioteca.com.br/reader/books/9788582606599/.
}{% BIBLIOGRAFIA COMPLEMENTAR
    KREUZBERGER, Dominik; KÜHL, Niklas; HIRSCHL, Sebastian. Machine Learning Operations (MLOps): overview, definition, and architecture. IEEE Access, v. 11, p. 31866–31879, 2023. DOI: 10.1109/ACCESS.2023.3262138. Disponível em: https://ieeexplore.ieee.org/document/10081336.
\newline \newline 
SYMEONIDIS, Georgios; NERANTZIS, Evangelos; KAZAKIS, Apostolos; PAPAKOSTAS, George A. MLOps – definitions, tools and challenges. In: IEEE 12th Annual Computing and Communication Workshop and Conference (CCWC), 2022, Las Vegas. Anais [...]. [S.l.]: IEEE, 2022. p. 453–460. DOI: 10.1109/CCWC54503.2022.9720902. Disponível em: https://ieeexplore.ieee.org/document/9720902.
\newline \newline 
TESTI, Matteo; BALLABIO, Michela; FRONTONI, Eleonora; IANNELLO, Guido; MOCCIA, Stefano; SODA, Pasquale; VESSIO, Giuseppe. MLOps: a taxonomy and a methodology. IEEE Access, v. 10, p. 63 606–63 618, 2022. DOI: 10.1109/ACCESS.2022.3181730. Disponível em: https://ieeexplore.ieee.org/document/9792270. 
}


\DetalhesDisciplina{
    <23> - Percepção e Ação Robótica
}{% EMENTA
    Tipos de sensores para robótica;
Percepção a partir de visão computacional;
Calibração de câmeras e visão estereoscópica;
Odometria visual;
Fusão sensorial e técnicas de localização robótica;
Localização e mapeamento simultâneo (SLAM);
Planejamento de rotas e navegação.
}{% BIBLIOGRAFIA BÁSICA
    NISE, Norman S.. Engenharia de Sistemas de Controle, 7a edição. São
Paulo: LTC, 2017. E-book. Disponível em: https://integrada.minhabiblioteca.
com.br/\#/books/9788521634379. Acesso em: 30 de maio de 2025.
\newline \newline 
Ben Coppin. 2004. Artificial Intelligence Illuminated. Jones and Bartlett Publishers, Inc., USA.
\newline \newline 
RUSSELL, Stuart J.; NORVIG, Peter. Inteligência artificial. Rio de Janeiro: Elsevier, Campus, 2013. 988 p. ISBN 9788535237016 (broch.).

}{% BIBLIOGRAFIA COMPLEMENTAR
    ROMERO, Roseli Aparecida F.; PRESTES, Edson; OSÓRIO, Fernando
et al. Robótica Móvel. São Paulo: LTC, 2014. E-book. Disponível em: https://integrada.minhabiblioteca.com.br/\#/books/978-85-216-2642-8. Acesso em: 30 de maio 2025.
\newline \newline 
MATARIC, Maja J.. Introdução á robótica. São Paulo: Editora Blucher, 2014. E-book. Disponível em: https://integrada.minhabiblioteca.com.br/\#/books/9788521208549. Acesso em: 30 de maio de 2025.
\newline \newline 
ANDRADE-CETTO, Juan; SANFELIU, Alberto. Environment learning for indoor mobile robots: a stochastic state estimation approach to simultaneous localization and map building.. Berlim; New York, NY: Springer, 2006. (Springer tracts in advanced robotics ; 23). Disponível em: https://link.springer.com/book/10.1007/11418382. Acesso em 30 de maio de 2025.
\newline \newline 
JAZAR, Reza N; SPRINGERLINK (ONLINE SERVICE). Theory of Applied Robotics : Kinematics, Dynamics, and Control (2nd Edition) . Springer eBooks 2nd. XXIII, 883p. 400 illus., 200 illus. in color ISBN 9781441917508. 
\newline \newline 
ZHANG, Dan; SPRINGERLINK (ONLINE SERVICE). Parallel Robotic Machine Tools. Springer eBooks ISBN 9781441911179.
}


\DetalhesDisciplina{
    <24> - Pesquisa e Inovação
}{% EMENTA
    Pesquisa em acervos físicos e virtuais: títulos, base de dados, periódicos, patentes, marcas, desenhos industriais. Método científico. Conceituação, tipos e metodologia de pesquisa. Elaboração de artigos técnico científicos, projetos de pesquisa e projetos de inovação.

Tipologias da inovação: industrial, em serviços, tecnológica, organizacional, aberta e fechada. Sistema de inovação, políticas de ciência, tecnologia e inovação (CT&I)
no Brasil e no mundo.
}{% BIBLIOGRAFIA BÁSICA
    WAZLAWICK, Raul Sidnei. Metodologia de pesquisa para ciência da computação. Rio de Janeiro : Elsevier.
2008.
\newline \newline 
DEZ tipos de inovação: Dez tipos de inovação** a disciplina de criação de avanços de ruptura. São Paulo: DVS Editora, 2015. 263 p. ISBN 978-85-8289-084-4.
\newline \newline 
VALLE, Rogério ; OLIVEIRA, Saulo Barbará de (org.). Análise e modelagem de processos de negócio:
foco na notação BPMN (Business Process Modeling Notation). São Paulo: Atlas, 2013. 207 p. ISBN 978-85-224-5621-5.
\newline \newline 

\newline \newline 


}{% BIBLIOGRAFIA COMPLEMENTAR
    A guide to the project management body of knowledge (PMBOK GUIDE)** A GUIDE to the project management body of knowledge (PMBOK GUIDE). 4rd. ed. Newtown Square, Pa: Project Management Institute, c2008. xxvi, 467 p. ISBN 9781933890517 (broch.).
\newline \newline 
FERREIRA, Valdinéia Barreto. E-science e políticas públicas para ciência, tecnologia e inovação no Brasil. Salvador: EDUFBA, 2018. 250 p. ISBN 978-85-2321707-5.
\newline \newline 
KALBACH, James. Mapeamento de experiências: um guia para criar valor por meio de jornadas, blueprints e diagramas. Rio de Janeiro: Alta Books, 2017. 362 p. ISBN 978-85-508-0061-5. Classificação: 658.4063 K21m (BCQ) Ac.217555
\newline \newline 
CHESBROUGH, Henry William. Modelos de negócios abertos: como prosperar no novo cenário da inovação. Porto Alegre: Bookman, 2012. xvi, 220 p. ISBN 978-85-7780-955-4.
\newline \newline 
EDLER, Jakob; WALZ, Rainer. Systems and Innovation Research in Transition: Research Questions and Trends in Historical Perspective. Springer Nature, 2024. DOI: https://doi.org/10.1007/978-3-031-66100-6 Acesso em: 3 de jun, 2025.
}


\DetalhesDisciplina{
    QXD0181 - Pesquisa Operacional
}{% EMENTA
    Introdução à Pesquisa Operacional e aos Sistemas de Apoio à Decisão.
Programação linear. Modelos de programação linear. Método simplex.
Problema do transporte. Dualidade. Técnicas avançadas em Pesquisa
Operacional.
}{% BIBLIOGRAFIA BÁSICA
    HILLIER, Frederick S.; LIEBERMAN, Gerald J.. Introdução à Pesquisa
Operacional. São Paulo: AMGH, 2013. E-book. Disponível em: https://
integrada.minhabiblioteca.com.br/\#/books/9788580551198. Acesso
em: 3 de out. 2022.
\newline \newline 
PASSOS, Eduardo J. P. F.. Programação linear como instrumento da
pesquisa operacional. Atlas, 2008. 451p. ISBN 9788522448395
\newline \newline 
COLIN, Emerson C.. Pesquisa Operacional -  170 Aplicações em Estratégia, Finanças, Logística, Produção, Marketing e Vendas, 2a edição. São Paulo: Atlas, 2017. E-book. Disponível em: https://integrada.
minhabiblioteca.com.br/\#/books/9788597014488. Acesso em: 3 de
out. 2022.
}{% BIBLIOGRAFIA COMPLEMENTAR
    ANDRADE, Eduardo Leopoldino de. Introdução à Pesquisa Operacional
- Método e Modelos para Análise de Decisões, 5a edição. São Paulo: LTC,
2015. E-book. Disponível em: https://integrada.minhabiblioteca.
com.br/\#/books/978-85-216-2967-2. Acesso em: 29 de set. 2022.
\newline \newline 
VIANA, Gerardo Valdisio Rodrigues. Meta-heurísticas e programação paralela em otimização combinatória. Fortaleza: Edições UFC, 1998. 250p.
\newline \newline 
ISBN 8572820396
CHAN, Alan H. S; AO, Sio-Iong SPRINGERLINK (ONLINE SERVICE).
Advances in Industrial Engineering and Operations Research. Springer
eBooks Boston, MA: Springer Science+Business Media, LLC, 2008. (Lecture Notes in Electrical Engineering ; 5) ISBN 9780387749051. Disponível
em: http://dx.doi.org/10.1007/978-0-387-74905-1. Acesso em : 21 set.
2010
\newline \newline 
QUDRAT-ULLAH, H; DAVIDSEN, P.I; SPECTOR, J.M SPRINGERLINK (ONLINE SERVICE). Complex Decision Making : Theory and
Practice . Springer e-books Berlin, Heidelberg: springer, 2008. (Understanding Complex Systems,) ISBN 9783540736653. Disponível em:
http://dx.doi.org/10.1007/978-3-540-73665-3. Acesso em: 21 set. 2010.
\newline \newline 
GOLDBARG, Marco Cesar. Otimização combinatória e programação linear: modelos e algoritmos. Rio de Janeiro: Elsevier: Campus, 2005. xvi,
518 p. : ISBN 9788535215205 (broch.).
}


\DetalhesDisciplina{
    <15> - Planejamento Automatizado
}{% EMENTA
    Modelos e linguagens de planejamento. (I) Planejamento Clássico: planejamento como satisfatibilidade (SATPLAN); grafo de planejamento (GRAPHPLAN); planejamento como busca heurística; construção de heurísticas independentes de domínio. (II) Variações e extensões de planejamento clássico: planejamento não-determinístico; busca num grafo AND-OR (LAO*); planejamento como verificação de modelos (baseado na lógica temporal CTL); planejamento contingente e conformante.
}{% BIBLIOGRAFIA BÁSICA
    RUSSELL, Stuart J.; NORVIG, Peter. Inteligência Artificial: Uma Abordagem Moderna. 4. ed. Rio de Janeiro: GEN LTC, 2022. E-book. p.Capa. ISBN 9788595159495. 
\newline \newline 
COPPIN, B Inteligência artificial. LTC, 2010. ISBN: 9788521617297
\newline \newline 
CORMEN, Thomas H.; LEISERSON, Charles E.; Ronald L. Rivest; et al. Algoritmos. 4. ed. Rio de Janeiro: GEN LTC, 2024. E-book. p.Capa. ISBN 9788595159914. 
}{% BIBLIOGRAFIA COMPLEMENTAR
    CLARKE, E. M. Model checking. Cambridge: MIT Press, 1999. 314 p. ISBN 9780262032704 (enc.).
\newline \newline 
BAIER, Christel; KATOEN, Joost-Pieter. Principles of model checking. Cambridge, Massachusetts: The Mit Press, 2008. xvii, 975 p. ISBN 9780262026499.
\newline \newline 
SILVA, Flávio Soares Corrêa da; FINGER, Marcelo; MELO, Ana Cristina Vieira de. Lógica para computação - 2ª edição. 2. ed. Porto Alegre: +A Educação - Cengage Learning Brasil, 2018. E-book. p.63. ISBN 9788522127191.
\newline \newline 
Ronald Brachman and Hector Levesque. 2004. Knowledge Representation and Reasoning. Morgan Kaufmann Publishers Inc., San Francisco, CA, USA.
\newline \newline 
Huth, Michael, and Mark Ryan. ""Lógica em ciência da computação: modelagem e argumentação sobre sistemas."" Rio de Janeiro, 2a edição edition. tradução e revisão técnica Valéria de Magalhães Iório (2008).
}

"\DetalhesDisciplina{
    QXD0109 - Pré-Cálculo
}{% EMENTA
    Funções no espaço contínuo: estudo de sinal, raízes, polinomiais, racionais, exponenciais, logarítmicas,trigonométricas.
}{% BIBLIOGRAFIA BÁSICA
    DEMANA, Franklin D. et al. Pré-cálculo . São Paulo, SP: Addison-Wesley, 2009. xv, 380 p. ISBN9788588639379 (broch.).
\newline \newline 
IEZZI, Gelson; MURAKAMI, Carlos. Fundamentos de matemática elementar: 1 : conjuntos, funções . 8. ed.,. SãoPaulo, SP: Atual, 2004. 374 p. ISBN 8570562705 (broch.) .
\newline \newline 
SILVA, Sebastiao Medeiros da; SILVA, Elio Medeiros da; SILVA, Ermes Medeiros da. Matemática básica paracursos superiores . São Paulo: Atlas, 2002. 227p.
}{% BIBLIOGRAFIA COMPLEMENTAR
    LEITHOLD, Louis. O Cálculo com geometria analítica. 3. ed. São Paulo: Harbra, c1994. 2 v. ISBN 8529400941v.1 (broch.)
\newline \newline 
IEZZI, Gelson. Fundamentos de matemática elementar, 6: complexos, polinômios, equações . 8. ed. São Paulo,SP: Atual, 2013. 250 p. ISBN 9788535717525 (broch.)
\newline \newline 
IEZZI, Gelson. Fundamentos de matemática elementar: 7 : geometria analítica. 5. ed. São Paulo, SP: Atual,2005. 282 p. ISBN 8535705465 (broch.)
\newline \newline 
IEZZI, Gelson. Fundamentos de matemática elementar, 3: trigonometria: 123 exercícios resolvidos, 385 exercícios propostos com resposta, 236 testes de vestibulares com resposta. 7. ed. São Paulo, SP: Atual, 1993.303 p. ISBN 8570562691
\newline \newline 
CALDEIRA, André Machado; SILVA, Luiza Maria Oliveira da,; MACHADO, Maria Augusta Soares. Pré-cálculo. 3.ed. rev. e ampl. São Paulo, SP: Cengage Learning; c2014. xv, 558p. ISBN 9788522116126 (broch.)
\newline \newline 
IEZZI, GELSON et al; Fundamentos de matemática elementar v.2: Logaritmos, 8 ed., 2004, Saraiva
}"

\DetalhesDisciplina{
    QXD0112 - Probabilidade e Estatística
}{% EMENTA
    Fundamentos de análise combinatória. Conceito de probabilidade e seus
teoremas fundamentais. Variáveis aleatórias. Distribuições de probabilidade.
Estatística descritiva. Noções de amostragem. Distribuições amostrais: discreta
e contínua. Inferência estatística: teoria da estimação, intervalos de confiança
e testes de hipóteses. Regressão linear simples. Correlação.
}{% BIBLIOGRAFIA BÁSICA
    LARSON, Ron; FARBER, Betsy. Estatística aplicada. 4. ed. São Paulo, SP:
Pearson/ Prentice Hall, 2010. xiv,637 p. ISBN 9788576053729 (broch.).
\newline \newline 
BARBETTA, Pedro Alberto; REIS, Marcelo Menezes; BORNIA, Antonio Cezar.
Estatística para cursos de engenharia e informática. 3. ed. São Paulo, SP:
Atlas, 2010. 410 p.
\newline \newline 
WALPOLE, Ronald E. Probabilidade e estatística: para engenharia e ciências.
8. ed. São Paulo, SP: Pearson/ Prentice Hall, 2009. xiv, 491 p. ISBN
9788576051992 (broch.).
\newline \newline 
MORETTIN, Luiz Gonzaga. Estatística básica:: probabilidade e inferência /
volume único. São Paulo, SP: Pearson Educational do Brasil. 2010. ISBN
8576053705 ISBN-13 9788576053705 (broch.).
}{% BIBLIOGRAFIA COMPLEMENTAR
    HAZZAN, Samuel. Fundamentos de matemática elementar, 5: combinatória,
probabilidade : 43 exercícios resolvidos, 439 exercícios propostos com
resposta, 155 testes de vestibulares com resposta. 7. ed. São Paulo, SP: Atual,
2004. 184 p. ISBN 8535704612 (broch.).
\newline \newline 
DANCEY, Christine P.; REIDY, John. Estatística sem matemática para
psicologia: usando SPSS para Windows. 3. ed. Porto Alegre, RS: Artmed,
2006. 608 p. (Biblioteca Artmed. Métodos de Pesquisa) ISBN 8536306882
(broch.).
\newline \newline 
HARPER, Brian D.; MERIAM, J. L; KRAIGE, L. G. Solving statistics problems in
\newline \newline 
MATLAB: engineering mechanics: statics. 6th ed. Massachusetts, [Estados
Unidos]: J. Wiley \& Sons, 2007. 139 p. ISBN 9780470099254 (broch.).
\newline \newline 
TRIOLA, Mario F. Introdução à estatística- Atualização da Tecnologia. 11. ed.
Rio de Janeiro, RJ:LTC, 2013. 740 p. ISBN 9788521622062 (broch.).
\newline \newline 
IEZZI, Gelson; HAZZAN, Samuel; DEGENSZAJN, David Mauro.Fundamentos
de matemática elementar: 11 : matemática comercial, matemática financeira,
estatística descritiva . São Paulo, SP: Atual, 2004. 232 p.
}


\DetalhesDisciplina{
    <25> - Processamento de Áudio e Voz
}{% EMENTA
    Modelos acústicos tradicionais, misturas de gaussianas e cadeias ocultas de markov. Representações de sinais para áudio e voz. Arquiteturas de redes neurais para reconhecimento de fala. Modelos de encoder, vocoder e arquiteturas para sintetização de voz. Arquiteturas para aplicações em música.
}{% BIBLIOGRAFIA BÁSICA
    HAYKIN, Simon S.. Redes neurais: princípios e prática. 2. ed. Porto Alegre: Bookman, 2001. xvii, 900 p.
ISBN 9788573077186.
\newline \newline 
Paulo S. R. Diniz; Eduardo A. B. Silva; Sergio L. Netto. Processamento digital de sinais. 
Grupo A. 2ª edição. Disponível em:
https://app.minhabiblioteca.com.br/reader/books/9788582601242 
\newline \newline 
Caseli, H.M.; Nunes, M.G.V. (org.) Processamento de Linguagem Natural: Conceitos, Técnicas e Aplicações em Português. 3 ed. BPLN, 2024. Disponível em: https://brasileiraspln.com/livro-pln/3a-edicao.
}{% BIBLIOGRAFIA COMPLEMENTAR
    RUSSELL, Stuart J.; NORVIG, Peter. Inteligência artificial. Rio de Janeiro: Elsevier, Campus, 2013. 988 p. ISBN 9788535237016 (broch.).
\newline \newline 
Júlio Serafim Martins; Maikon Lucian Lenz; Michel Bernardo Fernandes Da Silva; et al. Grupo A. Processamentos de Linguagem Natural. Disponível em: https://app.minhabiblioteca.com.br/reader/books/9786556900575
\newline \newline 
PAAß, Gerhard; GIESSELBACH, Sven. Foundation models for natural language processing: pre-trained language models integrating media. Cham: Springer, 2023. (Artificial Intelligence: Foundations, Theory, and Algorithms). DOI: https://doi.org/10.1007/978-3-031-23190-2. 
\newline \newline 
LIU, Zhiyuan; LIN, Yankai; SUN, Maosong (Ed.). Representation learning for natural language processing. 2. ed. Singapore: Springer, 2023. DOI: https://doi.org/10.1007/978-981-99-1600-9. 
\newline \newline 
AWAD, Mariette; KHANNA, Rahul. Efficient learning machines: theories, concepts, and applications for engineers and system designers. Berkeley, CA: Apress, 2015. DOI: https://doi.org/10.1007/978-1-4302-5990-9. 
}


\DetalhesDisciplina{
    <10> - Processamento de Dados em Larga Escala
}{% EMENTA
    Arquiteturas e paradigmas para manipulação e análise de grandes volumes de dados. Big Data. Arquiteturas e paradigmas para processamento de dados distribuído. Processamento de dados em lote (batch) e em fluxo (stream). Data Lakes e Data Warehouses. Estratégias de particionamento, tolerância a falhas e escalabilidade. Práticas de ETL, transformação e persistência de dados em ambientes distribuídos. Bancos de dados NoSQL e NewSQL. Computação em nuvem para Big Data. Otimização de performance e gerenciamento de recursos. Desenvolvimento de pipelines de dados escaláveis e eficientes.
}{% BIBLIOGRAFIA BÁSICA
    HUYEN, Chip. Projetando sistemas de Machine Learning: processo interativo para aplicações prontas para produção. Rio de Janeiro: Editora Alta Books, 2024. E-book. p.4. ISBN 9788550819648. Disponível em: https://app.minhabiblioteca.com.br/reader/books/9788550819648/.
\newline \newline 

PPADILHA, Juliana; SOARES, Juliane A.; ALVES, Nicolli S R.; et al. Analytics para big data. Porto Alegre: SAGAH, 2022. E-book. p.2. ISBN 9786556903477. Disponível em: https://app.minhabiblioteca.com.br/reader/books/9786556903477/.
\newline \newline 
SANTOS, Roger R.; BORDIN, Maycon V.; NUNES, Sergio E.; et al. Fundamentos de Big Data. Porto Alegre: SAGAH, 2021. E-book. p.Capa. ISBN 9786556901749. Disponível em: https://app.minhabiblioteca.com.br/reader/books/9786556901749/.

}{% BIBLIOGRAFIA COMPLEMENTAR
    ERL, Thomas; MONROY, Eric B. Computação em Nuvem: Conceitos, Tecnologia, Segurança e Arquitetura. 2. ed. Porto Alegre: Bookman, 2024. E-book. p.i. ISBN 9788582606599. Disponível em: https://app.minhabiblioteca.com.br/reader/books/9788582606599/.
\newline \newline 
RAGKOULIS, Marios; CARBONE, Paris; KALAVRI, Vasiliki; KATSIFODIMOS, Asterios. A survey on the evolution of stream processing systems. The VLDB Journal, [S.l.], v. 33, n. 2, p. 507–541, mar. 2024. DOI: 10.1007/s00778-023-00819-8. Disponível em: https://doi.org/10.1007/s00778-023-00819-8. 
\newline \newline 
LOURENÇO, João Ricardo; CABRAL, Bruno; CARREIRO, Paulo; VIEIRA, Marco; BERNARDINO, Jorge. Choosing the right NoSQL database for the job: a quality attribute evaluation. Journal of Big Data, [S.l.], v. 2, art. 18, p. 1–26, ago. 2015. DOI: 10.1186/s40537‑015‑0025‑0. Disponível em: https://journalofbigdata.springeropen.com/articles/10.1186/s40537-015-0025-0.
}


\DetalhesDisciplina{
    QXD0188 - Processamento de Imagens
}{% EMENTA
    Introdução ao processamento digital de imagens, Fundamentos sobre imagens
digitais, Formação de Imagens, Áreas de Aplicação. Transformações de
intensidade e filtragem espacial, Filtragem no domínio da frequência,
Amostragem e Quantização. Classificação de Imagens, Restauração e
reconstrução de imagens, Processamento de imagens coloridas, ​Wavelets e
processamento multiresolução, Compressão de imagens, Processamento
morfológico de imagens, Segmentação de imagens, Representação e
descrição, Reconhecimento de objetos.
}{% BIBLIOGRAFIA BÁSICA
    GONZALEZ, Rafael C.; WOODS, Richard E. Processamento digital de
imagens. 3. ed. Pearson, 2010. xv,624 p. ISBN 9788576054016 (broch.).
\newline \newline 
RUSS, John C. The Image Processing Handbook. Taylor \& Francis. 6 ed.
2010. ISBN 1439840458
\newline \newline 
PARKER, J. R. Algorithms for image Processing and Computer Vision. John
Wiley. 2 ed. 2010. ISBN 0470643854
}{% BIBLIOGRAFIA COMPLEMENTAR
    CONCI, Aura; AZEVEDO, Eduardo. Computação Gráfica, Volume 1 - Geração
de Imagens. Publicado por Elsevier. 2003. ISBN: 9788535212525, 384 páginas
GOMES, J. M.; VELHO, L. Fundamentos de computação gráfica. IMPA. 2008.
ISBN: 8524402008
\newline \newline 
BISWAS, Sambhunath; LOVELL, Brian C SPRINGERLINK (ONLINE
SERVICE). Bezier and Splines in Image Processing and Machine Vision.
Springer eBooks London: Springer-Verlag London Limited, 2008. ISBN
9781846289576. Disponível em : 
<http://dx.doi.org/10.1007/978-1-84628-957-6>. Acesso em : 21 set. 2010.
\newline \newline 
HUTCHISON, David; ELMOATAZ, Abderrahim; KANADE, Takeo; KITTLER,
Josef; KLEINBERG, Jon M; LEZORAY, Olivier; MAMMASS, Dris; MATTERN,
Friedemann; MITCHELL, John C; NAOR, Moni; NIERSTRASZ, Oscar;
\newline \newline 
NOUBOUD. Image and Signal Processing : 3rd International Conference,
ICISP 2008 Cherbourg-Octeville, France, July 1-3, 2008 Proceedings . Springer
eBooks Berlin, Heidelberg: Springer-Verlag Berlin Heidelberg, 2008. (Lecture
Notes in Computer Science, 5099) ISBN 9783540699057. Disponível em :
<http://dx.doi.org/10.1007/978-3-540-69905-7>. Acesso em : 21 set. 2010.
\newline \newline 
SARFRAZ, M SPRINGERLINK (ONLINE SERVICE). Interactive Curve
Modeling : With Applications to Computer Graphics, Vision and Image
Processing . Springer e-books London: Springer-Verlag London Limited, 2008.
ISBN 9781846288715. Disponível em :
<​http://dx.doi.org/10.1007/978-1-84628-871-5​>. Acesso em : 21 set. 2010.
\newline \newline 
CONCI, Aura; AZEVEDO, Eduardo.; LETA, Fabiana R. Computação gráfica,
v.2: teoria e prática. Rio de Janeiro, RJ: Elsevier: Campus, 2008. 407 p., [8] p.
de estampas + 1 CD-ROM ISBN 97885352232193 (broch).
\newline \newline 
JAIN, A.K. Fundamentals of Digital Image Processing. 1 ed. Prentice-Hall,
Addison-Wesley, 1988. ISBN-10: 0133361659 ISBN-13: 978-0133361650.
\newline \newline 
BOVIK, Al (ed.). Handbook of Image and Video Processing. 2 ed. Academic
Press, 2005. ISBN-10: 0121197921 ISBN-13: 9780121197926.
\newline \newline 
LIM, J. S. Two-dimensional Signal and Image Processing. Prentice Hall Press,
1990. ISBN-10: 0139353224 ISBN-13: 978-0139353222.
\newline \newline 
PETROU, Costas; PETROU, Maria. Image Processing: the fundamentals. 2 ed.
Wiley, 2010. ISBN-10: 047074586X ISBN-13: 9780470745861
\newline \newline 
JAHNE, Bernd. Practical Handbook on Image Processing for Scientific
Applications. 2 ed. CRC Press, 2004. ISBN-10: 0849319005 ISBN-13:
9780849319006.
}


\DetalhesDisciplina{
    <11> - Processamento de Linguagem Natural
}{% EMENTA
    Fundamentos e aplicações do Processamento de Linguagem Natural (PLN). Pré-processamento de texto. Representações textuais (bag-of-words, TF-IDF e representações distribuídas de palavras (word embeddings). Modelos de linguagem estatísticos e neurais. Arquiteturas baseadas em transformadores e mecanismos de atenção. Modelos pré-treinados. Classificação textual, análise de sentimentos, reconhecimento de entidades nomeadas (NER), extração de informações, tradução automática, sumarização de texto, sistemas de perguntas e respostas, chatbots e sistemas de diálogo. Métodos de avaliação de desempenho.
}{% BIBLIOGRAFIA BÁSICA
    Caseli, H.M.; Nunes, M.G.V. (org.) Processamento de Linguagem Natural: Conceitos, Técnicas e Aplicações em Português. 3 ed. BPLN, 2024. Licença Creative Commons. Disponível em: https://brasileiraspln.com/livro-pln/3a-edicao.
\newline \newline 
MCROY, Susan. Principles of Natural Language Processing. Milwaukee: Susan McRoy, 24 jul. 2021. 264 p. (Open textbook). ISBN 978‑1‑7376595‑0‑1. Disponível em: https://wisconsin.pressbooks.pub/naturallanguage/.
\newline \newline 
RUSSELL, Stuart J.; NORVIG, Peter. Inteligência Artificial: Uma Abordagem Moderna. 4. ed. Rio de Janeiro: GEN LTC, 2022. E-book. p.Capa. ISBN 9788595159495. Disponível em: https://app.minhabiblioteca.com.br/reader/books/9788595159495/. 
}{% BIBLIOGRAFIA COMPLEMENTAR
    HARIOM, Tatsat,; SAHIL, Puri,; BRAD, Lookabaugh,. Blueprints de aprendizado de máquina e ciência de dados para finanças: desenvolvendo desde estratégias de trades até robôs Advisors com Python. Rio de Janeiro: Editora Alta Books, 2024. E-book. p.i. ISBN 9788550821726. Disponível em: https://app.minhabiblioteca.com.br/reader/books/9788550821726/.
\newline \newline 
FACELI, Katti; LORENA, Ana C.; GAMA, João; AL, et. Inteligência Artificial - Uma Abordagem de Aprendizado de Máquina. 2. ed. Rio de Janeiro: LTC, 2021. E-book. p.Capa. ISBN 9788521637509. Disponível em: https://app.minhabiblioteca.com.br/reader/books/9788521637509/.
\newline \newline 
LIU, Zhiyuan; LIN, Yankai; SUN, Maosong (eds.). Representation Learning for Natural Language Processing. Singapore: Springer, 2023. eBook ISBN 978‑981‑99‑1600‑9; impresso ISBN 978‑981‑99‑1599‑6. Disponível em: https://link.springer.com/book/10.1007/978-981-99-1600-9. 

}


\DetalhesDisciplina{
    QXD0114 - Programação Funcional
}{% EMENTA
    Visão geral e motivação. Recursão sobre listas, números naturais, árvores, e
outros dados definidos recursivamente. Uso de funções como dados.
Expressões lambda. Avaliação preguiçosa. Prática de programação em
linguagem deste paradigma. Questões práticas como ​I/O​, depuração e
persistência de estruturas de dados.
}{% BIBLIOGRAFIA BÁSICA
    WAMPLER, Dean. Programação Funcional Para Desenvolvedores Java :
Ferramentas para Melhor Concorrência, Abstração e Agilidade. Novatec. 1a
ed., 2012. (ISBN 9788575223161)
\newline \newline 
SÁ, Claudio Cesar de. Haskell : uma abordagem prática. São Paulo, SP:
Novatec, 2006. 287 p.
\newline \newline 
SEIBEL, Peter. Practical common lisp . Berkeley, Ca: Apress, 2005. xxv, 499 p.
(The Expert’s voice in programming languages).
\newline \newline 
MICHAELSON, Greg. An Introduction to Functional Programming Through
Lambda Calculus . Dover Publications, 2011. ISBN: 0486478831
\newline \newline 
COUSINEAU, Guy; MAUNY, Michel; CALLAWAY, K. The Functional Approach
to Programming. Cambridge University Press; English edition, 1998. ISBN-10:
0521576814
}{% BIBLIOGRAFIA COMPLEMENTAR
    CHEN, Yinong; TSAI, Wei Tek. Introduction to programming languages:
programming in C, C++, Scheme, Prolog, C\#, and SOA. 2nd ed. xii, 383 p.
\newline \newline 
SEBESTA, Robert W. Conceitos de linguagens de programação. 9. ed. Porto
Alegre, RS: Bookman, 2011. ix, 792 p.
\newline \newline 
TUCKER, Allen B.; NOONAN, Robert. Linguagens de programação: princípios
e paradigmas. São Paulo, SP: McGraw Hill, 2009. xxi, 599p.
\newline \newline 
LEE, Kent SPRINGERLINK (ONLINE SERVICE). Programming Languages :
An Active Learning Approach . Springer eBooks Boston, MA: Springer-Verlag
US, 2008. ISBN 9780387794228. Disponível em :
<http://dx.doi.org/10.1007/978-0-387-79421-1>. Acesso em : 21 set. 2010.
\newline \newline 
GABBRIELLI, Maurizio; MARTINI, Simone; SPRINGERLINK 
(ONLINE SERVICE). Programming Languages: Principles and Paradigms . Springer
eBooks: Springer-Verlag London, 2010. ISBN 978-1-84882-914-5. Disponível em : <http://link.springer.com/book/10.1007/978-1-84882-914-5>. Acesso em :
12 jan. 2016.
\newline \newline 
EMERICK, Chas; CARPER, Brian; GRAND, Christophe. Clojure Programming.
O'Reilly Media; 1 edition, 2011. ISBN: 1449394701
\newline \newline 
LIPOVACA, Miran. Learn You a Haskell for Great Good!: A Beginner's Guide.
O'Reilly; 1 edition, 2011. ISBN: 1593272839
\newline \newline 
PETRICEK, Tomas; SKEET, Jon. Real-World Functional Programming: With
Examples in F\# and C\#. Manning Publications; 2010. ISBN-10: 1933988924
\newline \newline 
DYBVIG, R. Kent. The Scheme Programming Language, MIT Press; fourth
edition, 2009. ISBN-10: 026251298X
}


\DetalhesDisciplina{
    QXD0007 - Programação Orientada a Objetos
}{% EMENTA
    Introduzir o paradigma de Programação Orientada a Objetos (OO), juntamente
com seus conceitos de classes, objetos, herança, encapsulamento e
polimorfismo, além dos conceitos de Interfaces e exceções que são inerentes
às linguagens de programação orientadas a objetos. Desenvolvimento de um
pequeno sistema baseados no paradigma de programação OO.
}{% BIBLIOGRAFIA BÁSICA
    DEITEL, H. M. Java: como programar. 8 ed. Prentice Hall, 2010. ISBN:
9788576055631
\newline \newline 
HORSTMANN, Cay S. Core Java: volume I - fundamentos. 8. ed. São Paulo,
SP: Pearson, 2009. xiii, 383 p. ISBN 9788576053576
\newline \newline 
MCLAUGHLIN, Brett; POLLICE, Gary; WEST, David.Use a cabeça:análise e
projeto orientado ao objeto. Rio de Janeiro, RJ: Alta Books, 2007. xxviii, 441 p.
}{% BIBLIOGRAFIA COMPLEMENTAR
    BARNES, D.J.; KOLLING, M. Programação orientada a objetos com Java :
uma introdução prática usando BLUEJ. 4 ed. Prentice Hall Brasil, 2009. ISBN:
9788576051879
\newline \newline 
SIERRA, Kathy; BATES, Bert. Use a cabeça! Java. Rio de Janeiro: Alta Books,
2007. 470 p. ISBN 0596009208.
\newline \newline 
DEITEL, H. M. C++ como programar. 5 ed. Prentice Hall, 2006. ISBN:
8576050560.
\newline \newline 
BLAHA, Michael; RUMBAUGH, James. Modelagem e projetos baseados em
objetos com UML 2. 2.ed. rev. e atual. Rio de Janeiro, RJ: Campus; Elsevier,
2006. xvii, 496 p. ISBN 9788535217537 (broch.).
\newline \newline 
MANZANO, José Augusto N. G.; COSTA Jr., Roberto Affonso da. Java 7 -
Programação de Computadores - Guia Prático de Introdução, Orientação e
Desenvolvimento. 1. ed. Editora Érica, 2011. ISBN: 9788536503745.
\newline \newline 
MEYER, Bertrand. Object-Oriented Software Construction , Ed. Prentice Hall
PTR, 1997.
}


\DetalhesDisciplina{
    QXD0110 - Projeto de Pesquisa Científico-Tecnológico
}{% EMENTA
    O problema da pesquisa e sua formulação. Métodos e Técnicas de Pesquisa.
O planejamento da pesquisa. Elaboração de projeto de pesquisa referente
ao Trabalho de Conclusão de Curso.
}{% BIBLIOGRAFIA BÁSICA
    WAZLAWICK, Raul Sidnei. Metodologia de pesquisa para ciência da computação. Rio de Janeiro : Elsevier. 2008.
\newline \newline 
LAVILLE, Christian; Dionne, Jean. A Construção do Saber: Manual de
Metodologia da Pesquisa em Ciências Humanas. Porto Alegre: Artmed,
Belo Horizonte: Editora UFMG, 2008. 340 p.
\newline \newline 
MARCONI, Marina de Andrade; Lakatos, Eva Maria. Fundamentos de
Metodologia Científica.7. ed. São Paulo, SP: Atlas, 2010.
}{% BIBLIOGRAFIA COMPLEMENTAR
    YIN, Robert K. Estudo de Caso - Planejamento e Métodos. 4aed, Porto
Alegre : Bookman, 2010.
\newline \newline 
CERVO, A.; BERVIAN, P.A.; SILVA, R. Metodologia Científica. 6. ed.
2007. ISBN 8576050471
\newline \newline 
COOPER, D. R.; Schindler, Pamela S. Métodos de Pesquisa em Administração. Porto Alegre 7a ed Bookman. 2008.
\newline \newline 
FOWLER, F.J. Pesquisa de Levantamento. Porto Alegre: Pearson, 2011.
\newline \newline 
FREIRE, P. Extensão ou Comunicação.13. ed. Paz e Terra, 2006..
CHAUÍ, Marilena. Convite à Filosofia. 14 ed. Ática, 2011.
}


\DetalhesDisciplina{
    QXD0041 - Projeto e Análise de Algoritmos
}{% EMENTA
    Noções de análise de algoritmos: análise assintótica de pior caso e caso
médio; notação big-O, little-o, ômega e teta; principais classes de complexidade; medida empírica de performance; análise de algoritmos recursivos
utilizando relações de recorrência. Projeto de algoritmos: força bruta; gulosos; divisão e conquista; programação dinâmica. Algoritmos em grafos: grafos não-direcionados e direcionados; árvores; conectividade; árvores/florestas geradoras; ordenação topológica; caminho mais curto. NPcompletude: definição das classes P e NP; teorema de Cook; principais
problemas NP-completos; técnicas de redução.
}{% BIBLIOGRAFIA BÁSICA
    CORMEN, Thomas. Algoritmos - Teoria e Prática. São Paulo: GEN LTC,
2012. E-book. Disponível em: https://integrada.minhabiblioteca.
com.br/\#/books/9788595158092. Acesso em: 3 de out. 2022.
\newline \newline 
DASGUPTA, Sanjoy; PAPADIMITRIOU, Christos; VAZIRANI, Umesh.
Algoritmos.. São Paulo: AMGH, 2009. E-book. Disponível em: https://
integrada.minhabiblioteca.com.br/\#/books/9788563308535. Acesso
em: 3 de out. 2022.
\newline \newline 
ZIVIANI, Nivio. Projeto de Algoritmos: com implementações em JAVA
e C++. São Paulo: Cengage Learning Editores SA de CV, 2012. Ebook. Disponível em: https://integrada.minhabiblioteca.com.br/\#/
books/9788522108213. Acesso em: 3 de out. 2022.
}{% BIBLIOGRAFIA COMPLEMENTAR
    KLEINBERG, Jon; TARDOS, Éva. Algorithm design. Boston, Massachusetts: Pearson/Addison Wesley, c2006. 838 p. ISBN 0321295358.
\newline \newline 
GERSTING, Judith L.. Fundamentos Matemáticos para a Ciência da
Computação. São Paulo: LTC, 2016. E-book. Disponível em: https://
integrada.minhabiblioteca.com.br/\#/books/9788521633303. Acesso
em: 3 de out. 2022.
\newline \newline 
MENEZES, Paulo Blauth. Matemática Discreta para Computação e Informática - V16 - UFRGS. São Paulo: Bookman, 2013. E-book. Disponível em: https://integrada.minhabiblioteca.com.br/\#/books/
9788582600252. Acesso em: 3 de out. 2022.
\newline \newline 
ROSEN, Kenneth H.. Matemática Discreta e suas Aplicações. São
Paulo: ArtMed, 2010. E-book. Disponível em: https://integrada.
minhabiblioteca.com.br/\#/books/9788563308399. Acesso em: 3 de
out. 2022.
\newline \newline 
GOLDBARG, Marco Cesar; GOLDBARG, Elizabeth. Grafos: conceitos,
algoritmos e aplicações. Rio de Janeiro, RJ: Elsevier, 2012. 622 p. ISBN
9788535257168
\newline \newline 
TOSCANI, Laira V.; VELOSO, Paulo A. S. Complexidade de algoritmos:
análise, projeto e métodos. 3. ed. Porto Alegre: Sagra Luzzato, 2012. 262
p. (Serie Livros Didáticos Informática UFRGS ; 13). ISBN 9788540701380
(broch.).
}


\DetalhesDisciplina{
    <19> - Projeto Extensionista I
}{% EMENTA
    Introdução à extensão universitária e à política nacional de extensão. Diretrizes Nacionais da Extensão Universitária. Formação cidadã e transformação social. Planejamento e execução de ações em comunidades
de saberes.
}{% BIBLIOGRAFIA BÁSICA
    FREIRE, Paulo. Extensão ou comunicação. 13. ed. Rio de Janeiro: Paz e Terra, 2006. 93 p. (O mundo,
hoje, 24). ISBN 8521904274.
\newline \newline 
DEUS, Sandra de. Extensão universitária: trajetórias e desafios. Santa Maria, RS. Ed. PRE-UFSM, 2020.
96 p. ISBN: 978-65-87668-01-7
\newline \newline 
FORPROEX - FÓRUM DE PRÓ-REITORES DE EXTENSÃO DAS UNIVERSIDADES PÚBLICAS
BRASILEIRAS. Política Nacional de Extensão Universitária. Manaus, 2012. Disponível em:
<https://proex.ufsc.br/files/2016/04/Política-Nacional-de-Extensão-Universitária-e-book.pdf> Acesso em:
maio de 2025.
}{% BIBLIOGRAFIA COMPLEMENTAR
    COELHO, Geraldo Ceni. O papel pedagógico da extensão universitária. Revista Em Extensão,
Uberlândia, v. 13, n. 2, p. 11–24, 2015. DOI: 10.14393/REE-v13n22014\_art01. Disponível em:
https://seer.ufu.br/index.php/revextensao/article/view/26682. Acesso em: 12 maio. 2025.
\newline \newline 
LEMOS, Ronaldo.; DI FELICE, Massimo. A vida em rede. Campinas: Papirus, 2014. 142 p. (Papirus
debates). ISBN 9788561773618 (broch.).
\newline \newline 
YIN, Robert K.. Estudo de caso: planejamento e métodos. 4. ed. Porto Alegre: Bookman, 2010. xviii, 248 p.
ISBN 978-85-7780-655-3.
\newline \newline 
UNIVERSIDADE FEDERAL DO CEARÁ. Glossário: fundamentos da extensão universitária.
Fortaleza: Pró-Reitoria de Extensão, 2024. Disponível em: https://prex.ufc.br/wpcontent/uploads/2024/08/glossario-fundamentos-da-extensao-universitaria-240819-180754.pdf. Acesso em:
12 maio 2025.
\newline \newline 
ZITKOSKI, Jaime J.; STRECK, Danilo R.; REDIN, Euclides. Dicionário Paulo Freire. 2. ed. São Paulo:
Autêntica Editora, 2008. E-book. p.Cover. ISBN 9788582178089. Disponível em:
https://app.minhabiblioteca.com.br/reader/books/9788582178089/. Acesso em: 12 mai. 2025.
}


\DetalhesDisciplina{
    <20> - Projeto Extensionista II
}{% EMENTA
    Interdisciplinaridade e interprofissionalidade nas práticas extensionistas. Interação dialógica com comunidades de saberes. Formação discente integrada a contextos extensionistas. Transformação social. Indissociabilidade entre ensino, pesquisa e extensão. Práticas extensionistas.
}{% BIBLIOGRAFIA BÁSICA
    FREIRE, Paulo. Extensão ou comunicação. 13. ed. Rio de Janeiro: Paz e Terra, 2006. 93 p. (O mundo,
hoje, 24). ISBN 8521904274.
\newline \newline 
DEUS, Sandra de. Extensão universitária: trajetórias e desafios. Santa Maria, RS. Ed. PRE-UFSM, 2020.
96 p. ISBN: 978-65-87668-01-7
\newline \newline 
FORPROEX - FÓRUM DE PRÓ-REITORES DE EXTENSÃO DAS UNIVERSIDADES PÚBLICAS
BRASILEIRAS. Política Nacional de Extensão Universitária. Manaus, 2012. Disponível em:
<https://proex.ufsc.br/files/2016/04/Política-Nacional-de-Extensão-Universitária-e-book.pdf> Acesso em:
maio de 2025.
}{% BIBLIOGRAFIA COMPLEMENTAR
    ""COELHO, Geraldo Ceni. O papel pedagógico da extensão universitária. Revista Em Extensão,
Uberlândia, v. 13, n. 2, p. 11–24, 2015. DOI: 10.14393/REE-v13n22014\_art01. Disponível em:
https://seer.ufu.br/index.php/revextensao/article/view/26682. Acesso em: 12 maio. 2025.
\newline \newline 
LEMOS, Ronaldo.; DI FELICE, Massimo. A vida em rede. Campinas: Papirus, 2014. 142 p. (Papirus
debates). ISBN 9788561773618 (broch.).
\newline \newline 
YIN, Robert K.. Estudo de caso: planejamento e métodos. 4. ed. Porto Alegre: Bookman, 2010. xviii, 248 p.
ISBN 978-85-7780-655-3.
\newline \newline 
UNIVERSIDADE FEDERAL DO CEARÁ. Glossário: fundamentos da extensão universitária.
Fortaleza: Pró-Reitoria de Extensão, 2024. Disponível em: https://prex.ufc.br/wpcontent/uploads/2024/08/glossario-fundamentos-da-extensao-universitaria-240819-180754.pdf. Acesso em:
12 maio 2025.
\newline \newline 
ZITKOSKI, Jaime J.; STRECK, Danilo R.; REDIN, Euclides. Dicionário Paulo Freire. 2. ed. São Paulo:
Autêntica Editora, 2008. E-book. p.Cover. ISBN 9788582178089. Disponível em:
https://app.minhabiblioteca.com.br/reader/books/9788582178089/. Acesso em: 12 mai. 2025.""
}


\DetalhesDisciplina{
    QXD0126 - Psicologia e Percepção
}{% EMENTA
    O ser humano em sua relação com o mundo. Processos psicológicos
relacionados à percepção, sensação e cognição. Aprendizagem, memória,
motivação, emoção e linguagem.
}{% BIBLIOGRAFIA BÁSICA
    ARNHEIM, Rudolf. Arte e percepção visual: uma psicologia da visão criadora:
nova versão. São Paulo: Cengage Learning, 2017. 509 p. ISBN
9788522126002 (broch.)
\newline \newline 
GOMES FILHO, João. Gestalt do objeto: sistema de leitura visual da forma.
9.ed. São Paulo: Escrituras, 2009. ISBN: 9788586303579
\newline \newline 
MORRIS, Charles G. e MAISTO, Albert A. Introdução à psicologia. 6ª edição.
São Paulo, 2004. ISBN: 9788587918680
}{% BIBLIOGRAFIA COMPLEMENTAR
    FELDMAN, Robert S. Introdução à Psicologia. Porto Alegre: Grupo A, 2015.
E-book. ISBN 9788580554892. Disponíví el em:
https://app.minhabiblioteca.com.br/\#/books/9788580554892/. Acesso em:
26 set. 2022.
\newline \newline 
BOCK, Ana Mercês B.; FURTADO, Odair; TEIXEIRA, Maria de Lourdes T.
Psicologias: uma introdução ao estudo de psicologia. São Paulo: Editora
Saraiva, 2018. E-book. ISBN 9788553131327. Disponíví el em:
https://app.minhabiblioteca.com.br/\#/books/9788553131327/. Acesso em:
26 set. 2022.
\newline \newline 
DA MAIA, Gabriela Felten; FORECHI, Marcilene; LOPES, Daiane D.; et al.
Comunicação e Psicologia. Porto Alegre: Grupo A, 2020. E-book. ISBN
9786581492960. Disponíví el em:
https://app.minhabiblioteca.com.br/\#/books/9786581492960/. Acesso em:
26 set. 2022.
\newline \newline 
FREUD, Sigmund. Psicologia das massas e análise do eu e outros textos
(1920-1923). São Paulo: Companhia das Letras, 2011. 343 p. (Obras
completas ; v. 15). ISBN 9788535918717 (broch.).
\newline \newline 
BARBOSA, Livia. Sociedade de consumo. 2. ed. Rio de Janeiro: Zahar, 2008.
((Ciências Sociais Passo-a-Passo ; 49).). ISBN 9788571108134.
}


\DetalhesDisciplina{
    <26> - Raciocínio sob Incerteza
}{% EMENTA
    Tipos de incerteza; Teoria de probabilidade (revisão), raciocínio probabilístico, redes bayesianas, lógica probabilística, lógica fuzzy; Raciocínio Não-Monotônico, programação em lógica, raciocínio abdutivo; Problemas de decisão sob incerteza.
}{% BIBLIOGRAFIA BÁSICA
    LARSON, Ron; FARBER, Betsy. Estatística aplicada. 4. ed. São Paulo, SP: Pearson/ Prentice Hall, 2010. xiv,637 p. ISBN 9788576053729 (broch.).
\newline \newline 
COPPIN, Ben. Inteligencia artificial. Rio de Janeiro, RJ: LTC, 2010. 636 p. ISBN 9788521617297.
\newline \newline 
RUSSELL, Stuart J.; NORVIG, Peter. Inteligência artificial. Rio de Janeiro: Elsevier, Campus, 2013. 988 p. ISBN 9788535237016 (broch.).
}{% BIBLIOGRAFIA COMPLEMENTAR
    Lukasiewicz T. Probabilistic description logic programs. International Journal of Approximate Reasoning. 2007 Jul 1;45(2):288-307. Disponível em: https://www.sciencedirect.com/science/article/pii/S0888613X06000648    
\newline \newline 
LARSON, Ron; FARBER, Betsy. Estatística aplicada. 4. ed. São Paulo, SP: Pearson/ Prentice Hall, 2010. xiv,637
p. ISBN 9788576053729 (broch.)                                                                
\newline \newline 
BARBETTA, Pedro Alberto; REIS, Marcelo Menezes; BORNIA, Antonio Cezar. Estatística para cursos de
engenharia e informática. 3. ed. São Paulo, SP: Atlas, 2010. 410 p    
\newline \newline 
MORETTIN, Luiz Gonzaga. Estatística básica:: probabilidade e inferência / volume único. São Paulo, SP:
Pearson Educational do Brasil. 2010. ISBN 8576053705 ISBN-13 9788576053705 (broch.)                                                                       
\newline \newline 
WALPOLE, Ronald E. Probabilidade e estatística: para engenharia e ciências. 8. ed. São Paulo, SP: Pearson/
Prentice Hall, 2009. xiv, 491 p. ISBN 9788576051992 (broch.)                            
}


\DetalhesDisciplina{
    QXD0177 - Recuperação de Informação
}{% EMENTA
    Avaliação em recuperação de Informação; Modelos clássicos de recuperação (Booleano, Vetorial e Probabilístico); Operações sobre texto, indexação e consultas (pré-processamento); Mineração de texto; Indexação de textos na Web; Web Search; Web Crawling; Algoritmos baseados em links.
}{% BIBLIOGRAFIA BÁSICA
    BAEZA-YATES, Ricardo; RIBEIRO-NETO, Berthier. Recuperação de Informação. 2. ed. Porto Alegre: Bookman, 2013. E-book. p.1. ISBN 9788582600498. Disponível em: https://app.minhabiblioteca.com.br/reader/books/9788582600498/.
\newline \newline 
MANNING, Christopher D.; RAGHAVAN, Prabhakar; SCHÜTZE, Hinrich. Introduction to information retrieval. New York, NY: Cambridge at the University Press, 2009. xxi, 482 p. ISBN 9780521865715.
\newline \newline 
ELMASRI, R.; NAVATHE, S. B. Sistemas de banco de dados. 6 ed. Pearson/Addison-Wesley, 2011. ISBN: 9788579360855.

}{% BIBLIOGRAFIA COMPLEMENTAR
    SILBERSCHATZ, A.; SUDARSHAN, S. Sistema de banco de dados. Campus, 2006. ISBN: 9788535211078.
\newline \newline 
LARSON, Ron; FARBER, Betsy. Estatística aplicada. 4. ed. São Paulo, SP: Pearson/ Prentice Hall, 2010. xiv,637 p. ISBN 9788576053729 (broch.).
\newline \newline 
RAMAKRISHNAN, R.; GEHRKE, J.Sistemas de gerenciamento de banco de dados. McGrawHill, 2008. DOMINICH, Sendor SPRINGERLINK (ONLINE SERVICE). 
\newline \newline 
The Modern Algebra of Information Retrieval. Springer eBooks Berlin, Heidelberg: Springer-Verlag Berlin Heidelberg, 2008. (The Information Retrieval Series, 24) ISBN 9783540776598. Disponível em : <http://dx.doi.org/10.1007/978-3-540-77659-8>. Acesso em : 21 set. 2010.
\newline \newline 
HARPER, Brian D.; MERIAM, J. L; KRAIGE, L. G. Solving statistics problems in MATLAB: engineering mechanics: statics. 6th ed. Massachusetts, [Estados Unidos]: J. Wiley \& Sons, 2007. 139 p. ISBN 9780470099254 (broch.).
\newline \newline 
BARBETTA, Pedro Alberto; REIS, Marcelo Menezes; BORNIA, Antonio Cezar. Estatística para cursos de engenharia e informática. 3. ed. São Paulo, SP: Atlas, 2010. 410 p.
\newline \newline 
GROSSMAN, David A.; FRIEDER, Ophir. Information Retrieval: Algorithms and Heuristics. Springer, 2ª Ed., 2004. ISBN: 1402030045.
\newline \newline 
BUETTCHER, Stefan; Clarke, Charles L. A.; Cormack, Gordon V. Information Retrieval: Implementing and Evaluating Search Engines. MIT Press; 1ª Ed., 2010.
\newline \newline 
MORETTIN, Luiz Gonzaga. Estatística básica: probabilidade e inferência / volume único. São Paulo, SP: Pearson Educational do Brasil. 2010. ISBN 8576053705 ISBN-13 9788576053705.
}


\DetalhesDisciplina{
    QXD0021 - Redes de Computadores
}{% EMENTA
    Organização das redes de computadores. Modelos de referência OSI e
TCP/IP. Padrões de rede. Meios físicos de transmissão. Protocolos de acesso
ao meio. Interconexão de redes. Algoritmos e protocolos de roteamento.
Protocolos de redes. Protocolos de transporte TCP e UDP. Protocolos de
aplicação. Projeto e Dimensionamento de Redes.
}{% BIBLIOGRAFIA BÁSICA
    COMER, D. Redes de computadores e a internet. 4 ed. Bookman, 2007.ISBN:
9788560031368.
\newline \newline 
KUROSE, J.; ROSS, K. W. Redes de computadores e a Internet: uma
abordagem top-down, 6 ed. Pearson, 2013. ISBN: 8581436773.
\newline \newline 
FOROUZAN, Behrouz A. Comunicação de dados e redes de computadores. 4.
ed. São Paulo, SP: McGraw-Hill, 2008. xxxiv, 1134 p.
\newline \newline 
FOROUZAN, B. A.; FIROUZ, M. Redes de Computadores: Uma Abordagem
Top-down. 1a ed. 2013. McGraw-Hill. ISBN: 9788580551686
}{% BIBLIOGRAFIA COMPLEMENTAR
    TANENBAUM, Andrew S.; WETHERALL, D. Redes de computadores. 5. ed.
São Paulo: Pearson Prentice Hall, c2011. xvi, 582 p. ISBN 9788576059240.
\newline \newline 
COMER, D. Interligação de redes com TCP/IP.1. 5 ed.(vol.1). Campus, 2006.
ISBN: 9788535220179.
\newline \newline 
OLIVEIRA, Gorki Starlin da Costa. Redes de computadores comunicações de
dados TCP/IP : conceitos, protocolos e usos. Alta Books, 2004. ISBN :
8576080567
\newline \newline 
TORRES, G. Redes de computadores. Nova Terra, 2009.
\newline \newline 
MORIMOTO, C.E. Redes: guia prático. GDH Press, 2009. ISBN
9788599593097 (broch.).
\newline \newline 
PETERSON, Larry L.; DAVIE, Bruce S. Computer networks: a systems
approach . 5th ed. Amesterdam: Elsevier, c2012. xxxi, 884 p. ISBN
9780123850591
}


\DetalhesDisciplina{
    QXD0246 - Relações Étnico-Raciais e Africanidades
}{% EMENTA
    Negritude e pertencimento étnico. Conceitos de africanidades e afrodescendência. Cosmovisão africana: valores civilizatórios africanos presentes
na cultura brasileira. Ancestralidade e ensinamentos das religiosidades
tradicionais africanas nas diversas dimensões do conhecimento no Brasil.
Introdução à geografia e história da África. As origens africanas e as nações africanas representadas no Brasil. O sistema escravista no Brasil e no
Ceará. Aportes dos africanos à formação social e cultural do Brasil e do
Ceará. Personalidades africanas, afrodescendentes e da diáspora negra que
se destacaram em diferentes áreas do conhecimento. Contexto das Ações
Afirmativas hoje. Atualização do legado africano no Brasil. Desconstrução de preconceitos e desdobramentos teórico-práticos para a atuação do
profissional na sua área de inserção no mercado de trabalho.
}{% BIBLIOGRAFIA BÁSICA
    VAZZOLER por MUNANGA, Kabengele. Negritude - Nova Edição.
[Digite o Local da Editora]: Grupo Autêntica, 2019. E-book. ISBN 9788551306529. Disponível em: https://app.minhabiblioteca.com.
br/\#/books/9788551306529. Acesso em: 23 set. 2022.
OLIVA, Anderson R.; CHAVES, Marjorie N.; FILICE, Renísia Cristina G.;
\newline \newline 
NASCIMENTO, Wan. Tecendo redes antirracistas. [Digite o Local da Editora]: Grupo Autêntica, 2019. E-book. ISBN 9788551304877. Disponível
em: https://app.minhabiblioteca.com.br/\#/books/9788551304877.
Acesso em: 23 set. 2022.
\newline \newline 
ALMEIDA, Guilherme Assis de; CHRISTMANM, Martha Ochsenhofer.
Ética e Direito: Uma Perspectiva Integrada, 3a edição. São Paulo:
Grupo GEN, 2009. E-book. Disponível em: https://integrada.
minhabiblioteca.com.br/\#/books/9788522467150. Acesso em: 3 de
out. 2022.
}{% BIBLIOGRAFIA COMPLEMENTAR
    BRASIL. Síntese de indicadores Sociais: Uma análise das condições de vida
da população brasileira. Ministério do Planejamento, Orçamento e Gestão.
IBGE. Rio de Janeiro, 2013, 266p. Disponível em: https://biblioteca.
ibge.gov.br/visualizacao/livros/liv66777.pdf
\newline \newline 
CUNHA JUNIOR, H. Abolição inacabada e a educação dos afrodescendentes. Revista Espaço Acadêmico, no. 89, 2008. Disponível em:
http://www.espacoacademico.com.br/089/89cunhajr.pdf
\newline \newline 
ROMÃO, J. História da educação do negro e outras histórias. Brasília: Ministério da Educação, Secretaria de Educação Continuada, Alfabetização
e Diversidade. 2005, 278p. Disponível em: http://etnicoracial.mec.
gov.br/images/pdf/publicacoes/historia\\_educacao\\_negro.pdf
\newline \newline 
BRASIL. Constituição da República Federativa do Brasil. Biblioteca Digital da Câmara dos Deputados. 35a ed., 2012, 446p.
\newline \newline 
BRASIL. CNE. Parecer no. 03 de 10 de março de 2004. Dispõe sobre as diretrizes curriculares nacionais para a educação das relações
étnico-raciais e para o ensino de história e cultura afrobrasileira e africana. Relatora: Petronilha Beatriz Gonçalves e Silva. Ministério da
Educação. Brasília, julho de 2004. Disponível em: https://www.gov.
br/inep/pt-br/centrais-de-conteudo/acervo-linha-editorial/
publicacoes-diversas/temas-interdisciplinares/
diretrizes-curriculares-nacionais-para-a-educacao-das-relacoes-etnico-raciais-e-para-o-ensino-de-historia-e-cultura-afro-brasileira-e-africana
\newline \newline 
SANTIAGO, R. A história da educação do negro no Brasil: interdição institucional à escolarização pelo poder e seus reflexos no século XXI. Revista da ABPN. v.5, n.10, p.196-203, 2013. Disponível em: https://www.ipea.gov.br/portal/images/stories/PDFs/
politicas\\_sociais/bps\\_19\\_cap08.pdf
}


\DetalhesDisciplina{
    <04> - Representação do Conhecimento e Raciocínio
}{% EMENTA
    Lógica de primeira ordem (revisão); Representação de conhecimento em FOL: vocabulário/léxico, fatos simples, fatos compostos, consequências lógicas, fechamento dedutivo. Problema do Quadro. Cláusulas de Horn. Resolução. Árvores de Objetivos. Introdução ao Raciocínio Não-Monotônico.
}{% BIBLIOGRAFIA BÁSICA
    Ronald Brachman and Hector Levesque. 2004. Knowledge Representation and Reasoning. Morgan Kaufmann Publishers Inc., San Francisco, CA, USA.
\newline \newline 
COPPIN, Ben. Inteligencia artificial. Rio de Janeiro, RJ: LTC, 2010. 636 p. ISBN 9788521617297.
\newline \newline 
RUSSELL, Stuart J.; NORVIG, Peter. Inteligência artificial. Rio de Janeiro: Elsevier, Campus, 2013. 988 p. ISBN 9788535237016 (broch.).
}{% BIBLIOGRAFIA COMPLEMENTAR
    SILVA, Flávio Soares Corrêa da; FINGER, Marcelo; MELO, Ana Cristina Vieira de. Lógica para computação.
São Paulo, SP: Thomson Learning, 2006. 234 p. ISBN 8522105170 (broch)                                                                                                        
\newline \newline 
HUTH, Michael; RYAN, Mark. . Lógica em ciência da computação: modelagem e argumentação sobre sistemas .
2. ed. Rio de Janeiro: LTC, 2008. 322 p. ISBN 9788521616108 (broch.)
\newline \newline 
ALENCAR FILHO, Edgard de. Iniciação à lógica matemática. São Paulo: Nobel, [2002]. 203 p. ISBN
852130403X (broch)                                                                      
\newline \newline 
ENDERTON, Herbert B. A mathematical introduction to logic. 2nd ed. San Diego, California:
Harcourt/Academic Press, c2001. 317 p. ; ISBN 0122384520 (enc.)  
\newline \newline 
GERSTING, Judith L. Fundamentos matemáticos para a ciência da computação: um tratamento moderno
de matemática discreta . 5. ed. Rio de Janeiro: Livros Técnicos e Científicos, 2004. 597 p. ISBN 8521614225
(broch.).
}


\DetalhesDisciplina{
    QXD0069 - Segurança
}{% EMENTA
    Ameaças. Segurança como atributo qualitativo de projeto de software.
Autenticação. Autorização. Integridade. Confidencialidade. Criptografia (chaves
simétricas e assimétricas). Infraestrutura de chaves públicas brasileiras
(ICP-Brasil). Certificados digitais. Assinaturas digitais. Desenvolvimento de
software seguro. Noções de auditoria de sistemas. Norma NBR 27002.
}{% BIBLIOGRAFIA BÁSICA
    IMONIANA, Joshua Onome. Auditoria de sistemas de informação. 2. ed. São
Paulo: Atlas, 2008. 207 p. ISBN 9788522450022 (broch.).
\newline \newline 
STALLINGS, William. Criptografia e segurança de redes: princípios e práticas.
4. ed. São Paulo: Pearson/ Prentice Hall, 2008. 492 p. ISBN 9788576051190
(broch.).
\newline \newline 
BEAL, Adriana. Segurança da informação: princípios e melhores práticas para
a proteção dos ativos de informação nas organizações. São Paulo, SP: Atlas,
2008. 175 p. ISBN 9788522440856 (broch.).
}{% BIBLIOGRAFIA COMPLEMENTAR
    DASWANI, Neil; KERN, Christoph; KESAVAN, Anita. Foundations of security:
what every programmer needs to know . Berkeley, Ca: Apress, 2007. 290 p.
(The Expert’s voice in security) ISBN 9781590597842 (broch.).
\newline \newline 
KUROSE, James F.; ROSS, Keith W. Redes de computadores e a Internet:
uma abordagem top-down. 5. ed. São Paulo: Pearson Addison Wesley, 2010.
xxii, 614 p. ISBN 9788588639973 (broch.).
\newline \newline 
NAKAMURA, Emilio Tissato; GEUS, Paulo Lício de. Segurança de redes em
ambientes cooperativos. São Paulo: Novatec, c2007. ISBN 9788575221365
(broch.).
\newline \newline 
STATO FILHO, André. Linux: controle de redes. Florianópolis: Visual Books,
2009. 352 p. ISBN 9788575022443 (broch.).
\newline \newline 
ULBRICH, Henrique Cesar; DELLA VALLE, James. Universidade H4CK3R:
desvende todos os segredos do submundo dos hackers . 6. ed. São Paulo:
Digerati Books, 2009. 348p. (Série Universidade) ISBN 9788578730529
(broch).
\newline \newline 
ASSOCIAÇÃO BRASILEIRA DE NORMAS TÉCNICAS. ABNT NBR ISO/IEC
27001- Tecnologia da informação - técnicas de segurança - sistemas de
gestão de segurança da informação - requisitos. Rio de Janeiro, RJ, 2006. 34
p.
\newline \newline 
ASSOCIAÇÃO BRASILEIRA DE NORMAS TÉCNICAS. ABNT NBR ISO/IEC
27002- Tecnologia da informação - técnicas de segurança - código de prática
para a gestão da segurança da informação. Rio de Janeiro, RJ, 2005. 120 p.
ISBN 9788507006480.
}

\DetalhesDisciplina{
    QXD0043 - Sistemas Distribuídos
}{% EMENTA
    Introdução: caracterização de sistemas de computação distribuída; aplicações distribuídas (caracterização e aspectos de projeto); objetivos básicos de sistemas distribuídos (transparência, abertura, escalabilidade, etc.).
Modelos de sistemas distribuídos: sistemas cliente/servidor e sistemas multicamadas; sistemas peer-to-peer. Objetos distribuídos: interface versus
implementação; objetos remotos; chamadas de métodos remotos (RMI).
Processos em sistemas distribuídos: threads e seu uso em sistemas distribuídos; processos clientes e processos servidores; noções de código móvel
e agentes de software. Sincronização e Coordenação: o conceito de tempo
em sistemas distribuídos; consenso; exclusão mútua distribuída; eleição.
}{% BIBLIOGRAFIA BÁSICA
    MONTEIRO, Eduarda Rodrigues; JUNIOR, Ronaldo C. Mengato; LIMA,
Bruno Santos de et al. Sistemas Distribuídos. São Paulo: SAGAH, 2020.
E-book. Disponível em: https://integrada.minhabiblioteca.com.br/
\#/books/9786556901978. Acesso em: 29 de set. 2022.
\newline \newline 
COULOURIS, G. F.; DOLLIMORE, J.; KINDBERG, T. Sistemas distribuídos: conceitos e projetos. 4 ed. Bookman, 2007. ISBN: 9788560031498.
\newline \newline 
TANENBAUM, A.; STEEN, V. M. Sistemas distribuídos: princípios e
paradigmas. 2 ed. Prentice Hall, 2007. ISBN: 9788576051428.
}{% BIBLIOGRAFIA COMPLEMENTAR
    WHITE, T. Hadoop: the definitive guide. O’Reilly Media, 2009.
\newline \newline 
TANENBAUM, A. S. Sistemas operacionais modernos. 2 ed. Prentice
Hall, 2003.
\newline \newline 
ALONSO, G.; CASATI, F.; KUNO, K.; MACHIRAJU, V. Web Services: Concepts, Architectures and Applications. Springer, 2004. ISBN:
9783540440086
\newline \newline 
NAKAMURA, E. T.; GEUS, P.L. Segurança de redes em ambientes cooperativos. Novatec, 2007. ISBN: 9788575221365.
\newline \newline 
ERL, THOMAS. SOA. Princípios de design de serviço. Prentice Hall, 2009.
ISBN: 9788576051893.
}


\DetalhesDisciplina{
    QXD0076 - Sistemas Multiagentes
}{% EMENTA
    Agentes inteligentes: Conceitos, modelos e arquiteturas; Agentes reativos;
Agentes Deliberativos. Fundamentos da Inteligência Artificial Distribuída e, em
especial, dos Sistemas Multiagentes. Aspectos de comportamento emergente,
comunicação, negociação e coordenação entre agentes. Metodologias de
Desenvolvimento e Arquiteturas de Sistemas Multiagentes. Apresentação de
aplicações existentes e Utilização de Plataformas para o desenvolvimento de
Sistemas Multiagentes.
}{% BIBLIOGRAFIA BÁSICA
    WOOLDRIDGE, Michael J. An introduction to multiagent systems. 2. ed. New
York: J. Wiley \& Sons, 2009. 461 p. ISBN 9780470519462 (broch.).
\newline \newline 
RUSSELL, Stuart J.; NORVIG, Peter. Inteligência artificial. Rio de Janeiro:
Elsevier, Campus, 2013. 988 p. ISBN 9788535237016 (broch.).
\newline \newline 
FOWLER, Martin. UML essencial: um breve guia para a linguagem-padrão de
modelagem de objetos . 3. ed. Porto Alegre: Bookman, 2005. 160 p. ISBN
8536304545 (broch.)
}{% BIBLIOGRAFIA COMPLEMENTAR
    DEITEL, Paul J.; DEITEL, Harvey M. Java: como programar. 8. ed. São Paulo:
Pearson Prentice Hall, 2010. 1144 p. ISBN 9788576055631 (broch.).
\newline \newline 
BLAHA, Michael; RUMBAUGH, James. Modelagem e projetos baseados em
objetos com UML 2. 2.ed. rev. e atual. Rio de Janeiro, RJ: Campus; Elsevier,
2006. xvii, 496 p. ISBN 9788535217537.
\newline \newline 
FIPA. FIPA Communicative Act Library Specification. Foundation for Intelligent
Physical Agents (FIPA), 2002. Disponível em:
<www.fipa.org/specs/fipa00037/XC00037H.pdf>. Acesso em: 24 jan. 2013.
\newline \newline 
SHOHAM, Yoav; LEYTON-BROWN, Kevin. Multiagent systems: algorithmic,
game-theoretic, and logical foundations . New York, NY: Cambridge at the
University Press, 2009. xx, 483 p.
\newline \newline 
MEYER, Bertrand. Object - oriented software construction. 2nd. ed. New
Jersey: Prentice Hall PTR, 1997. 1254p ISBN 0136291554.
}


\DetalhesDisciplina{
    QXD0013 - Sistemas Operacionais
}{% EMENTA
    O histórico, o conceito e os tipos de sistemas operacionais. A estrutura de
sistemas operacionais. Conceito de processo. Gerência de processador:
escalonamento de processos, Concorrência e sincronização de processos.
Alocação de recursos e deadlocks. Gerenciamento de memória. Memória
virtual. Gerenciamento de arquivos. Gerenciamento de dispositivos de
entrada/saída.
}{% BIBLIOGRAFIA BÁSICA
    TANENBAUM, A. S. Sistemas Operacionais Modernos. São Paulo: 3a Edição,
Prentice Hall, 2010.
\newline \newline 
SILBERSCHATZ, Abraham; GAGNE, Greg; GALVIN, Peter; Fundamentos de
Sistemas Operacionais. Rio de Janeiro. 8a Edição. LTC. 2010.
\newline \newline 
OLIVEIRA, Rômulo Silva de; CARISSIMI, Alexandre da Silva; TOSCANI,
Simão Sirineo; UFRGS. Sistemas operacionais. 4. ed. Porto Alegre: Bookman,
2010. 374 p. (Livros didáticos. 11).
}{% BIBLIOGRAFIA COMPLEMENTAR
    FERREIRA, R. E. Linux: guia do administrador do sistema. 2 ed. Novatec,
2008. ISBN: 9788575221778.
\newline \newline 
MACHADO, F.B. Arquitetura de sistemas operacionais. 4 ed. LTC, 2007.
\newline \newline 
SILBERSCHATZ, A.; GALVIN, P.B.; GAGNE, G. Sistemas Operacionais com
Java. 7 ed. ​Campus,​ 2008. ISBN:9788535224061.
\newline \newline 
TANENBAUM, A. S. Sistemas operacionais: projeto e implementação. 3 ed.
Prentice Hall, 2008.
\newline \newline 
MENASCÉ, D.; ALMEIDA, V. Planejamento de capacidade para serviços na
web. ​Campus​, 2002.
}


\DetalhesDisciplina{
    QXD0162 - Sociedade, Culturas e Tecnologia
}{% EMENTA
    A Sociedade da informação. Tecnologias e sociabilidade. As tecnologias da
informação e a sociedade: aspectos econômicos, polítíicos e sociais. Internet:
informação, diversão e interação. Ciberespaço e interatividade. Novas tecnologias, novas mídí ias. Copyleft. Indústria cultural e indústria de jogos
eletrônicos.
}{% BIBLIOGRAFIA BÁSICA
    JENKINS, Henry. FORD, Sam. GREEN, Joshua. Cultura da Conexão: criando
valor e significado por meio da mídí ia propagável. 2014. São Paulo. Ed. Aleph.
ISBN: 8576571625.
\newline \newline 
RECUERO, Raquel. A conversação em rede: comunicação mediada pelo
computador e redes sociais na internet. 2012. Florianópolis. Ed. Sulina. ISBN
852050650X
\newline \newline 
SANTAELLA, Lucia. Comunicação Ubíqí ua. 2013. São Paulo. Ed. Paulus. ISBN
8534936374.
}{% BIBLIOGRAFIA COMPLEMENTAR
    MALDONADO, Tomás. Cultura, sociedade e técnica. São Paulo: Editora
Blucher, 2012. E-book. ISBN 9788521206521. Disponíví el em:
https://app.minhabiblioteca.com.br/\#/books/9788521206521/. Acesso em:
26 set. 2022.
\newline \newline 
BATISTA, Sueli Soares dos S.; FREIRE, Emerson. Sociedade e Tecnologia na Era Digital. São Paulo: Editora Saraiva, 2014. E-book. ISBN 9788536522531.
Disponíví el em:
https://app.minhabiblioteca.com.br/\#/books/9788536522531/. Acesso em:
26 set. 2022.
\newline \newline 
LIPOVETSKY, Gilles. A sociedade da sedução: democracia e narcisismo na
hipermodernidade liberal. Barueri: Editora Manole, 2019. E-book. ISBN
9788520459317. Disponíví el em:
https://app.minhabiblioteca.com.br/\#/books/9788520459317/. Acesso em:
26 set. 2022.
\newline \newline 
VIGLIAR, José Marcelo M. LGPD e a Proteção de Dados Pessoais na
Sociedade em Rede. [Digite o Local da Editora]: Grupo Almedina (Portugal),
2022. E-book. ISBN 9786556276373. Disponíví el em:
https://app.minhabiblioteca.com.br/\#/books/9786556276373/. Acesso em:
26 set. 2022.
\newline \newline 
SHIRKY, Clay. A cultura da participação : criatividade e generosidade no
mundo conectado. Rio de Janeiro: Zahar, 2001. ISBN 9788537805183.
Disponívíel em: <http://lectio.com.br/dashboard/midia/detalhe/1374>
}


\DetalhesDisciplina{
    QXD0046 - Teoria da Computação
}{% EMENTA
    Modelos computacionais universais: máquinas de Turing e funções recursivas. Tese de Church-Turing. Computabilidade. Problemas indecidíveis.
Introdução à complexidade computacional de problemas: complexidade de
tempo e espaço. Tópicos especiais em Teoria da Computação.
}{% BIBLIOGRAFIA BÁSICA
    SIPSER, M. Introdução à teoria da computação. 2 ed. Thompson Learning, 2007. ISBN: 9788522104994.
\newline \newline 
HOPCROFT, John E. Introdução à teoria de autômatos, linguagens e
computação. Rio de Janeiro: Elsevier,2003. ISBN 8535210725
\newline \newline 
DIVERIO, T.A. Teoria da computação: máquinas universais e computabilidade . 3 ed. Bookman, 2011. ISBN:9788577808243.
}{% BIBLIOGRAFIA COMPLEMENTAR
    CARNIELLI, W.; EPSTEIN, R. L. Computabilidade, Funções Computáveis, Lógica e os Fundamentos da Matemática. UNESP,2005. ISBN:
9788571398979.
\newline \newline 
ENDERTON, Herbert B. A mathematical introduction to logic. 2nd ed.
San Diego, California: Harcourt/Academic Press, c2001. xii, 317 p. ; ISBN
0122384520 (enc.) – ISBN 0122384520/9780122384523.
\newline \newline 
PAPADIMITROU, C. Computational complexity. Addison Wesley, 1994.
ISBN: 9780201530827.
\newline \newline 
MENEZES, Paulo Blauth. Linguagens formais e autômatos. 6. ed. Porto
Alegre: Bookman, 2011. 215 p. (Livros didáticos ; n.3 Série Livros Didáticos ; 3) ISBN: 9788577807659
\newline \newline 
BARKER-PLUMMER, Dave; BARWISE, Jon; ETCHEMENDY, John.
Language, proof and logic. CSLI, 2011. xiii 606 p. + 1 CD-ROM + 1
manual (CSLI lecture notes ; v 23) ISBN 9781575866321.
}


\DetalhesDisciplina{
    QXD0152 - Teoria dos Grafos
}{% EMENTA
    Grafos, subgrafos, grafos orientados, famílias de grafos. Árvores, caminhos,
ciclos. Conexidade. Grafos eulerianos. Grafos hamiltonianos.
Emparelhamento. Cliques e Conjuntos Independentes. Coloração de arestas.
Coloração de vértices. Grafos Perfeitos. Grafos planares.
}{% BIBLIOGRAFIA BÁSICA
    GOLDBARG, Marco Cesar; GOLDBARG, Elizabeth. ​Grafos: ​conceitos,
algoritmos e aplicações. Rio de Janeiro, RJ: Elsevier, 2012. 622 p. ISBN
9788535257168 (broch.).
\newline \newline 
BOAVENTURA NETTO, Paulo Oswaldo. ​Grafos: ​teoria, modelos, algoritmos.
5. ed. rev. e ampl. São Paulo, SP: Edgard Blücher, 2012. xiii, 310 p. ISBN
9788521206804 (broch.).
\newline \newline 
ROSEN, Kenneth H. ​Matemática discreta e suas aplicações. ​6. ed. São
Paulo: McGraw-Hill, c2009. xxi, 982 p. ISBN 9788577260362 (broch.).
}{% BIBLIOGRAFIA COMPLEMENTAR
    GERSTING, Judith L. Fundamentos matemáticos para a ciência da
computação: um tratamento moderno de matemática discreta . 5. ed. Rio de
Janeiro: Livros Técnicos e Científicos, c2004. xiv, 597 p. ISBN 8521614225
(broch.).
\newline \newline 
LIPSCHUTZ, Seymour; LIPSON, Marc. Teoria e problemas de matemática
discreta. 2. ed. Porto Alegre, RS: Bookman, 2004. 511p. (Coleção Schaum)
ISBN 9788536303611.
\newline \newline 
CORMEN, Thomas H. Algoritmos: teoria e prática. Rio de Janeiro: Elsevier,
2002. xvii , 916 p. ISBN 8535209263 (broch.).
\newline \newline 
DASGUPTA, Sanjoy; PAPADIMITRIOU, Christos H.; VAZIRANI, Umesh.
Algoritmos. São Paulo: McGraw-Hill, c2009. xiv, 320 p. ISBN 9788577260324
(broch.).
\newline \newline 
JUNGNICKEL, Dieter SPRINGERLINK (ONLINE SERVICE). Graphs, Networks
and Algorithms. Springer e-books Berlin, Heidelberg: Springer-Verlag Berlin
Heidelberg, 2008. (Algorithms and Computation in Mathematics, 5) ISBN
9783540727804. Disponível em:
<http://dx.doi.org/10.1007/978-3-540-72780-4>. Acesso em: 21 set. 2010.
}


\DetalhesDisciplina{
    <27> - Teoria dos Jogos
}{% EMENTA
    Introdução à teoria de decisão. Jogos. Equilíbrio de Nash. Jogos de soma zero. Jogos sequenciais. Jogos dinâmicos. Jogos com incerteza.

}{% BIBLIOGRAFIA BÁSICA
    HILLIER, Frederick S.; LIEBERMAN, Gerald J. Introdução à pesquisa operacional. 9. ed. Porto Alegre: AMGH, 2013. xxii, 1005 p. ISBN 978-85-8055-118-1.
\newline \newline 
OSBORNE, Martin J; RUBINSTEIN, Ariel. A course in game theory. Cambridge: Mit Press, c1994. 352 p. ISBN 9780262150415.
\newline \newline 
ROSEN, Kenneth H. Matemática discreta e suas aplicações. 6. ed. São Paulo: McGraw-Hill, c2009. xxi, 982 p. ISBN 9788577260362 (broch.).
}{% BIBLIOGRAFIA COMPLEMENTAR
    Prisoner’s Dilemma: John Von Neumann, Game Theory and the Puzzle of the
Bomb, by William Poundstone (1994 – MIT Press)
\newline \newline 
IEZZI, Gelson; MURAKAMI, Carlos. Fundamentos de matemática elementar: 1: conjuntos, funções. 9. ed. São Paulo: Atual, 2013. 410 p. ISBN 978-85-357-1680-1.
\newline \newline 
COPPIN, Ben. Inteligencia artificial. Rio de Janeiro, RJ: LTC, 2010. 636 p. ISBN 9788521617297.
\newline \newline 
RUSSELL, Stuart J.; NORVIG, Peter. Inteligência artificial. Rio de Janeiro: Elsevier, Campus, 2013. 988 p. ISBN 9788535237016 (broch.).
\newline \newline 
FIANI, Ronaldo. Teoria dos Jogos. 4. ed. Rio de Janeiro: GEN Atlas, 2015. E-book. p.i. ISBN 9788595156388. Disponível em: https://app.minhabiblioteca.com.br/reader/books/9788595156388/. Acesso em: 06 jun. 2025.
}


\DetalhesDisciplina{
    <28> - Tópicos Especiais em Aprendizado de Máquina
}{% EMENTA
    Estudos de tópicos especiais abordando Aprendizado Profundo OU Aprendizado por Reforço OU Aprendizado Não Supervisionado OU Aprendizado Auto-Supervisionado OU Modelos generativos OU Interpretação de Modelos de Aprendizado de Máquina OU Aprendizado de Máquina Automatizado OU Grandes Modelos de Linguagem OU MLOps
}{% BIBLIOGRAFIA BÁSICA
    COPPIN, Ben. Inteligência artificial. Rio de Janeiro: LTC, c2010. xxv, 636 p. ISBN 9788521617297.
HAYKIN, Simon. 
\newline \newline 
HAYKIN, Simon S.. Redes neurais: princípios e prática. 2. ed. Porto Alegre: Bookman, 2001. xvii, 900 p.
ISBN 9788573077186.
\newline \newline 
RUSSELL, Stuart J.; NORVIG, Peter. Inteligência artificial. Rio de Janeiro: Elsevier, Campus, 2004. 1021
p. ISBN 85-3521-177-2.

}{% BIBLIOGRAFIA COMPLEMENTAR
    LARSON, Ron; FARBER, Betsy. Estatística aplicada. 4. ed. São Paulo, SP: Pearson/ Prentice Hall, 2010. xiv,637 p. ISBN 9788576053729 (broch.).
\newline \newline 
DASGUPTA, Sanjoy; PAPADIMITRIOU, Christos H.; VAZIRANI, Umesh. Algoritmos. São Paulo:
McGraw-Hill, c2009. xiv, 320 p. ISBN 978-85-7726-032-4.
\newline \newline 
HARIOM, Tatsat,; SAHIL, Puri,; BRAD, Lookabaugh,. Blueprints de aprendizado de máquina e ciência de dados para finanças: desenvolvendo desde estratégias de trades até robôs Advisors com Python. Rio de Janeiro: Editora Alta Books, 2024. E-book. p.i. ISBN 9788550821726. Disponível em: https://app.minhabiblioteca.com.br/reader/books/9788550821726/. Acesso em: 30 mai. 2025.
\newline \newline 
HUTTER, Frank; KOTTHOFF, Lars; VANSCHOREN, Joaquin. Automated machine learning: methods, systems, challenges. Springer Nature, 2019. Dispnível em: https://link.springer.com/book/10.1007/978-3-030-05318-5 Acesso en: 30 mai. 2025.
\newline \newline 
AWAD, Mariette; KHANNA, Rahul. Efficient learning machines: theories, concepts, and applications for engineers and system designers. Berkeley, CA: Apress, 2015. DOI: https://doi.org/10.1007/978-1-4302-5990-9. 
}

\DetalhesDisciplina{
    <29> - Tópicos Especiais em Ciência de Dados
}{% EMENTA
    Tópicos especiais em Sistemas de Bancos de Dados OU Processamento de Dados OU Governança de Dados OU  Visualização de dados OU Mineração de Dados OU Engenharia de Dados e Big Data OU Detecção de Anomalias e Dados raros 
}{% BIBLIOGRAFIA BÁSICA
    TAN, Pang-Ning; STEINBACH, Michael; KUMAR, Vipin. Introdução ao DATAMINING:/ mineração de dados. Rio de Janeiro: Ciência Moderna, 2009. xxi, 900 p. ISBN 978-85-7393-761-9.
\newline \newline 
CARVALHO, André C. P. L. F de; MENEZES, Angelo G.; BONIDIA, Robson P. Ciência de Dados - Fundamentos e Aplicações. Rio de Janeiro: LTC, 2024. E-book. p.Capa. ISBN 9788521638766. Disponível em: https://app.minhabiblioteca.com.br/reader/books/9788521638766/ . Acesso em: 30 mai. 2025.
\newline \newline 
BEHRMAN, Kennedy R. Fundamentos de Python para ciência de dados. Porto Alegre: Bookman, 2023. E-book. p.i. ISBN 9788582605974. Disponível em: https://app.minhabiblioteca.com.br/reader/books/9788582605974/ . Acesso em: 30 mai. 2025.

}{% BIBLIOGRAFIA COMPLEMENTAR
    LARSON, Ron; FARBER, Elizabeth. Estatística aplicada. 4. ed. São Paulo: Pearson Prentice Hall, 2010.
xiv,637 p. ISBN 978-85-7605-372-9
\newline \newline 
MILANI, Alessandra M P.; SOARES, Juliane A.; ANDRADE, Gabriella L.; et al. Visualização de Dados. Porto Alegre: SAGAH, 2020. E-book. p.159. ISBN 9786556900278. Disponível em: https://integrada.minhabiblioteca.com.br/reader/books/9786556900278/ . Acesso em: 03 jun. 2025.
\newline \newline 
HAYKIN, Simon S.. Redes neurais: princípios e prática. 2. ed. Porto Alegre: Bookman, 2001. xvii, 900 p. ISBN 9788573077186.
\newline \newline 
RAMAKRISHNAN, Raghu; GEHRKE, Johannes. Database management systems. 3rd ed. Boston:
McGraw-Hill, 2003. xxxii, 1065 p. ISBN 13:978-0-07-246563-1.
\newline \newline 
MACHADO, Felipe Nery Rodrigues. Tecnologia e projeto de data warehouse: uma visão multidimensional.
5. ed. São Paulo, SP: Érica, 2010. 314 p. ISBN 9788536500126 (broch.).
\newline \newline 
CURRY, Edward et al. Technologies and Applications for Big Data Value. In: Technologies and Applications for Big Data Value. Cham: Springer International Publishing, 2022. p. 1-15. DOI: https://doi.org/10.1007/978-3-030-78307-5. Acesso em 3 de jun, 2025.
}


\DetalhesDisciplina{
    <30> - Tópicos Especiais em Filosofia, Ética e Humanidades
}{% EMENTA
    Conceito de ética;
O direito autoral no treinamento de IAs;
Relação entre a IA e a sociedade;
A ética do desenvolvedor e a ética da IA;
Consciência e inteligência;
O problema difícil da Inteligência Artificial.
}{% BIBLIOGRAFIA BÁSICA
    ALMEIDA, G. A.; CHRISTMANN, M. O. Ética e direito: uma perspectiva integrada.3 ed. Atlas, 2009. ISNN:9788522455072
\newline \newline 
MASIERO, P. C. Ética em computação. EDUSP, 2008 ISBN: 9788531405754
\newline \newline 
BARGER, R. N. Ética na computação: UMA ABORDAGEM BASEADA EM CASOS. LTC, 2011. ISBN: 9788521617761
}{% BIBLIOGRAFIA COMPLEMENTAR
    LIMBERGER, T. O direito à intimidade na era da informática. Livraria do Advogado, 2007.
\newline \newline 
CARBONI, G. C. O direito do autor na multimídia. Quartier Latim, 2003.
\newline \newline 
PAESANI, L. O direito na sociedade da informação III. Atlas, 2013.
\newline \newline 
LEVY, P. Tecnologias da inteligência: o futuro do pensamento na era da informática . 2. ed. Rio de Janeiro: Editora 34, 2010.
\newline \newline 
EPSTEIN, Richard G. The case of the killer robot: stories about the professional, ethical, and societal dimensions of computing. New York, NY: John Wiley \& Sons, 1997.
}


\DetalhesDisciplina{
    <31> - Tópicos Especiais em Inteligência Artificial
}{% EMENTA
    Estudos sobre Busca OU Otimização OU Problemas de Satisfação de Restrições OU Representação do Conhecimento OU Planejamento Automatizado
OU Raciocínio sob Incerteza OU  Processamento da Linguagem Natural OU Visão Computacional OU Robótica OU Sistemas Multiagentes.
}{% BIBLIOGRAFIA BÁSICA
    RUSSELL, Stuart J.; NORVIG, Peter. Inteligência Artificial: Uma Abordagem Moderna. 4. ed. Rio de Janeiro: GEN LTC, 2022. E-book. p.Capa. ISBN 9788595159495. Disponível em: https://app.minhabiblioteca.com.br/reader/books/9788595159495/. Acesso em: 22 mai. 2025.
\newline \newline 
COPPIN, B Inteligência artificial. LTC, 2010. ISBN: 9788521617297
\newline \newline 
MARTINS, Júlio S.; LENZ, Maikon L.; SILVA, Michel Bernardo Fernandes Da; et al. Processamentos de Linguagem Natural. Porto Alegre: SAGAH, 2020. E-book. p.Capa. ISBN 9786556900575. Disponível em: https://app.minhabiblioteca.com.br/reader/books/9786556900575/. Acesso em: 23 mai. 2025.
}{% BIBLIOGRAFIA COMPLEMENTAR
    Huth, Michael, and Mark Ryan. ""Lógica em ciência da computação: modelagem e argumentação sobre sistemas."" Rio de Janeiro, 2a edição edition. tradução e revisão técnica Valéria de Magalhães Iório (2008) 
\newline \newline 
SZELISKI, Richard. Computer Vision: Algorithms and Applications. Springer, 2010. ISBN 1848829345
\newline \newline 
BISHOP, CHRISTOPHER M.; Pattern Recognition and Machine Learning. SPRINGER VERLAG, 2006. (ISBN: 0387310738)
\newline \newline 
ROMERO, Roseli Aparecida F.; PRESTES, Edson; OSÓRIO, Fernando et al. Robótica Móvel. São Paulo: LTC, 2014. E-book. Disponível em: https://integrada.minhabiblioteca.com.br/\#/books/
978-85-216-2642-8. Acesso em: 26 de set. 2022.
\newline \newline 
SHOHAM, Y. Multiagent systems: algorithms, game theoretic. Cambridge University, 2009. ISBN: 9780521899437
}


\DetalhesDisciplina{
    <32> - Tópicos Especiais em Matemática Computacional
}{% EMENTA
    Programação linear ou Programação Inteira ou Pesquisa Operacional ou Cálculo Numérico e estudo de artigos e outras publicações atuais.
}{% BIBLIOGRAFIA BÁSICA
    BARROSO, Leônidas Conceição et al. Cálculo numérico: (com aplicações). 2.
ed. São Paulo, SP: Harbra, c1987. 367 p. ISBN 8529400895 (broch.).
\newline \newline 
GOLDBARG, Marco Cesar. Otimização combinatória e programação linear:
modelos e algoritmos. Rio de Janeiro: Elsevier: ​Campus,​ 2005. xvi, 518 p. :
ISBN 9788535215205 (broch.)
\newline \newline 
HILLIER, FREDERICK S.; LIEBERMAN, GERALD J. Introdução a Pesquisa
Operacional. MCGRAW HILL. 9a edição. (ISBN: 8580551188)
}{% BIBLIOGRAFIA COMPLEMENTAR
    RUGGIERO, Marcia A. Gomes; LOPES, Vera Lucia da Rocha. Cálculo
numérico: aspectos teóricos e computacionais . 2. ed. Pearson, c1997. ISBN
8534602042.
\newline \newline 
PASSOS, Eduardo José Pedreira Franco dos. Programação linear como
instrumento da pesquisa operaciona: Eduardo José Pedreira Franco dos
Passos. São Paulo, SP: Atlas, 2008. xii, 451p. ISBN 9788522448395 (broch.).
\newline \newline 
GERSTING, Judith L. Fundamentos matemáticos para a ciência da
computação: um tratamento moderno de matemática discreta . 5. ed. Rio de
Janeiro: Livros Técnicos e Científicos, c2004. xiv, 597 p. ISBN: 8521614225.
\newline \newline 
CHENEY, Ward; KINCAID, David (David Ronald). Numerical mathematics and
computing. 3rd. ed. Pacific Grove, CA: Books/Cole, c1994. 578p. ISBN
0534201121
\newline \newline 
SPERANDIO, Décio; MENDES, João Teixeira; SILVA, Luiz Henry Monken e.
Cálculo numérico: características matemáticas e computacionais dos métodos
numéricos. São Paulo, SP: Prentice Hall, 2003. ix, 354 p. ISBN 8587918745
(broch.).
}


\DetalhesDisciplina{
    QXD0182 - Visão Computacional
}{% EMENTA
    Introdução à visão computacional. Ferramentas de apoio. Formação da
imagem, dispositivos de captura e representação. Cor e textura.
Pré-processamento e Filtros. Segmentação. Rastreamento. Reconhecimento e
Classificação. Avaliação de desempenho de algoritmos de visão 
computacional. Aplicações e tópicos avançados.
}{% BIBLIOGRAFIA BÁSICA
    GONZALEZ, Rafael C.; WOODS, Richard E. Processamento digital de
imagens. 3. ed. São Paulo, SP: Pearson Education do Brasil, 2010. xv,624 p.
ISBN 9788576054016.
\newline \newline 
CONCI, A.; AZEVEDO, E.; LETA, F. Computação Gráfica, Volume 2 -
Processamento e Análise de Imagens Digitais. Publicado por Elsevier. 2007.
ISBN 97885352232193
\newline \newline 
SZELISKI, Richard. Computer Vision: Algorithms and Applications. Springer,
2010. ISBN 1848829345
\newline \newline 
PRINCE, Simon J. Computer Vision - Models, Learning and Inference. William
Morrow, 2012. ISBN 1107011795
\newline \newline 
BRADSKI, Gary; KAEHLER, Adrian. Learning OpenCV: Computer Vision in
C++ with the OpenCV Library. 2 ed. O’Reilly Media, 2012. ISBN 1449314651
}{% BIBLIOGRAFIA COMPLEMENTAR
    SARFRAZ, M SPRINGERLINK (ONLINE SERVICE). Interactive Curve
Modeling : With Applications to Computer Graphics, Vision and Image
Processing . Springer e-books London: Springer-Verlag London Limited, 2008.
ISBN 9781846288715. Disponível em :
<http://dx.doi.org/10.1007/978-1-84628-871-5>. Acesso em : 21 set. 2010.
\newline \newline 
HUTCHISON, David; ELMOATAZ, Abderrahim; KANADE, Takeo; KITTLER,
Josef; KLEINBERG, Jon M; LEZORAY, Olivier; MAMMASS, Dris; MATTERN,
Friedemann; MITCHELL, John C; NAOR, Moni; NIERSTRASZ, Oscar;
\newline \newline 
NOUBOUD. Image and Signal Processing : 3rd International Conference,
ICISP 2008 Cherbourg-Octeville, France, July 1-3, 2008 Proceedings . Springer
eBooks Berlin, Heidelberg: Springer-Verlag Berlin Heidelberg, 2008. (Lecture
Notes in Computer Science, 5099) ISBN 9783540699057. Disponível em :
<http://dx.doi.org/10.1007/978-3-540-69905-7>. Acesso em : 21 set. 2010.
\newline \newline 
CONCI, A.; AZEVEDO, E.; Computação Gráfica, Volume 1 - Geração de
Imagens. Publicado por Elsevier. 2003. ISBN: 9788535212525, 384 páginas
\newline \newline 
GOMES, J. M.; VELHO, L. Fundamentos de computação gráfica. IMPA. 2008.
ISBN: 8524402008
\newline \newline 
LIDWELL, W.; HOLDEN, K.; BUTLER, J. Princípios universais do design. Porto
Alegre: Bookman. 2010. ISBN 9788577807383
\newline \newline 
RUSS, John C. The Image Processing Handbook. Taylor \& Francis. 6 ed.
2010. ISBN 1439840458
\newline \newline 
PARKER, J. R. Algorithms for Image Processing and Computer Vision. John
Wiley. 2 ed. 2010. ISBN 0470643854
\newline \newline 
ANGEL, E. Interactive Computer Graphics: A Top-Down Approach Using
OpenGL. Edition: 5th. Published by Addison-Wesley. 2009. ISBN-10:
0321535863, ISBN-13: 9780321535863, 864 pages
}


\DetalhesDisciplina{
    <33> - Visualização de Dados
}{% EMENTA
    Introdução à Visualização de Dados. Marcas e canais. Abstração de Dados. Abstração de Tarefas. Visualização de dados tabulares. Visualização de dados espaciais. Visualização de redes e árvores. Mapeamento de cor. Estratégias para lidar com´complexidade em visualizações.
}{% BIBLIOGRAFIA BÁSICA
    MILANI, Alessandra M P.; SOARES, Juliane A.; ANDRADE, Gabriella L.; et al. Visualização de Dados. Porto Alegre: SAGAH, 2020. E-book. p.159. ISBN 9786556900278. Disponível em: https://integrada.minhabiblioteca.com.br/reader/books/9786556900278/ . Acesso em: 03 jun. 2025.
\newline \newline 
KYRAN, Dale,. Visualização de dados com Python e JavaScript: raspe, limpe, explore e transforme seus dados – Tradução da 2ª edição. Rio de Janeiro: Editora Alta Books, 2024. E-book. p.i. ISBN 9788550821801. Disponível em: https://integrada.minhabiblioteca.com.br/reader/books/9788550821801/ . Acesso em: 03 jun. 2025.
\newline \newline 
WILKE, Claus O. Fundamentals of Data Visualization: A Primer on Making Informative and Compelling Figures. 1 ed., O’Reilly, 2019. ISBN-13: 978-1492031086, ISBN-10: 1492031089. Disponível online: https://clauswilke.com/dataviz/  Acesso em: 03 jun. 2025.
}{% BIBLIOGRAFIA COMPLEMENTAR
    
AIGNER, Wolfgang et al. Visualization of time-oriented data. London: Springer, 2011. DOI: : https://doi.org/10.1007/978-1-4471-7527-8 Acesso em: 03 ju. 2025.
\newline \newline 
BURGHARDT, Dirk; DEMIDOVA, Elena; KEIM, Daniel A. Volunteered Geographic Information: Interpretation, Visualization and Social Context. Springer Nature, 2024.DOI:  https://doi.org/10.1007/978-3-031-35374-1 Acesso em: 03 ju. 2025.
\newline \newline 
Starbuck, C. (2023). Data Visualization. In: The Fundamentals of People Analytics. Springer, Cham. DOI: https://doi.org/10.1007/978-3-031-28674-2\_15
\newline \newline 

FERREIRA, Rafael G C.; MIRANDA, Leandro B. A de; PINTO, Rafael A.; et al. Preparação e Análise Exploratória de Dados. Porto Alegre: SAGAH, 2021. E-book. p.Capa. ISBN 9786556902890. Disponível em: https://app.minhabiblioteca.com.br/reader/books/9786556902890/  . Acesso em: 03 jun. 2025.
\newline \newline 
CURRY, Edward et al. Technologies and Applications for Big Data Value. In: Technologies and Applications for Big Data Value. Cham: Springer International Publishing, 2022. p. 1-15. DOI: https://doi.org/10.1007/978-3-030-78307-5 . Acesso em 3 de jun, 2025.
}
		% \apendice{Tabelas ...}
\label{ap:B}
	% 	% \apendice{Códigos-fontes utilizados para...}
% \label{ap:C}

% \lstinputlisting[language=C++,caption={Hello World em C++}]{figuras/main.cpp}


% \begin{lstlisting}[language=Java,caption={Hello World em Java}]
% public class HelloWorld {
% 	public static void main(String[] args) {
% 		System.out.println("Hello World!");
% 	}
% }
% \end{lstlisting}

	% 	% \apendice{\textit{IEEE CEFC 2016}}
% \label{ap:D}

% \textit{Digest} submetido ao \textit{The 17th Biennial Conference on Eletromagnetic Field Computation, Miami FL - NOV 13-16, 2016, USA}.

% %Código fonte para inserir um arquivo em PDF
% \includepdf[pages={-}]{3-pos-textuais/apendices/PID4416093.pdf}
	\imprimiranexos
		% Adicione aqui os anexos do seu trabalho
		\anexo{Manual de Normatização de Atividades Complementares}
\label{an:manual-atividades-complementares}

% Um anexo é um documento que não foi elaborado pelo autor, ou seja, o autor apenas anexa. Anexos podem ser tabelas, mapas, diagramas, \textit{datasheets}, manuais e etc. 



		\anexo{Manual de Normatização de Estágio Supervisionado}
\label{an:manual-estagio-supervisionado}

% O autor pode anexar um \gls{PDF}, traduzido como formato portátil de documento. Veja o código fonte utilizado para anexar o arquivo ``Sikasil.pdf'' que foi colocado dentro da pasta ``anexos'' que por sua vez está dentro da pasta ``elementos-pos-textuais''. Tenha muita atenção na hora de especificar o local do arquivo. Recomenda-se não utilizar caracteres especiais para nomear pastas e, principalmente, arquivos. 

% Pode-se fazer uma descrição sucinta do arquivo anexado.

% %Comando para incluir um arquivo em PDF:
% \includepdf[pages={-}]{3-pos-textuais/anexos/Sikasil.pdf}

		\anexo{Manual de Normatização da Extensão }
\label{an:manual-extensao}
		\anexo{Manual de Normatização de Trabalho de Conclusão de Curso}
\label{an:manual-tcc}
		\anexo{Fluxograma de Integralização Curricular}
\label{an:fluxograma}

\begin{center}
  \includegraphics[height=0.99\textwidth,angle=90]{"Fluxograma IA - 2025-01.jpg"}
\end{center}

\begin{center}
  \includegraphics[height=0.99\textwidth,angle=90]{"Fluxograma IA - 2025-02.jpg"}
\end{center}

% \begin{center}
% \resizebox{.75\textwidth}{!}{\rotatebox{90}{%
% \begin{tikzpicture}[
%   node distance=0.6cm and 1.5cm,
%   box/.style={draw, rounded corners, fill=blue!10, text width=3.7cm, minimum height=2.5cm, align=center},
%   opt/.style={draw, rounded corners, fill=purple!20, text width=3.7cm, minimum height=2.5cm, align=center},
%   label/.style={font=\bfseries, anchor=east}]

% % Labels dos semestres
% \node[label] at ( 1, 2) {Semestre 1};
% \node[label] at ( 6, 2) {Semestre 2};
% \node[label] at (11,2) {Semestre 3};
% \node[label] at (16,2) {Semestre 4};
% \node[label] at (21,2) {Semestre 5};
% \node[label] at (26,2) {Semestre 6};
% \node[label] at (31,2) {Semestre 7};
% \node[label] at (36,2) {Semestre 8};

% % Disciplinas (linha a linha)
% % Linha 1
% \node[box] at ( 0, 0) {<02> \\ Introdu\c{c}\~ao \`a Intelig\^encia Art.};
% \node[box] at ( 5, 0) {QXD0007 \\ Programa\c{c}\~ao Orientada a Obj.};
% \node[box] at (10, 0) {QXD0011 \\ Fundamentos de Banco de Dados};
% \node[box] at (15, 0) {QXD0178 \\ Minera\c{c}\~ao de Dados};
% \node[box] at (20, 0) {<08> \\ \'{E}tica em Inteligência Artificial};
% \node[box] at (25, 0) {<10> \\ Processamento de Dados em Larga Escala};
% \node[box] at (30, 0) {<13> \\ Intelig\^encia Artif. Generativa};
% \node[opt] at (35, 0) {Optativa};

% % Linha 2
% \node[box] at ( 0, -3) {QXD0001 \\ Fundamentos de Programa\c{c}\~ao};
% \node[box] at ( 5, -3) {QXD0010 \\ Estrutura de Dados};
% \node[box] at (10, -3) {QXD0115 \\ Estrutura de Dados Avan\c{c}ada};
% \node[box] at (15, -3) {QXD0041 \\ Projeto e An\'alise de Algoritmos};
% \node[box] at (20, -3) {<06> \\ Aprendizado Supervisionado};
% \node[box] at (25, -3) {QXD0182 \\ Vis\~ao Computacional};
% \node[box] at (30, -3) {<12> \\ Aprendizado por Refor\c{c}o};
% \node[opt] at (35, -3) {Optativa};

% % Linha 3
% \node[box] at ( 0, -6) {QXD0109 \\ Pr\'e-C\'alculo};
% \node[box] at ( 5, -6) {QXD0017 \\ L\'ogica para Computa\c{c}\~ao};
% \node[box] at (10, -6) {QXD0112 \\ Probabilidade e Estat\'istica};
% \node[box] at (15, -6) {<04> \\ Representação do Conhecimento e Racioc\'inio};
% \node[box] at (20, -6) {<05> \\ Aprendizado N\~ao Supervisionado};
% \node[box] at (25, -6) {<11> \\ Processamento de Linguagem Natural};
% \node[box] at (30, -6) {<14> \\ MLOps};
% \node[opt] at (35, -6) {Optativa};

% % Linha 4
% \node[box] at ( 0, -9) {QXD0056 \\ Matem\'atica B\'asica};
% \node[box] at ( 5, -9) {QXD0008 \\ Matem\'atica Discreta};
% \node[box] at (10, -9) {QXD0116 \\ \`Algebra Linear};
% \node[box] at (15, -9) {QXD0019 \\ Engenharia de Software};
% \node[box] at (20, -9) {<07> \\ Busca e Otimiza\c{c}\~ao};
% \node[box] at (25, -9) {<09> \\ Aprendizado Profundo};
% \node[box] at (30, -9) {<15> \\ Planejamento Automatizado};
% \node[opt] at (35, -9) {Optativa};

% % Linha 5
% \node[box] at ( 0, -12) {QXD0162 \\ Sociedade,\\Cultura e Tec.};
% \node[box] at ( 5, -12) {QXD0006 \\ C\'alculo Diferencial e Integral I};
% \node[box] at (10, -12) {<03> \\ Agentes\\Inteligentes};
% \node[box] at (15, -12) {QXD0076 \\ Sistemas\\Multiagentes};
% \node[opt] at (20, -12) {Optativa};
% \node[opt] at (25, -12) {Optativa};
% \node[opt] at (30, -12) {Optativa};
% \node[box] at (35, -12) {<01> \\ Atividades \\ Complementares};

% % Linha 6
% \node[box] at (20, -15) {<19> \\ Projeto\\Extensionista I};
% \node[box] at (25, -15) {<20> \\ Projeto\\Extensionista II};
% \node[box] at (30, -15) {QXD0110 \\ Projeto de Pesq. Científico-Tec.};
% \node[box] at (35, -15) {<16> \\ Est\'agio Curricular \\ Supervisionado};

% % Linha 7
% \node[box] at (30, -18) {<17> \\ Trabalho de\\Concl. de Curso I};
% \node[box] at (35, -18) {<18> \\ Trabalho de\\Concl. de Curso II};

% \end{tikzpicture}
% }}
% \end{center}

		
	\imprimirindice

\end{document}